\documentclass[11pt]{article}

\usepackage[margin=1.02in]{geometry}
\usepackage{amsmath,amssymb,amsthm,mathtools,mathrsfs}
\usepackage{bm}
\usepackage{enumitem}
\usepackage{booktabs,array,tabularx,longtable}
\usepackage{xcolor}
\usepackage{microtype}
\usepackage[colorlinks=true,linkcolor=blue,citecolor=blue,urlcolor=blue]{hyperref}
\hypersetup{bookmarksdepth=2}
\usepackage{aliascnt}
\usepackage[nameinlink,capitalise]{cleveref}
\usepackage{tikz-cd}
\setlength{\emergencystretch}{4em}
\raggedbottom
\hbadness=10000

% -----------------------------------------------------------------------------
% theorem environments
% -----------------------------------------------------------------------------
\newtheorem{theorem}{Theorem}[section]
\numberwithin{equation}{section}

\newaliascnt{proposition}{theorem}
\newtheorem{proposition}[proposition]{Proposition}
\aliascntresetthe{proposition}

\newaliascnt{lemma}{theorem}
\newtheorem{lemma}[lemma]{Lemma}
\aliascntresetthe{lemma}

\newaliascnt{corollary}{theorem}
\newtheorem{corollary}[corollary]{Corollary}
\aliascntresetthe{corollary}

\theoremstyle{definition}
\newaliascnt{definition}{theorem}
\newtheorem{definition}[definition]{Definition}
\aliascntresetthe{definition}

\newaliascnt{construction}{theorem}
\newtheorem{construction}[construction]{Construction}
\aliascntresetthe{construction}

\newaliascnt{assumption}{theorem}
\newtheorem{assumption}[assumption]{Assumption}
\aliascntresetthe{assumption}

\theoremstyle{remark}
\newaliascnt{remark}{theorem}
\newtheorem{remark}[remark]{Remark}
\aliascntresetthe{remark}

\crefname{theorem}{Theorem}{Theorems}
\Crefname{theorem}{Theorem}{Theorems}
\crefname{proposition}{Proposition}{Propositions}
\Crefname{proposition}{Proposition}{Propositions}
\crefname{lemma}{Lemma}{Lemmas}
\Crefname{lemma}{Lemma}{Lemmas}
\crefname{corollary}{Corollary}{Corollaries}
\Crefname{corollary}{Corollary}{Corollaries}
\crefname{definition}{Definition}{Definitions}
\Crefname{definition}{Definition}{Definitions}
\crefname{construction}{Construction}{Constructions}
\Crefname{construction}{Construction}{Constructions}
\crefname{assumption}{Assumption}{Assumptions}
\Crefname{assumption}{Assumption}{Assumptions}
\crefname{remark}{Remark}{Remarks}
\Crefname{remark}{Remark}{Remarks}

% -----------------------------------------------------------------------------
% notation
% -----------------------------------------------------------------------------
\newcommand{\N}{\mathbb N}
\newcommand{\Z}{\mathbb Z}
\newcommand{\R}{\mathbb R}
\newcommand{\C}{\mathbb C}
\newcommand{\one}{\mathbf 1}
\newcommand{\id}{\mathrm{id}}
\newcommand{\dd}{\mathrm d}
\newcommand{\Tr}{\operatorname{Tr}}
\newcommand{\tr}{\operatorname{tr}}
\newcommand{\rank}{\operatorname{rank}}
\newcommand{\Ker}{\operatorname{Ker}}
\newcommand{\Ran}{\operatorname{Ran}}
\newcommand{\Span}{\operatorname{span}}
\newcommand{\Hom}{\operatorname{Hom}}
\newcommand{\supp}{\operatorname{supp}}
\newcommand{\spec}{\operatorname{spec}}
\newcommand{\diag}{\operatorname{diag}}
\newcommand{\norm}[1]{\lVert#1\rVert}
\newcommand{\abs}[1]{\lvert#1\rvert}
\newcommand{\ip}[2]{\left\langle #1,#2\right\rangle}
\newcommand{\Grand}{\mathrm{Grand}}
\newcommand{\GT}{\boldsymbol{\mathbb T}_{\Grand}}
\newcommand{\GTX}{\boldsymbol{\mathbb T}_{\Grand,X}}
\newcommand{\Hist}{\mathcal A^{\mathrm{hist}}}
\newcommand{\Sym}{\operatorname{Sym}}
\newcommand{\Alt}{\operatorname{Alt}}
\newcommand{\Cl}{\mathrm{Cl}}
\newcommand{\HS}{\mathrm{HS}}
\newcommand{\op}{\mathrm{op}}
\newcommand{\CP}{\mathrm{CP}}
\newcommand{\UCP}{\mathrm{UCP}}
\newcommand{\vol}{\mathrm{vol}}
\newcommand{\Ric}{\mathrm{Ric}}
\newcommand{\Ein}{\mathrm{Ein}}
\newcommand{\Krein}{\mathcal K}
\newcolumntype{L}[1]{>{\raggedright\arraybackslash}p{#1}}


\title{\textbf{Renewal Geometry and the Emergence of Lorentzian Spacetime}}

\author{Aurelien Pelissier$^{1,\dagger}$\\
{\small $^{\dagger}$Correspondence: \href{mailto:aurelien.pelissier.38@gmail.com}{aurelien.pelissier.38@gmail.com}}}
\date{}

\begin{document}
\maketitle

\begin{abstract}
We develop a finite operational framework in which Lorentzian spacetime and vacuum Einstein dynamics are reconstructed from renewal data. Quotienting histories by equality of their conditional futures yields a canonical predictive carrier, a positive history algebra, and an intrinsic word-Gram hierarchy; together these define the \emph{Renewal Grand Tensor}, from which relational geometry is reconstructed. Within specified finite relational classes, two complementary selection principles identify three spatial dimensions, while one additional protected signed direction yields Lorentzian signature. Calibrated renewal statistics then reconstruct the spatial metric, lapse, and shift.
%
A variational structure is reconstructed directly from the finite dynamics. On noncollapsed continuum branches satisfying explicit stationarity, compactness, and curvature-propagation conditions, its local curvature-linear representative is of Palatini--Holst type and the limiting geometry satisfies the vacuum Einstein equation. We also construct a cutoff-uniform nonsymmetric three-dimensional class, stable across an open finite-dimensional family of renewal laws sharing the same Einstein continuum limit.
\end{abstract}


\medskip
\noindent\textbf{Keywords.} Emergent spacetime; Lorentzian geometry; predictive dynamics; renewal processes; variational dynamics; Palatini--Holst gravity; ADM decomposition.
\setcounter{tocdepth}{1}
\begingroup
\tableofcontents
\endgroup
\clearpage


% =============================================================================

\section{Introduction}
\label{sec:introduction}
% =============================================================================

One of the oldest ambitions of fundamental physics is to explain why spacetime
has the structure that it does, rather than to take that structure as the stage
on which all other physics occurs.  General relativity transformed gravity
from a force propagating through spacetime into the dynamics of spacetime
itself, but it still begins with a differentiable manifold carrying a
Lorentzian metric.  Most subsequent approaches to quantum gravity likewise
replace some part of the classical geometric description while retaining
another part as primitive.  Causal-set theory begins with a locally finite
causal order \cite{bombelli-sorkin,surya}; graph, simplicial, and group-field
approaches begin with combinatorial incidence data whose collective phases are
expected to approximate geometry
\cite{konopkamarkopouloseverini2008,ambjornjurkiewiczloll2009,oriti2009};
and noncommutative geometry replaces the underlying space by spectral and
operator-algebraic data \cite{Connes1994}.  Other programmes seek gravitational
dynamics from thermodynamics, entropy, or quantum information
\cite{jacobson1995,verlinde2011,vanraamsdonk2010}.  These approaches differ
profoundly in their microscopic assumptions, yet they share a central
difficulty: the mathematical language used to reconstruct spacetime often
already contains some analogue of geometry, locality, causal order, dimension,
or temporal structure.

The question pursued here is more primitive.  We ask whether the structures
normally used to define spacetime can instead arise from the organization of
\emph{predictive histories}.  The microscopic object is not a point of a
manifold, a simplex, a causal element, or a prescribed spectral triple.  It is
an operational renewal event: an intervention occurs, a finite outcome is
registered, some distinctions remain capable of influencing later operations,
and the process may subsequently renew.  At this level a history is simply a
sequence of operational events.  It has no assigned position, distance, metric,
connection, curvature, or spacetime dimension, and its ordering is not yet
interpreted as motion through a pre-existing geometry.

\paragraph{Prediction before geometry.}
This choice of starting point is motivated by a simple observation.  Physical
description is never sensitive to an arbitrary microscopic history in its full
detail.  What matters operationally is which distinctions in that history can
still affect a future experiment.  Two past histories that produce exactly the
same statistics under every admissible continuation are therefore physically
indistinguishable for all subsequent purposes.  Identifying such histories is
a predictive rather than a geometric reduction.  It is closely related in
spirit to causal-state constructions in computational mechanics, where pasts
are quotiented by equality of their conditional futures
\cite{shalizicrutchfield,TraversCrutchfield2014}, and to continuous-time
predictive representations of stochastic processes
\cite{MarzenCrutchfield2017,JurgensCrutchfield2021}.  The difference is that
the present construction retains the complete operational intervention data
and the positive process structure required for the subsequent algebraic and
geometric reconstruction.

Because future equivalence is preserved under admissible continuation, the
quotient is not merely a compressed description of the past.  It carries a
canonical successor structure and therefore becomes the minimum persistent
predictive state of the process.  Persistence is thus reconstructed from the
organization of future distinguishability rather than supplied as an
independent state space, graph adjacency, or background coordinate system.

\paragraph{Renewal as a microscopic principle.}
This perspective leads naturally to renewal theory.  Renewal processes are
usually introduced as stochastic models of systems that repeatedly reset,
regenerate, or lose part of their past.  They underlie classical results in
probability and statistical mechanics \cite{feller}, while non-Markovian
renewal processes include heavy-tailed and infinite-mean regimes with
qualitatively different long-time behavior
\cite{garsialamperti1962,erickson1970,caravennadoney2019}.  In conventional
applications, however, renewal occurs \emph{inside} an already given time and
state space.  Waiting times are measured by an external clock, and the process
evolves on a previously specified collection of states.

Here we reverse that logical order.  Renewal is not placed in spacetime;
renewal belongs to the microscopic organization from which spacetime is
reconstructed.  The microscopic law may retain history through waiting-time,
record, or return statistics and need not be Markovian.  A local Markovian
evolution can nevertheless arise after predictive compression and coarse
graining, once unresolved renewal memory has collapsed onto a sufficient
future state.  In this sense the familiar locality of continuum evolution is
naturally interpreted as a \emph{Markovian scaling regime of a more general
renewal dynamics}, rather than as the fundamental form of microscopic
evolution.

There is a useful conceptual analogy with statistical mechanics and
renormalization.  Thermodynamic variables emerge after most microscopic
information has been discarded, while renormalization retains degrees of
freedom that remain relevant at longer scales.  Here the reduction is instead
organized by \emph{future distinguishability}: microscopic information is
discarded precisely when no admissible future can read it.  What survives is
therefore not an arbitrary coarse graining but a canonical predictive quotient.
The geometry developed below is built on that quotient.  Loosely speaking,
spacetime is the stable organization of those relational distinctions that
prediction cannot forget.

\paragraph{Main result.}
The central result is a finite-to-continuum reconstruction and certification
theorem.  Complete conditional future statistics first determine a canonical
minimum predictive carrier and a positive history algebra, whose intrinsic
word-Gram hierarchy together form the \emph{Renewal Grand Tensor}.  The finite relational dimension is selected in two independent ways.  Exact
endpoint--cycle source-rank balance among connected simple regular relations
selects the tetrahedral cell $K_4$ without an alternating-form hypothesis;
within the stronger pair-homogeneous simplex category, minimal faithful
alternating completion selects the same cell.  Its standard module is the
three-dimensional $A_3$ root space.  Calibrated root-race statistics then
reconstruct the ten ADM variables.  The positive Grand Tensor cannot itself
supply a negative direction: an independently occurring protected binary
orientation is an irreducible signed datum, and adjoining one signed direction
to the positive three-dimensional spatial form fixes Lorentzian inertia
$(1,3)$ before a Clifford realization is chosen.

The dynamical part is reconstructed from the complete accepted physical
cylinder rather than imposed as a continuum action.  A marked renewal pressure
recovers the full action slope, finite curl and period tests determine whether
the increments integrate to one scalar common action, and a positive
Fisher--Bregman gap, with its physical test-lift and cutoff constants retained,
certifies stationarity of the accepted field dynamics for that independently
reconstructed action.  Finite Gibbs/Kullback--Leibler identities separately
compare the accepted probability rows with their reconstructed costs.  After resolved memory is removed,
Lorentz covariance and the finite bivector commutant classify the first-order
gravitational score as Palatini--Holst plus determinant volume.  A positive
spatial cut margin supplies regulator-uniform Sobolev, Poincar\'e, counting,
and compact-screen bounds.  Continuous and fully discrete propagation estimates
give a spacetime $L^2$ curvature bound from initial energy and explicit
growth, transfer and differentiated-source budgets; a sharp action-current
comparison can supply the source budget.  Connection curvature can be identified either through a separate
connection compactness certificate with a finite Hodge--temporal sufficient criterion, or, on the literal-link branch, directly from anisotropic spatial connection energy together with the measured electric and magnetic link records; the latter route requires no time derivative of the temporal connection coefficient.  Interface-complete Cartan comparisons remain available for piecewise reconstructions.  On the branch where these finite and cofinal
certificates pass, the weak--strong Palatini limit gives
\begin{equation}
\boxed{G_{\mu\nu}+\Lambda g_{\mu\nu}=0}
\end{equation}
distributionally on the determining test core.  If a required condition fails,
the construction instead returns the corresponding finite or cofinal
obstruction.  The result is therefore an Einstein-or-obstruction theorem,
rather than a claim that every microscopic renewal process flows to general
relativity.  A failed sufficient certificate does not exclude an Einstein
subsequence obtainable by a different argument.

The positive branch is nonempty and is not confined to symmetry-reduced
kinematics.  In addition to the homogeneous and Gowdy sectors described
below, a full three-dimensional local harmonic renewal writer admits a
cutoff-independent nonlinear Cauchy theorem in a fixed Sobolev ball for all
ten metric components and all spatial lattice modes.  More strongly, this is
not an isolated engineered update: the same estimates and Einstein limit hold
uniformly on a fixed open $55$-dimensional family of local differential laws,
with no convergence of the law parameter required along the cutoff sequence.
Around finite realizations of these laws, conditional physical residual
moments and a profile-transport metric define relatively open neighborhoods
of history-dependent transition rows whose successful histories have the same
Einstein limit in probability and, along the declared odd refinement, eventually
almost surely.  On the cofinally consistent initial-action branch, compatible limiting
initial records can also be obtained from convergent finite preparations by solving the original reconstructed action's complete lapse--shift
constraint rows from unconstrained small free tensor records; no continuum
vacuum metric is required as an input to that preparation.

These population results and the general certification theorem have different
logical roles.  \Cref{thm:main-open-3plus1,thm:main-open-law-basin} construct
nonsymmetric $3+1$ classes and reach the Einstein equation by cutoff-uniform
evolution, consistency and finite preparation estimates, whereas
\Cref{thm:main-renewal-einstein} is a recognition theorem for arbitrary
complete gravitational renewal cylinders.  Its five certificates are not the
definition of the constructive population.  Nor is the stable law family
silently identified with the raw exact finite Euler flow of the reconstructed
action.  That stronger route is analyzed independently in
\Cref{thm:main-exact-action-response,thm:supp-exact-initial-layer} and
Appendix~\ref{supp:exact-action-audit}.  The exact-action analysis contains an all-cutoff regular transverse source
block, an exact scalar--trace response decomposition, and source-neutral finite
preparations.  On the amplitude-graded three-site chart the boundary argument is rank-free, and a same-event weak-Poisson certificate gives a five-dimensional weak kernel and the exact rapid split $(8,8,13)$.  Under the complete graded-chart conditions, the unchanged finite action has exactly constrained boundary-selected histories through the certified prepared event with $O(a)$ metric and link curvature on $0\le t\le ca^2$, where $a$ is field amplitude and the spatial cutoff is fixed.  The selector changes only higher preparation coefficients; it is not a forward Cauchy theorem for arbitrary prepared initial records.  Exact rapid and slow normal factorizations on their stated invariant chart give necessary local-residence widths $O(ae^{-\mu_rT/a^3})$ and $O(ae^{-\mu_sT/a})$ for a fixed physical horizon $T$.  On the intermediate scale $t=a\tau$, bounded rescaled rates require a five-component compatibility map on a $297$-dimensional leading critical sheet.  Exact initialization automatically satisfies the three constant-shift/Fredholm rows, while the remaining two initial rows are an explicit projected acceleration defect computable from the unchanged finite action.  For actual Lipschitz bounded-rate slow limits, exact elimination of the two generalized-zero rows yields a quadratic harmonic-slope gate.  On the literal record-preserving branch the harmonic quartic columns vanish, so this slope map factors through the two generalized-zero rows and has Jacobian rank at most two; a fixed cokernel condition precedes a two-equation quadratic root problem.  Positive-time nonlinear propagation of a compatible root branch and cutoff-uniform cofinal control remain open.

The homogeneous isotropic sector gives a complementary exact variational
realization: arbitrary positive renewal variables yield the full lapse-varied
vacuum Friedmann equations, so the flat and de~Sitter rate profiles arise as
stationary solutions rather than being prescribed in advance.  A finite
positive accept/reject process reconstructs the same reduced action directly
from measured trial probabilities and converges quantitatively to the
stationary branch.  Exact conservative time blocking of this quadratic class
preserves the scale-free action--law defect, so it does not generate an open
basin toward the matched surface.  Beyond homogeneity, a nonlinear unpolarized
Gowdy $T^3$ regulator preserves an exact finite momentum invariant and, on its
finite offset alphabet, has pathwise $O(h)$ convergence of the full Riemann
tensor on smooth bounded slabs.  These constructions provide complementary
open, homogeneous, and nonlinear symmetry-reduced populations without
asserting generic microscopic attraction to the certified branch.

The precise operational entrance is intentionally narrower than a claim that
all geometric information is generated from positive probabilities alone.
No predictive state space or persistent relational carrier is postulated
independently: histories are identified exactly when every admissible
continuation has the same conditional law, and the resulting
future-equivalence quotient is the unique minimum persistent predictive
carrier through which all future distinctions factor.  The geometric branch
also retains two kinds of information that this positive predictive quotient
cannot manufacture.  A protected binary orientation is an actual two-valued
operational distinction, but no interpretation as past and future,
positive and negative norm, or Lorentzian signature is assumed microscopically.
Likewise, normalized race probabilities determine only relative rates, so
absolute waiting-time and measure calibrations are required when those rates
are interpreted geometrically.

A further distinction is essential.  The future-equivalence class is derived
as an abstract predictive state, but spacetime reconstruction must physically
read that class at a cut and correlate it on the same history with the
occurring root event, scheduler-separated fast path, orientation, calibration,
and other declared marks.  These data form a minimal faithful readable relation
packet.  The packet does not introduce a second predictive state: its state
coordinate is the already derived future-equivalence class.  Its additional
content is physical occurrence and same-cut readability.  The enrichment is
conservative: forgetting the added marks recovers the original event cylinders
exactly, while future minimization removes every unread refinement.  Thus
predictive persistence is derived, the physical information needed for the
geometric reconstruction must actually occur together, and metric, connection,
curvature, and spacetime dimension remain downstream outputs.

\paragraph{The Renewal Grand Tensor and Lorentzian reconstruction.}
The first major structural output is the object we call the
\emph{Renewal Grand Tensor}.  The name does not refer to a tensor field on a
pre-existing manifold.  Predictively indistinguishable histories are first
identified, producing the unique minimum persistent predictive carrier.
Tomographically complete intervention statistics then reconstruct the finite
positive process tensors, while support-minimal completely positive
realization removes unobservable environmental degrees of freedom
\cite{Stinespring1955,Choi1975}.  The surviving mixed histories generate a
finite $C^*$-algebra, and their intrinsic positive overlaps form a word-Gram
hierarchy.  Together these structures record which distinctions survive into
the future, how renewal operations compose, and how the resulting histories
overlap.  The Grand Tensor is therefore not geometry imposed on a renewal
process; it is the canonical algebraic and positive structure retained by the
process after every future-invisible distinction has been removed.

This places the construction near noncommutative and spectral geometry, but
with a different starting point.  In ordinary spectral geometry an algebra,
Hilbert-space representation, and Dirac operator are supplied as geometric
data \cite{Connes1994}.  Lorentzian variants enrich this structure with an
indefinite inner product, fundamental symmetry, temporal datum, or twist
\cite{strohmaier,franco,franco-eckstein,
vandendungenpaschkerennie2013,vandendungen2016,
DevastatoFarnsworthLizziMartinetti2018,
DevastatoFilaciMartinettiSingh2020}.  Recent work has also examined the
emergence of temporal and Lorentzian structures within twisted
almost-commutative settings \cite{nieuviarts2026a,nieuviarts2026b}.  The
present construction operates one level upstream.  Neither spectral geometry
nor a Krein structure is taken as the initial object.  The positive history
algebra is reconstructed from operational renewal statistics, while the
protected binary orientation becomes an indefinite/Krein structure only after
it acts on the predictively reduced relational carrier.  Its Clifford
realization then distinguishes one timelike direction.

Predictive collapse does more than compress a finite history.  Because future
equivalence is preserved under admissible continuation, the minimum predictive
state propagates canonically from one event to the next and therefore supplies
a persistent relational carrier.  The local dimension is then tested by two
logically distinct finite selectors.  The first uses no alternating structure:
for a connected simple regular relation, exact equality of the faithful
endpoint-source and cycle-source ranks has the unique nontrivial solution
$K_4$.  The second works in the more restrictive pair-homogeneous simplex
category and asks for a nondegenerate alternating temporal--spatial coefficient
pairing together with minimum faithful source innovation; it independently
selects the same $K_4$.  Thus the alternating form is not the sole route to
three spatial dimensions.  It remains a branch axiom of the second selector
and supplies only an antisymmetric coefficient pairing, not a symplectic phase
space, Hamiltonian flow, metric, or Lorentzian sign.  Neither selector is a
universality theorem over arbitrary higher-arity microscopic systems.

The metric appears one layer later.  The relational carrier supports a finite
family of elementary root transitions.  Their occurrence rates and waiting
times provide observable first- and second-order statistics.  The second-order
statistics determine the spatial kinetic form, the directed first moment
determines the shift, and a physical speed density fixes the remaining scale.
These quantities invert uniquely to the ADM spatial metric, lapse, and shift.
The metric coefficients are therefore not microscopic labels but collective
statistical quantities reconstructed from the renewal process.

The geometric clock is likewise not identified with the slower renewal clock.
A nondegenerate continuum bracket requires a scheduler-separated root race
with total rate of order $h^{-2}$.  Its winner and waiting-time statistics
supply the absolute kinetic scale, whereas root-direction probabilities alone
leave the common rate scale undetermined.  The protected binary orientation then acquires its geometric meaning.  The
positive Grand-Tensor and root-race forms remain positive and therefore cannot
produce an indefinite direction.  The orientation Read supplies one independent signed line through the explicit minimal signed-extension
condition; leaving the reconstructed three-dimensional spatial form positive, the minimal nondegenerate signed extension has inertia
$(-,+,+,+)$, up to overall sign convention.  Only after this inertia is fixed
do we choose a minimal Clifford representation.  The gamma matrices therefore
realize a signature already determined by the one-sign extension; they do not
select the signature.  Thus the Lorentzian sign is neither obtained by analytic
continuation of a positive diffusion tensor nor inserted through a Dirac-matrix
ansatz.

\paragraph{From kinematics to Einstein dynamics.}
Producing Lorentzian kinematics is not yet a derivation of gravitational
dynamics.  In particular, prescribing a renewal-rate history whose
reconstructed metric is de~Sitter would establish only realizability of one
Einstein solution.  The dynamical part therefore begins instead from the
complete accepted physical cylinder.  Future refinement reconstructs the
minimum determining-field quotient and the accepted transition kernel acting
on it.  On the same resolved history, a protected vector reward together with
physical duration determines an analytic renewal pressure.  Its first
derivative restores the affine action anchor lost by a centered probability
score, while its higher jets reconstruct the mixed action Gram.  Finite
face-curl and period tests then determine whether these increments integrate
to a single scalar common action and reconstruct that action when they do.

The actually accepted transition is next compared with the update generated
by this independently reconstructed action.  On the interior variational
branch, the Fisher--Bregman gap controls the squared represented first
variation, and explicit physical test-lift bounds retain the cutoff dependence
needed in the continuum passage.  Exact Gibbs/Kullback--Leibler identities
provide a separate comparison of finite probability rows.  The field-action
comparison, rather than a cost fitted after observing the transition law,
supplies the stationarity test.

The gravitational handoff is likewise formulated as a reconstruction and
certification problem.  The protected cycle reward entering this step is
initially a generic measured face reward: no Palatini, Holst,
Einstein--Hilbert, or cosmological Lagrangian is named in the entrance data.
After resolved cut memory is removed, binary renewal subdivision and stable
additivity produce a first-order face functional without assuming a derivative
at a zero face.  Lorentz covariance and the residual bivector commutant then
leave only the identity and internal Hodge star, so only at this stage is the
local first-order gravitational score classified as Palatini--Holst plus
determinant volume.

Connection stationarity yields zero torsion on the spinless branch.  Global
spatial compactness is not introduced as a separate geometric axiom: the
physical endpoint/capacity record determines an exact cut constant, and a
uniform positive margin gives Sobolev, Poincar\'e, low-energy-counting, and
finite-rank compact-screen bounds.  Curvature control is supplied by a separate
positive propagation estimate.  A bounded initial curvature energy together
with regulator-uniform observer/transport and differentiated-source budgets
yields a spacetime $L^2$ curvature bound, also for a fully discrete curvature
record.  Strong coframe convergence supplies the nonlinear product passage;
strong connection convergence and interface-complete consistency separately
identify its curvature.  These are the two analytic inputs to the
weak--strong Palatini passage.  On the positive certificate branch,
the limiting metric variation is
\begin{equation}
\chi\int_K\sqrt{-g}\,
(G_{\mu\nu}+\Lambda g_{\mu\nu})k^{\mu\nu},
\end{equation}
and determining stationarity therefore gives the vacuum Einstein equation
distributionally.

These finite reconstructions assemble into a certification-or-obstruction
alternative.  For a complete gravitational renewal cylinder, the
determining-field quotient, accepted kernel, renewal pressure, common-action
exactness, action gap, spatial cut constant, face linearization, memory short,
connection residual, curvature-propagation budgets, and Cartan compatibility
are finite or cofinally testable.  If the corresponding positive margins
persist, the Einstein continuum follows.  If they do not, the construction
returns the failed object itself: a nonzero common-action curl or period, a
positive Bregman/KL row, a macroscopic cut/potential witness, a
face-linearity or memory residual, a torsion/connection defect, an unbounded
curvature-propagation budget, a Cartan mismatch, or a cutoff compactness
defect.

The homogeneous isotropic sector addresses a narrower question not settled by
the abstract Einstein handoff.  Instead of prescribing a de~Sitter rate
profile, we begin with arbitrary positive root-symmetric variables
$\kappa(t)$ and $\varrho(t)$.  The reconstructed ADM map converts them
bijectively to the scale factor and lapse.  Rewriting the classified vacuum
action in $(\kappa,\varrho)$ and varying both variables before fixing the lapse
yields the Hamiltonian constraint and scale equation, equivalently the full
spatially flat vacuum Friedmann system.  The exponential de~Sitter profile
therefore appears only after the stationary equations have been solved.

The same distinction persists at finite cutoff.  A midpoint variational action
retains independent lapse-weighted intervals, and its stationary recurrence
converges at second order to the continuum flat/de~Sitter branch.  A finite
positive accept/reject law reconstructs this action directly from its proposal
distribution, accepted row, and total acceptance probability.  A measured rank
defect determines whether the local low-noise update and the cosmological
parameter of the reconstructed action agree.  On the compatible branch, finite
stochastic cylinders have vanishing Euler residual and converge quantitatively
to the stationary de~Sitter solution.

The finite evolution mechanism is not confined to homogeneous trajectories.
In the unpolarized Gowdy $T^3$ vacuum class both wave polarizations and a
spatially varying nonlinear frame are retained.  A
source/characteristic/source splitting preserves an exact staggered momentum
invariant, while the finite offset support and shared Hermite readout give
pathwise second-jet control and hence $O(h)$ convergence of the full Riemann
tensor on fixed smooth slabs.  The independently varied unreduced action has
the same vanishing continuum first variations.  This symmetry-reduced theorem
does not imply generic attraction to de~Sitter or unrestricted four-dimensional
regularity, but it provides a nonlinear inhomogeneous class in which full
curvature compactness is derived rather than assumed.

\paragraph{Scope and organization.}
The emergence claim is relative to a precise operational entrance.  Actual
occurrence, complete contextual prediction, protected orientation, and
absolute physical calibration are irreducible rows of the declared finite
language.  For gravitational dynamics one additionally needs the complete
accepted field/action cylinder to occur physically.  Once these records are
present, however, the predictive carrier, determining field, accepted kernel,
action slope, common scalar action, ADM geometry, compact spatial screen, and
stationarity residual are reconstructed or certified rather than supplied as
independent continuum geometry.  Vanishing of the action gap and persistence
of the positive noncollapse and curvature-propagation margins remain genuine
physical conditions that can fail.

The dimensional claim has the same controlled scope.  The $K_4/A_3$ result
does not state that every possible microscopic theory must have three spatial
dimensions.  Rather, two independent finite classifications select the same
cell: endpoint--cycle source-rank balance does so without an alternating-form
assumption, while the alternating/minimal-source simplex selector does so in a
stronger category.  Their agreement removes the alternating form from the role
of sole dimensional selector, while still falling short of universality over
arbitrary higher-arity systems.  Likewise the protected binary orientation is
not a microscopic Lorentzian metric.  The positive operational rows cannot
reconstruct it, as shown by explicit separating models; once the $A_3$ spatial
form is positive, this single independent sign fixes the minimal indefinite
extension to Lorentzian inertia $(1,3)$.  The subsequent Clifford matrices are
a representation of that quadratic form, not the origin of its signature.

The conceptual proposal is therefore that renewal may be more fundamental
than its conventional role in probability theory suggests.  From the usual
point of view, renewal describes memory and resetting \emph{within} time.  From
the present point of view, retention and forgetting in operational history
determine which distinctions survive at all, while predictive equivalence
reconstructs the persistent state on which relational geometry is built.
Effective Markovian locality then appears not as the microscopic starting
principle but as a limiting organization of predictive memory.  In this sense,
the usual hierarchy
\begin{equation}
\boxed{\text{spacetime}\longrightarrow\text{stochastic dynamics}}
\end{equation}
is replaced by
\begin{equation}
\boxed{
\text{renewal dynamics}
\longrightarrow
\text{predictive structure}
\longrightarrow
\text{relational geometry}
\longrightarrow
\text{Lorentzian spacetime}.}
\end{equation}
The claim is not that renewal theory in its classical form already contains
general relativity.  It is that when renewal is treated operationally,
predictively, and with a minimal protected orientation, persistence and
relational structure can be reconstructed before ordinary spacetime geometry
is defined.

Sections~\ref{sec:grand-tensor-main} and~\ref{sec:adm-main} reconstruct the
Renewal Grand Tensor and the calibrated Lorentzian geometry.
Section~\ref{sec:einstein-dynamics} reconstructs the common action,
compares it with the accepted dynamics, and classifies its local
gravity representative. Section~\ref{sec:einstein-limit} supplies the
noncollapse and curvature estimates and proves the general
renewal-to-Einstein certification theorem.
Section~\ref{sec:constructive-main} then constructs a nonsymmetric
cutoff-uniform Einstein class, establishes its law-space robustness, and
states the compatibility with a common chronological action process.
Its final subsection summarizes the separate exact-action question;
the response, boundary-realization and slow-compatibility analyses are
proved in Appendices~\ref{supp:exact-action-audit} and~\ref{supp:exact-action-slow}.
Section~\ref{sec:vacuum-main} treats homogeneous and Gowdy sectors,
including their operational realizations and the conservative-blocking
obstruction. Section~\ref{sec:discussion} discusses interpretation and
open problems. The remaining appendices follow the same order as the
main argument; the appendix guide identifies the proof location of each
principal result.


\section{Predictive reconstruction and the Renewal Grand Tensor}
\label{sec:grand-tensor-main}
% =============================================================================

\subsection{Operational histories and the predictive quotient}

At a finite cutoff $X$, the operational entrance is a finite typed event
protocol with a declared finite horizon, finite intervention and outcome
alphabets, and normalized conditional probabilities for every reachable
history. The remaining horizon is included in the cut type, so admissible
continuations are compared only at compatible cuts. An unbounded-horizon
protocol is also allowed when a finite-index predictive congruence is
separately certified; a finite alphabet alone does not imply this property.  The physical intervention letters include tomographically complete
frames on each exposed finite operator port; their multitime probabilities are
affine under randomization, additive under coarse graining, causal under
deterministic discard, and compatible on common prefixes.  A distinguished
binary orientation is retained as an actual protected operational distinction.
No predictive state space, persistent relational carrier, metric, scheduler,
connection, or spacetime dimension is part of this entrance.

For a reachable history $h$ ending at cut type $x$, let $\mathsf F_{X,x}$ be
the set of all admissible future words within the declared horizon
starting at $x$ (or all finite words on the certified finite-index branch).  Write
$p_X(w\mid h)$ for the corresponding conditional cylinder probability.  Define
\begin{equation}
 h\sim_X h'
 \quad\Longleftrightarrow\quad
 p_X(w\mid h)=p_X(w\mid h')
 \quad\text{for every }w\in\mathsf F_{X,x}.
 \label{eq:future-equivalence-main}
\end{equation}
Histories of different cut type are never identified.  If additional actual
terminal Reads or bounded readout queries are present, they are included in
the future word alphabet; thus \eqref{eq:future-equivalence-main} means equality
of the complete declared future signature.

Future equivalence is a typed right congruence.  Indeed, if $h\sim_Xh'$ and an
event $a:x\to y$ is admissible, then the one-letter future gives
$p_X(a\mid h)=p_X(a\mid h')$.  On a positive-probability branch, the chain rule
then gives
\[
 p_X(v\mid ha)
 =\frac{p_X(av\mid h)}{p_X(a\mid h)}
 =\frac{p_X(av\mid h')}{p_X(a\mid h')}
 =p_X(v\mid h'a)
\]
for every continuation $v$.  Hence $ha\sim_Xh'a$.

\paragraph{The minimum persistent predictive carrier.}

Define
\begin{equation}
 Z_{X,x}^{\min}:=\mathsf H_{X,x}/\!\sim_X,
 \qquad q_{X,x}:\mathsf H_{X,x}\twoheadrightarrow Z_{X,x}^{\min}.
 \label{eq:predictive-state-main}
\end{equation}
For an admissible event $a:x\to y$, the right-congruence property defines
\begin{equation}
 T_{X,a}([h])=[ha],
 \qquad
 \kappa_{X,a}([h])=p_X(a\mid h).
 \label{eq:predictive-update-main}
\end{equation}
Thus the event law itself reconstructs a deterministic predictive state
machine.  If $A_1,A_2,\ldots$ is the realized event sequence,
$H_n=A_1\cdots A_n$, and $Z_n=q_X(H_n)$, then
\begin{equation}
 Z_{n+1}=T_{X,A_{n+1}}(Z_n),
 \qquad
 \mathfrak p_{X,n}^{\min}
   =(Z_n,A_{n+1},Z_{n+1}).
 \label{eq:derived-relation-packet-main}
\end{equation}
The persistent relation is therefore the transport of a future-equivalence
class under event composition, not an independently supplied state variable.

The quotient is universal.  Every reachable deterministic predictive
realization $S_{X,x}$ of the same conditional event law admits a unique
surjection
\begin{equation}
 \pi_{X,x}:S_{X,x}\twoheadrightarrow Z_{X,x}^{\min}
 \label{eq:predictive-universal-main}
\end{equation}
intertwining branch probabilities and updates.  Consequently every declared
downstream readout depending only on the complete future signature factors
uniquely through $Z_X^{\min}$.  Generated ledgers, hidden-state labels, and
other deterministic predictive presentations are therefore representations of
the canonical quotient rather than additional foundational coordinates.

The physical geometric branch considered below requires more than this
predictive reduction.  In particular, same-cut readability of the predictive
class, the protected orientation, and the absolute winner--waiting calibration
of the fast root race are actual physical rows.  They are not determined by
the positive predictive quotient.  The purpose of
\eqref{eq:predictive-state-main}--\eqref{eq:derived-relation-packet-main} is to
remove independently postulated persistence and every future-invisible state
refinement.

The quotient determines the state distinctions that prediction retains.
The geometric construction also needs the corresponding state and its
calibrations to be readable on the same physical history. The next
completion separates this occurrence requirement from state minimization.

\subsection{Readable relational completion}

On a spacetime-capable branch, let $O_{X,n}$ denote the occurring root-bearing
event at the $n$th slow cut, let $\Gamma_{X,n}$ denote the scheduler-separated
fast root path between slow cuts, and let $\Xi_{X,n}$ collect the remaining
same-history physical marks used by the geometric reconstruction, including the
protected orientation and absolute calibrations.  A raw readable relation
packet is
\begin{equation}
 \mathfrak p_{X,n}^{\rm rel}
 =\bigl(
 Z_{X,n}^{\min},\Gamma_{X,n},O_{X,n},\Xi_{X,n},Z_{X,n+1}^{\min}
 \bigr).
 \label{eq:readable-relation-packet-main}
\end{equation}
The state coordinates in \eqref{eq:readable-relation-packet-main} are not new
hidden variables: they are same-cut physical realizations of the predictive
classes already reconstructed in \eqref{eq:predictive-state-main}.  Two raw
packets are identified when every admitted continuation, readout query, and
terminal Read has the same conditional law.  Denote the resulting finite
future-minimal quotient by $\mathfrak P_X^{\rm rel,min}$.

\begin{proposition}[Conservative readable completion]
\label{prop:readable-relational-completion}
Suppose the additional readable marks in
\eqref{eq:readable-relation-packet-main} are attached to the underlying finite
operational process by normalized same-history conditional kernels.  Then:
\begin{enumerate}[label=\textup{(\roman*)},leftmargin=2.8em]
\item forgetting $\Gamma$, $\Xi$, and the explicit same-cut realization of the
      predictive class recovers every cylinder probability and branch map of
      the underlying operational process exactly;
\item future-signature refinement terminates at
      $\mathfrak P_X^{\rm rel,min}$, and every other finite readable packet
      sufficient for the same geometric reconstruction has a unique surjection onto
      it;
\item unread refinements add neither predictive content nor source dimension;
      metric, coframe, connection, curvature, and Einstein data remain
      downstream readout outputs and are not coordinates of the relation
      packet.
\end{enumerate}
Thus physical readability is a conservative enrichment of a derived
predictive carrier, not a second postulated state space.
\end{proposition}

The proof is given in Appendix~\ref{supp:readable-relational-completion}.  This
proposition is the interface between predictive reconstruction and the
geometric branch: the first determines which state distinctions exist, while
the second records which of those distinctions and physical calibrations
actually occur together on one history.

Finite exact reconstruction, its horizon convention and the presentation
category are treated in Appendix~\ref{supp:grand-reconstruction}.
The equality tests concern the supplied probability table, not arbitrary
finite samples of an unknown real-valued law.

\subsection{The Renewal Grand Tensor}

The predictive quotient fixes persistence; the positive multitime table
fixes how the retained operations compose and overlap. These structures
are assembled into one object before any metric is introduced.

The operational multitime table on the physical intervention frames
reconstructs a positive causal process tensor in the quantum-comb sense
\cite{Chiribella2009,Pollock2018}.  Its prefix-support-minimal square-root
purification supplies canonical sequential memories, unique up to
record-intertwining memory unitaries in the purifying category defined in
\Cref{thm:supp-comb-tomography}.  The represented minimal-state projections,
mixed forward maps, and their Hilbert--Schmidt reverses generate a finite
history algebra $\mathcal A_X^{\mathrm{hist}}$.  If $[w]$ denotes the
represented history word and $\widehat\tau_{\mathrm{reg}}$ is the normalized
regular trace, define
\begin{equation}
 \mathbb K_{X,r}(v,w)
 =\widehat\tau_{\mathrm{reg}}\bigl([v]^*[w]\bigr),
 \qquad |v|,|w|\le r.
 \label{eq:word-gram-main}
\end{equation}
The positive structural Renewal Grand Tensor is
\begin{equation}
 \boxed{
 \GTX
 =\left(
 \mathcal A_X^{\mathrm{hist}},
 \{\mathbb K_{X,r}\}_{r\ge0},
 Z_X^{\min}
 \right).}
 \label{eq:grand-tensor-main}
\end{equation}
The term \emph{Grand Tensor} denotes this algebra-valued positive moment
object, not a tensor field on a manifold. Each word-Gram form is a covariant
sesquilinear tensor on the finite word space. Under independent composition
with independently accessible operations, the predictive quotient, history
algebra and word moments factor as in
\Cref{prop:supp-grand-monoidal}; thus the tensor terminology also has a
precise product interpretation on that category.  Finite dimensionality forces the represented word hierarchy to
become flat after a finite depth, at which point one further word layer
reconstructs the complete represented algebra.

For the spacetime branch, \eqref{eq:grand-tensor-main} is used together with
the minimal faithful readable packet $\mathfrak P_X^{\rm rel,min}$ of
\Cref{prop:readable-relational-completion}.  The distinction is deliberate:
$\GTX$ records the canonical positive predictive-history structure, whereas
the readable packet records the actual same-history occurrence of orientation,
calibration, scheduler, and root-path data required to interpret that
structure geometrically.  Forgetting the latter returns the positive
structural tensor exactly.

\begin{theorem}[Canonical predictive Grand-Tensor reconstruction]
\label{mt:grand-tensor}
Let a finite typed operational renewal process have a finite event grammar,
normalized prefix-compatible conditional probabilities, tomographically
complete physical intervention frames, finite operational positivity, causal
trace recursion, and a protected binary orientation as described above. On an unbounded-horizon
branch, assume also a common finite-dimensional sequential realization of
the complete represented history operations. Then:
\begin{enumerate}[label=\textup{(\roman*)},leftmargin=2.8em]
\item future equivalence \eqref{eq:future-equivalence-main} is a typed right
      congruence, and $Z_X^{\min}$ is the unique coarsest reachable
      deterministic predictive realization of the event law;
\item the successor maps \eqref{eq:predictive-update-main} are canonical and
      the persistent event packet \eqref{eq:derived-relation-packet-main} is
      derived from composition rather than independently postulated;
\item the state-conditioned Hankel quotient is the unique reachable,
      future-separated minimum linear predictor and determines the executable
      forward algebra;
\item the physical probability table reconstructs unique positive causal
      prefix combs whose square-root supports determine a canonical
      prefix-support-minimal sequential isometric class, unique up to
      record-intertwining memory unitaries; this is a purifying realization,
      not a minimization over arbitrary mixed-state memory implementations;
\item the contextual history quotient, normalized regular trace, and positive
      word-Gram hierarchy are canonical, and finite flatness reconstructs the
      complete represented monoidal $*$-algebra;
\item every deterministic predictive presentation and every future-operational
      compiler factors uniquely through $Z_X^{\min}$; consequently
      \eqref{eq:grand-tensor-main} is canonical without an independently
      supplied ledger, hidden state, persistent state space, preferred Kraus
      frame, primitive metric, or preselected history algebra;
\item any finite same-history readable enrichment by normalized
      conditional kernels admits the minimal faithful quotient
      $\mathfrak P_X^{\rm rel,min}$ of
      \eqref{eq:readable-relation-packet-main}; its forgetful image is the
      underlying operational law, so same-cut readability does not alter the
      predictive reconstruction or introduce new geometric coordinates.
\end{enumerate}
\end{theorem}

The detailed proof is given in Appendix~\ref{supp:grand-reconstruction}.  Three minimizations
remain logically distinct: predictive state minimization, linear Hankel
minimization, and contextual $C^*$-history separation.  Conflating them either
retains future-invisible spectator structure or erases distinctions visible to
an actual operational continuation.

The independence of orientation and absolute calibration, and the exact
remaining calibration freedom, are established in
\Cref{thm:supp-relative-primitive-floor,thm:supp-calibration-fibre}.
The next section uses these additional physical records to reconstruct
geometry on the canonical carrier; it does not alter the positive
history object just constructed.


\section{Relational dimension and Lorentzian geometry}
\label{sec:adm-main}
% =============================================================================

\subsection{Endpoint--cycle selection and a secondary simplex criterion}

The primary finite dimensional statement is endpoint--cycle source balance:
inside connected simple regular relations it uniquely selects $K_4$, without
an alternating form. Its simplex representation-theoretic counterpart is the
orientation-twisted equivalence of the complete cycle and endpoint sources.
A secondary criterion in the pair-homogeneous simplex category combines
alternating parity with an explicit minimum-source rule. That criterion
selects the same cell, but its minimality requirement is a selection
hypothesis, not a consequence of predictive compression.

\begin{proposition}[Endpoint--cycle balance selects $K_4$ without an alternating form]
\label{prop:main-regular-graph-k4}
Let $G$ be a connected simple $k$-regular relation graph on $v>1$ endpoints.
If its faithful endpoint source and cycle source have equal rank, then
\begin{equation}
 v-1=\frac{kv}{2}-v+1,
 \qquad\hbox{so}\qquad
 (4-k)v=4.
 \label{eq:main-regular-graph-balance}
\end{equation}
The unique nontrivial connected simple regular solution is
\begin{equation}
 v=4,\qquad k=3,\qquad G\cong K_4.
 \label{eq:main-regular-graph-k4}
\end{equation}
Thus the associated faithful endpoint module has dimension $v-1=3$.  This
selector contains no temporal--spatial alternating-form hypothesis.
\end{proposition}

\begin{proof}
For a connected graph the endpoint-source rank is $v-1$, while the cycle
rank is $|E|-v+1=kv/2-v+1$. Equality gives $(4-k)v=4$.
For integers $v>1$ and $k\ge1$ with $k\le v-1$, the only possibility
is $v=4$, $k=3$; a simple graph with these parameters is $K_4$.
\end{proof}

For the simplex selector, consider a local relation event on the minimum
persistent predictive carrier of \Cref{mt:grand-tensor}, with one
future-readable source and one future-readable target.  Physical reversal
turns it into a signed edge.  Future minimization removes parallel labels that
no future can distinguish, while pair-homogeneous recurrence completes the
orbit of one nonzero relation.  The relation-colour-free skeleton is therefore
$K_N$.  Writing $\R^N=\R u_0\oplus W_N$, with
$u_0=N^{-1/2}\one_N$, the signed edge space decomposes as
\begin{equation}
 C_1(K_N)=u_0\wedge W_N\oplus\Lambda^2W_N.
 \label{eq:cut-cycle-main}
\end{equation}
The first summand is the endpoint or cut source and the second the complete
cycle or interference source, of dimensions $N-1$ and
$\binom{N-1}{2}$ respectively.

\begin{proposition}[Alternating/minimal-source simplex selector]
\label{prop:main-alternating-k4}
Assume in addition that the temporal--spatial coefficient space
$\R t\oplus W_N$ carries a nondegenerate alternating form, that the endpoint
and complete signed-edge sources are faithful, that the cycle source is
nonzero, and that the future-readable source innovation is minimal among
conservative completions satisfying these queries.  Then
\begin{equation}
 N=4,\qquad G_{\mathrm{rel}}=K_4,
 \qquad W_{\mathrm{sp}}=W_4\cong\R^3.
 \label{eq:k4-main}
\end{equation}
Moreover, after complexification,
\begin{equation}
 \Hom_{S_N}(\Lambda^2W_N,W_N\otimes\operatorname{sgn})
 =\begin{cases}\C,&N=4,\\0,&N\ne4,\end{cases}
 \label{eq:main-oriented-copy}
\end{equation}
so only at $N=4$ is the complete cycle representation one oriented copy of the
spatial standard module rather than an additional inequivalent source type.
\end{proposition}

The proof is given in Appendix~\ref{supp:dimension-adm}.  The first conclusion
uses parity, nonzero interference and the stipulated source minimum. The
representation-theoretic statement is a separate algebraic characterization
of $N=4$ inside the simplex family, not evidence of generic microscopic selection.

\begin{remark}[Exact scope of the dimensional claim]
\label{rem:main-dimension-scope}
The endpoint--cycle balance theorem is the primary finite classification;
the alternating/minimal-source theorem is a secondary classification with
different hypotheses.  The first
shows that the alternating form is not the sole origin of the three-dimensional
spatial module.  The second does not become a hidden symplectic or Hamiltonian
assumption: its alternating form supplies only an antisymmetric coefficient
pairing and no metric, lapse, shift, connection, curvature, or Lorentzian sign.
Neither theorem is a classification of arbitrary graphs, hypergraphs, or
higher-arity microscopic relation systems.
\end{remark}

\paragraph{Source equivalence.}
Within the faithful simplex category the complete cycle source is one
orientation-twisted copy of the endpoint source exactly at $N=4$;
\Cref{prop:supp-source-equivalence} proves uniqueness up to scale and
orientation. This interprets the rank balance within that category, without
asserting that an arbitrary operational law satisfies either selector.


\subsection{The \texorpdfstring{$A_3$}{A3} root race}

Choose an anchored basis of $W_4$ in which the six unoriented $A_3$ roots are
\begin{equation}
 r_1=e_1,\quad r_2=e_2,\quad r_3=e_3,\quad
 r_4=e_2-e_1,\quad r_5=e_3-e_1,\quad r_6=e_3-e_2.
 \label{eq:a3-roots-main}
\end{equation}
Let $k_a^\pm$ be the two directed rates associated with $\pm r_a$, and let $h$
be the mesh scale.  The predictable second and first moments are
\begin{equation}
 \mathcal B_h
 =h^2\sum_{a=1}^{6}(k_a^++k_a^-)r_ar_a^{\mathsf T},
 \qquad
 \beta_h
 =h\sum_{a=1}^{6}(k_a^+-k_a^-)r_a.
 \label{eq:bracket-drift-main}
\end{equation}
The six dyadics $r_ar_a^{\mathsf T}$ form a basis of $\Sym_3$, so the rate sums
reconstruct all six spatial inverse-metric coefficients.  The rate differences
map with rank three onto $\beta_h$; their remaining three coordinates are
circulating root currents.  Winner probabilities alone lose the common rate
scale.  Observing both the winning root and its waiting time gives a faithful
score frame for all directed rates.

Let $\varrho_h>0$ be the physical speed density.  The spatial metric and lapse
are reconstructed by
\begin{equation}
 \boxed{
 g_h=\varrho_h^{2/3}(\det\mathcal B_h)^{1/3}\mathcal B_h^{-1},
 \qquad
 N_h=\varrho_h^{1/3}(\det\mathcal B_h)^{1/6}.}
 \label{eq:adm-inversion-main}
\end{equation}
Equivalently,
\begin{equation}
 \sqrt{\det g_h}=\varrho_h,
 \qquad
 \mathcal B_h=N_h^2g_h^{-1}.
 \label{eq:adm-forward-main}
\end{equation}
Together with the shift $\beta_h$, these are exactly the ten ADM coefficients.
The formula is an analytic bijection on the positive cone.  A nondegenerate
bracket requires total root rate of order $h^{-2}$; identifying the geometric
scheduler with a slower $O(h^{-1})$ renewal clock would make
$\mathcal B_h\to0$ and collapse the metric.

\paragraph{Linear tomography versus positive realizability.}
Spanning $\Sym_3$ is a linear statement. Positive rates in a fixed root frame
realize only the cone
\begin{equation}
 \mathcal C_{A_3}:=\left\{\sum_{a=1}^6s_ar_ar_a^{\mathsf T}:s_a>0\right\},
 \qquad s_a=h^2(k_a^++k_a^-),
 \label{eq:main-root-positive-cone}
\end{equation}
not every positive-definite matrix in that same coordinate frame.
The component inequalities and the positive directed lift of a bounded
shift are recorded in Appendix~\ref{supp:dimension-adm}. Thus the ADM
map is an analytic bijection on positive bracket--density pairs, whereas
realization by the declared race also requires an interior cone margin.


\subsection{One signed direction and Lorentzian inertia}

The Grand Tensor reconstructed in the previous section is positive: its
canonical word-Gram hierarchy is a Hilbert-space Gram structure, and the
root-race bracket is a positive quadratic form on the three-dimensional
spatial module.  Neither can contain a negative direction.  The additional
physical datum is a protected two-valued operational Read
\begin{equation}
 \epsilon_X\in\{+1,-1\},
 \label{eq:orientation-read-main}
\end{equation}
with both values separated by the declared Read policy.  In the support-minimal
Hilbert realization it has orthogonal readout projections $P_+$ and $P_-$,
$P_++P_-=I$, and therefore
\begin{equation}
 J:=P_+-P_-,\qquad J=J^*=J^{-1},
 \qquad [u,v]_J:=\langle u,Jv\rangle.
 \label{eq:krein-main}
\end{equation}
This is the operator representation of an already occurring $\mathbb Z_2$
mark.  The microscopic record itself carries no norm, cone, lapse, temporal
coordinate, or indefinite metric.

The sign is not reconstructible from the positive packet.  The separating
models of \Cref{thm:supp-relative-primitive-floor} keep the complete positive
operational data fixed while reversing the protected orientation, and the
obstruction survives common conservative refinement.  The signed row is
therefore irreducible relative to the declared operational language rather than
a reparametrization of the positive Grand Tensor.

\begin{definition}[Minimal signed-extension condition]
\label{def:main-signed-extension}
A minimal signed extension of the positive spatial space $(W,g)$ consists of
one independent real line $L$, a negative-definite quadratic form $q_L$ on
$L$, and the orthogonal extension $q=q_L\oplus g$. Calibration fixes the
absolute coefficient of $q_L$. A protected orientation selects one of the
two directions on $L$; it does not, by itself, determine the sign of $q_L$.
\end{definition}

A binary Read has a representation $J=P_+-P_-$ on its Hilbert carrier, but
the dimensions of the eigenspaces of $J$ need not be one and three and are
not the dimension of the reconstructed tangent space. The passage from this
Read to a single negatively signed geometric line is precisely
\Cref{def:main-signed-extension}. Neither positivity of a word Gram nor a
choice of orientation alone implies it. Once this condition is imposed,
Sylvester inertia determines the Lorentzian type without choosing gamma
matrices. In inverse-metric conventions the calibrated line is
$-N^{-2}\tau^2$; the corresponding cotangent extension with shift is
\eqref{eq:dirac-square-main}.


\begin{proposition}[Minimal one-sign extension has Lorentzian inertia]
\label{prop:main-signature-minimality}
Let $(W,g)$ be the positive three-dimensional spatial quadratic space
reconstructed from the $K_4/A_3$ root race.  Adjoin one independent real line
$\R t$ with the negative quadratic form specified by the minimal
signed-extension condition, and leave
$g|_W$ unchanged.  Then every minimal nondegenerate orthogonal signed extension
is, up to multiplication of the whole quadratic form by $-1$,
\begin{equation}
 \boxed{q(\tau t+w)=-N^{-2}\tau^2+g(w,w),
 \qquad \operatorname{inertia}(q)=(1,3).}
 \label{eq:main-signature-inertia}
\end{equation}
No choice of Clifford representation enters this conclusion.
\end{proposition}

The proof and the corresponding independence statement are given in
Appendix~\ref{supp:dimension-adm}.  The point is that the inertia conclusion
precedes the Clifford realization.

The signed-extension condition, rather than positivity or orientation alone,
is the extra geometric hypothesis. Its role is now explicit, and the
Clifford layer only represents the resulting quadratic form.

\paragraph{Clifford realization.}
Once the signed quadratic form is fixed, the explicit matrices and
principal-symbol calculation in Appendix~\ref{supp:dimension-adm}
give a faithful minimal Clifford realization. No representation-theoretic
choice enters the inertia conclusion. We can now collect the reconstruction
steps in a single theorem.

\begin{theorem}[Minimal faithful $3+1$ ADM reconstruction from either finite selector]
\label{mt:adm}
Let the minimum persistent predictive carrier of \Cref{mt:grand-tensor} admit a
geometric branch with one source and one target, physical reversal, faithful
endpoint data, and a local relation selected as $K_4$ by either
\Cref{prop:main-regular-graph-k4} or the alternating/minimal-source simplex
criterion of \Cref{prop:main-alternating-k4}.  Assume an occurring
scheduler-separated root-race calibration supplies directed winners, their
waiting times, and a positive physical speed density.  Assume also that the
protected orientation packet supplies the Read
\eqref{eq:orientation-read-main}, independent of the positive packet, with one
nontrivial signed direction in the sense of
\Cref{def:main-signed-extension}.  Then:
\begin{enumerate}[label=\textup{(\roman*)},leftmargin=2.8em]
\item the selected local relation cell is $K_4$ and its spatial module is the
      $A_3$ root space $W_4\cong\R^3$;
\item the winner--waiting statistics of the scheduler-separated root race form
      a faithful coordinate system for all directed rates;
\item the moments \eqref{eq:bracket-drift-main} and the speed density invert
      uniquely to the ADM spatial metric, lapse, and shift through
      \eqref{eq:adm-inversion-main};
\item a nondegenerate continuum bracket requires an $O(h^{-2})$ fast scheduler;
\item positivity of the reconstructed spatial form together with exactly one
      independent signed direction fixes Lorentzian inertia $(1,3)$ up to
      overall sign, before any Clifford representation is chosen;
\item the Clifford cell \eqref{eq:clifford-main} faithfully realizes that
      already determined quadratic form, and the principal symbol has square
      \eqref{eq:dirac-square-main};
\item under summably correctable mixed-Gram and refinement defects, the
      dimension, source ranks, pair gap, ADM packet, and signed inertia persist
      on the canonical cofinal limit quotient.
\end{enumerate}
No persistent predictive state, metric, coframe, connection, curvature tensor,
or spacetime dimension is independently supplied.  Persistence is reconstructed
by \Cref{mt:grand-tensor}; the protected orientation and absolute
winner--waiting calibration remain irreducible physical input rows on this
branch.  The alternating form is required only on the second dimensional
selector and is not a prerequisite of the first.
\end{theorem}

The residual status of the last two rows is not merely a limitation of the
construction: \Cref{thm:supp-relative-primitive-floor} gives explicit
separating models showing that neither the protected orientation nor the
absolute calibration is functorially reconstructed from the remaining
positive operational data, even after common cofinal refinement.

The detailed proof is given in Appendix~\ref{supp:dimension-adm}.  The local algebraic audit is
finite: on the free $A_3$ accumulator, all equal-endpoint word relations are
generated by inverse digons, root triangles, and commutator squares of depth at
most four.  Global topology enters separately through a finite period bank.


\section{Reconstruction of the gravitational action}
\label{sec:einstein-dynamics}

The preceding construction is kinematical. We now reconstruct the
variational object to which a continuum limit can be applied. The order
of the argument is fixed: the complete accepted record determines the
field dynamics; marked renewal statistics determine action increments;
finite exactness integrates those increments; and an independent
comparison tests stationarity. Only then is the local gravitational
representative classified.

\paragraph{Three kinds of hypotheses.}
\begin{center}
\small
\begin{tabularx}{\linewidth}{@{}p{.20\linewidth}X@{}}
\toprule
Input or condition & Mathematical role\\
\midrule
Operational records & Complete conditional futures, readable root/waiting
and density data, a signed-line extension, and the same-history accepted,
reward and duration records determine the finite objects being compared.\\
Finite structural tests & The relational selector, action curl and periods,
local degree and covariance class, and nonzero calibrated Palatini coefficient
specify the branch. They are not automatic consequences of positivity.\\
Cofinal analytic tests & Scaled physical stationarity, coframe compactness,
connection compactness, curvature propagation and full Cartan consistency
permit the nonlinear variational limit. Each is imposed on one common
realized regulator sequence.\\
\bottomrule
\end{tabularx}
\end{center}

This distinction will also govern the limit theorem in
Section~\ref{sec:einstein-limit}. Its analytic hypotheses are not
assumed in order to reconstruct the finite objects below.
\subsection{The determining field and accepted kernel}

At cutoff $X$, let $\widetilde\Theta_X$ denote the finite same-history record
carrying the reconstructed geometric variables used by the gravitational
sector, generic protected reward and duration coordinates, and the resolved
accepted-branch labels.  At this stage the reward coordinates are not named as
Palatini, Holst, Einstein--Hilbert, or cosmological terms; they are simply actual
finite Reads (or fixed functions of actual protected record coordinates) whose
variation content will be reconstructed and classified below.  Two values
$x,y\in\widetilde\Theta_X$ are equivalent
when every admitted future word and every declared determining coordinate has
the same value from $x$ and $y$, including branch-domain indicators.  Put
\begin{equation}
 \Theta_{{\rm fld},X}^{\min}
 :=\widetilde\Theta_X/\!\sim_{{\rm fld},X}.
 \label{eq:main-determining-field}
\end{equation}
For resolved accepted subbranches $a$, let $p_{a,X}(\theta)$ be their physical
probabilities and $u_a^{\min}$ their descended updates.  On the positive
acceptance domain define
\begin{equation}
 K_X^{\rm acc}(\theta'\mid\theta)
 =\frac{1}{p_{{\rm acc},X}(\theta)}
   \sum_{\substack{a:\,u_a^{\min}(\theta)=\theta'}}p_{a,X}(\theta).
 \label{eq:main-accepted-kernel}
\end{equation}

\begin{theorem}[Determining-field and accepted-kernel reconstruction]
\label{thm:main-determining-kernel}
The quotient \eqref{eq:main-determining-field} is the unique coarsest finite
record through which every declared determining-field future and every
primitive update factors.  Finite partition refinement terminates, and
\eqref{eq:main-accepted-kernel} is the unique table-native accepted transition
reproducing all selected post-accepted field futures.
\end{theorem}

Thus the accepted field dynamics is not selected after the fact once the
complete accepted cylinder is present.  What cannot be reconstructed from the
marginal accepted/not-accepted bit is the complete field cylinder itself; this
is a physical occurrence requirement rather than an implicit state-space
assumption.  The proof and the corresponding non-identifiability boundary are
given in Appendix~\ref{supp:dynamics}.

\subsection{Renewal pressure and common-action exactness}

Let $(\Omega_X,p_X)$ be a finite faithful marked first-return law, let
$\tau_X:\Omega_X\to(0,\infty)$ be the protected physical duration, and let
$E_X$ be the finite variation space.  A protected cycle-reward synthesis is a
linear map
\begin{equation}
 \mathscr G_X:E_X\longrightarrow L^2(\Omega_X,p_X),
 \qquad v\longmapsto g_{v,X},
 \label{eq:main-reward-synthesis}
\end{equation}
whose values are actual Reads or fixed functions of protected record
coordinates.  Define
\begin{equation}
 \mathcal L_X(q,r)
 =\mathbb E_X\!\left[e^{-g_{q,X}-r\tau_X}\right].
 \label{eq:main-reward-laplace}
\end{equation}
Since $\partial_r\mathcal L_X(0,0)=-\mathbb E_X\tau_X<0$, there is a unique
analytic pressure $\mathcal P_X$ near the origin with
\begin{equation}
 \mathcal P_X(0)=0,\qquad
 \mathcal L_X(q,\mathcal P_X(q))=1.
 \label{eq:main-renewal-pressure}
\end{equation}
The normalized tilted law is
$p_{X,q}(\omega)=p_X(\omega)
\exp[-g_{q,X}(\omega)-\mathcal P_X(q)\tau_X(\omega)]$.
Let $S_Xv=D\log p_{X,q}|_{q=0}[v]$, let
$\pi_X=D\mathcal P_X(0)$, and let $T_Xc=c\tau_X$.

\begin{theorem}[Renewal-pressure action reconstruction]
\label{thm:main-action-reconstruction}
The marked renewal table reconstructs $\mathcal P_X$ and all of its finite
jets.  If $\bar\tau_X=\mathbb E_X\tau_X$, then
\begin{equation}
 \pi_X(v)=-\frac{\mathbb E_X g_{v,X}}{\bar\tau_X},
 \qquad
 \mathscr G_X=-S_X-T_X\pi_X.
 \label{eq:main-action-slope-reconstruction}
\end{equation}
Hence the centered probability score together with physical duration recovers
the complete uncentered action slope, including its affine anchor.  The pressure
Hessian and duration cross-jets reconstruct the full marginal and mixed action
Gram whenever all reward sectors occur on the same marked cycle record.
\end{theorem}

The remaining question is whether the reconstructed increments belong to one
scalar action.  Let $\mathscr V_X$ be the finite protected variation complex and
let $\alpha_X\in C^1(\mathscr V_X;\R)$ be the reconstructed action increment.

\begin{theorem}[Finite common-action exactness]
\label{thm:main-common-action}
There is a scalar common action $\mathscr S_X$ satisfying
\begin{equation}
 \boxed{\alpha_X=d_0\mathscr S_X}
 \label{eq:main-common-action}
\end{equation}
if and only if $d_1\alpha_X=0$ on every declared $2$-cell and the periods of
$\alpha_X$ vanish on one integral basis of the free part of $H_1(\mathscr V_X;\Z)$.  After one
additive anchor is fixed on each connected component, $\mathscr S_X$ is unique and is reconstructed by path
integration.  Moreover, if $P_{\rm har}$ is the harmonic projection and
$\sigma_{\rm coex}>0$ is the least singular value of $d_1$ on the coexact
range, then
\begin{equation}
 \operatorname{dist}(\alpha_X,\Ran d_0)^2
 \le \|P_{\rm har}\alpha_X\|^2
 +\sigma_{\rm coex}^{-2}\|d_1\alpha_X\|^2.
 \label{eq:main-common-action-obstruction}
\end{equation}
Thus failure of exactness supplies a finite obstruction to a common action.
\end{theorem}

\begin{remark}[Non-circularity of the gravitational reward]
\label{rem:main-gravitational-reward-noncircular}
The synthesis $\mathscr G_X$ is fixed by the protected cycle record before it is
compared with the accepted kernel $K_X^{\rm acc}$.  Neither the transition
probabilities nor a desired continuum equation are used to fit a gravitational
cost after the fact.  Moreover, no Palatini or Holst decomposition is assumed at
this stage.  That decomposition appears only after derivative-free face
linearization, complete resolved-memory subtraction, and the Lorentz-equivariant
commutant classification below.  The subsequent Bregman/KL test is therefore an
independent question: it asks whether the actually accepted dynamics is
stationary for the previously reconstructed action.
\end{remark}

Exactness identifies which scalar action the recorded increments represent.
It does not yet relate that action to the actually accepted evolution.
The next estimate makes this comparison at the level of field variations,
with all cutoff-dependent test-lift factors retained.

\subsection{Stationarity of the reconstructed action}

On a compact convex chart $\Theta\subset\R^d$ of the reconstructed determining
field, let $\Psi$ and the reconstructed common action $\mathcal A$ be $C^2$ on
an open neighbourhood of $\Theta$.  The operational potential is strictly
convex, with Fisher metric
\begin{equation}
 G(\theta)=\nabla^2\Psi(\theta),
 \qquad mI\preceq G(\theta)\preceq MI,
 \label{eq:main-fisher-metric}
\end{equation}
and, for $\lambda,L\ge0$ and $\eta>0$,
\begin{equation}
 -\lambda G(\theta)\preceq D^2\mathcal A(\theta)
 \preceq L G(\theta),\qquad \eta\lambda<1.
 \label{eq:main-action-hessian}
\end{equation}
For the Bregman divergence~\cite{Bregman1967}
\begin{equation}
 D_\Psi(y,\theta)=\Psi(y)-\Psi(\theta)
       -\langle\nabla\Psi(\theta),y-\theta\rangle,
 \label{eq:main-bregman}
\end{equation}
put
\begin{equation}
 \Phi_\theta(y)=\mathcal A(y)+\eta^{-1}D_\Psi(y,\theta),
 \qquad y_\theta^*=\arg\min_{y\in\Theta}\Phi_\theta(y).
 \label{eq:main-proximal-action}
\end{equation}
The minimum is unique.  On the branch used for unrestricted first variations,
$y_\theta^*$ lies in the interior of $\Theta$ for every supported source
$\theta$.  This interiority is important: a constrained minimum at a chart
boundary need not have zero ordinary action gradient.

For an actually accepted transition $\theta\to y$, define
\begin{equation}
 \Delta_\Psi(\theta,y)
       =\Phi_\theta(y)-\Phi_\theta(y_\theta^*)\ge0.
 \label{eq:main-action-gap}
\end{equation}
If $K=(1-\vartheta)I+\vartheta K^{\rm acc}$, with state-independent
$0<\vartheta\le1$, has stationary law $\nu$, set
\begin{equation}
 \begin{aligned}
 \overline\Delta_\Psi
  &=\int\nu(\dd\theta)K^{\rm acc}(\dd y\mid\theta)
                          \Delta_\Psi(\theta,y),\\
 \mathfrak S(\theta)
  &=\nabla\mathcal A(\theta)^*G(\theta)^{-1}\nabla\mathcal A(\theta).
 \end{aligned}
 \label{eq:main-stationarity-gap}
\end{equation}

\begin{theorem}[Accepted action gap controls interior stationarity]
\label{thm:main-action-stationarity}
Under the preceding hypotheses, including interiority of the supported
proximal minimizers, an explicit finite constant
$C_{\Psi,\mathcal A,\eta}$ satisfies
\begin{equation}
 \int\mathfrak S(\theta)\,\nu(\dd\theta)
       \le C_{\Psi,\mathcal A,\eta}\,\overline\Delta_\Psi.
 \label{eq:main-stationarity-bound}
\end{equation}
In particular, zero gap forces the represented first variation of the common
action to vanish $\nu$-almost surely.
\end{theorem}

The proof in Appendix~\ref{supp:dynamics} keeps the convexity and chart
constants explicit.  To pass to physical metric tests at varying cutoffs, let
$v_{X,K}(k;\theta)$ be the actual represented lift of a fixed determining test
$k$ on a compact cylinder $K$, and put
\begin{equation}
 B_{X,K}(k)^2:=\sup_{\theta\in\supp\nu_X}
 \|v_{X,K}(k;\theta)\|_{G_X(\theta)}^2.
 \label{eq:main-test-lift-budget}
\end{equation}
Duality then gives the quantitative physical-test estimate
\begin{equation}
 \mathbb E_{\nu_X}
       |\delta\mathcal A_X[v_{X,K}(k)]|^2
       \le B_{X,K}(k)^2C_X\overline\Delta_{\Psi,X},
 \qquad C_X=C_{\Psi_X,\mathcal A_X,\eta_X}.
 \label{eq:main-physical-stationarity}
\end{equation}
Thus the right-hand side, rather than the unweighted gap alone, is the
quantity required to vanish.  On a countable determining core this gives
convergence of the first variations in mean square and an almost-sure
vanishing subsequence under any common coupling; the deterministic continuum
handoff is applied to the resulting realized sequences.  The construction and
proof are given in \Cref{prop:supp-physical-test-stationarity}.

For a finite probability row with reference law $p_\theta(y)>0$ and
same-history action cost $c_\theta(y)$, the entropy-proximal row is
\begin{equation}
 q_\theta^*(y)=\frac{p_\theta(y)e^{-\eta c_\theta(y)}}
                    {\sum_zp_\theta(z)e^{-\eta c_\theta(z)}},
 \label{eq:main-gibbs-row}
\end{equation}
and its row variational gap is exactly
$\eta^{-1}D_{\rm KL}(K^{\rm acc}(\cdot\mid\theta)\Vert q_\theta^*)$.
This is an independent likelihood test against the reconstructed cost.
Equality of two finite-temperature probability rows does not by itself make
every sampled field configuration a critical point of $\mathcal A$; the
field-action estimate and its physical test lifts remain the stationarity
criterion used in the Einstein passage.

The preceding tests reconstruct an action and measure its stationarity,
without choosing a gravitational Lagrangian. We now restrict the local
incidence and degree class and use Lorentz covariance to identify its
possible bulk terms. This algebraic step is separate from the analytic
limit in the next section.

\subsection{Classification of the gravitational representative}

Let $e^I$ be the reconstructed coframe, $\omega^{IJ}$ a metric-compatible
reconstructed Cartan connection, $B^{IJ}=e^I\wedge e^J$, and $T^I=De^I$.
Before any continuum gravitational Lagrangian is named, the generic protected
face reward is subjected to two finite operations described in
Appendix~\ref{supp:dynamics}: binary renewal subdivision, which produces a unique
first-order linear functional on its exact or summably stable branch, and
subtraction of the complete resolved cut-memory log-partition.  Denote the
resulting generic first-order gravitational score by $A^{\rm grav}_{X,\rm sh}$.
For later reference, write the three Lorentz-natural first-order densities (up
to the overall normalization convention) as
\begin{align}
 L_{\rm Palatini}
 &=\frac12\epsilon_{IJKL}e^I\wedge e^J\wedge R^{KL},\\
 L_{\rm Holst}
 &=e^I\wedge e^J\wedge R_{IJ},\\
 L_{\rm vol}
 &=\frac1{4!}\epsilon_{IJKL}
 e^I\wedge e^J\wedge e^K\wedge e^L.
 \label{eq:main-phv-densities}
\end{align}

\begin{definition}[Local curvature-linear gravitational class]
\label{def:main-gravitational-class}
The gravitational class considered here consists of local polynomial
four-form densities in an oriented coframe $e$ and a metric-compatible
Lorentz connection $\omega$. The curvature-dependent part has bidegree
$(2,1)$ in $(e,R)$, and the curvature-independent part has degree four in
$e$. Coefficients are spacetime-constant real scalars, and the only internal
contractions are the Lorentz metric, its oriented volume tensor, and the
represented scalar multiplicity commutant. No inverse coframe, independent
torsion polynomial, derivative of curvature, additional background tensor,
or matter field is admitted in this class. Densities are compared modulo
exact terms whose variations vanish on the declared compactly supported
bulk tests; boundary functionals are retained separately when boundary tests
are allowed.
\end{definition}

This definition specifies the incidence and degree of the reconstructed
score. Here \emph{curvature linear} is more restrictive than merely first
order in derivatives of an independent connection. In particular, it is not
a classification of every generally covariant first-order gravity action.


\begin{theorem}[Lorentz-natural classification of the generic gravitational score]
\label{thm:main-gravitational-classification}
Assume $A^{\rm grav}_{X,\rm sh}$ belongs to the polynomial local class of
\Cref{def:main-gravitational-class}, is real and Lorentz natural, has scalar
residual typed multiplicity commutant on the bivector carrier, and represents
determinant coframe volume.
Then, up to coefficient normalization, its only first-order Lorentz-natural
bulk terms are
\begin{equation}
 A^{\rm grav}_{X,\rm sh}
 \in\Span_{\R}\{L_{\rm Holst},L_{\rm Palatini},L_{\rm vol}\}.
 \label{eq:main-phv-classification}
\end{equation}
Equivalently, there are real coefficients $\alpha_X,\beta_X,\lambda_X$ such that
\begin{equation}
 A^{\rm grav}_{X,\rm sh}
 =\alpha_XL_{\rm Holst}+\beta_XL_{\rm Palatini}+\lambda_XL_{\rm vol}.
 \label{eq:main-phv-coefficients}
\end{equation}
The hypothesis specifies the incidence and degree class, not its three
coefficients or the name of a gravitational Lagrangian. The classification
is not asserted for all first-order connection actions.
\end{theorem}

The classification follows because both the coframe bivector and curvature
transform in $\Lambda^2\R^{1,3}$ and
$\operatorname{End}_{SO^+(1,3)}(\Lambda^2\R^{1,3})
=\Span_{\R}\{I,\star\}$; the scalar residual commutant excludes any additional
internal tensor, and determinant volume is the only independent alternating
four-coframe scalar at this order.  The full proof is in
Appendix~\ref{supp:dynamics}.

\paragraph{Compatibility at the level of physical first variations.}
Classification of a face score is not, by itself, convergence of the first
variation of the complete physical action. Let $S_{{\rm H},X}$,
$S_{{\rm P},X}$ and $V_X$ denote the actual common-cylinder realizations of
the three classified densities, with their physical cell measures and
boundary conventions. On a represented determining variation $v_X$ define
\begin{equation}
 \begin{split}
 \mathfrak r_X[v_X]:={}&\delta\mathcal A_X[v_X]
  -\alpha_X\delta S_{{\rm H},X}[v_X]
  -\beta_X\delta S_{{\rm P},X}[v_X]
  -\lambda_X\delta V_X[v_X].
 \end{split}
 \label{eq:main-action-variation-remainder}
\end{equation}
The gravitational reduction used below requires, on the same metric and
connection determining cores,
\begin{equation}
 \mathfrak r_X[v_{X,K}(k)]\longrightarrow0.
 \label{eq:main-first-variation-reduction}
\end{equation}
At each cutoff these are comparisons of actual represented test rows; the
finite banks exhaust the fixed core with the stated test-interpolation
control. This condition includes the variation of any retained subdivision
error, the complete resolved-memory subtraction, and any boundary remainder.
The memory log-partition must be the contribution represented on that same
physical record, not an unrelated convex function chosen for subtraction.

Condition \eqref{eq:main-first-variation-reduction} is automatic when the
physical action agrees exactly, up to the specified boundary completion,
with its classified representative on the chart. More generally, a uniform
$C^1$ bound on the remainder over the physically lifted test bank is
sufficient. Convergence of action values alone is not sufficient: the smooth
functions $A_n(x)=n^{-1}\sin(nx)$ converge uniformly to zero, although
$A_n'(0)=1$. Thus the Einstein handoff uses variational consistency, not an
unjustified interchange of differentiation and a cutoff limit.


On the spinless nondegenerate branch, the connection variation controls
torsion by \Cref{thm:supp-palatini-torsion}. Classification and stationarity
are therefore in place. To obtain an Einstein metric, it remains to pass
these variations to one common noncollapsed limit and to identify its
nonlinear connection curvature.


\section{The renewal-to-Einstein limit}
\label{sec:einstein-limit}

This section proves the conditional continuum theorem. A positive spatial
cut margin supplies spatial compactness estimates. Initial-energy propagation
controls curvature size, while a separate connection argument identifies
the weak curvature limit. The final variational passage uses these estimates
on the same realized regulator sequence.

The literal-link criterion is the principal connection route presented
here. Fully discrete energy propagation, the differentiated action-current
estimate, a Hodge--temporal alternative and interface comparisons are
retained in Appendix~\ref{supp:continuum}. They are alternative or supporting
certificates, not additional conclusions inferred from a curvature norm alone.
\subsection{Global noncollapse and compact spatial screens}

Let the physical endpoint quotient of a global regulator carry masses
$m_{v,X}$, symmetric spatial conductances $c_{e,X}$, physical cell scale
$h_X$, and provenance-safe efficient capacities $u_{e,X}$.  The Dirichlet
form and its nonnegative self-adjoint Laplacian are
\begin{equation}
 \mathcal E_X^{\rm sp}(f)
 =\sum_{\{u,v\}\in E_X}c_{uv,X}|f(u)-f(v)|^2,
 \qquad
 \mathcal E_X^{\rm sp}(f)=\langle f,\Delta_X^{\rm sp}f\rangle_{L^2(\mu_X)}.
 \label{eq:main-spatial-dirichlet}
\end{equation}
The endpoint/action packet also contains a finite capacity--conductance
comparison.  On the positive comparison branch there are cutoff-independent
constants $0<a_-\le a_+<\infty$ such that
\begin{equation}
 a_-h_Xc_{e,X}\le u_{e,X}\le a_+h_Xc_{e,X}
 \qquad(e\in E_X).
 \label{eq:main-capacity-conductance}
\end{equation}
This is an explicit same-history comparison between two reconstructed physical
rows; it is not an identification made after taking the continuum limit.
Define the efficient and analytic cut constants
\begin{equation}
 I_X^{\rm eff}:=
 \min_{0<\mu_X(A)\le\mu_X(V_X)/2}
 \frac{u_X(\partial A)}{\mu_X(A)^{2/3}},
 \qquad
 I_X^{\rm sp}:=
 \min_{0<\mu_X(A)\le\mu_X(V_X)/2}
 \frac{h_Xc_X(\partial A)}{\mu_X(A)^{2/3}}.
 \label{eq:main-cut-constant}
\end{equation}
Then
\begin{equation}
 a_- I_X^{\rm sp}\le I_X^{\rm eff}\le a_+I_X^{\rm sp}.
 \label{eq:main-cut-comparison}
\end{equation}
Thus a positive efficient-capacity margin is exactly the operational input
needed to obtain a positive spatial isoperimetric margin on this branch.


\begin{theorem}[Cut margin and compact-screen alternative]
\label{thm:main-spatial-screen}
Assume the total physical volume is uniformly bounded,
$\mu_X(V_X)\le V_*$, the weighted degree satisfies
$\sum_uc_{uv,X}\le D_*h_X^{-2}m_{v,X}$, and
\eqref{eq:main-capacity-conductance} holds.  If
$\inf_X I_X^{\rm eff}=I_*^{\rm eff}>0$, then
$I_X^{\rm sp}\ge I_*^{\rm eff}/a_+=:I_*^{\rm sp}>0$ and every mean-zero $f$
obeys regulator-uniform Sobolev and Poincar\'e bounds
\begin{equation}
 \|f\|_{L^6(\mu_X)}^2\le C_{\rm S}\mathcal E_X^{\rm sp}(f),
 \qquad
 \|f\|_{L^2(\mu_X)}^2\le C_{\rm P}\mathcal E_X^{\rm sp}(f),
 \label{eq:main-sobolev-poincare}
\end{equation}
with cutoff-independent constants $C_{\rm S},C_{\rm P}$ given explicitly in
Appendix~\ref{supp:continuum}.
The low-energy counting function satisfies
\begin{equation}
 N_X(R)\le1+4e^{3/2}V_*(C_{\rm S}R)^{3/2}
 \qquad(R>0),
 \label{eq:main-weyl-count}
\end{equation}
and for $P_{X,R}=\one_{[0,R]}(\Delta_X^{\rm sp})$,
\begin{equation}
 \|(I-P_{X,R})f\|_{L^2(\mu_X)}^2
 \le R^{-1}\mathcal E_X^{\rm sp}(f).
 \label{eq:main-screen-tail}
\end{equation}
Hence the spectral projections form uniformly exhausted finite-rank compact
screens.  If instead $I_{X_j}^{\rm eff}>0$ tends to zero, the minimizing
cuts $A_j$ give normalized potentials $\phi_j$ for the demand
\eqref{eq:main-cut-demand}, defined in Appendix~\ref{supp:continuum}, with
\begin{equation}
 \sum_eu_{e,X_j}|\dd\phi_j(e)|=1,
 \qquad
 \left|\sum_v b_{A_j,X_j}(v)\phi_j(v)\right|
 =(I_{X_j}^{\rm eff})^{-1}\longrightarrow\infty,
 \label{eq:main-neck-obstruction}
\end{equation}
which is an explicit macroscopic-neck obstruction.  If a cut has zero
capacity already at finite cutoff, that disconnected-support cut is retained
directly, without dividing by its capacity.
\end{theorem}

The discrete Sobolev--Nash--heat-kernel mechanism behind this theorem is
standard in spirit; see, for example, Chung~\cite[Chs.~10--11]{Chung1997} and
Saloff-Coste~\cite{SaloffCoste2002}.  Because the regulator normalization is
part of the present result, Appendix~\ref{supp:continuum} gives the complete
finite proof, including the constants, the $L^1\!\to L^\infty$ heat-semigroup
bound, and the derivation of the $R^{3/2}$ counting law.

Thus compact spatial control is derived on the positive cut branch and is
replaced by a concrete obstruction on the collapsed branch.  Uniform spatial
ranks and tails alone do not establish spacetime compactness.  The
common-cylinder Hodge/action packet also retains convergence of transported
finite screens and temporal control of the field coefficients.
Appendix~\ref{supp:continuum} gives the compact-screen upgrade from this
packet to the strong convergence needed for quadratic Palatini insertions.

\subsection{Curvature propagation from initial energy}

The cut estimate does not control curvature. We therefore retain a separate
positive energy for the actual electric and magnetic records and propagate
it from the initial cut. The cancellation below is checked before taking
norms, so the large skew principal operator does not enter the growth
constant; compare \cite{ACY1998,KRS2015}.

On a compact foliated cylinder $K=[0,T_*]\times\Sigma$, let
$y_X(t)=(E_X(t),B_X(t))$ be the actual retained electric/magnetic curvature
coefficients in the reconstructed observer frame, with positive physical Gram
$M_X(t)=\operatorname{diag}(M_{E,X},M_{B,X})\succ0$.  Put
\begin{equation}
 z_X(t)^2:=y_X(t)^TM_X(t)y_X(t).
 \label{eq:main-curvature-energy}
\end{equation}
The finite Bianchi/transport propagation writer is required to occur in the
form
\begin{equation}
 \dot y_X=(K_X+L_X)y_X+f_X,
 \qquad M_XK_X+K_X^TM_X=0,
 \label{eq:main-curvature-propagation-writer}
\end{equation}
where $K_X$ is the positive-Gram skew principal curl block, $L_X$ contains
lower-order observer, lapse, connection, and transport terms, and $f_X$ is the
complete differentiated curvature source, including any incoming boundary-energy
contribution when the cylinder has a timelike spatial boundary.  Define
\begin{equation}
 a_X(t):=\frac12\max\left\{0,
 \lambda_{\max}\!\left[
 M_X^{-1/2}(\dot M_X+M_XL_X+L_X^TM_X)M_X^{-1/2}
 \right]\right\}.
 \label{eq:main-curvature-growth}
\end{equation}
On the gravitational shared-link branch, $f_X$ contains the physically
transferred, volume-normalized Palatini metric row together with the
differentiated Cartan and finite Bianchi defects; Appendix~\ref{supp:dynamics}
keeps these terms separate.

The principal cancellation is tested on the actual record.  For example, a
recorded curl $C_X:V_{B,X}\to V_{E,X}$ has positive-Gram adjoint
$C_X^\dagger=M_{B,X}^{-1}C_X^TM_{E,X}$, and the block
$\left(\begin{smallmatrix}0&C_X\\-C_X^\dagger&0\end{smallmatrix}\right)$
is exactly $M_X$-skew.  Defining this adjoint does not prove that an actual
geometric second curl equals $-C_X^\dagger$: their discrepancy is retained
in the source or symmetric-growth record.  Consequently no hidden
$h_X^{-1}$ principal norm is discarded in the estimate.

\begin{theorem}[Initial-energy curvature propagation certificate]
\label{thm:main-curvature-propagation}
Assume \eqref{eq:main-curvature-propagation-writer}.  Then
\begin{equation}
 z_X(t)\le e^{\int_0^t a_X}
 \left[z_X(0)+\int_0^t e^{-\int_0^s a_X}
 \|f_X(s)\|_{M_X(s)}\,\dd s\right].
 \label{eq:main-curvature-energy-propagation}
\end{equation}
Suppose moreover that the actual curvature field entering the Palatini test has
a uniformly stable interpolation
\begin{equation}
 R_X=\mathcal I_Xy_X+\zeta_X,
 \qquad
 \|\mathcal I_Xy_X(t)\|_{L^2(\Sigma_t)}\le C_Iz_X(t),
 \qquad
 \|\zeta_X\|_{L^2(K)}\le C_\zeta,
 \label{eq:main-curvature-incidence}
\end{equation}
and that the lapse is bounded by $N_*$.  If, uniformly in $X$,
\begin{equation}
 z_X(0)\le Z_*,\qquad
 \int_0^{T_*}a_X(t)\,\dd t\le A_*,\qquad
 \int_0^{T_*}\|f_X(t)\|_{M_X(t)}\,\dd t\le F_*,
 \label{eq:main-curvature-budgets}
\end{equation}
then
\begin{equation}
 \|R_X\|_{L^2(K)}
 \le C_I\sqrt{N_*T_*}\,e^{A_*}(Z_*+F_*)+C_\zeta.
 \label{eq:main-derived-l2-curvature}
\end{equation}
Thus the curvature-size requirement needed below follows with $p=2>3/2$.
\end{theorem}


The fully discrete version is \Cref{thm:main-discrete-curvature};
\Cref{prop:main-action-current} gives the sharply scaled sufficient
source estimate. Both are proved in Appendix~\ref{supp:continuum}.
Neither estimate alone identifies the curvature of a limiting connection.
\subsection{Gauge-free anisotropic compactness from literal links}

\paragraph{Periodic-link notation.}
Let $N=2m+1$, $h=N^{-1}$, and
$\Lambda_h=(h\mathbb Z/\mathbb Z)^3$.  With shifts $S_i$ set
\begin{equation}
 D_i^+=(S_i-I)/h,\qquad D_i^-=(I-S_i^{-1})/h,\qquad
 D_i^0=(D_i^++D_i^-)/2.
 \label{eq:main-open-differences}
\end{equation}
The spatial $H_h^1$ norm is the sum of the discrete $L^2$ norm and the
$L^2$ norms of these forward differences; coefficient norms use a fixed
positive matrix inner product. The nonlinear evolution in the next section
uses this same regulator.

\label{subsec:main-literal-link-compactness}

We next identify the nonlinear curvature from the original link record.
Unlike a compactness assumption on the full connection, this route requires
strong convergence only of the spatial coefficients; the temporal
coefficient may remain weakly convergent. On the periodic branch let
\begin{equation}
 U_{i,h}=\exp(hA_{i,h}),\qquad
 Z_{i,h}=\frac{U_{i,h}-I}{h},
 \label{eq:main-link-coordinate}
\end{equation}
and let $E_{i,h}$ and $F_{ij,h}$ denote respectively the normalized literal
electric link and spatial plaquette-curvature records in the same common
trivialization.  All coefficient norms below use a fixed positive Euclidean
matrix norm; no indefinite Lorentz norm is used in the compactness estimate.

\begin{theorem}[Gauge-free anisotropic literal-link compactness]
\label{thm:main-literal-link-compactness}
On an odd periodic spatial regulator, suppose the literal links remain in one
dimensionless identity chart,
\begin{equation}
 h|A_{i,h}|\le\delta_*,
 \label{eq:main-link-chart}
\end{equation}
with $\delta_*$ fixed and sufficiently small, and suppose on a compact time
interval
\begin{equation}
 \sup_h\left(
 \|A_h\|_{L^\infty_tL_h^2}+\|A_h\|_{L^2_tH_h^1}
 +\|A_{0,h}\|_{L^2}+\|E_h\|_{L^2}+\|F_h\|_{L^2}\right)<\infty.
 \label{eq:main-link-energy}
\end{equation}
Then, after extraction in the original common trivialization,
\begin{equation}
 \mathcal I_hA_h\to A\ \hbox{strongly in }L^2,
 \qquad
 \mathcal I_hA_{0,h}\rightharpoonup A_0\ \hbox{weakly in }L^2,
 \label{eq:main-link-convergence}
\end{equation}
and the weak $L^2$ limits of the literal electric and magnetic records are
exactly the components of
\begin{equation}
 R(\omega)=\dd\omega+\omega\wedge\omega,
 \qquad \omega=A_0\,\dd t+\sum_iA_i\,\dd x^i.
 \label{eq:main-link-curvature-limit}
\end{equation}
No Coulomb condition, gauge smallness, temporal-mean bound,
$D^+A_{0,h}$ bound, or $\partial_tA_{0,h}$ bound is required.  The
identity-link chart and the anisotropic energy bound
\eqref{eq:main-link-energy} remain explicit hypotheses.
\end{theorem}

The proof is given in Appendix~\ref{supp:continuum}.  Its key point is to
compactify $Z_{i,h}$ rather than differentiate the logarithm in time: the
literal electric identity gives a cutoff-uniform $H_h^{-2}$ time-derivative
bound for $Z_h$, while $Z_h-A_h\to0$ strongly.  Weak--strong passage in the
exact electric equation and the second-order plaquette expansion then
identify every component of \eqref{eq:main-link-curvature-limit}.


A second route assumes compactness of the full connection and verifies the
Cartan and interface residuals explicitly; see
\Cref{ass:main-connection-compactness,thm:main-hodge-connection,thm:supp-curvature-compactness}.
The limit theorem allows either route, with the respective hypotheses
unchanged.
\subsection{The conditional Einstein limit}

The finite and analytic steps can now be assembled. The hypotheses below
specify which records occur, which finite structural tests pass, and which
bounds persist on a common cofinal sequence. The constructive existence
question is treated separately in Section~\ref{sec:constructive-main}.

Call a cofinal regulator a \emph{complete gravitational renewal cylinder} when
the geometric packet, complete accepted determining-field cylinder, generic
protected reward-duration record, and physical endpoint, propagation and
Cartan records occur on compatible restrictions of one operational process.
The records include the physical test lifts, action calibration, actual
curvature interpolation and common-cylinder comparison maps used above.
The word \emph{gravitational} labels the sector being tested; it does not mean
that an Einstein Lagrangian is stored in the entrance record.


\begin{theorem}[General renewal-to-Einstein certification]
\label{thm:main-renewal-einstein}
For a complete gravitational renewal cylinder, consider a cofinal
sequence of Lorentzian relational regulators of \Cref{mt:adm}.  On every
compact spacetime cylinder $K$, define the following five certification
conditions on compatible restrictions of the same operational process.
\begin{enumerate}[label=\textup{(E\arabic*)},leftmargin=3em]
\item The complete accepted cylinder and generic protected reward-duration record
      reconstruct the determining kernel and scalar common action.  Its actual
      first variations tend to zero on a fixed determining metric-test core.
      The scaled bound \eqref{eq:main-physical-stationarity}, with the realized
      subsequence of \Cref{prop:supp-physical-test-stationarity}, is a sufficient
      operational route to this condition.  The physical coframe lifts of these
      tests have uniformly bounded, convergent coefficient tensors in the
      curvature and volume insertions.
\item Complete resolved-memory subtraction and exact or summably stable binary
      face subdivision give the first-order score classified by       \Cref{thm:main-gravitational-classification}, with the physical
       first-variation remainder condition
       \eqref{eq:main-first-variation-reduction} on the same metric and
       connection tests. Pointwise convergence of action values is not used
       in place of this condition.
\item In a fixed physical action calibration the spacetime-constant
      coefficients converge,
      $(\alpha_X,\beta_X,\lambda_X)\to(\alpha,\beta,\lambda)$, with
      $\beta\ne0$ the coefficient of $L_{\rm Palatini}$.  The spinless
      connection Euler residual tends to zero and the Cartan map is uniformly
      nondegenerate on the retained coframe chart.
\item The endpoint packet satisfies the hypotheses of
      \Cref{thm:main-spatial-screen}, with
      $\liminf_X I_X^{\rm eff}>0$.  Its common-cylinder compact-screen packet,
      including transported-screen convergence and temporal control, gives a
      nondegenerate limit with $e_X\to e$ strongly in local $L^2$ and
      $\sup_X\|e_X\|_{L^6(K)}<\infty$.
\item The actual curvature records satisfy either the continuous propagation
      certificate of \Cref{thm:main-curvature-propagation} or the fully discrete
      certificate of \Cref{thm:main-discrete-curvature}, with uniform initial,
      growth, transfer and source budgets and stable curvature interpolation.
      The scaled action-current estimate
      \eqref{eq:main-action-source-budget} may supply the source budget.       For the connection/curvature handoff, either the connection satisfies the
       separate compactness certificate \Cref{ass:main-connection-compactness},
       with \Cref{thm:main-hodge-connection} as a sufficient Hodge--temporal
       realization and the two Cartan identification residuals in
       \eqref{eq:main-identified-curvature-hypotheses}, or the periodic literal-link
       branch satisfies \Cref{thm:main-literal-link-compactness}, which directly
       supplies strong spatial connection compactness and identifies the full
       nonlinear curvature from the electric and magnetic records.
\end{enumerate}
If conditions \textup{(E1)--(E5)} hold on a cofinal tail, then a subsequence
has zero limiting torsion and, for each determining compactly supported
inverse-metric test $k^{\mu\nu}$, the metric first variation converges to
\begin{equation}
 \chi\int_K\sqrt{-g}\,
       (G_{\mu\nu}+\Lambda g_{\mu\nu})k^{\mu\nu},\qquad \chi\ne0.
 \label{eq:main-einstein-variation}
\end{equation}
Consequently
\begin{equation}
 \boxed{G_{\mu\nu}+\Lambda g_{\mu\nu}=0}
 \label{eq:main-vacuum-einstein}
\end{equation}
holds distributionally on the determining test core, and on its closure in
the declared test topology.
\end{theorem}
\begin{proof}[Proof outline]
The calibrated connection equation and the uniform Cartan inverse give
vanishing limiting torsion. The cut and transport estimates supply the
strong coframe limit; the propagation estimate supplies the curvature bound;
and the selected connection criterion identifies its weak limit.
Strong interpolation then passes the coframe bivectors against the weak
curvature, while the physical first-variation remainder tends to zero.
On the torsion-free branch the Holst coframe variation vanishes and the
nonzero Palatini coefficient leaves the Einstein insertion displayed above.
Stationarity on the determining core proves the conclusion.
The full argument, including the volume term and the common subsequence,
is \Cref{thm:supp-renewal-palatini,cor:supp-renewal-einstein} in
Appendix~\ref{supp:continuum}.
\end{proof}

\begin{corollary}[Einstein limit or failed certificate]
\label{thm:main-einstein-alternative}
For a complete gravitational renewal cylinder, apply the finite
reconstructions and conditions \textup{(E1)--(E5)} on a selected cofinal
sequence. If the conditions pass on a cofinal tail, the sequence has the
Einstein subsequential limit of \Cref{thm:main-renewal-einstein}.
Otherwise the failing algebraic test, residual, calibrated comparison,
positive margin, propagation budget or compactness tail is retained on
the same physical records. The complete witness list and its cofinal
quantifiers are given in Appendix~\ref{subsec:supp-failed-certificates}.
Failure is an obstruction to this sufficient certification procedure,
not a proof that every Einstein subsequence is impossible.
\end{corollary}
\begin{proof}
The positive conclusion is \Cref{thm:main-renewal-einstein}.
The complementary witnesses are obtained by negating the corresponding
finite or cofinal tests, as detailed in
Appendix~\ref{subsec:supp-failed-certificates}.
\end{proof}

\begin{remark}[Certification is not the definition of the constructive branch]
\label{rem:main-certification-not-tautology}
Conditions \textup{(E1)--(E5)} recognize limits of general complete
cylinders. They do not define the constructive law of
\Cref{thm:main-open-3plus1}, whose estimates are proved directly. Conversely,
the recognition theorem is not restricted to that law.
\end{remark}

The nonzero Palatini coefficient is distinct from invertibility of the
Palatini--Holst connection operator.  A nonzero pure Holst coefficient can
eliminate torsion, but its torsion-free metric variation does not supply the
Einstein equation.  Likewise the propagated curvature bound supplies size,
whereas \eqref{eq:main-identified-curvature-hypotheses} identifies the nonlinear
connection curvature.  These distinctions prevent either condition from being
hidden inside a generic ``positive branch.''


\section{A constructive nonsymmetric \texorpdfstring{$3+1$}{3+1} Einstein class}
\label{sec:constructive-main}

\subsection{The local harmonic law and uniform realization}
\label{subsec:main-open-3plus1}

The preceding theorem recognizes Einstein limits from operational
certificates. We now address a different question: whether a nonempty
nonsymmetric class admits cutoff-uniform evolution and convergence.
For the declared local harmonic law, these bounds are proved directly on
a fixed Sobolev neighbourhood; the five recognition conditions are not
assumed as its definition. The law is not identified with the raw finite
Euler flow of the reconstructed action.

We use the odd periodic regulator and the forward, backward and centred
differences of \eqref{eq:main-open-differences}.
For $g=\eta+q$ in a small Lorentzian chart, write the harmonic normal row as
\begin{equation}
 a(g)q_{tt}+2b^i(g)\partial_iq_t-c^{ij}(g)\partial_i\partial_jq
             -F(g,\partial g)=0,
 \qquad a=-g^{00},\ b^i=-g^{0i},\ c^{ij}=g^{ij},
 \label{eq:main-harmonic-normal-row}
\end{equation}
where $F$ is the analytic first-jet source of the harmonic-reduced vacuum
Ricci equation, equivalently of its trace-reversed reduced Einstein system. This identifies the target equation of the constructive
realization explicitly. Its selection is not derived again from the root
race, nor is consistency of this writer used as a substitute for the
independent action classification and stationarity certificates.  Put
\begin{equation}
 \mathsf K_b=\sum_i\{M_{b^i}D_i^0+D_i^0M_{b^i}\},
 \label{eq:main-open-skew-transport}
\end{equation}
and, with $Y_h=(v,D_1^+q,D_2^+q,D_3^+q)$,
\begin{equation}
 \mathsf G_h(q,v)=F(g,Y_h)
 -\sum_{i,j}Dc^{ij}(g)[D_i^+q]D_j^+q
 +\sum_iDb^i(g)[D_i^+q]v.
 \label{eq:main-open-compensator}
\end{equation}
The local evolution law is
\begin{equation}
 q_t=v,\qquad
 a(g)v_t=\sum_{i,j}D_i^-\!\bigl(c^{ij}(g)D_j^+q\bigr)
           -\mathsf K_bv+\mathsf G_h(q,v).
 \label{eq:main-open-writer}
\end{equation}
The divergence and skew placements are fixed by the positive quadratic
principal energy.  The compensator removes exactly the first-derivative terms
introduced by these placements, so the smooth continuum limit is
\eqref{eq:main-harmonic-normal-row}.  A parity obstruction of the bare
$A_3\cong D_3$ execution graph requires one additional odd-parity execution
incidence for a cutoff-uniform Cartesian Sobolev realization; this incidence
is used only to execute the writer and is not added to the root race entering
\eqref{eq:bracket-drift-main}.  The complete proof is in
Appendix~\ref{supp:open-3plus1}.  In the main theorem we use the discrete
Sobolev norms
\begin{equation}
 \|u\|_{H_h^r}^2=\sum_{|\alpha|\le r}\|D^\alpha u\|_h^2,
 \qquad
 \|X\|_{\mathcal X_h^r}^2=\|q\|_{H_h^{r+1}}^2+\|v\|_{H_h^r}^2,
 \label{eq:main-open-sobolev-norms}
\end{equation}
where $D^\alpha$ is formed from the commuting forward differences.

For an integer $s\ge11$ define the compatible continuum data class
\begin{equation}
 \begin{split}
 \mathcal D_s(\varepsilon)=\bigl\{(\gamma,K):\;&
 \gamma\in H^{s+1}(\mathbb T^3),\ K\in H^s(\mathbb T^3),\\
 &R(\gamma)+(\tr_\gamma K)^2-|K|_\gamma^2=0,\qquad
 \nabla_\gamma^iK_{ij}-\partial_j\tr_\gamma K=0,\\
 &\|\gamma-I\|_{H^{s+1}}+\|K\|_{H^s}<\varepsilon\bigr\}.
 \end{split}
 \label{eq:main-open-data-class}
\end{equation}
This is relative openness in the vacuum constraint set; it is not an ambient
open ball of fixed nonzero constraint violations.

\begin{theorem}[Constructive nonsymmetric $3+1$ renewal--Einstein realization]
\label{thm:main-open-3plus1}
There are $\varepsilon_s,T_s,h_s,C_s>0$, independent of the cutoff, with the
following properties.
\begin{enumerate}[label=\textup{(O\arabic*)},leftmargin=3em]
\item For every $h<h_s$, every real ten-component initial record
      $X_h(0)=(q_h(0),v_h(0))$ in the fixed ball
      $\|q_h(0)\|_{H_h^{s+1}}+\|v_h(0)\|_{H_h^s}\le\varepsilon_s$
      has a unique solution of \eqref{eq:main-open-writer} on $[0,T_s]$, and
      \begin{equation}
       \sup_{t\le T_s}
       \bigl(\|q_h(t)\|_{H_h^{s+1}}+\|v_h(t)\|_{H_h^s}\bigr)
       \le C_s\|X_h(0)\|_{\mathcal X_h^s}.
       \label{eq:main-open-uniform-time}
      \end{equation}
      No symmetry, Fourier-support restriction, analyticity, or
      cutoff-shrinking amplitude is imposed.
\item Every sufficiently small $(\gamma,K)\in\mathcal D_s(\varepsilon_s)$
      has a harmonic finite preparation for which the whole cutoff sequence
      converges on $[0,T_s]\times\mathbb T^3$ to a unique metric $g$ with the
      prescribed initial data.  Writing
      $\varepsilon_0=\|\gamma-I\|_{H^{s+1}}+\|K\|_{H^s}$, one has
      \begin{equation}
       \boxed{G(g)=0,\qquad
       \sup_{t\le T_s}\left(
       \|g_h-g\|_{H^{s-1}}+\|\partial_tg_h-\partial_tg\|_{H^{s-2}}
       +\|\partial_t^2g_h-\partial_t^2g\|_{H^{s-3}}
       \right)\le C_sh\varepsilon_0.}
       \label{eq:main-open-einstein-rate}
      \end{equation}
      Here $g_h$ denotes the canonical interpolation of $\eta+q_h$.
\item For the compatible preparations in item~\textup{(O2)}, the harmonic
      subsidiary field and Einstein residual obey, at the corresponding lower orders,
      \begin{equation}
       \sup_{t\le T_s}
       \bigl(\|c(g_h)\|_{H^{s-2}}+\|\partial_tc(g_h)\|_{H^{s-3}}
                  +\|G(g_h)\|_{H^{s-3}}\bigr)
       \le C_sh\varepsilon_s,
       \label{eq:main-open-subsidiary}
      \end{equation}
      while the positive $A_3$ rate lift remains in an interior cone and its
      spatial cut margin stays uniformly positive.  In particular the
      reconstructed metric is uniformly noncollapsed and its metric curvature
      is uniformly bounded on the common slab and converges to the curvature
      of $g$ at the lower Sobolev order supplied by
      \eqref{eq:main-open-einstein-rate}.
\item The differential writer has a genuinely finite renewal realization.
      With segment duration $\tau_h\asymp h^2$, fixed-degree local
      Taylor--Hermite steps and Sobolev-scaled finite precision produce
      committed successor words whose complete acceleration defect is
      $O(\varepsilon h^4)$ in $C_tH_h^s$ and whose differentiated defect is
      $O(\varepsilon h^2)$ in $C_tH_h^{s-3}$ on the common slab;
      \Cref{thm:supp-open-finite-words} gives stronger rates for the explicit construction.
      With the compatible finite initialization specified there, every such
      certified history inherits the bounds above.  This statement
      concerns finite-range information flow; it does not identify the
      arithmetic service time with the physical root clock.
\end{enumerate}
Consequently the positive Einstein population is not confined to homogeneous
or one-dimensional inhomogeneous reductions: it contains a relatively open
small-data class with arbitrary three-dimensional dependence.  The next
result shows that the differential writer itself is also not isolated.
\end{theorem}

\begin{remark}[Scope of the constructive law]
\label{rem:main-open-3plus1-scope}
This is a cutoff-uniform small-data realization, not a universal attraction
result or an identification with exact finite Euler flow. The independent
exact-action question is summarized at the end of this section and analyzed
in Appendices~\ref{supp:exact-action-audit} and~\ref{supp:exact-action-slow}.
\end{remark}


\subsection{Structural stability of renewal laws and action-prepared finite seeds}
\label{subsec:main-open-law-basin}

The fixed writer of \Cref{thm:main-open-3plus1} is the centre of a larger
finite law class.  Let
\begin{equation}
 \Lambda_h=-\sum_{i=1}^3D_i^-D_i^+,
 \qquad
 \mathcal B=\{B=B^T\in\mathbb R^{10\times10}:\|B\|_{\op}<b_0\},
 \label{eq:main-law-ball}
\end{equation}
where $b_0>0$ is fixed below the spatial coercivity margin of the common small
metric chart.  Thus $\mathcal B$ is a $55$-dimensional open ball.  For
$B\in\mathcal B$ define
\begin{equation}
 q_t=v,\qquad
 a(g)v_t=\sum_{i,j}D_i^-\!\bigl(c^{ij}(g)D_j^+q\bigr)
       -\mathsf K_bv+\mathsf G_h(q,v)-h^2B\Lambda_h^2q.
 \label{eq:main-open-law-family}
\end{equation}
The added fourth-order stencil has fixed execution radius.  It need not be
small in the top ultraviolet vector-field norm; its incidence and sign are
instead controlled by the same positive shifted energy.

For a finite chronological realization of \eqref{eq:main-open-law-family},
let $I_j$ have duration $\tau_j\asymp h^2$.  A candidate physical letter $e$
at a live prefix $\pi$ carries a piecewise-Hermite segment $Q_e$ with actual
acceleration defect
\begin{equation}
 f_e=Q_e''-V_{B,h}(Q_e,Q_e'),
 \label{eq:main-law-residual}
\end{equation}
and, with unit cost assigned to a grammar or execution failure, define
\begin{equation}
 c_{h,j}(\pi,e)=\varepsilon^{-2}
 \int_{I_j}\bigl(\|f_e\|_{H_h^s}^2
       +\|\partial_tf_e\|_{H_h^{s-3}}^2\bigr)\,\dd t,
 \qquad C_h=\sum_jc_{h,j}.
 \label{eq:main-law-cost}
\end{equation}
The cost is determined by the physical letter and its prefix, not by its
probability.

Finally let $\mathscr H_h(\gamma,\pi,N,\beta)$ be the canonical Hamiltonian
obtained from the reconstructed phase-compatible action after stationary
connection and Legendre elimination, and set
\begin{equation}
 \mathscr C_h(X)=
 \left(-\nabla_N\mathscr H_h,\ \nabla_\beta\mathscr H_h\right)
       \big|_{N=1,\,\beta=0}.
 \label{eq:main-action-initial-constraints}
\end{equation}
This is the original four-row finite initial constraint map; it is not the
residual of the local evolution writer. Finite zero-set preparation and its
transport to continuum Cauchy data are separate assertions. The latter uses
comparison of the actual stationary/Legendre maps, made explicit in
\Cref{ass:supp-initial-consistency,lem:supp-initial-elimination-consistency}.

\begin{theorem}[Structurally stable renewal laws with action-prepared finite seeds]
\label{thm:main-open-law-basin}
After decreasing the common small-data radius if necessary, the following
statements hold with constants independent of the cutoff.
\begin{enumerate}[label=\textup{(L\arabic*)},leftmargin=3em]
\item For every $B\in\overline{\mathcal B}$ the law
      \eqref{eq:main-open-law-family} has the same fixed-radius nonlinear
      lifespan as \Cref{thm:main-open-3plus1}.  There is an energy
      $\mathscr E^B_{s,h}$ uniformly equivalent to
      $\|X\|_{\mathcal X_h^s}^2$ such that
      \begin{equation}
       (\mathscr E^B_{s,h})'
       \le C_s\mathscr E^B_{s,h}
          +C_s(\mathscr E^B_{s,h})^{3/2}.
       \label{eq:main-law-energy}
      \end{equation}
\item Let $B_h\in\overline{\mathcal B}$ be arbitrary, with no convergence of
      $B_h$ required.  For the compatible data of
      \eqref{eq:main-open-data-class}, the whole sequence of solutions of
      \eqref{eq:main-open-law-family} converges to the same harmonic vacuum
      development $g$ as the central writer, and
      \begin{equation}
       \sum_{j=0}^3
       \|q_{B_h,h}^{(j)}-q_{0,h}^{(j)}\|_{C_tH_h^{s-1-j}}
          \le C_sh\|B_h\|_{\op}\varepsilon.
       \label{eq:main-law-jet-comparison}
      \end{equation}
      Hence the metric, subsidiary, noncollapse and curvature conclusions of
      \Cref{thm:main-open-3plus1} are uniform on the entire coefficient ball.
\item Let $K_{h,j}(e\mid\pi)$ be arbitrary history-dependent finite transition
      rows satisfying, at every live prefix,
      \begin{equation}
       \sum_eK_{h,j}(e\mid\pi)c_{h,j}(\pi,e)<M\tau_jh^4.
       \label{eq:main-law-row-moment}
      \end{equation}
      No independence, invariant starting law, or per-letter small-defect
      condition is imposed.  Then
      \begin{equation}
       \mathbb EC_h\le MTh^4,
       \qquad \mathbb P\{C_h>h^2\}\le MTh^2.
       \label{eq:main-law-probability}
      \end{equation}
      On the complementary event the finite history shadows a member of the
      differential-law family closely enough to have the same Einstein limit.
      For a fixed finite grammar containing a sufficiently accurate completion
      at every live source, the strict conditions
      \eqref{eq:main-law-row-moment} form nonempty relatively open convex
      subsets of the product of row simplexes and may contain strictly
      positive rows. The construction in \Cref{thm:supp-open-finite-words}
      supplies the required strict residual margin.  Along the odd refinement sequence the exceptional
      probabilities are summable, so under any common coupling the Einstein
      conclusion holds eventually almost surely.
\item For every fixed sufficiently small nonzero preparation amplitude $t$,
      the exact map \eqref{eq:main-action-initial-constraints} has a finite,
      nonsymmetric analytic zero-set chart generated from arbitrary small free
      tensor records $W$ after removal of four fixed balancing modes, with
      the free-record bound taken in $\mathcal P_h^{s+2}$ as defined in the appendix. Its
      prepared records $\mathfrak I_h(t,W)$ satisfy
      \begin{equation}
       \mathscr C_h(\mathfrak I_h(t,W))=0,\qquad
       \|\mathfrak I_h(t,W)\|\le C_s|t|.
       \label{eq:main-action-prepared-zero}
      \end{equation}
      The chart is relatively open in the complete finite zero set, has
      codimension $4N^3$, and the linearized constraint map has a right inverse
      obeying
      \begin{equation}
       \|\mathscr R_{h,t}f\|
       \le C_s\bigl(\|P_\perp f\|+|t|^{-1}|P_0f|\bigr).
       \label{eq:main-action-prepared-inverse}
      \end{equation}
      No CMC, conformal-flatness, spatial symmetry, finite Fourier support, or
      continuum-compatible metric is imposed on the free record $W$.
      On the cofinally consistent initial-action branch of
      \Cref{ass:supp-initial-consistency}, every cofinal sequence of these
      records with $\mathcal I_hW_h$ convergent in $\mathcal P^{s+1}$
      has a harmonic completion converging to the vacuum development of
      its limiting prepared data, for arbitrary
      $B_h\in\overline{\mathcal B}$. For merely bounded free records
      this conclusion is subsequential, not whole-sequence convergence
      to one fixed metric. The original four lapse--shift rows vanish
      identically at the initial cut independently of the subsequently
      chosen acceleration law.
\end{enumerate}
The law-space statement is deliberately quantitative. Here structural
stability means a common continuum limit in the stated Sobolev and residual
topologies; it does not mean attraction under a renormalization flow.  It gives a fixed open
ball of differential laws and relatively open finite probability-row
neighborhoods in a physically scaled residual topology; it does not assert a
fixed unscaled total-variation ball of arbitrary microscopic kernels or
attraction from arbitrary leading-order forces.
\end{theorem}

\begin{remark}[The robustness class is nontrivial but not arbitrary]
\label{rem:main-law-basin-scope}
The family in \Cref{thm:main-open-law-basin} can differ from the central
writer by a nonvanishing amount at the highest lattice frequencies.  The law
parameter is, however, retained during a physical history.  Unrestricted
rapid switching between two individually stable coefficient marks can
parametrically amplify a resolved high-frequency mode.  Likewise, a row may
be close in ordinary total variation while accumulating order-one failure
probability over $O(h^{-2})$ stages.  These boundaries explain why the
neighborhood is formulated by energy-compatible coefficients and complete
physical residual profiles rather than by an unscaled norm on arbitrary
kernels.  The proofs are in Appendix~\ref{supp:open-law-basin}.
\end{remark}


\subsection{Compatibility on one chronological process}
\label{subsec:main-same-cylinder}

The stability theorem controls a class of finite evolution laws. Its
compatibility with the independently reconstructed action is a further
question, because equality of the two finite dynamics has not been assumed.
The following chronological correction preserves the specified physical
marginals while retaining both the evolution-residual and stationarity scales.  Let $\nu_{h,j}$ be the physical marginal at
stage $j$ and let $Q_{h,j}(y\mid x)$ be the independently reconstructed
common-action comparator on the admitted transition support.  Put
\begin{equation}
 R_{h,j}(x,y)=\nu_{h,j}(x)Q_{h,j}(y\mid x),\qquad
 \rho_{h,j+1}=\nu_{h,j}Q_{h,j},\qquad
 d_{h,j}=\nu_{h,j+1}-\rho_{h,j+1}.
 \label{eq:main-same-cylinder-reference}
\end{equation}
For the source and target incidence maps
$Ah(x)=\sum_yh(x,y)$ and $Bh(y)=\sum_xh(x,y)$, define the source-shorted
operator
\begin{equation}
 \mathcal L_{h,j}
 =\diag(\rho_{h,j+1})
  -Q_{h,j}^{\mathsf T}\diag(\nu_{h,j})Q_{h,j}.
 \label{eq:main-thomson-operator}
\end{equation}
When $d_{h,j}\in\Ran\mathcal L_{h,j}$ and the source-preserving
Thomson correction has leverage at most $1/2$, the explicit formula
\eqref{eq:main-thomson-correction} in Appendix~\ref{supp:same-cylinder}
defines a nonnegative corrected coupling and its stochastic kernel.
It has exactly the source and target marginals specified above.
$(Q_{h,j}c)(x)=\sum_yQ_{h,j}(y\mid x)c(x,y)$.  Its occurrence energy is
\begin{equation}
 \mathcal E_{{\rm occ},h,j}
 =\langle d_{h,j},\mathcal L_{h,j}^{\dagger}d_{h,j}\rangle.
 \label{eq:main-occurrence-energy}
\end{equation}

\begin{theorem}[Same-cylinder action-occurrence Einstein compiler]
\label{thm:main-same-cylinder-einstein}
Consider the open nonsymmetric $3+1$ chart of
\Cref{thm:main-open-3plus1,thm:main-open-law-basin}.  Suppose at each
chronological stage the preceding Thomson correction exists on the same
support as the common-action comparator, and let $c_{h,j}\ge0$ be the
stable-writer residual cost with duration $\tau_{h,j}$.  Assume, uniformly,
\begin{equation}
 \begin{split}
 m_{h,j}:=\mathbb E_{R_{h,j}}c_{h,j}&\le M_0\tau_{h,j}h^4,\\
 \mathcal E_{{\rm occ},h,j}&\le M_1\tau_{h,j}h^4,\\
 \mathcal V_{h,j}:=\mathbb E_{R_{h,j}}
 |c_{h,j}-Q_{h,j}c_{h,j}|^2&\le V_0\tau_{h,j}h^4,
 \end{split}
 \qquad \sum_j\tau_{h,j}\le T.
 \label{eq:main-same-cylinder-writer-scale}
\end{equation}
Let $\Delta_{h,j}$ be the Bregman field-action gap of the same reconstructed
common action, on its supported interior proximal chart.  For the compared
source and target marginals retain the actual positive action drop
\begin{equation}
 d^+_{\mathcal A,h,j}=
 \left[\mathbb E_{\nu_{h,j}}\mathcal A_{h,j}
             -\mathbb E_{\nu_{h,j+1}}\mathcal A_{h,j}\right]_+.
 \label{eq:main-chronological-action-drop}
\end{equation}
Assume the corrected test-lift bound
\eqref{eq:supp-same-cylinder-physical-scale}, including this drop, tends
to zero on the fixed determining core.  Equal marginals are a sufficient
special case for removal of the drop; prescribed chronological marginals
need not be equal.  Finally assume the physical first-variation reduction
and calibrated coefficient/torsion conditions \textup{(E2)--(E3)} of
\Cref{thm:main-renewal-einstein}, including nonzero Palatini normalization,
and the stated noncollapse and calibration conditions on the constructive
small metric chart.  Then
\begin{enumerate}[label=\textup{(\roman*)},leftmargin=2.8em]
\item the corrected finite process has exactly the prescribed chronological
      marginals and remains on the common-action support;
\item for $C_h=\sum_jc_{h,j}(X_j,X_{j+1})$,
      \begin{equation}
       \mathbb EC_h\le
       (M_0+\sqrt{M_1V_0})Th^4,
       \qquad \Pr\{C_h>h^2\}=O(h^2);
       \label{eq:main-same-cylinder-probability}
      \end{equation}
\item on $\{C_h\le h^2\}$ the constructive metric-jet bounds imply the
      literal-link budget \eqref{eq:main-link-energy}; hence
      \Cref{thm:main-literal-link-compactness} makes connection compactness
      and nonlinear curvature identification outputs of the successful
      history rather than independent assumptions;
\item the represented first variation of the same reconstructed action tends
      to zero in mean square on the determining core;
\item every continuum cluster point on a realized cofinal subsequence where
      the residual and stationarity conclusions hold satisfies
      \begin{equation}
       G_{\mu\nu}+\Lambda g_{\mu\nu}=0
       \label{eq:main-same-cylinder-einstein}
      \end{equation}
      distributionally on that core.
\end{enumerate}
Along any common coupling with $\sum_nh_n^2<\infty$, the successful residual
branch occurs eventually almost surely; no independence between cutoffs or
between chronological stages is required.
\end{theorem}

This theorem is a compatibility result, not an identification of two
finite dynamics.  The common action is reconstructed before the chronological
correction is formed, while the writer-residual moments independently test
whether that action-supported process remains in the constructive Einstein
basin.

\subsection{Exact finite-action dynamics: realization and remaining obstructions}
\label{sec:main-exact-action}

The constructive and chronological results do not identify their laws with
the unchanged finite Euler flow. That stronger question has its own finite
constraint geometry and its own limiting scales. We state the established
conclusions together here, while retaining their different branch hypotheses.
The detailed response and boundary arguments are in
Appendix~\ref{supp:exact-action-audit}; slow compatibility and obstruction
theory are in Appendix~\ref{supp:exact-action-slow}.

\begin{theorem}[Summary of exact-action realization and compatibility]
\label{thm:main-exact-action-response}
\label{thm:main-exact-action-summary}
For the original finite action on the respective branches specified in
Appendices~\ref{supp:exact-action-audit} and~\ref{supp:exact-action-slow}:
\begin{enumerate}[label=\textup{(X\arabic*)},leftmargin=3em]
\item The regular finite response admits lossless transverse elimination and
an exact positive-reproduction/scalar--trace decomposition. Nontrivial
source-neutral prepared families and the stated all-cutoff active families
exist. These static results do not assert invariance under the exact flow.
\item At the certified event on the three-site periodic regulator, the weak
Poisson block has rank twenty-four and a five-dimensional kernel; the rapid
generator has stable, unstable and centre dimensions eight, eight and
thirteen. Under the complete amplitude-graded chart assumptions, the
unchanged action admits an exactly constrained boundary-selected history
on an interval of length proportional to the square of the amplitude.
The first field and multiplier tilt are preserved, and both metric and
literal-link curvatures are bounded by a constant times the amplitude.
The constants here are fixed-cutoff constants.
\item On the separate invariant normal-factor chart, residence through a
fixed physical time requires the rapid and slow exponentially small normal
widths of \Cref{thm:main-exact-residence}. This is a necessary local-residence
condition, not a curvature-blowup statement.
\item Under the stated initialization hypotheses, the three constant-shift
compatibility rows are automatic. For slow limits with bounded rescaled rates, the remaining two
initial rows are measured by the original-action acceleration defect.
On the record-preserving branch, the subsequent necessary quadratic
harmonic-slope map factors through a two-dimensional space. Its fixed
cokernel obstruction must vanish before the two-equation root problem
can be solved; a root alone does not establish a positive-time trajectory.
\end{enumerate}
The assertions use their stated, potentially different, charts. They do not
jointly supply a cutoff-uniform exact-action Cauchy theorem.
\end{theorem}
\begin{proof}
The response and preparations are proved in
\Cref{thm:supp-exact-transverse,thm:supp-exact-scalar-reconstruction,thm:supp-exact-active-regular,thm:supp-exact-neutral-preparation}.
The fixed-event spectrum and boundary history are
\Cref{thm:supp-certified-weak-rank,thm:supp-certified-rapid-spectrum,thm:main-exact-boundary};
\Cref{thm:main-exact-residence} proves the normal-width assertion under its
separate assumptions. Finally,
\Cref{thm:supp-exact-initial-gate,thm:supp-exact-slope-gate,thm:supp-exact-rank-two-factorization,thm:supp-exact-two-equation-slope}
give the initial defect and the necessary slow-slope conditions.
\end{proof}

Boundary selection fixes unstable coordinates at a terminal cut and is
therefore not a forward microscopic selection rule. The remaining problems
are positive slow-time propagation, physical forward selection and uniform
control under spatial refinement. They do not enter the constructive
Einstein theorem above.


\section{Explicit vacuum sectors and operational diagnostics}
\label{sec:vacuum-main}

The general limit and the nonsymmetric realization answer the principal
continuum questions. The following sectors serve complementary purposes:
the homogeneous model makes lapse variation and action--law matching
explicit, whereas the Gowdy model retains nonlinear spatial variation and
an exact discrete momentum invariant. Their special hypotheses are not
substitutes for the general certification conditions.
\subsection{Homogeneous renewal coordinates and the full lapse constraint}
\label{subsec:homogeneous-renewal-main}

On the $A_3$ root-symmetric branch write
\begin{equation}
 k_{\alpha,h}(t)=\frac{\kappa(t)}{8h^2},
 \qquad
 m_h(t)=\varrho(t)h^3,
 \qquad \kappa(t)>0,\quad\varrho(t)>0.
 \label{eq:main-general-homogeneous-rates}
\end{equation}
The root identity gives $\mathcal B=\kappa I_3$ and reversal symmetry gives
$\beta=0$.  If $g_{ij}=a^2\delta_{ij}$, the ADM inversion
\eqref{eq:adm-inversion-main} is therefore exactly
\begin{equation}
 a=\varrho^{1/3},\qquad
 N=\varrho^{1/3}\kappa^{1/2},\qquad
 \varrho=a^3,\qquad
 \kappa=\frac{N^2}{a^2}.
 \label{eq:main-isotropic-renewal-map}
\end{equation}
Thus $\kappa$ is not the scale factor: together with the physical density it
contains the lapse information that must remain independent until after
variation.

On a compact interval and torus of coordinate volume $V_0$, the torsion-free
Palatini--volume branch of \Cref{thm:main-renewal-einstein}, equivalently its
Einstein--Hilbert representative, reduces after the standard
Gibbons--Hawking--York Dirichlet boundary completion
\cite{York1972,GibbonsHawking1977} and division by the nonzero constant
$\chi V_0$, to
\begin{equation}
 S[a,N]=\int\left[-\frac{6a\dot a^2}{N}-2\Lambda Na^3\right]\dd t.
 \label{eq:main-homogeneous-action-an}
\end{equation}
Here the induced spatial metric is fixed at the initial and final slices, so
$\delta a(t_0)=\delta a(t_1)=0$.  The lapse remains an independent bulk
variation and is not gauge-fixed before the Euler equations are taken.  The
four-dimensional reduction and the exact cancellation of the second-derivative
endpoint term by the boundary completion are written out in
Appendix~\ref{supp:homogeneous-dynamics}.

\begin{theorem}[Vacuum Friedmann equations in renewal coordinates]
\label{thm:main-homogeneous-friedmann}
Let $u=\log\kappa$, $v=\log\varrho$, and
$\mathcal H=\dot a/(Na)$, with $D_\tau=N^{-1}\dd/\dd t$.  Under
\eqref{eq:main-isotropic-renewal-map}, the action
\eqref{eq:main-homogeneous-action-an} becomes
\begin{equation}
 S[\kappa,\varrho]
 =\int\left[-\frac{2\dot\varrho^2}
 {3\varrho^{4/3}\sqrt\kappa}
 -2\Lambda\varrho^{4/3}\sqrt\kappa\right]\dd t.
 \label{eq:main-renewal-rate-action}
\end{equation}
Its Euler expressions, defined by
$\delta S=\int(E_u\delta u+E_v\delta v)\dd t$, are
\begin{equation}
 E_u=Na^3(3\mathcal H^2-\Lambda),\qquad
 E_v=Na^3\left[4D_\tau\mathcal H+
 \frac83(3\mathcal H^2-\Lambda)\right].
 \label{eq:main-renewal-euler}
\end{equation}
Consequently stationarity in the two independent positive renewal variables is
equivalent to
\begin{equation}
 3\mathcal H^2=\Lambda,\qquad D_\tau\mathcal H=0,
 \label{eq:main-vacuum-friedmann}
\end{equation}
which is the complete vacuum Einstein equation within the spatially flat
homogeneous isotropic ansatz, including the lapse constraint.
\end{theorem}

\begin{proof}
The map $(u,v)\mapsto(\log a,\log N)=(v/3,u/2+v/3)$ has nonzero
Jacobian, so no homogeneous variation is lost.  Varying
\eqref{eq:main-homogeneous-action-an} before fixing $N$ gives
\[
 E_N=2a^3(3\mathcal H^2-\Lambda),\qquad
 E_a=6Na^2(2D_\tau\mathcal H+3\mathcal H^2-\Lambda).
\]
The chain rule gives $E_u=(N/2)E_N$ and
$E_v=(a/3)E_a+(N/3)E_N$, proving
\eqref{eq:main-renewal-euler}.  The two equations are therefore equivalent to
\eqref{eq:main-vacuum-friedmann}.  Appendix~\ref{supp:homogeneous-dynamics}
checks directly that these are exactly the $00$ and spatial Einstein equations
for the reconstructed metric, rather than relying on a symmetry-reduction
principle.
\end{proof}

\begin{corollary}[Flat and de Sitter rates are stationary renewal solutions]
\label{cor:main-derived-desitter-rates}
For $\Lambda>0$ put $H_0=\sqrt{\Lambda/3}$ and let
$\tau(t)=\int^tN(s)\dd s$.  Every positive connected stationary solution of
\eqref{eq:main-renewal-rate-action} has
\begin{equation}
 a(\tau)=a_0e^{\sigma H_0(\tau-\tau_0)},\qquad
 \varrho=a^3,\qquad
 \kappa=N^2/a^2,
 \qquad \sigma\in\{-1,1\}.
 \label{eq:main-stationary-homogeneous-solutions}
\end{equation}
After variation, choose proper-time gauge $N=1$.  On the expanding branch,
\begin{equation}
 k_{\alpha,h}(t)=
 \frac{a_0^{-2}e^{-2H_0(t-t_0)}}{8h^2},\qquad
 m_h(t)=a_0^3e^{3H_0(t-t_0)}h^3.
 \label{eq:desitter-regulator-main}
\end{equation}
For $\Lambda=0$, $a$ is constant; for $\Lambda<0$ there is no positive real
solution in this spatially flat vacuum sector.
\end{corollary}

The exponential factor in \eqref{eq:desitter-regulator-main} is therefore an
output of the stationary renewal equations on this branch, not an input used to
obtain those equations.  The result does not select $\Lambda$, the expanding
versus contracting sign, or the initial scale $a_0$.

\subsection{A finite lapse-varied renewal action}
\label{subsec:finite-homogeneous-action-main}

The same distinction can be kept at finite cutoff.  The construction is a
midpoint variational discretization in the standard sense of discrete
mechanics~\cite{MarsdenWest2001}, but every equation needed below is derived
explicitly.  Put
$q=a^{3/2}=\sqrt\varrho$ and let
$s_j=N_j\Delta t_j>0$ be independent lapse-weighted intervals.  With
$\bar q_j=(q_j+q_{j+1})/2$, $A_0=8/3$, and $B_0=2\Lambda$, define
\begin{equation}
 S_{\rm d}(q,s)
 =-\sum_{j=0}^{M-1}\left[
 A_0\frac{(q_{j+1}-q_j)^2}{s_j}
 +B_0s_j\bar q_j^2\right].
 \label{eq:main-discrete-homogeneous-action}
\end{equation}
This discretizes the boundary-completed reduced action, with fixed
endpoint amplitudes and independently varied lapse-weighted intervals.
The full endpoint convention and the conversion back to renewal variables
are given in Appendix~\ref{supp:homogeneous-dynamics}.

\begin{theorem}[Finite stationary classification and de Sitter convergence]
\label{thm:main-finite-homogeneous-stationarity}
Let $v_j=(q_{j+1}-q_j)/s_j$.  The full Euler equations of
\eqref{eq:main-discrete-homogeneous-action} are
\begin{align}
 A_0v_j^2-B_0\bar q_j^2&=0,\label{eq:main-discrete-lapse-constraint}\\
 2A_0(v_j-v_{j-1})
 -B_0(s_{j-1}\bar q_{j-1}+s_j\bar q_j)&=0.
 \label{eq:main-discrete-scale-euler}
\end{align}
For $\Lambda>0$, put $\omega=\sqrt{B_0/A_0}=\sqrt{3\Lambda/4}$.
Every positive stationary solution has one common sign
$\sigma\in\{-1,1\}$ and
\begin{equation}
 q_{j+1}=q_j
 \frac{1+\sigma\omega s_j/2}{1-\sigma\omega s_j/2},
 \qquad 0<\omega s_j/2<1.
 \label{eq:main-finite-stationary-recurrence}
\end{equation}
Conversely every such recurrence is stationary.  If
$\tau_j=\sum_{i<j}s_i$, $T=\sum_js_j$, $h_t=\max_js_j$, and
$\omega h_t/2\le b<1$, then
\begin{equation}
 \max_j\left|\log\frac{q_j}{q_0}-\sigma\omega\tau_j\right|
 \le \frac{\omega^3T}{12(1-b^2)}h_t^2.
 \label{eq:main-finite-profile-convergence}
\end{equation}
Hence $a_j=q_j^{2/3}$, $\varrho_j=q_j^2$, and in proper-time gauge
$\kappa_j=q_j^{-4/3}$ converge with second-order logarithmic error to the
stationary flat/de~Sitter profiles of
\Cref{cor:main-derived-desitter-rates}.
\end{theorem}

The proof, including the lapse variation and the exact rational recurrence, is
in Appendix~\ref{supp:homogeneous-dynamics}.  In particular, fixing $N=1$
before variation would lose \eqref{eq:main-discrete-lapse-constraint} and admit
spurious non-Einstein homogeneous trajectories.

\subsection{Measured acceptance law and action--dynamics matching}
\label{subsec:operational-homogeneous-main}

To relate the finite variational action directly to an operational probability
law, we next exhibit a finite positive class in which the action is recovered
from the trial probabilities themselves.  For $A,B>0$, $C\in\R$, and $s>0$, consider
\begin{equation}
 E_s(x,y)=A\frac{(y-x)^2}{s}
 +Bs\left(\frac{x+y}{2}\right)^2+C(y^2-x^2),
 \qquad
 \begin{pmatrix}A&C\\ C&B\end{pmatrix}\succeq0.
 \label{eq:main-measured-quadratic-cost}
\end{equation}
Within this declared finite parametric class, three measured costs identify
$A,B,C$.  Define the directly measurable rank defect
\begin{equation}
 \mathfrak D=AB-C^2\ge0.
 \label{eq:main-rank-defect}
\end{equation}
For a finite proposal alphabet $F$, let $p_s(y\mid x)>0$ be a normalized
proposal row and accept a proposal with probability
\begin{equation}
 a_s(x,y)=e^{-E_s(x,y)/\varepsilon}.
 \label{eq:main-operational-acceptance}
\end{equation}
Write
\begin{equation}
 Z_s(x)=\sum_{y\in F}p_s(y\mid x)a_s(x,y),\qquad
 Q_s(y\mid x)=\frac{p_s(y\mid x)a_s(x,y)}{Z_s(x)}.
 \label{eq:main-accepted-row-normalizer}
\end{equation}
Both $Q_s$ and the total acceptance probability $Z_s$ are entries of the same
complete trial table.

\begin{theorem}[Measured acceptance reconstructs the finite action]
\label{thm:main-operational-action-matching}
The complete trial table reconstructs the unnormalized cost exactly:
\begin{equation}
 E_s(x,y)=-\varepsilon\log
 \frac{Z_s(x)Q_s(y\mid x)}{p_s(y\mid x)}.
 \label{eq:main-observed-cost}
\end{equation}
For a path with fixed endpoints, the observed writer
\begin{equation}
 \mathscr A_{\rm obs}
 =\varepsilon\sum_j\left[
 \log\frac{Q_{s_j}(q_{j+1}\mid q_j)}{p_{s_j}(q_{j+1}\mid q_j)}
 +\log Z_{s_j}(q_j)\right]
 +C(q_M^2-q_0^2)
 \label{eq:main-observed-homogeneous-action}
\end{equation}
is identically
\begin{equation}
 \mathscr A_{\rm obs}
 =-\sum_j\left[A\frac{(q_{j+1}-q_j)^2}{s_j}
 +Bs_j\bar q_j^2\right]
 \label{eq:main-offshell-action-match}
\end{equation}
off shell on the parametric domain.  Thus every interior-amplitude and lapse
first variation agrees exactly with the reconstructed finite action.

The one-step minimizing skeleton has proper-time amplitude drift
$c=-C/A$.  The bulk action in
\eqref{eq:main-offshell-action-match} has
$\Lambda_{\rm act}=4B/(3A)$, whereas the limiting constant-Hubble geometry of
the minimizing skeleton has
$\Lambda_{\rm geom}=4C^2/(3A^2)$.  Consequently
\begin{equation}
 G[g]+\Lambda_{\rm act}g
 =\frac{4\mathfrak D}{3A^2}g.
 \label{eq:main-measured-einstein-mismatch}
\end{equation}
The local operational update and the vacuum parameter of its measured bulk
action therefore agree exactly if and only if $\mathfrak D=0$.
\end{theorem}

The total acceptance probability in
\eqref{eq:main-observed-cost} is essential: the accepted row $Q_s$ alone is
invariant under addition of a source-dependent cost and therefore does not
identify the lapse action.  The complete trial record removes this ambiguity.
On the rank-one branch with $A=A_0$, $B=A_0\omega^2$, and
$C=-\sigma A_0\omega$, the cost is the square
\begin{equation}
 E_s(x,y)=\frac{A_0}{s}
 \left[y-x-\sigma\omega s\frac{x+y}{2}\right]^2,
 \label{eq:main-rank-one-cost}
\end{equation}
so its zero-cost recurrence is exactly
\eqref{eq:main-finite-stationary-recurrence}.

\begin{theorem}[Finite positive operational family with an Einstein limit]
\label{thm:main-explicit-operational-family}
Fix $T,\omega,q_0>0$ and a sign $\sigma$.  There exists a sequence of genuinely
finite accept/reject cylinders with $M$ time stages, finite positive proposal
alphabets, calibrated step $s=T/M$, rank-one costs
\eqref{eq:main-rank-one-cost}, and a finite retry cap such that, for constants
independent of $M$,
\begin{equation}
 \mathbb E\mathcal E_M\le CM^{-4},\qquad
 \mathbb E\max_j|q_j-q_0e^{\sigma\omega\tau_j}|\le CM^{-2},
 \label{eq:main-stochastic-profile}
\end{equation}
where
$\mathcal E_M=A_0\sum_jR_j^2/s$ and
$R_j=q_{j+1}-q_j-\sigma\omega s\bar q_j$.  The expected absolute first
variation of \eqref{eq:main-discrete-homogeneous-action} against every fixed
smooth scale/lapse test is $O(M^{-2})$, including the lapse constraint.  The
retry-failure probability is at most $M^{-6}$.  Under any common coupling of
cutoffs, failures occur only finitely often almost surely and the successful
piecewise-affine paths converge strongly in $H^1$ and uniformly to the
stationary solution.  Attaching the $A_3$ root race with
\begin{equation}
 \varrho_j=q_j^2,\qquad \kappa_j=q_j^{-4/3},\qquad
 k_{\alpha,h_x,j}=\frac{\kappa_j}{8h_x^2},\qquad
 m_{h_x,j}=\varrho_jh_x^3
 \label{eq:main-stochastic-root-rates}
\end{equation}
therefore gives positive renewal regulators converging to a vacuum metric with
$\Lambda=4\omega^2/3$.
\end{theorem}

Appendix~\ref{supp:homogeneous-dynamics} gives the finite grids, temperature
schedule, entropy--energy estimate, retry bound, and deterministic residual
estimates used in this theorem.  The construction proves nonemptiness of a
matched operational class; it does not assert that an arbitrary Grand Tensor
satisfies the quadratic duration law or the rank-one condition.

\paragraph{Flat and de Sitter regulators.}
The stationary profiles above give explicit geometric regulators.
\Cref{mt:vacuum} in Appendix~\ref{supp:vacuum} collects their generator,
noncollapse, curvature and summable exact-link estimates. In particular,
flat and de Sitter geometries arise after solving the lapse-varied equations,
rather than being inserted as prescribed microscopic rate histories.

\subsection{Conservative blocking preserves the action--law defect}
\label{subsec:homogeneous-blocking-main}

The preceding construction raises a distinct stability question: can ordinary
time coarse graining drive a neighboring quadratic law toward the matched
surface?  For the conservative interval elimination of the measured quadratic
cost, the answer is no.  Assume $0<s\sqrt{B/A}<2$ and set
\begin{equation}
 \mu=\sqrt{AB},\qquad
 \theta=2\operatorname{artanh}\!\left(\frac{s}{2}\sqrt{B/A}\right),\qquad
 \rho=\frac{C}{\mu},\qquad
 \Delta_{\rm rel}=1-\rho^2=\frac{AB-C^2}{AB}.
 \label{eq:main-relative-rank-defect}
\end{equation}
Then the same cost can be written
\begin{equation}
 \mathcal E_{\mu,\theta,C}(x,y)
 =\mu\coth\theta\,(x^2+y^2)-2\mu\operatorname{csch}\theta\,xy
   +C(y^2-x^2).
 \label{eq:main-hyperbolic-cost}
\end{equation}
For two successive intervals define conservative blocking by
\begin{equation}
 (E_1\star_0E_2)(x,y)=\inf_{z>0}\{E_1(x,z)+E_2(z,y)\}.
 \label{eq:main-conservative-blocking}
\end{equation}

\begin{theorem}[No transverse matching attraction under conservative blocking]
\label{thm:main-no-blocking-attractor}
For every $\theta_1,\theta_2>0$,
\begin{equation}
 \mathcal E_{\mu,\theta_1,C}\star_0
 \mathcal E_{\mu,\theta_2,C}
 =\mathcal E_{\mu,\theta_1+\theta_2,C}.
 \label{eq:main-hyperbolic-blocking}
\end{equation}
Hence $AB=\mu^2$, $C$, $\mathfrak D=AB-C^2$, and the scale-free mismatch
$\Delta_{\rm rel}$ are invariant under every finite sequence of homogeneous
blockings.  For two identical midpoint intervals of duration $s$, rewritten as
one interval of duration $2s$,
\begin{equation}
 A'=A+\frac{Bs^2}{4},\qquad
 B'=\frac{AB}{A+Bs^2/4},\qquad C'=C.
 \label{eq:main-homogeneous-block-map}
\end{equation}
Consequently $\Delta_{\rm rel}>0$ remains strictly positive: the matched surface
$\Delta_{\rm rel}=0$ has no open basin of attraction for this conservative
blocking operation, even after changes of time/amplitude units or common
positive rescaling of the action.  On finite positive amplitude alphabets with a symmetric untwisted trial
kernel, the endpoint twist $C$ is retained exactly under finite blocking, and a controlled
low-temperature/grid limit recovers the same nonzero relative defect.
\end{theorem}

This is a no-go result for a specified natural coarse-graining operation, not a
no-go theorem for every possible renewal renormalization.  In particular, a
driven or conditioned law may define a different physical flow.  The theorem
shows that generic Einstein matching cannot be inferred merely from conservative
elimination of intermediate amplitudes.

The same blocking operation also gives an exact consistency test:
local minimizers compose across intervals precisely on the matched branch.
\Cref{prop:main-optimizer-composition} in
Appendix~\ref{supp:homogeneous-dynamics} proves this characterization.
It tests matching; it does not generate an attracting flow.
\subsection{A nonlinear inhomogeneous two-polarization vacuum sector}
\label{subsec:gowdy-main}

The unpolarized Gowdy $T^3$ class provides a complementary nonlinear test:
a finite splitting preserves an exact momentum invariant, and the readout
controls the full curvature pathwise. This is a symmetry-reduced result
on a smooth slab \cite{Moncrief1981}. Write
$[t_0,t_1]\times\mathbb T^3$, $t_0>0$, write
\begin{equation}
 g=t^{-1/2}e^{\lambda/2}(-\dd t^2+\dd\theta^2)
 +t\left[e^P(\dd x+Q\dd y)^2+e^{-P}\dd y^2\right].
 \label{eq:main-gowdy-metric}
\end{equation}
Set
\begin{equation}
 a=P_t,\qquad b=P_\theta,\qquad
 c=e^P Q_t,\qquad d=e^P Q_\theta,
 \label{eq:main-gowdy-frame-vars}
\end{equation}
and
\begin{equation}
 e=\frac12(a^2+b^2+c^2+d^2),\qquad f=ab+cd.
 \label{eq:main-gowdy-ef}
\end{equation}
The first-order vacuum system is
\begin{align}
 a_t&=b_\theta-a/t+c^2-d^2, & b_t&=a_\theta,\nonumber\\
 c_t&=d_\theta-c/t-ac+bd, & d_t&=c_\theta+ad-bc,\nonumber\\
 P_t&=a, & Q_t&=e^{-P}c, & \lambda_t&=2te,
 \label{eq:main-gowdy-frame-system}
\end{align}
with coordinate and momentum constraints
\begin{equation}
 P_\theta=b,\qquad e^P Q_\theta=d,\qquad
 \lambda_\theta=2tf.
 \label{eq:main-gowdy-constraints}
\end{equation}
The identities
\begin{equation}
 e_t-f_\theta=-\frac{a^2+c^2}{t},\qquad
 f_t-e_\theta=-\frac ft
 \label{eq:main-gowdy-stress-identities}
\end{equation}
propagate the last constraint and preserve the zero mean of $tf$ required for
periodic $\lambda$.

For the variational test the areal gauge is imposed only after variation.  With
independent fields $(R,\lambda,P,Q,N,\beta)$, $R,N>0$, and
$v_u=\dot u-\beta u'$, use the boundary-marked ADM density
\begin{align}
 \mathscr L_{\rm G}={}&
 -\frac{v_R(v_\lambda-4\beta')}{2N}
 +\frac R{2N}(v_P^2+e^{2P}v_Q^2)\nonumber\\
 &+N\left[-2R''+\frac{R'\lambda'}2
 -\frac R2\{(P')^2+e^{2P}(Q')^2\}\right].
 \label{eq:main-gowdy-adm-density}
\end{align}
Its lapse and shift equations are independent constraints.  The solution of
\eqref{eq:main-gowdy-frame-system}--\eqref{eq:main-gowdy-constraints} lifts to
those equations with $R=t$, $N=1$, $\beta=0$; hence the finite action test below
is not obtained by first gauge-fixing away the lapse and shift variations.

\begin{theorem}[Nonlinear two-polarization renewal regulator]
\label{thm:main-gowdy-regulator}
Let $X_*(t)$ be a smooth solution of
\eqref{eq:main-gowdy-frame-system}--\eqref{eq:main-gowdy-constraints} on a fixed
slab $[t_0,t_1]\times\mathbb T^1$, with the spatial derivatives required below
uniformly bounded and with the reconstructed metric uniformly nondegenerate.
Let $h=\ell$ be the synchronized temporal and spatial mesh.  Then there is a
nearest-neighbour symmetric source/transport/source splitting with the following
properties.
\begin{enumerate}[label=\textup{(G\arabic*)},leftmargin=3em]
\item With
      $U=(u_1,u_2,u_3,u_4)=\sqrt t\,(a,b,c,d)$,
      $\mathcal J(U)=u_1u_2+u_3u_4$,
      $D_+=(S_+-I)/\ell$, and $\mathsf A_+=(I+S_+)/2$, the quantity
      \begin{equation}
       \mathcal K_\ell
       =D_+\lambda-2\mathsf A_+\mathcal J(U)
       \label{eq:main-gowdy-staggered-momentum}
      \end{equation}
      is preserved exactly by every deterministic full step.  The map is
      second-order consistent on the stated smooth chart.
\item If
      $X_{j+1}=\Phi_h(t_j,X_j)+r_j$ and
      \begin{equation}
       \mathcal E_h=\sum_j\frac{\|r_j\|_{2,\ell}^2}{2h},\qquad
       \mathcal F_h=\ell^{-(2r+1)}\mathcal E_h,
       \label{eq:main-gowdy-residual-energy}
      \end{equation}
      then, under the stated smooth-chart bootstrap or an enforced offset
      envelope,
      \begin{equation}
       \max_j\|X_j-\mathsf S_\ell X_*(t_j)\|_{r,\infty,\ell}
       \le C\bigl(\epsilon_0+h^2+\sqrt{\mathcal F_h}\bigr).
       \label{eq:main-gowdy-state-error}
      \end{equation}
      For $r\ge1$, the accumulated staggered-constraint error is bounded by its
      initial value plus $C\sqrt{\mathcal F_h}$.
\item On the finite offset alphabet with
      $\|\epsilon_j\|_\infty\le c_0h^4$ and initial $C_h^1$ error $O(h^2)$,
      the shared tensor-product cubic Hermite readout is pathwise $O(h^2)$ in
      $W^{1,\infty}$ and $O(h)$ in its second derivatives.  Consequently, on
      compact subslabs,
      \begin{equation}
       \|g_h-g_*\|_{W^{1,\infty}}
       +\|D^2g_h-D^2g_*\|_{L^\infty}
       +\|\operatorname{Riem}(g_h)-\operatorname{Riem}(g_*)\|_{L^\infty}
       \le Ch,
       \label{eq:main-gowdy-full-curvature}
      \end{equation}
      where the first term is in fact $O(h^2)$.  Thus
      $\sup_h\|\operatorname{Riem}(g_h)\|_{L^2(K)}<\infty$ and
      $\|\operatorname{Ric}(g_h)\|_{L^\infty(K)}\le Ch$ pathwise.  The
      Levi--Civita curvature of the limiting metric is therefore identified
      directly.
\item An independent midpoint/centered discretization of
      \eqref{eq:main-gowdy-adm-density}, with $N$ and $\beta$ varied
      independently before restriction to areal gauge, has first-variation
      residual tending to zero against every fixed smooth local test.  A finite
      positive offset/acceptance realization can be chosen with the support cap
      in item~\textup{(G3)} and obeys the exact Gibbs entropy--energy bound; its
      successful histories therefore inherit the pathwise curvature estimate.
\end{enumerate}
\end{theorem}

The construction and proof are given in Appendix~\ref{supp:gowdy}.  It is a
nonlinear and spatially inhomogeneous population result, but its scope remains
symmetry reduced and local to a smooth positive slab.  Unlike the general
certificate of \Cref{thm:main-curvature-propagation}, the curvature bound in
\eqref{eq:main-gowdy-full-curvature} is obtained directly from pathwise
second-jet control and requires no separately assumed all-time curvature budget
on this class.  It does not imply an unrestricted four-dimensional nonlinear
Cauchy theorem.


The finite proposal-access estimate, including reset and retry assumptions,
is \Cref{prop:main-safe-access} in Appendix~\ref{supp:gowdy}. It determines
when the measured process reaches the admissible histories; it is separate
from the deterministic curvature estimate.


\section{Interpretation and discussion}
\label{sec:discussion}
% =============================================================================

The construction developed in this paper separates several structures that
are usually introduced simultaneously.  Predictive persistence is obtained
first by quotienting operational histories according to their complete
conditional futures.  Relational geometry is then reconstructed on this
future-minimal carrier from physically occurring incidence, calibration, and
renewal statistics.  Lorentzian signature requires one additional protected
signed datum, and gravitational dynamics enters only after a common
variational structure has been reconstructed and its continuum behaviour
controlled.

The resulting picture is therefore not a direct identification of stochastic
dynamics with spacetime geometry.  It is a sequence of reconstruction and
compatibility problems.  Each stage has its own finite data, its own
obstructions, and its own continuum requirements.  This separation is
important both conceptually and mathematically: the predictive quotient does
not already contain a Lorentzian metric, Lorentzian kinematics does not
already imply Einstein dynamics, and the existence of a reconstructed action
does not imply that an arbitrary microscopic evolution is stationary for
that action.


% -----------------------------------------------------------------------------
\subsection{Predictive persistence and Lorentzian geometry}
% -----------------------------------------------------------------------------

The first structural point is that persistence need not be supplied as a
microscopic state space.  Histories that have identical statistics under
every admissible future continuation are operationally indistinguishable and
are therefore identified.  Because this equivalence is preserved by
continuation, the quotient carries a canonical successor structure.  The
minimum predictive carrier is thus reconstructed from the process itself
rather than introduced as a hidden state, lattice site, or geometric point.

The Renewal Grand Tensor records the corresponding finite positive
history structure: the minimum predictive carrier, represented history
algebra, and positive word-Gram hierarchy.  Its role is not that of a metric
tensor on a pre-existing space.  It records which distinctions survive
predictive reduction, how represented historical operations compose, and
which positive overlaps remain between them.  Geometry becomes meaningful
only after this predictive structure is combined with the additional
same-history information required by the geometric branch.

\paragraph{Dimension and relational structure.}

The selection of three spatial dimensions is likewise conditional and
structural rather than universal.  Two independent finite classifications
select the tetrahedral relation cell $K_4$.  The endpoint--cycle balance
criterion does so among connected simple regular relations without invoking
an alternating form, while the stronger pair-homogeneous simplex criterion
selects the same cell from minimal faithful alternating completion.  Both
lead to the three-dimensional $A_3$ standard module.

The endpoint--cycle balance argument is the primary result and does not
depend on the parity argument. The orientation-twisted source equivalence
provides its representation-theoretic interpretation; the alternating
criterion is secondary and retains its explicit minimization hypothesis.  At the same time,
neither theorem claims that every possible microscopic relation system must
belong to these categories.  Higher-arity interactions, enriched incidence
structures, or different primitive relational questions can lead outside the
classified branch.  The result should therefore be read as a pair of finite
uniqueness statements within explicitly stated structural classes.

\paragraph{Metric reconstruction and calibration.}

Once the $A_3$ relational carrier has been selected, the metric is
reconstructed from observable renewal statistics.  The symmetric second
moments of the directed root race determine the spatial kinetic form, the
antisymmetric first moments determine the shift, and the physical speed
density together with the waiting-time scale fixes the remaining ADM
normalization.  The spatial metric, lapse, and shift are consequently
collective quantities derived from the finite renewal process rather than
labels attached to microscopic events.

This reconstruction also identifies what cannot be obtained from normalized
probabilities alone.  Winner probabilities determine only relative rates;
an absolute waiting-time calibration is required to recover their common
scale.  More generally, the calibration fibre described above separates
physically meaningful scale information from conventional changes of units.
Absolute calibration and same-cut readability are therefore genuine physical
inputs of the geometric branch, even though the metric itself is not.

\paragraph{The Lorentzian sign.}

The Lorentzian signature occupies a particularly clear position in this
hierarchy.  Both the Grand-Tensor Gram structure and the root-race spatial
form are positive.  They cannot generate a negative direction by themselves.
The protected binary orientation is an additional operational distinction,
and the separating constructions above show that it is independent of the
remaining positive data within the declared language.

Once the positive spatial module has dimension three, adjoining one
independent nondegenerate signed direction yields Lorentzian inertia
$(1,3)$, up to overall convention.  This conclusion precedes the Clifford
representation.  The gamma matrices therefore realize an already determined
Lorentzian quadratic form; they are not responsible for selecting its
signature.  In particular, the Lorentzian direction is neither obtained by
analytic continuation of a positive diffusion form nor inserted through a
choice of Dirac matrices.


% -----------------------------------------------------------------------------
\subsection{From reconstructed action to Einstein dynamics}
% -----------------------------------------------------------------------------

Lorentzian reconstruction alone is a kinematical statement.  The dynamical
problem begins with the complete accepted physical cylinder.  The determining
field is obtained from future-equivalence on that cylinder, while a marked
renewal pressure reconstructs the affine action slope from protected reward
and duration data.  Finite curl and period tests then determine whether the
resulting increments integrate to a single scalar action.

This ordering makes the variational statement noncircular.  The action is
reconstructed before the accepted dynamics is compared with it.  The
Fisher--Bregman and related finite probability identities subsequently test
whether the observed evolution is stationary for that independently
reconstructed action; they do not define the action by fitting a cost to the
observed transition kernel.

After resolved memory has been removed, Lorentz covariance restricts the
first-order gravitational score to the Palatini--Holst--volume family.
Connection stationarity gives the torsion-free branch, while the positive
spatial cut condition supplies regulator-uniform noncollapse and compact
screens.  Curvature control is obtained separately from initial energy and
propagation estimates and is identified with the curvature of the limiting
connection through the Cartan comparison.  These ingredients allow the
weak--strong variational passage to
\begin{equation}
\boxed{
G_{\mu\nu}+\Lambda g_{\mu\nu}=0
}
\end{equation}
on the certified continuum branch.

The point of this construction is therefore not merely that particular
renewal histories can be chosen whose reconstructed metrics happen to solve
Einstein's equation.  Geometry, action provenance, stationarity, noncollapse,
and curvature compactness are tested independently and are joined only at the
continuum handoff.

\paragraph{Constructive and certification results.}

There are two complementary theorem levels.  The general
renewal-to-Einstein theorem is a recognition result: a complete gravitational
renewal cylinder satisfying the stated finite and cofinal certificates has an
Einstein subsequential limit.  Its hypotheses deliberately separate action
stationarity, first-order gravitational provenance, noncollapse, and
curvature propagation.

The open $3+1$ population theorem addresses the converse concern that these
certificates might describe an empty or highly fine-tuned class.  For an
explicit local writer it proves a cutoff-independent nonlinear Cauchy interval
with arbitrary three-dimensional dependence in a fixed Sobolev
neighbourhood.  The surrounding open family of local differential laws has
the same continuum Einstein limit, and the corresponding finite
history-dependent transition rows admit relatively open residual-controlled
neighbourhoods.  The same-cylinder compiler further shows that, under explicit
occurrence-energy and variance scalings, a chronological correction of the
independently reconstructed action comparator preserves both the writer
residual scale and the field-action stationarity scale, while the literal-link
compactness theorem supplies the connection handoff on the successful branch.

The homogeneous and Gowdy constructions play a different role.  The
homogeneous sector makes the lapse variation and action--dynamics matching
fully explicit and shows that the flat and de~Sitter profiles are stationary
outputs rather than prescribed rate histories.  The unpolarized Gowdy sector
provides a nonlinear inhomogeneous example with exact finite momentum
propagation and direct full-curvature convergence.  Together with the open
$3+1$ population, these examples show that the positive branch is not
restricted to a single symmetric trajectory.

\paragraph{Einstein matching is not an automatic attractor.}

The robustness results should not be confused with a universality theorem.
For the conservative homogeneous blocking operation studied above, the
dimensionless action--law mismatch is invariant.  Coarse graining of that
specific type therefore does not drive a generic microscopic law toward the
Einstein-matched surface.  The optimizer-composition identity instead gives
an exact consistency characterization of the matched branch.

This distinction is useful more generally.  Stability of an already
controlled Einstein branch and dynamical selection of that branch from a
larger microscopic theory are different questions.  The former is
established for the finite law classes treated here; the latter remains a
problem of microscopic selection and renormalization.

\paragraph{The exact-action route.}

A further distinction concerns the exact Euler flow of the reconstructed
finite action.  The constructive local writer used in the open population
theorem is action compatible, but it is not assumed to coincide with the raw
finite Euler flow at every cutoff.  The exact-action analysis shows that this
stronger identification is nontrivial.

The finite response problem has substantial regular structure: transverse
source blocks can be eliminated exactly, the signed response separates into
positive source reproduction and a scalar--trace correction, and nontrivial
source-neutral prepared families exist.  On the certified fixed-cutoff event
the weak first-Poisson block has exact rank twenty-four, the five-dimensional
weak kernel splits into two generalized-zero and three constant-shift
directions, and the rapid generator has the semisimple $(8,8,13)$
stable--unstable--centre decomposition.  Under the full graded-chart conditions, the rank-free nonlinear two-point
construction supplies an exactly constrained original-action
history through that event with $O(a)$ curvature on an $O(a^2)$ physical
interval while changing only higher preparation coefficients.  This does not
give forward Cauchy stability: rapid and slow normal expansion on the separate
normal-factor chart impose distinct exponential preparation scales for
residence through a fixed physical time.  At the initial cut exact initialization automatically satisfies the three
constant-shift/Fredholm rows, while the other two rows are measured by an
explicit acceleration defect rather than inferred from amplitude analyticity.
For bounded-rate positive-time limits, elimination of the generalized-zero
rows produces a quadratic harmonic-slope condition.  The literal harmonic
Ward structure forces its Jacobian to have rank at most two, yielding a
cokernel obstruction and, when that obstruction vanishes, a two-equation
root-curve problem.  The remaining problem is nonlinear propagation of a
compatible slow branch, physical selection without future endpoint data, and
regulator-uniform continuation under spatial refinement.


% -----------------------------------------------------------------------------
\subsection{Scope and open problems}
% -----------------------------------------------------------------------------

The results establish a finite reconstruction and continuum mechanism for a
specified class of operational renewal processes.  They do not imply that
every renewal law produces Lorentzian geometry or that every such geometry
enters the Einstein branch.  Actual occurrence of the required operational
rows, protected orientation, absolute calibration, noncollapse, action
stationarity, and curvature propagation remain genuine conditions on the
microscopic realization.

Likewise, the paper does not establish a universal derivation of three
spatial dimensions.  The $K_4/A_3$ result is conditional on the finite
relational categories classified above.  Its strength is that two
mathematically distinct selectors lead to the same spatial module, not that
all conceivable pregeometric theories are forced into that category.

\paragraph{Microscopic selection and universality.}

The most immediate open problem is to understand why a broad microscopic
class should enter the controlled Lorentzian--Einstein branch.  The present
results establish nonempty and relatively open families once the appropriate
physical topology is chosen, but they do not derive an attracting
renormalization flow from arbitrary renewal laws.  A useful next step would
be to identify natural conditioning, coarse-graining, or learning dynamics
under which the action, noncollapse, and propagation margins become stable
large-scale quantities.

The exact-action problem is a sharper version of the same question.  One
would like to determine whether the reconstructed finite action itself admits
a regular cofinal branch whose exact Euler flow preserves the required
constraint and slow compatibility structure on a cutoff-independent interval
and enters the same Einstein continuum class as the constructive writers.
The fixed-cutoff boundary construction and the two exponential residence
bounds do not settle this selection and propagation problem.

\paragraph{Matter and internal structure.}

The present analysis concerns the vacuum gravitational sector.  An extension
to matter requires more than appending additional fields to the limiting
Einstein equation.  Internal symmetry, matter incidence, fermionic
multiplicity, and their common action should be represented on the same
finite carrier as the reconstructed spacetime structure, and their nonlinear
stress must remain controlled in the continuum passage.  This raises both
finite operator-algebraic questions and independent compactness questions for
the coupled gravity--matter system.

\paragraph{Residual memory.}

A final question concerns the information lost in the local continuum
description.  The microscopic dynamics considered here may be genuinely
non-Markovian, whereas the controlled continuum branches are effectively
local after predictive compression.  It is therefore natural to ask which
subleading structures retain memory of the microscopic renewal process.
Controlled departures from the Markovian scaling regime could provide a
systematic language for nonlocal or history-dependent corrections to the
effective geometric dynamics.

Within the scope established here, the main conclusion is more limited and
more precise.  A persistent state, a spatial metric, Lorentzian signature,
and the vacuum Einstein equation need not all be supplied at the microscopic
level.  Predictive persistence can be reconstructed from future
distinguishability; a calibrated relational branch can reconstruct the ADM
geometry; one independent protected sign fixes the Lorentzian inertia; and,
when the independently reconstructed action, noncollapse, and curvature
conditions are satisfied, the corresponding continuum branch obeys vacuum
Einstein dynamics.  The remaining problem is to understand how broadly
microscopic physics selects and stabilizes that branch.


\section*{Data and formalization availability}
Source code for selected finite algebraic components is available at
\url{https://github.com/AI-MathPhys/renewal-geometry}.  The fixed-cutoff
exact-action rank and spectral statements above are finite computer-assisted
certificates: the proof uses outward interval bounds, exact rational matrix
checks, and finite-field/Pfaffian identities as described in
Appendices~\ref{supp:exact-action-audit} and~\ref{supp:exact-action-slow}.
The analytic estimates and conditional finite-chart theorems in this article
are given as mathematical proofs. In particular the exact-action boundary
theorem retains all of \Cref{ass:supp-graded-exact-chart}; its analytic
divisibility, source and readout inequalities are not consequences of the
rank or spectral certificate alone. Instantiation of a numerical finite-chart certificate requires its full
coefficient data, same-root conventions, and validated inequalities; ranks,
homogeneities and degree bounds alone are not executable coefficient data.


\clearpage
\appendix
\section*{Guide to the appendices}
\phantomsection
\addcontentsline{toc}{section}{Appendices: guide to proofs}
\addtocontents{toc}{\protect\setcounter{tocdepth}{0}}
The appendices retain the detailed constructions, quantitative constants
and alternative sufficient criteria. Each has one mathematical purpose;
all exact-action statements are separate from the general Einstein-limit
proof.
\begin{center}
\small
\begin{tabularx}{\linewidth}{@{}p{.10\linewidth}X@{}}
\toprule
Appendix & Proofs and dependencies\\
\midrule
\ref{supp:grand-reconstruction} & Predictive quotient, process reconstruction, history algebra and calibration.\\
\ref{supp:dimension-adm} & Relational selectors, root statistics, positive realizability and signed geometry.\\
\ref{supp:dynamics} & Determining field, common action, stationarity, local classification and torsion.\\
\ref{supp:continuum} & Noncollapse, propagation, connection alternatives, interfaces and the Einstein limit.\\
\ref{supp:open-3plus1} & Uniform harmonic evolution and genuinely finite renewal realization.\\
\ref{supp:open-law-basin} & Law-space robustness, finite constraint preparation and same-cylinder compatibility.\\
\ref{supp:exact-action-audit} & Exact-action response, preparations, fixed-cutoff boundary histories and residence.\\
\ref{supp:exact-action-slow} & Necessary slow compatibility, acceleration tests and the quadratic slope obstruction.\\
\ref{supp:sectors} & Homogeneous matching, conservative blocking, Gowdy and flat/de Sitter regulators.\\
\bottomrule
\end{tabularx}
\end{center}


\section{Operational reconstruction and positive history algebra}
\label{supp:grand-reconstruction}

This appendix proves \Cref{mt:grand-tensor} and
\Cref{prop:readable-relational-completion}. It uses only the finite
operational entrance from Section~\ref{sec:grand-tensor-main}; the final
subsection isolates the additional orientation and calibration information.

\subsection{Finite typed operational renewal datum}

\begin{definition}[Finite typed event grammar]
\label{def:supp-event-grammar}
At cutoff $X$, a finite typed event grammar consists of:
\begin{enumerate}[label=\textup{(E\arabic*)},leftmargin=3em]
\item a finite set of cut types $\mathsf O_X$;
\item finite sets $\Sigma_X(x,y)$ of primitive event letters $a:x\to y$;
\item a finite-horizon prefix-closed family $\mathsf H_{X,x}$ of reachable
      typed histories ending at each cut type $x$, with the remaining horizon
      included in the type, typed concatenation and empty words; alternatively,
      an unbounded family with a separately certified finite-index future
      congruence;
\item finite preparation, intervention-outcome, and terminal-Read alphabets,
      together with finite exposed operator ports;
\item a distinguished protected binary orientation record whose two values are
      separated by the declared physical Read policy.
\end{enumerate}
No state label, persistent relation object, metric, scheduler, connection, or
spacetime semantic is assigned to a history.
\end{definition}

For $h\in\mathsf H_{X,x}$ let $\Sigma_{X,h}$ be the event letters admissible
after $h$.

\begin{definition}[Finite conditional operational renewal law]
\label{def:supp-operational-datum}
A finite conditional operational renewal law on
\Cref{def:supp-event-grammar} assigns
\begin{equation}
 p_X(a\mid h)\in[0,1],
 \qquad
 \sum_{a\in\Sigma_{X,h}}p_X(a\mid h)=1,
 \label{eq:supp-one-step-law}
\end{equation}
for every reachable history.  For an admissible future word
$w=a_1\cdots a_k$ define
\begin{equation}
 p_X(w\mid h)
 =\prod_{j=1}^{k}
 p_X(a_j\mid ha_1\cdots a_{j-1}),
 \qquad p_X(\epsilon_x\mid h)=1.
 \label{eq:supp-chain-law}
\end{equation}
The physical intervention letters include, on every exposed Hermitian
Choi-coordinate space, a finite physical $\CP$ frame whose real span is the
full Hermitian space.  The induced multitime frame probabilities are affine
under randomization, additive under coarse graining, normalized on complete
instruments, causal under deterministic discard, compatible on common
prefixes, and satisfy the finite positivity conditions of
\Cref{thm:supp-comb-tomography}.  Actual bounded terminal Reads, marks, and
readout queries may be included as terminal event outcomes; equivalently they
may be appended to the future test bank.
\end{definition}

The distinguished renewal semantics used in the geometric branch is carried by
the admitted completion/return event family.  No ordinary reset state is
needed for the predictive construction below; if a stronger reset law is
imposed, it is an additional property of the same event process rather than a
state-space axiom.

\subsection{Future equivalence and the derived persistent carrier}

\begin{definition}[Complete predictive future signature]
\label{def:supp-future-signature}
For $h\in\mathsf H_{X,x}$ define
\begin{equation}
 \Sigma_X(h)=\bigl(p_X(w\mid h)\bigr)_{w\in\mathsf F_{X,x}},
 \label{eq:supp-future-signature}
\end{equation}
where $\mathsf F_{X,x}$ is the complete horizon-compatible future-test bank
starting at $x$; on an unbounded-horizon protocol it comprises all finite
admissible continuations.  Histories $h,h'\in\mathsf H_{X,x}$ are future equivalent when
\begin{equation}
 h\sim_Xh'
 \quad\Longleftrightarrow\quad
 p_X(w\mid h)=p_X(w\mid h')
 \text{ for every }w\in\mathsf F_{X,x}.
 \label{eq:supp-future-equivalence}
\end{equation}
Histories of different cut type are never identified.
\end{definition}

\begin{theorem}[History-level right congruence]
\label{thm:supp-right-congruence}
Future equivalence is a typed right congruence.  If $h\sim_Xh'$ and
$a:x\to y$ is admissible, then
\begin{equation}
 p_X(a\mid h)=p_X(a\mid h'),
 \label{eq:supp-one-letter-equality}
\end{equation}
and on the positive-probability branch
\begin{equation}
 ha\sim_Xh'a.
 \label{eq:supp-right-congruence}
\end{equation}
\end{theorem}

\begin{proof}
Use the one-letter future $a$ in
\eqref{eq:supp-future-equivalence}.  If the common probability is positive,
then for any continuation $v$ the chain rule gives
\[
 p_X(v\mid ha)
 =\frac{p_X(av\mid h)}{p_X(a\mid h)}
 =\frac{p_X(av\mid h')}{p_X(a\mid h')}
 =p_X(v\mid h'a).
\]
Hence the complete future signatures agree after appending $a$.
\end{proof}

\begin{definition}[Minimum predictive state]
\label{def:supp-predictive-state}
For every cut type $x$ put
\begin{equation}
 Z_{X,x}^{\min}:=\mathsf H_{X,x}/\!\sim_X,
 \qquad
 q_{X,x}:\mathsf H_{X,x}\twoheadrightarrow Z_{X,x}^{\min}.
 \label{eq:supp-predictive-state}
\end{equation}
For an admissible event $a:x\to y$ define
\begin{equation}
 T_{X,a}([h])=[ha],
 \qquad
 \kappa_{X,a}([h])=p_X(a\mid h).
 \label{eq:supp-predictive-transition}
\end{equation}
\end{definition}

\begin{corollary}[Derived predictive state machine]
\label{cor:supp-derived-state-machine}
The maps in \eqref{eq:supp-predictive-transition} are well defined.  If
$w=a_1\cdots a_k$, $z_0=q_X(h)$, and
$z_j=T_{X,a_j}(z_{j-1})$, then
\begin{equation}
 p_X(w\mid h)=\prod_{j=1}^{k}\kappa_{X,a_j}(z_{j-1}).
 \label{eq:supp-state-factorization}
\end{equation}
Thus the predictive state machine is reconstructed from the conditional event
law.
\end{corollary}

\begin{proof}
Well-definedness is \Cref{thm:supp-right-congruence}; the product formula is
\eqref{eq:supp-chain-law} written on future-equivalence classes.
\end{proof}

\begin{definition}[Deterministic predictive realization]
\label{def:supp-predictive-realization}
A reachable deterministic predictive realization consists of finite sets
$S_{X,x}$, history maps $s_x:\mathsf H_{X,x}\to S_{X,x}$, partial updates
$u_a:S_{X,x}\to S_{X,y}$, and branch functions
$r_a:S_{X,x}\to[0,1]$ such that on every reachable branch
\begin{equation}
 s_y(ha)=u_a(s_x(h)),
 \qquad
 p_X(a\mid h)=r_a(s_x(h)).
 \label{eq:supp-realization}
\end{equation}
\end{definition}

\begin{theorem}[Canonical minimum predictive realization]
\label{thm:supp-minimal-record}
The quotient $Z_X^{\min}$ is the unique coarsest reachable deterministic
predictive realization.  For every reachable realization $S_X$ there is a
unique surjection
\begin{equation}
 \pi_{X,x}:S_{X,x}\twoheadrightarrow Z_{X,x}^{\min}
 \label{eq:supp-universal-surjection}
\end{equation}
such that
\begin{equation}
 q_x=\pi_xs_x,
 \qquad
 \pi_yu_a=T_{X,a}\pi_x,
 \qquad
 r_a=\kappa_{X,a}\pi_x.
 \label{eq:supp-universal-intertwining}
\end{equation}
Consequently a generated ledger or hidden-state presentation, when present,
is only an intermediate predictive realization and has a unique quotient onto
$Z_X^{\min}$.
\end{theorem}

\begin{proof}
If $s_x(h)=s_x(h')$, induction on future-word length using
\eqref{eq:supp-realization} gives $p_X(w\mid h)=p_X(w\mid h')$ for every
future word $w$.  Hence $h\sim_Xh'$, so
$\pi_x(s_x(h)):=[h]$ is well defined.  Reachability makes $\pi_x$ unique and
surjective.  The intertwining identities follow from the definitions, and the
quotient itself is a realization by
\Cref{cor:supp-derived-state-machine}.
\end{proof}

\begin{theorem}[Persistence derived from composition]
\label{thm:supp-persistence-derived}
Let $A_1,A_2,\ldots$ be the realized event sequence, let
$H_n=A_1\cdots A_n$, and put $Z_n=q_X(H_n)$.  Then on every realized branch
\begin{equation}
 Z_{n+1}=T_{X,A_{n+1}}(Z_n).
 \label{eq:supp-persistence}
\end{equation}
Consequently
\begin{equation}
 \mathfrak p_{X,n}^{\min}=(Z_n,A_{n+1},Z_{n+1})
 \label{eq:supp-derived-packet}
\end{equation}
is a persistent relational event packet derived from the event process.  Every
other deterministic persistent predictive packet maps uniquely onto
\eqref{eq:supp-derived-packet}.
\end{theorem}

\begin{proof}
By definition,
$Z_{n+1}=q_X(H_nA_{n+1})=T_{X,A_{n+1}}q_X(H_n)$.  For any other predictive
realization, apply the universal map of
\Cref{thm:supp-minimal-record} at consecutive cuts.
\end{proof}

\begin{corollary}[Future-operational transfer]
\label{cor:supp-future-operational-transfer}
Let $\mathcal D$ be any downstream finite readout whose value is unchanged
whenever two candidate states have the same complete declared future
signature.  Then there is a unique $\overline{\mathcal D}$ on
$Z_X^{\min}$ such that
\begin{equation}
 \mathcal D=\overline{\mathcal D}\circ\pi_X.
 \label{eq:supp-compiler-factorization}
\end{equation}
Thus deleting every future-null state coordinate preserves such downstream
outputs exactly.
\end{corollary}

\begin{proof}
The readout is constant on the fibres of the universal surjection
\eqref{eq:supp-universal-surjection}; define $\overline{\mathcal D}$ on a
fibre by the common value.  Surjectivity gives uniqueness.
\end{proof}

\subsection{Conservative readable relational completion}
\label{supp:readable-relational-completion}

The predictive quotient above determines persistence but does not by itself
assert that the quotient class, scheduler path, orientation, and calibration
are jointly readable at one cut.  We now isolate the finite conservative
extension used by the geometric reconstruction.

Let $\Phi_{a,*}$ be the support-minimal Schr\"odinger-picture branch map for an
underlying event letter $a$.  Let $z\in Z_X^{\min}$ be its predictive source
class.  For a finite family of added same-history marks $m=(\gamma,\xi,z')$,
where $\gamma$ is a scheduler-path record, $\xi$ contains orientation and
calibration marks, and $z'$ is the next readable predictive class, let
\begin{equation}
 K_{X,a}(m\mid z)\ge0,
 \qquad
 \sum_m K_{X,a}(m\mid z)=1.
 \label{eq:supp-relational-kernel}
\end{equation}
On block states $\rho=(\rho_z)_z$, define
\begin{equation}
 [\widetilde\Phi_{a,m,*}(\rho)]_{z'}
 =\sum_z K_{X,a}(m\mid z)\,\Phi_{a,*}(\rho_z).
 \label{eq:supp-relational-refinement}
\end{equation}
The symbol in \eqref{eq:supp-relational-refinement} denotes a same-history
conditional extension, not a tensor-product decomposition of the state space.

The refined branch maps are completely positive because their coefficients
are nonnegative. Their normalized forgetful identity is
\begin{equation}
       \mathsf F_X\!\left(\sum_m\widetilde\Phi_{a,m,*}(\rho)\right)
       =\Phi_{a,*}(\mathsf F_X\rho);
       \label{eq:supp-relational-forgetful}
      \end{equation}
This also supplies the explicit realization used in
\Cref{prop:readable-relational-completion}.


\begin{proof}[Proof of \Cref{prop:readable-relational-completion}]
Complete positivity follows because every target block in
\eqref{eq:supp-relational-refinement} is a nonnegative scalar combination of
completely positive branch maps.  Summing
\eqref{eq:supp-relational-kernel} over the added marks gives one, and therefore
summing the refined branches proves
\eqref{eq:supp-relational-forgetful}.  Induction on word length gives exact
preservation of every underlying finite cylinder.

For future minimality, define two raw readable packets to be equivalent when
all declared future signatures agree.  Prefixing the same admissible event to
two equal future signatures leaves them equal, so the relation is a typed
right congruence.  Finite partition refinement terminates because every strict
round increases the number of classes.  The usual fibre argument then gives
its universal property: any other sufficient readable presentation identifies
only packets with equal complete future signatures and therefore surjects
uniquely onto the stable quotient.
\end{proof}


\subsection{Finite distinguishing depth and the presentation category}

\begin{proposition}[Finite predictive reconstruction with an explicit horizon convention]
\label{prop:supp-finite-prediction}
For a finite-horizon protocol, include the remaining horizon in the cut type
and compare all admissible continuations. The future quotient is finite and
is a right congruence on every nonterminal positive-probability branch.
For an unbounded-horizon protocol admitting a reachable finite presentation
with $m$ records, the coarsest predictive partition is obtained after at most
$m-1$ strict partition refinements. A stable finite partition is sufficient
for equality of all finite futures. Finite alphabets without either a horizon
bound or a finite-index hypothesis do not imply a finite quotient.
\end{proposition}
\begin{proof}
At a fixed horizon there are finitely many histories. The conditional chain
rule proves right congruence; inclusion of the remaining horizon in the type
ensures that the continuation $av$ and its suffix $v$ are both legitimate in
the corresponding comparison. For a finite presentation, initialize the
partition by type, immediate read values, branch domains and probabilities.
At each round, additionally require equal successor blocks for each letter
of positive probability. Induction shows that the partition after $r$ rounds
records equality of the immediate data and of all future tests with at most
$r$ further successor operations. Each strict refinement increases the
number of blocks, which is at most $m$. At a fixed point, induction on word
length gives equality of every finite future. Conversely, future-equivalent
records cannot be separated by any round. This proves termination and the
claimed quotient. Exact comparisons of the table values are part of this
algebraic statement.

For the last assertion, consider a binary renewal process whose probability
of the next renewal, after $n$ consecutive nonrenewals, is $1/(n+2)$.
Every such history is reachable and its immediate next-event law differs
from that at every other age. The alphabet is finite but the future quotient
has infinitely many states. Thus finiteness of the quotient is not a
consequence of alphabet size alone.
\end{proof}

\begin{proposition}[Terminality and naturality of predictive reduction]
\label{prop:supp-predictive-category}
Fix a typed operational law. Let a morphism $f:S\to S'$ of reachable
predictive presentations satisfy $s'=f\circ s$, preserve reads and branch
probabilities, and intertwine every positive-probability successor. Then
$Z_X^{\min}$ is terminal in this category. Isomorphisms of typed laws induce
isomorphisms of their predictive quotients, compatibly with composition.
\end{proposition}
\begin{proof}
If two histories have the same record in $S$, their complete futures agree,
so $f_S(s(h))=[h]$ is well defined. It preserves the immediate data and
successors by the chain rule. Reachability forces every morphism from $S$
to the quotient to have this value, proving terminality. An isomorphism of
laws carries equal future signatures to equal future signatures and therefore
descends to the quotient; the inverse descends as well. The descended maps
respect composition because they agree on every represented history.
\end{proof}


\subsection{The state-conditioned Hankel quotient}

Fix a minimum predictive state sector $(x,z)$.  Let
$\mathsf P_{x,z}$ and $\mathsf F_{x,z}$ denote the finite pasts and futures
meeting at that sector, and set
\begin{equation}
 \mathbb p_{x,z}(f,p)=\Pr(f\circ p).
 \label{eq:supp-Hankel-table}
\end{equation}
For each past $p$ define $h_p(f)=\mathbb p_{x,z}(f,p)$ and
\begin{equation}
 M_{x,z}=\Span_{\R}\{h_p:p\in\mathsf P_{x,z}\}.
 \label{eq:supp-hankel-space}
\end{equation}

\begin{theorem}[Canonical typed Hankel realization]
\label{thm:supp-hankel}
For every primitive event
$a:(x,z)\to(y,T_{X,a}z)$ there is a unique linear map
\begin{equation}
 A_a:M_{x,z}\longrightarrow M_{y,T_{X,a}z},
 \qquad A_ah_p=h_{a\circ p}.
 \label{eq:supp-hankel-transition}
\end{equation}
Every future $f$ defines $\lambda_f(m)=m(f)$ and
\begin{equation}
 \mathbb p(f,p)=\lambda_f(h_p),
 \qquad
 \lambda_fA_a=\lambda_{f\circ a}.
 \label{eq:supp-hankel-evaluation}
\end{equation}
The realization is reachable and future separated.  For any other linear
realization $N_{x,z}$,
\begin{equation}
 N_{x,z}^{\mathrm{reach}}/N_{x,z}^{\mathrm{unobs}}
 \cong M_{x,z};
 \label{eq:supp-hankel-minimality}
\end{equation}
therefore every linear realization has dimension at least $\dim M_{x,z}$,
with equality exactly on the reachable future-separated branch.
\end{theorem}

\begin{proof}
If $\sum_jc_jh_{p_j}=0$, then for every target future $f$,
\[
 \sum_jc_jh_{a\circ p_j}(f)
 =\sum_jc_j\mathbb p(f\circ a,p_j)=0.
\]
Hence appending $a$ respects every column relation and
\eqref{eq:supp-hankel-transition} is well defined.  Reachability is built into
\eqref{eq:supp-hankel-space}; every nonzero response function is detected by
some future.  For another realization, map its reachable state $n$ to the
function $f\mapsto\mu_f(n)$.  The image is $M_{x,z}$ and the kernel is exactly
the future-unobservable subspace, proving
\eqref{eq:supp-hankel-minimality}.
\end{proof}

This is the reachable/observable Hankel construction of rational-series
realization theory \cite{BerstelReutenauer}, applied separately to the typed
operational sectors. The maps $A_a$ generate the table-native executable forward algebra.  A
Hilbert adjoint is not introduced at this stage; it is selected only after the
positive process tensor has fixed the support-minimal Hilbert realization.

\subsection{Comb tomography and support-minimal sequential realization}

For each exposed Hermitian Choi space select a real basis
$\{F_\alpha\}$ from a spanning physical frame, and let $\{D^\alpha\}$ be
its trace-dual basis, so $\Tr(D^\alpha F_\beta)=\delta^\alpha_\beta$.
An overcomplete frame may instead be used with a dual reconstruction frame;
the Kronecker identity is required only for the selected basis.  For an $N$-step table define
\begin{equation}
 R^{(N)}
 =\sum_{\alpha_1,\ldots,\alpha_N}
 p(\alpha_1,\ldots,\alpha_N)
 D^{\alpha_1}\otimes\cdots\otimes D^{\alpha_N}.
 \label{eq:supp-comb-tomography}
\end{equation}

\begin{theorem}[Finite process-comb tomography]
\label{thm:supp-comb-tomography}
The tensor $R^{(N)}$ in \eqref{eq:supp-comb-tomography} is independent of the
chosen spanning physical frame and is the unique Hermitian tensor representing
all multilinear intervention probabilities.  It is a deterministic finite
quantum process exactly when
\begin{equation}
 R^{(N)}\succeq0,
 \qquad
 \Tr_{2k-1}R^{(k)}=I_{2k-2}\otimes R^{(k-1)},
 \qquad
 R^{(0)}=1.
 \label{eq:supp-comb-conditions}
\end{equation}
Here support minimality means that the dilation of every prefix channel
is a minimal Stinespring dilation; all purifying memory is retained until
the corresponding final discard. On this branch there are memories
$A_k\cong\overline{\supp R^{(k)}}$, $A_0=\C$, and isometries
\begin{equation}
 V_k:H_{2k-2}\otimes A_{k-1}
 \longrightarrow H_{2k-1}\otimes A_k
 \label{eq:supp-comb-isometry}
\end{equation}
whose sequential link followed by final memory discard realizes $R^{(N)}$.
Every other prefix-support-minimal sequential isometric realization satisfies
\begin{equation}
 V'_k=(I\otimes U_k)V_k(I\otimes U_{k-1}^*)
 \label{eq:supp-comb-unitary}
\end{equation}
for unique memory unitaries $U_k$ on these supports. Named operators
retained as part of the compared record realization are transported by
these same unitaries. The minimum predictive state remains a separate
classical quotient; its sharp physical readability is not inferred from
purification alone.
\end{theorem}

\begin{proof}
Products of the physical frames span the full Hermitian tensor space, so
the dual-frame formula is the unique representing tensor. Positivity and
the partial-trace recursion are necessary by complete positivity and causal
normalization. Their sufficiency follows from the sequential construction
below, in the comb realization framework \cite{Chiribella2009}. The explicit
support-minimal construction is also given in \cite[Theorem~2]{Bisio2011}.

Let $|\psi_k\rangle=|(R^{(k)})^{1/2}\rangle\!\rangle$ with purifying
space $A_k=\overline{\supp R^{(k)}}$.  Its ancillary Schmidt support is
all of $A_k$, and tracing $A_k$ gives $R^{(k)}$.  Suppose the prefix through
$k-1$ has been realized. Tensor $|\psi_{k-1}\rangle$ with the unnormalized
maximally entangled vector on the new input and a copy of that input.
This is a minimal purification of
$R^{(k-1)}\otimes I_{2k-2}$, with purifying space
$A_{k-1}\otimes H_{2k-2}'$. Viewed over the same base space,
$|\psi_k\rangle$ purifies exactly that operator because of
\eqref{eq:supp-comb-conditions}; its purifying space is
$H_{2k-1}\otimes A_k$. Schmidt decomposition therefore supplies a unique
isometry from the first purifying support into the second. Identifying the
input copy with $H_{2k-2}$ gives $V_k$ in
\eqref{eq:supp-comb-isometry}. Induction and final discard realize the full
comb.

The prefix purifying dimension cannot be below $\rank R^{(k)}$, by Schmidt
rank, and the construction attains that value. For two such realizations,
minimal-purification uniqueness at prefix $k$ gives a unique $U_k$ relating
their prefix vectors. Insert that relation at consecutive prefixes.
Contraction with all external prefix vectors spans
$H_{2k-2}\otimes A_{k-1}$, by prefix support minimality, and therefore
gives the operator identity \eqref{eq:supp-comb-unitary}. This proves the
stated sequential, rather than merely terminal, uniqueness. It does not
assert uniqueness or the same memory cost for nonisometric realizations
that discard or classically compress memory between rounds.
\end{proof}

\begin{remark}[Why support reduction is necessary]
An arbitrary nonminimal positive factorization may contain spectator
directions.  Taking ambient Hilbert adjoints before removing them can create
mixed history words absent from the operational process.  The support-minimal
square-root class is therefore the canonical boundary at which a history
$C^*$-algebra becomes well defined.
\end{remark}

\subsection{History algebra, contextual separation, and flat word Grams}

Choose one support-minimal realization.  Let $\mathcal A_X^{\mathrm{raw}}$ be
the finite-dimensional $*$-algebra generated by the represented projections
onto $Z_X^{\min}$, primitive branch maps, and their Hilbert--Schmidt adjoints.
A represented element is \emph{contextually null} when every admitted left and
right history context and every terminal Read has zero matrix coefficient on
it.  Let $\mathcal J_X^{\mathrm{ctx}}$ be the set of all such elements.

\begin{lemma}[Contextual nullity is a two-sided $*$-ideal]
\label{lem:supp-contextual-ideal}
$\mathcal J_X^{\mathrm{ctx}}$ is the largest two-sided $*$-ideal invisible in
every admitted history context.  Hence
\begin{equation}
 \mathcal A_X^{\mathrm{hist}}
 =\mathcal A_X^{\mathrm{raw}}/\mathcal J_X^{\mathrm{ctx}}
 \label{eq:supp-history-algebra}
\end{equation}
is the unique minimum context-preserving separated history envelope.
\end{lemma}

\begin{proof}
If $a$ is null, inserting any represented word on either side only enlarges the
tested context, and reversing the context shows that $a^*$ is null.  Thus the
null set is a two-sided $*$-ideal.  Every context-invisible ideal is contained
in it by definition, which gives the universal property.
\end{proof}

For a finite-dimensional unital $C^*$-algebra $\mathcal A$, define the
normalized regular trace
\begin{equation}
 \widehat\tau_{\mathrm{reg}}^{\mathcal A}(a)
 =\frac{\Tr_{\C}(L_a)}{\dim_{\C}\mathcal A},
 \qquad L_a(b)=ab.
 \label{eq:supp-regular-trace}
\end{equation}
If $\mathcal A\cong\bigoplus_\lambda M_{n_\lambda}(\C)$, then
$\Tr_{\C}(L_a)=\sum_\lambda n_\lambda\Tr(a_\lambda)$, so
$\widehat\tau_{\mathrm{reg}}$ is a faithful tracial state invariant under every
$*$-isomorphism.

Let $\mathcal W_{X,\le r}$ have basis the composable represented history words
of length at most $r$, and let $W_re_w=[w]$ in
$\mathcal A_X^{\mathrm{hist}}$.  Define
\begin{equation}
 \mathbb K_{X,r}(v,w)
 =\widehat\tau_{\mathrm{reg}}([v]^*[w]).
 \label{eq:supp-word-gram}
\end{equation}

\begin{theorem}[Finite flat word reconstruction]
\label{thm:supp-word-flatness}
Every $\mathbb K_{X,r}$ is positive.  There is a least $r_X$ such that
\begin{equation}
 \rank\mathbb K_{X,r_X}=\rank\mathbb K_{X,r_X+1}.
 \label{eq:supp-flat-stop}
\end{equation}
At this depth the represented word span is the whole history algebra, and the
panel through depth $r_X+1$ reconstructs the object units, minimum predictive
state algebra, multiplication, involution, predictive core algebras, and every
represented mixed history.
\end{theorem}

\begin{proof}
Positivity follows from
$\sum_{v,w}\bar c_vc_w\mathbb K(v,w)
=\widehat\tau_{\mathrm{reg}}(a^*a)\ge0$.  Faithfulness makes the Gram kernel
exactly the represented word-relation space.  The nested word spans stabilize
because the history algebra is finite-dimensional.  Equality of two
consecutive ranks therefore means equality of two consecutive spans; the
stabilized span contains the units and is invariant under left multiplication
by every generator, so it equals the full generated algebra.  The one-step
extension reconstructs the generator multiplication matrices and reversed
words reconstruct the involution.
\end{proof}

\subsection{Covariance of word moments and independent tensor composition}

\begin{proposition}[Tensorial covariance and independent composition]
\label{prop:supp-grand-monoidal}
A change of word-space coordinates $c=C\widetilde c$ transforms the Gram
matrix by $\widetilde{\mathbb K}=C^*\mathbb K C$. A record-preserving
unitary equivalence of the support-minimal realization induces a
$*$-isomorphism of history algebras preserving all word moments.

Suppose two independent operational processes have product preparations,
independent future tests and independently accessible local operations,
including identity operations. Assume joint reachability is the product of
local reachability. With the product contextual test family, their composite
has
\begin{equation}
 Z_{12}^{\min}\cong Z_1^{\min}\times Z_2^{\min},\qquad
 \mathcal A_{12}^{\rm hist}\cong
 \mathcal A_1^{\rm hist}\otimes\mathcal A_2^{\rm hist}.
 \label{eq:supp-grand-product}
\end{equation}
On a rectangular bank of independent word pairs,
\begin{equation}
 \mathbb K_{12}((v_1,v_2),(w_1,w_2))
 =\mathbb K_1(v_1,w_1)\mathbb K_2(v_2,w_2).
 \label{eq:supp-grand-gram-product}
\end{equation}
These identifications respect reassociation and interchange of the independent
factors. No factorization is asserted for correlated processes or for a
restricted diagonal operation grammar.
\end{proposition}
\begin{proof}
The coordinate rule follows by substituting $c=C\widetilde c$ in the
quadratic form. A realizing unitary conjugates every generator and its
adjoint, carries the contextual null ideal to the corresponding null ideal,
and preserves the normalized regular trace under the resulting isomorphism.
For independent processes, future probabilities factor. Equality of both
marginal future signatures gives equality of the joint signature; conversely,
a test on one factor with discard on the other detects any unequal marginal.
Product reachability therefore identifies the joint quotient with the product.

Independently accessible generators contain $a\otimes I$ and $I\otimes b$,
so they generate the tensor product of the two separated finite algebras.
Product contexts separate this algebra: the spans of the local context
functionals are the full dual spaces after contextual separation, and their
products span the joint dual. There is consequently no further contextual
ideal. Finally
$L_{a\otimes b}=L_a\otimes L_b$ and
$\dim(\mathcal A_1\otimes\mathcal A_2)=\dim\mathcal A_1\dim\mathcal A_2$,
so the normalized regular trace is multiplicative. This proves the Gram
identity. All identifications are specified on elementary histories and
therefore commute with the usual associativity and interchange maps.
\end{proof}


\subsection{Proof of Theorem~\ref{mt:grand-tensor}}

\begin{proof}[Proof of Theorem~\ref{mt:grand-tensor}]
Item~\textup{(i)} is
\Cref{thm:supp-right-congruence,thm:supp-minimal-record}.  Item~\textup{(ii)}
is \Cref{thm:supp-persistence-derived}.  Item~\textup{(iii)} is
\Cref{thm:supp-hankel}.  Item~\textup{(iv)} is
\Cref{thm:supp-comb-tomography}.  Item~\textup{(v)} is
\Cref{lem:supp-contextual-ideal,thm:supp-word-flatness}.  Item~\textup{(vi)}
follows from \Cref{thm:supp-minimal-record,cor:supp-future-operational-transfer}
and support-minimal unitary uniqueness.  Item~\textup{(vii)} is
\Cref{prop:readable-relational-completion}.  None of these constructions
introduces an independent predictive state space, metric, connection, or
spacetime dimension beyond the actual operational rows explicitly retained.
\end{proof}

\subsection{Relative minimality of the operational entrance}
\label{subsec:supp-relative-minimality}

For the purpose of the foundational audit, keep the typed potential event
grammar fixed and distinguish four semantic rows.  Let $\mathfrak O$ denote
the actual-occurrence interface selecting which potential events, writers, and
Reads are physically present; let $\mathbb P^+$ denote the complete contextual
predictive law, including the finite ancillary/matrix-order positivity and
causal-recursion conditions used in the comb reconstruction; let $\Theta$
denote the protected orientation section; and let $\Lambda$ denote the
absolute physical calibration needed when durations, rates, lengths, or
measures are interpreted in physical units.  This separation is an audit of
information content, not an enlargement of the operational model above.

\begin{theorem}[Relative primitive-basis minimality for the paper-level entrance]
\label{thm:supp-relative-primitive-floor}
Within the fixed category of finite typed operational-predictive models used in
this paper, the semantic entrance
\begin{equation}
 [\mathfrak G;\mathfrak O,\mathbb P^+,\Theta,\Lambda]
 \label{eq:supp-relative-floor}
\end{equation}
has the following status.
\begin{enumerate}[label=\textup{(R\arabic*)},leftmargin=3em]
\item The rows already assumed in the paper are sufficient for the foundational
      reconstruction actually used below: $\mathbb P^+$ and composition give
      the minimum future quotient and persistent predictive carrier, the
      contextual positive table gives the support-minimal causal process and
      history algebra, while $\Theta$ and $\Lambda$ supply precisely the signed
      and absolute-scale information required on the Lorentzian ADM branch.
\item Each of $\mathfrak O$, $\mathbb P^+$, $\Theta$, and $\Lambda$ is
      independent of the other three relative to the fixed typed grammar: for
      each forgotten row there are two anchored finite models with identical
      reducts and inequivalent full models.
\item These four independence obstructions survive arbitrary common
      conservative refinement and cofinal reindexing; hence no functorial
      cutoff limit can reconstruct a forgotten row from the other three.
\item Protected orientation can be absorbed into occurrence exactly when the
      occurrence packet contains a named actual Read together with a physical
      decoder selecting one orientation section.  This is a definitional
      bundling, not an informational reduction of the orientation choice.
\end{enumerate}
Consequently the displayed semantic entrance is a minimal known sufficient and
independent basis relative to the declared operational language and target
class.  No claim of absolute minimality among all conceivable foundational
languages is made.
\end{theorem}

\begin{proof}
Sufficiency for the predictive and positive-process part is exactly
\Cref{thm:supp-right-congruence,thm:supp-minimal-record,cor:supp-derived-state-machine,thm:supp-comb-tomography,lem:supp-contextual-ideal,thm:supp-word-flatness}.  On the geometric branch,
\Cref{lem:supp-race-score,thm:supp-adm-frame,lem:supp-fast-clock} show how the
physically calibrated root race supplies the absolute kinetic data used by the
ADM inversion, while
\Cref{lem:supp-positive-krein-distinct,lem:supp-clifford-cell,cor:supp-lorentzian-square} use the protected sign and do not derive it from
positive data.

For row-irredundancy it is enough to give one anchored separating pair for each
row.  For $\mathbb P^+$, fix one complete binary instrument with named outcomes
$a,b$, keep every other row fixed, and compare Bernoulli laws
$\Pr(a)=p$ and $\Pr(a)=q$ with $0<p\ne q<1$.  Forgetting the predictive row
makes the models identical, whereas any full anchored isomorphism fixes the
outcome names and therefore cannot identify the two probabilities.

For $\mathfrak O$, adjoin a potential deterministic writer $W=0$ whose
mathematical extension is unique.  In one model retain only the base event as
actual; in the other declare $W$ to be an actual readable writer as well.  The
structural and probabilistic data agree after the occurrence incidence is
forgotten, but the full models differ on whether $W$ belongs to the physical
interface.  Thus structural realizability, even when unique, does not imply
actual occurrence.

For $\Theta$, take the same anchored real two-plane with named basis
$e_1,e_2$ and positive Gram $I_2$ in both models, but choose the unit
alternating forms $\omega(e_1,e_2)=+1$ and
$\omega(e_1,e_2)=-1$.  Every positive Gram datum agrees.  An anchored
isomorphism fixes the ordered basis and cannot reverse the alternating form,
so the protected orientation is not selected by the positive packet.

For $\Lambda$, fix a projective race with winner law $p_a$ and an internal
waiting variable $S\sim\operatorname{Exp}(1)$ independent of the winner.  Two
models with the same grammar, occurrence, predictive law, and orientation but
physical durations $T=c_0S$ and $T=c_1S$, $c_0\ne c_1$, have identical
uncalibrated reducts and different absolute rates $k_a=p_a/c_i$.  Hence the
common physical scale is not determined by normalized/projective race data.
This is the same scale direction whose loss is visible in
\Cref{lem:supp-race-score} when waiting-time calibration is omitted.

To see cofinal stability, tensor either member of any separating pair with the
same strictly refining finite spectator process and let every cutoff map append
and later discard only spectator records.  The row-forgotten theories are then
naturally isomorphic at every cutoff.  If a full natural isomorphism appeared
on a cofinal tail, composing it with spectator discard would give an anchored
isomorphism of the original finite separating pair, a contradiction.  Thus
common refinement or passage to a cofinal limit cannot create the missing row.

Finally, if occurrence itself contains a named actual orientation Read and a
deterministic physical decoder into the two-element orientation torsor, the
decoded section and a separately named $\Theta$ determine one another.  If the
decoder is omitted, the two opposite sections above remain compatible with the
same positive occurrence statistics.  The former is therefore only a change
of packaging, while the latter retains genuine orientation independence.
\end{proof}

\paragraph{The missing calibration scale.}
\label{rem:main-calibration-fibre}
Winner probabilities lose the common rate scale; the following result
states the complete log-affine calibration freedom and the required anchors.

\begin{theorem}[Exact calibration fibre and minimum anchor bank]
\label{thm:supp-calibration-fibre}
Write all positive physical scales and additive origins in log-affine
coordinates $x\in\mathbb R^N$.  Suppose the physically established incidence
relations are
\begin{equation}
 Lx=b.
 \label{eq:supp-calibration-incidence}
\end{equation}
If the solution set is nonempty and $x_0$ is one solution, then:
\begin{enumerate}[label=\textup{(K\arabic*)},leftmargin=3em]
\item the complete calibration fibre is
      \begin{equation}
       x_0+\mathcal K_L,
       \qquad \mathcal K_L:=\Ker L,
       \label{eq:supp-calibration-fibre}
      \end{equation}
      of dimension $N-\rank L$;
\item exactly $\dim\mathcal K_L$ independent scalar anchors are necessary and
      sufficient to determine one calibration in this affine fibre;
\item if $U:\mathbb R^d\to\mathbb R^N$ encodes conventional changes of base
      units with $LU=0$, the intrinsic physical residual is
      \begin{equation}
       \mathcal K_L^{\rm phys}
       =\mathcal K_L/\operatorname{Im}U,
       \label{eq:supp-physical-calibration-fibre}
      \end{equation}
      and one anchor per independent quotient direction is necessary and
      sufficient;
\item winner probabilities for a root race instantiate a one-dimensional
      common-rate direction of this type: multiplying all directed rates by the
      same positive constant leaves the winner law unchanged, whereas the
      waiting-time scale separates that direction.
\end{enumerate}
\end{theorem}

\begin{proof}
If $Lx=Lx_0=b$, then $x-x_0\in\Ker L$, and conversely every
$x_0+k$, $k\in\Ker L$, solves \eqref{eq:supp-calibration-incidence}; this gives
\eqref{eq:supp-calibration-fibre} and its dimension by rank--nullity.  A basis
of $\mathcal K_L$ and the corresponding dual linear functionals give a
sufficient scalar anchor bank.  Any smaller bank has nontrivial common kernel
on $\mathcal K_L$ and is therefore insufficient.  Quotienting the subgroup of
pure unit conventions gives
\eqref{eq:supp-physical-calibration-fibre} and the same argument on the quotient
gives the intrinsic anchor count.  For the root race, the winner law is
$p_b=k_b/\sum_ck_c$ and is invariant under $k_b\mapsto ck_b$; the exponential
waiting time has total rate $\sum_ck_c$ and therefore fixes the missing common
scale.
\end{proof}


\section{Minimal relational dimension and ADM reconstruction}
\label{supp:dimension-adm}

This appendix supplies the algebraic and statistical details of
\Cref{mt:adm}. It assumes the carrier of Section~\ref{sec:grand-tensor-main},
keeps the two relational selectors separate, and then treats calibration,
positive realizability, signed extension and Clifford representation.

\subsection{Endpoint--cycle balance and the simplex comparison}

The primary regular-graph classification is proved directly in
\Cref{prop:main-regular-graph-k4}. The stronger simplex category is developed
next; its extra assumptions are not used in that proof.

\begin{assumption}[Minimal faithful predictive relation branch]
\label{ass:supp-relational-branch}
A local relation event has one future-readable source and one future-readable
target.  Physical reversal retains both orientations.  Future-minimality
identifies parallel relation labels exactly when every continuation identifies
them.  A nonempty relation set is invariant under a group acting transitively
on unordered endpoint pairs.
\end{assumption}

\begin{lemma}[Pair-homogeneous completion]
\label{lem:supp-pair-completion}
Under \Cref{ass:supp-relational-branch}, the relation-colour-free local skeleton
is the complete graph $K_N$ on the endpoint set.
\end{lemma}

\begin{proof}
The selected branch of the derived persistent predictive carrier contains
exactly one future-readable source and one future-readable target, so one
occurring relation gives one unoriented pair after physical reversal.  The
orbit of that pair under a group transitive on $\binom{X}{2}$ is the complete
set of unordered pairs.  Future-minimality removes only unread parallel labels,
not distinct endpoint pairs.  Hence the skeleton is $K_N$.
\end{proof}

Let $X=\{1,\ldots,N\}$, let $u_0=N^{-1/2}\one_N$, and write
$\R^N=\R u_0\oplus W_N$.  Identify the signed edge $[i\to j]$ with
$e_i\wedge e_j\in\Lambda^2\R^N$.  The boundary map is
$\partial(e_i\wedge e_j)=e_j-e_i$.

\begin{theorem}[Cut--cycle decomposition]
\label{thm:supp-cut-cycle}
There is an orthogonal decomposition
\begin{equation}
 C_1(K_N)=u_0\wedge W_N\oplus\Lambda^2W_N,
 \label{eq:supp-cut-cycle}
\end{equation}
with
\begin{equation}
 \partial(u_0\wedge w)=\sqrt N\,w,
 \qquad
 \Ker\partial=\Lambda^2W_N.
 \label{eq:supp-boundary-action}
\end{equation}
Consequently
\begin{equation}
 \dim\Ran\partial=N-1,
 \qquad
 \dim\Ker\partial=\binom{N-1}{2}.
 \label{eq:supp-cut-cycle-dim}
\end{equation}
\end{theorem}

\begin{proof}
The exterior square of
$\R^N=\R u_0\oplus W_N$ is
$u_0\wedge W_N\oplus\Lambda^2W_N$.  For $x,y\in\R^N$,
\[
 \partial(x\wedge y)
 =\left(\sum_ix_i\right)y-
   \left(\sum_iy_i\right)x.
\]
Taking $x=u_0$, $y=w\in W_N$ gives the first identity.  If both vectors lie in
$W_N$, both coordinate sums vanish, so the second summand is in the kernel.
Dimension counting shows that it is the whole kernel and gives
\eqref{eq:supp-cut-cycle-dim}.
\end{proof}

\begin{remark}[Metric tomography does not select $N$]
The root dyadics of a complete simplex define an isomorphism
\[
 \R^{\binom N2}\longrightarrow\Sym(W_N),
 \qquad
 (a_{ij})\longmapsto\sum_{i<j}a_{ij}
 (e_j-e_i)(e_j-e_i)^{\mathsf T}
\]
for every $N$.  Indeed, a vanishing weighted sum is a graph Laplacian whose
off-diagonal entries are $-a_{ij}$.  Thus the equality
$\binom N2=\dim\Sym(W_N)$ is universal and cannot by itself select three
spatial dimensions.  The paper uses two different non-tomographic selectors:
exact endpoint--cycle source-rank balance, and the stronger
alternating/minimal-source simplex criterion.
\end{remark}

\begin{remark}[Conditional character of the selector]
The alternating form in the next theorem is part of the definition of the
selected relational category.  It is not reconstructed from the positive word
Grams and does not, by itself, specify a metric, a Hamiltonian, a symplectic
phase space of continuum fields, or a Lorentzian sign.  Accordingly the theorem
asserts uniqueness only among faithful pair-homogeneous simplex cells satisfying
this alternating/minimal-source query.  It is the second dimensional selector,
not the sole derivation of $K_4$; the regular-graph source-balance result below
uses no alternating form.
\end{remark}

\begin{theorem}[Minimal faithful $K_4$ selector on the alternating simplex branch]
\label{thm:supp-k4-selector}
Assume, in addition to \Cref{ass:supp-relational-branch}, that:
\begin{enumerate}[label=\textup{(D\arabic*)},leftmargin=3em]
\item the temporal--spatial coefficient space is $\R t\oplus W_N$ and carries a
      nondegenerate alternating form;
\item the endpoint and complete signed-edge source Grams are faithful;
\item the cycle source is nonzero;
\item among conservative completions satisfying these queries, the
      future-readable source innovation is minimal.
\end{enumerate}
Then
\begin{equation}
 N=4,
 \qquad
 G_{\mathrm{rel}}=K_4,
 \qquad
 W_{\mathrm{sp}}=W_4\cong\R^3.
 \label{eq:supp-k4-result}
\end{equation}
The faithful endpoint and edge source ranks are
$r_0(N)=N-1$ and $r_1(N)=\binom N2$.
\end{theorem}

\begin{proof}
The alternating coefficient space has dimension
$1+\dim W_N=N$.  A nondegenerate alternating form exists only in even
dimension, so $N$ is even.  A nonzero cycle source requires
$\Lambda^2W_N\ne0$, hence $N\ge3$; together with evenness this gives $N\ge4$.
On the even integers $N\ge4$, both $N-1$ and $\binom N2$ are strictly
increasing.  The unique minimum of the pair of faithful source ranks is therefore
attained at $N=4$.  The graph conclusion follows from
\Cref{lem:supp-pair-completion}, and $\dim W_4=3$.
\end{proof}

\begin{remark}[The minimization category]
The minimum in \Cref{thm:supp-k4-selector} is taken over the declared
admissible simplex cells, with source cost ordered by the pair
$(N-1,\binom N2)$; equivalently one may use any cost strictly increasing in
both faithful ranks. It is not a theorem that a fixed law with more
future-distinguishable endpoints has a conservative four-endpoint
compression. Such a compression would contradict faithfulness unless the
additional distinctions were genuinely unread. The independent
endpoint--cycle balance condition is not inferred from this minimum rule.
\end{remark}

A representation-theoretic strengthening explains why the cycle packet is
particularly economical at $N=4$.  After complexification,
$W_N\cong S^{(N-1,1)}$,
$\Lambda^2W_N\cong S^{(N-2,1,1)}$, and
$W_N\otimes\operatorname{sgn}\cong S^{(2,1^{N-2})}$.  These Young diagrams
coincide exactly for $N=4$.  Hence
\begin{equation}
 \Hom_{S_N}(\Lambda^2W_N,W_N\otimes\operatorname{sgn})
 =\begin{cases}\C,&N=4,\\0,&N\ne4,\end{cases}
 \label{eq:supp-oriented-copy}
\end{equation}
on the complexified nontrivial simplex branch; the corresponding real
intertwiner space is one-dimensional over $\R$ at $N=4$. Thus only $K_4$ identifies the entire cycle
representation with one oriented spatial copy rather than an additional
inequivalent source type.


The source-equivalence condition is
\begin{equation}
 \Lambda^2W_N\cong W_N\otimes\operatorname{sgn},\qquad N=4.
 \label{eq:main-source-equivalence}
\end{equation}

\begin{proposition}[Faithful orientation-twisted source equivalence]
\label{prop:supp-source-equivalence}
For $N\ge3$, an $S_N$-equivariant isomorphism
$\Lambda^2W_N\to W_N\otimes\operatorname{sgn}$ exists if and only if
$N=4$. At $N=4$ its space is one dimensional. With the standard invariant
Euclidean forms, an isometric choice is unique up to sign.
\end{proposition}
\begin{proof}
An isomorphism requires
$(N-1)(N-2)/2=N-1$, which gives $N=4$. On the three-dimensional space
$W_4$, the Hodge identification $\Lambda^2W_4\to W_4$ satisfies
$*(gu\wedge gv)=\det(g)g*(u\wedge v)$. The determinant of the standard
permutation representation restricted to $W_4$ is $\operatorname{sgn}(g)$,
since its complementary constant line is fixed. Thus Hodge identification
is the required twisted intertwiner. The complexified standard
representation is irreducible and its commuting endomorphisms are scalar;
composing any two intertwiners proves uniqueness up to a real scalar.
Normalization of the invariant forms restricts that scalar to $\pm1$.
\end{proof}


\subsection{The \texorpdfstring{$A_3$}{A3} source frame}

In the anchored basis of \eqref{eq:a3-roots-main}, define
$Q_a=r_ar_a^{\mathsf T}$.

\begin{lemma}[Root dyadics span $\Sym_3$]
\label{lem:supp-root-frame}
The six matrices $Q_1,\ldots,Q_6$ form a basis of $\Sym_3$.
\end{lemma}

\begin{proof}
$Q_1,Q_2,Q_3$ are the three diagonal matrix units.  For $i<j$,
\[
 (e_j-e_i)(e_j-e_i)^{\mathsf T}-e_ie_i^{\mathsf T}-e_je_j^{\mathsf T}
 =-(e_ie_j^{\mathsf T}+e_je_i^{\mathsf T}).
\]
The three difference roots therefore recover the three symmetric off-diagonal
matrix units.  The six matrices are linearly independent and span the
six-dimensional space $\Sym_3$.
\end{proof}

Let the twelve oriented roots race with positive rates $k_a^\pm$.  Put
$K=\sum_a(k_a^++k_a^-)$, let $A$ be the winning oriented root, and let $T$ be
the first waiting time.

\begin{lemma}[Winner--waiting score is faithful]
\label{lem:supp-race-score}
For each oriented branch $b$, the log-rate score is
\begin{equation}
 s_b=\one_{\{A=b\}}-k_bT.
 \label{eq:supp-race-score}
\end{equation}
Its Gram matrix is diagonal and positive:
\begin{equation}
 \mathbb E[s_bs_c]=\delta_{bc}\frac{k_b}{K}.
 \label{eq:supp-race-score-gram}
\end{equation}
Winner probabilities alone have covariance
$\diag(p)-pp^{\mathsf T}$ and therefore lose the common rate scale.
\end{lemma}

\begin{proof}
For independent exponential clocks, $T\sim\operatorname{Exp}(K)$, the winner
has probability $p_b=k_b/K$, and $A$ and $T$ are independent.  Therefore
\[
 \begin{aligned}
 \mathbb E[s_bs_c]
 &={}\delta_{bc}p_b-p_bk_c\mathbb ET-k_bp_c\mathbb ET
      +k_bk_c\mathbb ET^2\\
 &={}\delta_{bc}\frac{k_b}{K}
   -2\frac{k_bk_c}{K^2}+2\frac{k_bk_c}{K^2},
 \end{aligned}
\]
which proves \eqref{eq:supp-race-score-gram}.  The covariance of the winner
indicator is singular in the all-ones rate direction because probabilities
are unchanged by common rescaling.
\end{proof}

\subsection{Positive-rate realizability}

The dyadic span does not make every positive matrix a positive rate sum.
For the cone \eqref{eq:main-root-positive-cone}, the explicit inversion is
\begin{equation}
 \begin{gathered}
 s_4=-\mathcal B_{12},\quad s_5=-\mathcal B_{13},\quad
 s_6=-\mathcal B_{23},\\
 s_1=\mathcal B_{11}+\mathcal B_{12}+\mathcal B_{13},\quad
 s_2=\mathcal B_{22}+\mathcal B_{12}+\mathcal B_{23},\quad
 s_3=\mathcal B_{33}+\mathcal B_{13}+\mathcal B_{23}.
 \end{gathered}
 \label{eq:main-root-cone-inverse}
\end{equation}
These six strict inequalities are a finite realizability test. The ADM
formula is an analytic bijection on abstract positive bracket--density
pairs; its realization by the declared positive race is restricted to
\eqref{eq:main-root-positive-cone}. In the orthonormal $D_3$ coordinates
used for the constructive small-data class, the isotropic bracket is an
interior point of the corresponding cone. A bounded shift also has a positive
directed lift at sufficiently fine cutoff: for a fixed right inverse $J_R$
of $R=(r_1\ \cdots\ r_6)$, choose
$d=h^{-1}J_R\beta$ and
$k_a^\pm=(s_a/h^2\pm d_a)/2$. Positivity holds whenever
$h|(J_R\beta)_a|<s_a$. Thus local positive realizability follows from an
interior margin, not from the dimension of the dyadic span alone.
\subsection{ADM inversion and scheduler scale}

Define $\mathcal B_h$ and $\beta_h$ by
\eqref{eq:bracket-drift-main}.

\begin{theorem}[Complete finite ADM frame]
\label{thm:supp-adm-frame}
The six rate sums map isomorphically to $\mathcal B_h\in\Sym_3$.  The six rate
differences map with rank three onto $\beta_h\in\R^3$ and have a
three-dimensional kernel, the tetrahedral circulation space.  For every
$\mathcal B_h\succ0$ and $\varrho_h>0$, the formula
\eqref{eq:adm-inversion-main} is the unique solution of
\begin{equation}
 \mathcal B_h=N_h^2g_h^{-1},
 \qquad
 \varrho_h=\sqrt{\det g_h}.
 \label{eq:supp-adm-characterization}
\end{equation}
Hence the root rates, waiting-time scale, and speed density reconstruct the six
spatial metric components, three shift components, and lapse.
\end{theorem}

\begin{proof}
The first assertion is \Cref{lem:supp-root-frame}.  The difference map has
image containing $r_1,r_2,r_3$, hence has rank three; rank--nullity gives a
three-dimensional kernel.  For the inversion, take determinants in
$\mathcal B=N^2g^{-1}$:
$\det\mathcal B=N^6/(\det g)=N^6/\varrho^2$.  Thus
$N=\varrho^{1/3}(\det\mathcal B)^{1/6}$.  Substituting back gives
$g=N^2\mathcal B^{-1}$, which is exactly
\eqref{eq:adm-inversion-main}.  Positivity makes the solution unique.
\end{proof}

\begin{lemma}[Fast-clock lower bound]
\label{lem:supp-fast-clock}
Suppose $\norm{r_a}\le R$ and $\mathcal B_h\succeq cI_3$.  Then
\begin{equation}
 \sum_a(k_a^++k_a^-)\ge\frac{3c}{R^2h^2}.
 \label{eq:supp-fast-clock-bound}
\end{equation}
In particular, a total rate $O(h^{-1})$ produces $\mathcal B_h=O(h)$ and a
collapsed continuum bracket.
\end{lemma}

\begin{proof}
Taking traces in \eqref{eq:bracket-drift-main},
\[
 3c\le\tr\mathcal B_h
 \le R^2h^2\sum_a(k_a^++k_a^-).
\]
This gives \eqref{eq:supp-fast-clock-bound}.  The final assertion follows by
substitution.
\end{proof}

\subsection{Finite Krein orientation from the protected sign}

Let $\epsilon_X\in\{+1,-1\}$ be the actual protected orientation Read of
\eqref{eq:orientation-read-main} on the finite support-minimal relational
carrier $\mathcal H_J$.  Functional calculus of this classical two-point Read
gives orthogonal projections $P_\pm$ with $P_++P_-=I$; define
$J=P_+-P_-$.  Thus
\begin{equation}
 \mathcal H_J=\mathcal H_J^+\oplus\mathcal H_J^-,
 \qquad
 \mathcal H_J^\pm=\Ran P_\pm,
 \qquad
 J|_{\mathcal H_J^\pm}=\pm I.
 \label{eq:supp-krein-splitting}
\end{equation}
Assume both summands are nonzero.  Define
\begin{equation}
 [u,v]_J=\langle u,Jv\rangle.
 \label{eq:supp-krein-form}
\end{equation}
Then $[\cdot,\cdot]_J$ is a nondegenerate Hermitian form with one positive and one negative sector in the orientation factor, hence a finite Krein structure.  This statement uses only the protected sign involution and the support-minimal Hilbert realization; no continuum metric is involved.

\begin{lemma}[Positive geometry and Krein signature are distinct layers]
\label{lem:supp-positive-krein-distinct}
The positive word-Gram hierarchy of \Cref{thm:supp-word-flatness} remains positive after adjoining the protected orientation record, while the indefinite sign is carried entirely by $J$.  Consequently no positive combination of the spatial root dyadics can replace the Krein involution, and the Lorentzian signature of the principal symbol requires the signed layer in addition to the $A_3$ bracket.
\end{lemma}

\begin{proof}
The word-Gram hierarchy is a Gram matrix for the faithful positive regular trace and is therefore positive semidefinite.  Likewise every spatial bracket is a positive sum $\sum_a c_a r_ar_a^{\mathsf T}$ with $c_a>0$.  Such a form has no negative direction.  By contrast, $J$ has both $+1$ and $-1$ spectral subspaces by \eqref{eq:supp-krein-splitting}; the form \eqref{eq:supp-krein-form} is therefore indefinite.  The two structures are independent in exactly the required sense: the spatial geometry remains positive while $J$ supplies the sign subsequently represented by the temporal Clifford generator.
\end{proof}

\begin{theorem}[Irreducible sign and Lorentzian inertia before Clifford realization]
\label{thm:supp-signature-minimality}
Let $W_4$ carry the positive spatial form reconstructed by
\Cref{thm:supp-adm-frame}.  The complete positive Grand-Tensor packet cannot
supply a negative direction, while the protected orientation row is independent
of that packet by \Cref{thm:supp-relative-primitive-floor}.  If the geometric
extension retains the positive restriction on $W_4$ and satisfies
\Cref{def:main-signed-extension} on exactly one independent line $\R t$,
then, up to overall sign,
\begin{equation}
 q(\tau t+w)=-a^2\tau^2+g(w,w),\qquad a>0,
 \qquad \operatorname{inertia}(q)=(1,3).
 \label{eq:supp-signature-inertia}
\end{equation}
The conclusion depends only on positivity of the spatial form, independence of
the signed row, and one-line minimality; it does not depend on a choice of gamma
matrices.
\end{theorem}
\begin{proof}
The positive word-Gram and root-race forms have no negative eigenvalues by
\Cref{lem:supp-positive-krein-distinct}.  The separating pair in the proof of
\Cref{thm:supp-relative-primitive-floor} shows that reversing the protected
orientation leaves all positive operational data unchanged, so its orientation is not
a derived eigenvalue of the positive packet. Assignment of a negative
quadratic form to the added line is the separate signed-extension condition.  A minimal extension introduces
one real line and no additional signed directions.  Nondegeneracy makes its
coefficient nonzero; choosing the relative sign opposite to the positive
spatial restriction gives one negative eigenvalue and the three positive
eigenvalues of $g$.  Sylvester's law of inertia makes this count invariant
under every real change of basis.  Multiplying the entire form by $-1$ reverses
the sign convention but not the Lorentzian type.  No Clifford representation
has been used.
\end{proof}

\subsection{Clifford realization of the signed quadratic form}

The inertia has already been fixed by
\Cref{thm:supp-signature-minimality}. We now give the explicit representation
used in the geometric reconstruction.

The matrix representation is now introduced only after
\Cref{prop:main-signature-minimality}.  A minimal depth-two realization uses
binary factors $A,B$ to represent the four Clifford axes.  The physical sign
content remains the single Read \eqref{eq:orientation-read-main}; the second
matrix factor is representation-theoretic multiplicity and is not a second
operational time-orientation choice.  Set
\begin{equation}
 \gamma^0=-iX_A\otimes I_B,
 \quad
 \gamma^1=Z_A\otimes Z_B,
 \quad
 \gamma^2=Z_A\otimes Y_B,
 \quad
 \gamma^3=Y_A\otimes I_B.
 \label{eq:clifford-main}
\end{equation}
They satisfy
$(\gamma^0)^2=-I$, $(\gamma^a)^2=I$, and
$\gamma^\mu\gamma^\nu+\gamma^\nu\gamma^\mu=0$ for $\mu\ne\nu$; together they
generate $M_4(\C)\cong\Cl_{1,3}(\C)$.  If $E_h^{\mathsf T}E_h=g_h^{-1}$, the
principal symbol
\begin{equation}
 \sigma_h(\xi)
 =\gamma^0N_h^{-1}(\xi_0-\beta_h\!\cdot\!\vec\xi)
  +\gamma^a(E_h)_a{}^i\xi_i
 \label{eq:dirac-symbol-main}
\end{equation}
has square
\begin{equation}
 \sigma_h(\xi)^2
 =\left[-N_h^{-2}(\xi_0-\beta_h\!\cdot\!\vec\xi)^2
        +g_h^{ij}\xi_i\xi_j\right]I.
 \label{eq:dirac-square-main}
\end{equation}
The Clifford cell therefore represents the Lorentzian form fixed in
\eqref{eq:main-signature-inertia}; it does not create its one-minus-three-plus
signature.

\begin{lemma}[Depth-two Lorentzian Clifford cell]
\label{lem:supp-clifford-cell}
The matrices $\gamma^\mu$ satisfy
\begin{equation}
 (\gamma^0)^2=-I,
 \qquad
 (\gamma^a)^2=I,
 \qquad
 \{\gamma^\mu,\gamma^\nu\}=0\quad(\mu\ne\nu),
 \label{eq:supp-clifford-relations}
\end{equation}
and generate $M_4(\C)\cong\Cl_{1,3}(\C)$.
\end{lemma}

\begin{proof}
The square relations follow from $X^2=Y^2=Z^2=I$.  Every mixed product contains
exactly one local Pauli anticommutation, so the mixed anticommutators vanish.
The generators contain $X_A\otimes I$ and $Y_A\otimes I$, hence also
$Z_A\otimes I$.  Multiplication by $\gamma^1$ and $\gamma^2$ then gives
$I\otimes Z_B$ and $I\otimes Y_B$, hence $I\otimes X_B$.  Both local matrix
factors are generated, so the total algebra is
$M_2(\C)\otimes M_2(\C)=M_4(\C)$.
\end{proof}

\begin{corollary}[Lorentzian principal square]
\label{cor:supp-lorentzian-square}
Let $E_h^{\mathsf T}E_h=g_h^{-1}$ and define $\sigma_h$ by
\eqref{eq:dirac-symbol-main}.  Then \eqref{eq:dirac-square-main} holds.
If $g_h\succ0$ and $N_h>0$, the characteristic quadratic form realizes the
already determined inertia $(1,3)$ in the convention $(-,+,+,+)$.
\end{corollary}

\begin{proof}
The temporal--spatial cross terms vanish by the Clifford anticommutation
relations.  The temporal square contributes
$-N_h^{-2}(\xi_0-\beta_h\cdot\vec\xi)^2$, while the spatial square is
$\delta^{ab}(E_h)_a{}^i(E_h)_b{}^j\xi_i\xi_j=g_h^{ij}\xi_i\xi_j$.
\end{proof}

\subsection{Cofinal transport and proof of \texorpdfstring{\Cref{mt:adm}}{the ADM reconstruction theorem}}

For adjacent cutoffs, let the complete old/detail positive form have block
matrix
\begin{equation}
 M_Y=\begin{pmatrix}A_{Y/X}&B_{Y/X}\\
                     B_{Y/X}^*&D_{Y/X}\end{pmatrix},
 \label{eq:supp-refinement-block}
\end{equation}
and Schur complement
$S_{Y/X}=A_{Y/X}-B_{Y/X}D_{Y/X}^{\dagger}B_{Y/X}^*$.  Define
\begin{equation}
 \sigma_{Y/X}=\norm{S_{Y/X}-M_X}_{\op},
 \qquad
 \theta_{Y/X}=\norm{B_{Y/X}D_{Y/X}^{\dagger}}_{\op},
 \qquad
 \mu_{Y/X}=\lambda_{\min}^+(D_{Y/X}).
 \label{eq:supp-refinement-defects}
\end{equation}
The harmonic lift
\begin{equation}
 \mathcal H_{Y/X}u
 =J_{Y/X}u-K_{Y/X}D_{Y/X}^{\dagger}B_{Y/X}^*u
 \label{eq:supp-harmonic-lift}
\end{equation}
minimizes the fine energy over detail completions and has energy
$\ip{u}{S_{Y/X}u}$.  This is ordinary completion of the square.

If
\begin{equation}
 \sum_X\bigl(\sigma_X+\theta_X+\mu_X^{-1}\bigr)<\infty,
 \label{eq:supp-summable-refinement}
\end{equation}
then the finite forms and their regular isolated spectral projections are
Cauchy after harmonic transport.  In particular, a separated threshold rank,
a positive pair gap, and a finite-dimensional Clifford source persist.  This
is the finite-dimensional core of the common-screen limit used here.

\begin{proof}[Proof of Theorem~\ref{mt:adm}]
The predictive carrier is supplied by \Cref{thm:supp-persistence-derived}.
For the regular-graph route, the local dimension follows directly from
\Cref{prop:main-regular-graph-k4}.  For the pair-homogeneous simplex route, the
local graph is complete by \Cref{lem:supp-pair-completion} and the same
$K_4/A_3$ module follows from \Cref{thm:supp-cut-cycle,thm:supp-k4-selector}.
Thus the alternating form is needed only on the second route.  The root sums and
differences, faithful rate coordinates, ADM inversion, and clock lower bound
are \Cref{lem:supp-root-frame,lem:supp-race-score,thm:supp-adm-frame,lem:supp-fast-clock}.
The positive/signed separation is \Cref{lem:supp-positive-krein-distinct};
\Cref{thm:supp-signature-minimality} fixes Lorentzian inertia before the
Clifford matrices are introduced, and
\Cref{lem:supp-clifford-cell,cor:supp-lorentzian-square} realizes that form.
Finally,
\eqref{eq:supp-summable-refinement} makes every corrected finite matrix packet
Cauchy in operator norm on each protected finite screen; spectral separation
therefore preserves the ranks and positive gaps, while the Schur mismatch and
leakage tails vanish.  The canonical limit quotient consequently retains the
$K_4/A_3$ module, the physically occurring temporal calibration, the ADM packet,
and the Lorentzian Clifford cell.
\end{proof}

\begin{remark}[Exact dimension scope]
The result applies to the derived future-minimal, one-source/one-target, reciprocal,
pair-homogeneous, faithful, and minimal faithful branch.  A deliberately
higher-arity boundary, a relation-coloured theory, a nonfaithful cycle quotient,
or an enriched cross-product packet changes the primitive query and is not a
counterexample.  In particular, a seven-dimensional vector cross product
$\Lambda^2W\to W$ has a fourteen-dimensional kernel and requires an additional
stable $G_2$ three-form; it is an enriched theory rather than an equal-cost
realization of the faithful $K_4$ branch.
\end{remark}


\section{Action reconstruction, stationarity and gravitational classification}
\label{supp:dynamics}

This appendix proves the finite action and stationarity statements of
Section~\ref{sec:einstein-dynamics}. The determining-field quotient,
renewal pressure and common action are reconstructed before comparison
with the accepted law; the final steps classify the restricted local
representative and control torsion. Continuum estimates are in
Appendix~\ref{supp:continuum}.
\subsection{Accepted determining-field quotient and kernel}

Let $\widetilde\Theta_X$ be the finite accepted ledger of
\Cref{eq:main-determining-field}.  Two states $x,y$ are future equivalent when
every declared determining coordinate and every admitted future word agrees
from $x$ and $y$, including all branch-domain indicators.  Let
$\Theta_{{\rm fld},X}^{\min}$ be the quotient.


\begin{proof}[Proof of \Cref{thm:main-determining-kernel}]
Future equality is a typed right congruence.  If $x\sim_{{\rm fld},X}y$, then
for every primitive branch $a$ and continuation $w$, every determining read
satisfies
\[
 d(u_wu_a(x))=d(u_{wa}(x))=d(u_{wa}(y))=d(u_wu_a(y)),
\]
so every primitive update descends to the quotient.  At each refinement round,
split two states whenever an immediate coordinate differs or one primitive
update reaches different blocks.  The $k$th partition identifies precisely the
states not separated by a future of length at most $k$.  Every strict round
increases the number of finite blocks, hence the refinement terminates.  Its
stable partition has the universal property by the usual fibre argument.

Accepted branch probabilities and target quotient classes are constant on the
stable fibres.  Summing the resolved accepted branches whose target lies in a
specified quotient class gives \eqref{eq:main-accepted-kernel}, and these sums
are exactly the probabilities of the corresponding post-accepted Reads.
Uniqueness follows because those Reads separate the quotient classes.
\end{proof}

\begin{remark}[Occurrence boundary]
The theorem reconstructs the dynamical state and kernel from the \emph{complete}
accepted cylinder.  It does not reconstruct that cylinder from an accepted
binary clock alone.  Indeed, an identity accepted kernel and a full-support
reset kernel can be coupled to the same state-independent acceptance bit while
producing different terminal field laws.  Thus occurrence of the complete
accepted field record is a genuine physical input row.
\end{remark}

\subsection{Renewal pressure and the common scalar action}

Let $(\Omega_X,p_X)$ be a faithful marked first-return law, let
$\tau_X>0$ be the protected physical duration, and let
$\mathscr G_X:E_X\to L^2(\Omega_X,p_X)$ be the reward synthesis of
\eqref{eq:main-reward-synthesis}.  Put
\[
 \mathcal L_X(q,r)=\mathbb E_X[e^{-g_{q,X}-r\tau_X}],
 \qquad
 \bar\tau_X=\mathbb E_X\tau_X,
 \qquad
 T_Xc=c\tau_X.
\]


For the score, differentiating the tilted law also gives
\begin{equation}
 S_Xv=-\mathscr G_Xv-\pi_X(v)\tau_X.
 \label{eq:supp-pressure-score}
\end{equation}
\begin{proof}[Proof of \Cref{thm:main-action-reconstruction}]
The transform is a finite sum of known probabilities times known protected
writer values.  At the origin, $\mathcal L_X(0,0)=1$ and
$\partial_r\mathcal L_X(0,0)=-\bar\tau_X<0$, so the analytic implicit-function
theorem gives the pressure root.  Differentiating
$\mathcal L_X(q,\mathcal P_X(q))=1$ in direction $v$ gives
\[
 0=-\mathbb E_Xg_{v,X}-\pi_X(v)\bar\tau_X,
\]
which proves the first identity.  Differentiating the normalized tilted law
gives the score formula and hence the synthesis reconstruction.
\end{proof}

\begin{corollary}[Pressure-jet reconstruction of mixed action blocks]
\label{cor:supp-pressure-jet}
With $\bar\tau_X(q)=\mathbb E_{X,q}\tau_X$,
$\delta_X=D\bar\tau_X(0)$, and
$m_{2,X}=\mathbb E_X\tau_X^2$,
\begin{equation}
 D^2\mathcal P_X(0)=\bar\tau_X^{-1}S_X^*S_X,\qquad
 \delta_X=T_X^*S_X,
 \label{eq:supp-pressure-jet-a}
\end{equation}
\begin{equation}
 \mathscr G_X^*\mathscr G_X
 =\bar\tau_XD^2\mathcal P_X(0)
 +\delta_X^*\pi_X+\pi_X^*\delta_X
 +m_{2,X}\pi_X^*\pi_X.
 \label{eq:supp-pressure-jet-b}
\end{equation}
Consequently one pressure surface reconstructs every marginal and mixed action
block represented on the same marked cycle record.
\end{corollary}

\begin{proof}
Differentiate the pressure normalization twice.  The first factors are the
centered scores $-S_Xu$ and $-S_Xv$, which gives the first identity.
Differentiating the tilted mean duration gives $\delta_X=T_X^*S_X$; expanding
$\mathscr G_X=-S_X-T_X\pi_X$ gives the second formula.
\end{proof}

Let $\mathscr V_X$ be the finite protected variation complex, with coboundaries
$d_0:C^0\to C^1$ and $d_1:C^1\to C^2$, $d_1d_0=0$, and let
$\alpha_X\in C^1(\mathscr V_X;\R)$ be the reconstructed physical action
increment.


\begin{proof}[Proof of \Cref{thm:main-common-action}]
If $\alpha_X=d_0\mathscr S_X$, telescoping gives zero face sums and zero periods.
Conversely, fix a base configuration in each connected component and define $\mathscr S_X(q)$ by summing
$\alpha_X$ along a path from the base to $q$.  Face closure makes the sum
invariant under elementary $2$-cell moves and the period conditions remove the
remaining homology class.  One anchor on each connected component removes its additive constant.  For the
estimate use the finite Hodge decomposition
$C^1=\Ran d_0\oplus\mathcal H^1\oplus\Ran d_1^*$ and the singular-value bound
$\|d_1\beta\|\ge\sigma_{\rm coex}\|\beta\|$ on the coexact range.
\end{proof}

\begin{remark}[Provenance is not stationarity]
Reconstructing a common action does not prove that the observed accepted
transition is generated by it.  The accepted/action comparison below is an
independent same-history test.  Conversely, fitting a cost to the observed
kernel after the fact would be tautological and has no stationarity content.
\end{remark}

\subsection{Fisher--Bregman stationarity and physical test lifts}

Let $\Theta$, $\Psi$, $G$, $\mathcal A$, and $\eta$ satisfy the convex-chart
and interiority hypotheses of \Cref{thm:main-action-stationarity}.  Define
$D_\Psi$, $\Phi_\theta$, $y_\theta^*$, $\Delta_\Psi$, and
$\overline\Delta_\Psi$ by
\eqref{eq:main-bregman}--\eqref{eq:main-stationarity-gap}.

\begin{theorem}[Quantitative accepted-action stationarity]
\label{thm:supp-action-stationarity}
Put
\begin{equation}
 \mu_\Phi=(\eta^{-1}-\lambda)m>0,\qquad
 \Lambda_\Phi=(\eta^{-1}+L)M,\qquad
 \ell=\max\{\lambda,L\}.
 \label{eq:supp-action-constants}
\end{equation}
Then
\begin{equation}
 \int D_\Psi(y,\theta)\,
       \nu(\dd\theta)K^{\rm acc}(\dd y\mid\theta)
       \le\eta\overline\Delta_\Psi,
 \label{eq:supp-bregman-bound}
\end{equation}
\begin{equation}
 \int\|y-\theta\|^2\,
       \nu(\dd\theta)K^{\rm acc}(\dd y\mid\theta)
       \le\frac{2\eta}{m}\overline\Delta_\Psi,
 \label{eq:supp-displacement-bound}
\end{equation}
and
\begin{equation}
 \int\mathfrak S(\theta)\,\nu(\dd\theta)
       \le C_{\Psi,\mathcal A,\eta}\overline\Delta_\Psi,
 \label{eq:supp-stationarity-bound}
\end{equation}
where one admissible constant is
\begin{equation}
 C_{\Psi,\mathcal A,\eta}
 =\frac{4\Lambda_\Phi^2}{m\mu_\Phi}
   +\frac{4M^2(1+\eta\ell)^2}{m^2\eta}.
 \label{eq:supp-stationarity-constant}
\end{equation}
\end{theorem}
\begin{proof}
State independence of $\vartheta>0$ and stationarity of $\nu$ imply
$\nu K^{\rm acc}=\nu$.  Since $y_\theta^*$ minimizes $\Phi_\theta$ and
$\Phi_\theta(\theta)=\mathcal A(\theta)$,
\[
 \mathcal A(y)-\mathcal A(\theta)+\eta^{-1}D_\Psi(y,\theta)
       \le\Delta_\Psi(\theta,y).
\]
The action difference cancels on averaging, proving
\eqref{eq:supp-bregman-bound}.  The lower Hessian bound gives
$D_\Psi(y,\theta)\ge m\|y-\theta\|^2/2$, hence
\eqref{eq:supp-displacement-bound}.

The Hessian of $\Phi_\theta$ lies between $\mu_\Phi I$ and
$\Lambda_\Phi I$ on the convex chart.  Interiority of $y_\theta^*$ gives
$\nabla\Phi_\theta(y_\theta^*)=0$.  Strong convexity and the upper Hessian
bound therefore imply
\[
 \|y-y_\theta^*\|^2\le\frac{2\Delta_\Psi(\theta,y)}{\mu_\Phi},
 \qquad
 \|\nabla\Phi_\theta(y)\|^2
       \le\frac{2\Lambda_\Phi^2}{\mu_\Phi}\Delta_\Psi(\theta,y).
\]
For
\[
 r_\Psi(\theta,y)=\eta\nabla\Phi_\theta(y)
       =\nabla\Psi(y)-\nabla\Psi(\theta)+\eta\nabla\mathcal A(y)
\]
we have
\[
 \eta\nabla\mathcal A(\theta)
 =r_\Psi-[\nabla\Psi(y)-\nabla\Psi(\theta)]
       +\eta[\nabla\mathcal A(\theta)-\nabla\mathcal A(y)].
\]
The last two terms together have norm at most
$M(1+\eta\ell)\|y-\theta\|$.  Squaring with
$\|a+b\|^2\le2\|a\|^2+2\|b\|^2$, using $G^{-1}\preceq m^{-1}I$, and
averaging yields \eqref{eq:supp-stationarity-bound} with
\eqref{eq:supp-stationarity-constant}.
\end{proof}

\begin{remark}[Boundary minima]
The first two inequalities do not require interiority, but the ordinary
gradient conclusion does.  For example, on $\Theta=[0,1]$ with
$\Psi(x)=x^2/2$ and $\mathcal A(x)=x$, the source $\theta=0$ has proximal
minimum $y_\theta^*=0$.  The stationary identity transition supported at zero
has zero gap although $\mathcal A'(0)=1$.  A boundary formulation must retain
the normal-cone or projected variational residual.  It cannot be substituted
for unrestricted determining metric variations without checking their
admissibility.
\end{remark}

\begin{proposition}[Action-gap control for nonstationary marginals]
\label{prop:supp-nonstationary-action-gap}
Let $\pi(\dd\theta,\dd y)$ be any coupling with source $\alpha$ and target
$\beta$ in the convex chart of \Cref{thm:supp-action-stationarity}, and
retain the interiority hypothesis.  Set
\[
 \bar\Delta_\pi=\mathbb E_\pi\Delta_\Psi,\qquad
 d_{\mathcal A}=\mathbb E_\alpha\mathcal A-
                         \mathbb E_\beta\mathcal A,
 \qquad
 C_1=\frac{4\Lambda_\Phi^2}{m\mu_\Phi},\quad
 C_2=\frac{4M^2(1+\eta\ell)^2}{m^2\eta}.
\]
Then
\begin{equation}
 \mathbb E_\pi D_\Psi(y,\theta)
 \le\eta(\bar\Delta_\pi+d_{\mathcal A}),\qquad
 \mathbb E_\alpha\mathfrak S
 \le(C_1+C_2)\bar\Delta_\pi+C_2d_{\mathcal A}.
 \label{eq:supp-nonstationary-gap}
\end{equation}
In particular the latter right side may be enlarged by replacing
$d_{\mathcal A}$ with its positive part.  For a represented physical test
with lift bound $B(k)$,
\begin{equation}
 \mathbb E_\alpha|D\mathcal A[v(k)]|^2
 \le B(k)^2\bigl[(C_1+C_2)\bar\Delta_\pi+C_2(d_{\mathcal A})_+\bigr].
 \label{eq:supp-nonstationary-test}
\end{equation}
No invariant law is required.  Equality of source and target marginals
recovers the stationary estimate.
\end{proposition}
\begin{proof}
The pointwise inequality in the stationary proof is
$\eta^{-1}D_\Psi(y,\theta)\le
\Delta_\Psi(\theta,y)+\mathcal A(\theta)-\mathcal A(y)$.
Averaging gives the first bound, and hence
$\mathbb E_\pi\|y-\theta\|^2\le
2\eta(\bar\Delta_\pi+d_{\mathcal A})/m$.
The pointwise gradient argument in
\Cref{thm:supp-action-stationarity} does not use stationarity.  Its
proximal-gradient term contributes $C_1\bar\Delta_\pi$, and its
displacement term contributes $C_2(\bar\Delta_\pi+d_{\mathcal A})$.
This proves the second bound.  Positive-metric duality gives the test
estimate.  A zero proximal gap at one chronological step need not force
zero gradient: it can instead describe genuine action descent, represented
by $d_{\mathcal A}>0$.
\end{proof}

\begin{proposition}[Cutoff-aware stationarity on the determining test core]
\label{prop:supp-physical-test-stationarity}
Let $k$ be a fixed physical test and let its represented lift satisfy
\eqref{eq:main-test-lift-budget}.  Then
\eqref{eq:main-physical-stationarity} holds.  If
\begin{equation}
 B_{X,K}(k)^2C_X\overline\Delta_{\Psi,X}\longrightarrow0
 \label{eq:supp-scaled-test-gap}
\end{equation}
for each test in a countable determining core, the corresponding first
variations converge to zero in mean square.  Under any common coupling there
is a deterministic subsequence of cutoffs on which all these first variations
converge to zero almost surely.
\end{proposition}
\begin{proof}
At each represented field value, duality for the positive $G_X$ pairing gives
\[
 |D\mathcal A_X[v]|^2
 \le\|D\mathcal A_X\|_{G_X^{-1}}^2\|v\|_{G_X}^2.
\]
Take the supported supremum of the second factor and apply
\Cref{thm:supp-action-stationarity}.  This proves the mean-square estimate.
Enumerate the determining core as $k_1,k_2,\ldots$.  Select increasing cutoffs
$X_j$ so that the right-hand side of
\eqref{eq:main-physical-stationarity} is at most $2^{-j}$ for each of
$k_1,\ldots,k_j$.  For a fixed $k_m$, the sum over $j\ge m$ of the expected
squared residuals is finite.  Tonelli's theorem implies that the corresponding
sum of squared residuals is finite almost surely, so the residual tends to
zero.  A countable intersection gives the assertion for the whole core.
No independence between cutoffs is used.  Uniform continuity of the limiting
variational functional on the declared test topology then extends zero from
the core to all determining tests.
\end{proof}

\begin{corollary}[Finite Gibbs/KL certificate]
\label{cor:supp-gibbs-gap}
For a finite reference row $p_\theta(y)>0$ and reconstructed same-history cost
$c_\theta(y)$, the row $q_\theta^*$ of \eqref{eq:main-gibbs-row} is the unique
entropy-proximal row and
\begin{equation}
 \Delta_{{\rm Gibbs},\theta}
       =\eta^{-1}D_{\rm KL}
       \bigl(K^{\rm acc}(\cdot\mid\theta)\Vert q_\theta^*\bigr)
 \label{eq:supp-gibbs-gap}
\end{equation}
is exactly its row variational gap.
\end{corollary}
\begin{proof}
For any probability row $q$,
\[
 \sum_yq(y)c_\theta(y)+\eta^{-1}D_{\rm KL}(q\Vert p_\theta)
 =\eta^{-1}D_{\rm KL}(q\Vert q_\theta^*)
       -\eta^{-1}\log\sum_zp_\theta(z)e^{-\eta c_\theta(z)}.
\]
Relative entropy is nonnegative and vanishes only at equality.  This proves a
statement about probability rows; it does not identify their sampled points
with interior critical points of the field action.
\end{proof}

\subsection{Derivative-free face linearization and complete memory shorting}

The first-order gravitational score is not assumed as a derivative at a zero
face.  It is obtained from renewal subdivision after all resolved cut-memory
coordinates have been retained.

\begin{theorem}[Derivative-free renewal face linearization]
\label{thm:supp-face-linearization}
Let $E$ be a finite-dimensional real normed space, $U\subset E$ a balanced
convex neighborhood of zero, and $A:U\to\R$ with $A(0)=0$.  If
\begin{equation}
 A(z)=2A(z/2)
 \label{eq:supp-face-divisibility}
\end{equation}
for every $z\in U$, $A$ is bounded near zero, and
\[
 |A(x+y)-A(x)-A(y)|\le H\|x\|\,\|y\|
\]
near zero, then $A$ is the restriction of a unique real linear functional.
More generally, if the renormalized scores
$L_m(z)=2^mA_m(2^{-m}z)$ have summable adjacent-cutoff defects and their
additivity defect tends to zero, then $L_m$ converges locally uniformly to one
continuous linear functional.
\end{theorem}

\begin{proof}
For the exact branch iterate \eqref{eq:supp-face-divisibility}:
\[
 A(x+y)-A(x)-A(y)
 =2^m[A(2^{-m}(x+y))-A(2^{-m}x)-A(2^{-m}y)].
\]
The right-hand side is bounded by $H2^{-m}\|x\|\|y\|$ and therefore vanishes as
$m\to\infty$.  Thus $A$ is additive; boundedness near zero gives continuity and
hence real homogeneity.  On the stable branch the adjacent-cutoff estimate makes
$L_m$ locally uniformly Cauchy; passing to the limit in the additivity defect
gives a continuous linear limit.
\end{proof}

\begin{theorem}[Complete resolved-memory subtraction]
\label{thm:supp-memory-short}
Let the resolved cut-memory state be
$\rho_{\rm cq}=\bigoplus_\omega p_\omega\rho_\omega$ with faithful blocks and
linear perturbation
$K_{\rm cq}(z)=\bigoplus_\omega K_\omega(z)$.  Put
\begin{equation}
 M_{\rm cq}(z)
 =\log\Tr\exp(\log\rho_{\rm cq}+K_{\rm cq}(z))-\bar m(z),
 \label{eq:supp-memory-logpartition}
\end{equation}
where $\bar m$ is the average linear slope.  Then
$M_{\rm cq}(0)=DM_{\rm cq}(0)=0$, its Hessian is the sum of a classical branch
covariance and a positive Kubo--Mori form, and
\begin{equation}
 A_{\rm sh}=A_{\rm coarse}-M_{\rm cq}
 \label{eq:supp-memory-short}
\end{equation}
is the canonical fully memory-shorted face score.  Thus unresolved memory is
removed at the complete log-partition level rather than only to second order.
\end{theorem}

\begin{proof}
The block exponential and trace split over the classical label.  Duhamel's
formula gives the first derivative in each faithful block and the Kubo--Mori
quadratic form as the second logarithmic derivative.  Differentiating the outer
finite log-sum-exp adds the covariance of the branch means.  Both terms are
positive.  Subtracting the complete log-partition therefore removes exactly the
resolved cut-memory contribution and leaves the displayed shorted score.
\end{proof}

\subsection{Local gravitational classification and variational consistency}

We use the densities \eqref{eq:main-phv-densities}, with the same
normalization and boundary convention.

\begin{proposition}[Classification in the prescribed incidence and degree class]
\label{prop:supp-restricted-classification}
Within \Cref{def:main-gravitational-class}, the real Lorentz-natural bulk
space is generated by $L_{\rm Holst}$, $L_{\rm Palatini}$ and $L_{\rm vol}$.
The assertion holds before torsion is set to zero.
\end{proposition}
\begin{proof}
Let $V=\mathbb R^{1,3}$. Antisymmetry of the exterior products identifies
the curvature-dependent coefficient with an invariant bilinear pairing
between two copies of $\Lambda^2 V$. The Lorentz metric identifies one copy
with its dual. On real bivectors the internal Hodge operator satisfies
$\star^2=-I$; after complexification its $+i$ and $-i$ eigenspaces are the
two nonisomorphic irreducible bivector representations of the connected
Lorentz group. Schur's lemma makes a commuting complex operator scalar on
each eigenspace. Reality makes the two scalars complex conjugate, leaving
precisely $uI+v\star$, $u,v\in\mathbb R$. These two pairings give
$e^I\wedge e^J\wedge R_{IJ}$ and
$\tfrac12\epsilon_{IJKL}e^I\wedge e^J\wedge R^{KL}$.
Scalar multiplicity excludes another coefficient operator on an internal
factor. The alternating part of a four-coframe coefficient is an invariant
in $\Lambda^4 V^*$, a one-dimensional space generated by the oriented
volume tensor. Symmetric metric contractions do not create an additional
alternating four-coframe invariant. This proves the asserted span. Exact
terms are dealt with by the boundary convention in the definition, not by
deleting physical boundary variations.
\end{proof}

\begin{remark}[Topological terms and torsion dependence]
\label{rem:supp-topological-scope}
Euler and Pontryagin densities are quadratic in curvature and do not belong
to this bidegree class. Their local bulk variations are exact, but their
boundary and global information need not vanish. The Nieh--Yan identity is
\begin{equation}
 \dd(e^I\wedge T_I)
   =T^I\wedge T_I-e^I\wedge e^J\wedge R_{IJ}.
 \label{eq:supp-nieh-yan}
\end{equation}
It follows from $De^I=T^I$ and $DT^I=R^I{}_J\wedge e^J$, with contraction
turning the exterior covariant derivative into $\dd$. Thus enlarging the
class to include $T^I\wedge T_I$ introduces a relation to the Holst term
modulo an explicitly retained boundary density; it is not justified by the
bivector commutant alone. More general inverse-coframe or torsion invariants
require a different classification. The distinction between bulk equations,
variational principles and boundary charges is essential in first-order
gravity \cite{CorichiReview}. Neither the Palatini coefficient nor its ratio
to the volume coefficient is fixed by Lorentz covariance.
\end{remark}

\begin{lemma}[Transfer of a classified variational representative]
\label{lem:supp-variation-transfer}
Suppose the classified curvature and volume variations converge on each
physical determining test, and \eqref{eq:main-first-variation-reduction}
holds on those tests. Then the complete action variation has the same limit.
If that complete variation tends to zero, the limiting classified variation
vanishes. The assertion also holds on one common realized subsequence of
stochastic records satisfying these hypotheses.
\end{lemma}
\begin{proof}
Use \eqref{eq:main-action-variation-remainder}, the convergence of the three
calibrated coefficients, and the vanishing remainder. All operations take
place on the same test and the same sequence. In the stochastic case the
argument is applied after the common almost-sure subsequence has been
selected; no independent subsequences are multiplied together.
\end{proof}


\subsection{Connection stationarity and torsion}

Let $B^{IJ}=e^I\wedge e^J$, $T^I=De^I$, and
$P_{\alpha,\beta}=\alpha I+\beta\star$, with the same
spacetime-constant coefficients as in \eqref{eq:main-phv-coefficients}.

\begin{theorem}[Connection stationarity and the torsion alternative]
\label{thm:supp-palatini-torsion}
If $(\alpha,\beta)\ne(0,0)$, then
\begin{equation}
 P_{\alpha,\beta}^{-1}
 =\frac{\alpha I-\beta\star}{\alpha^2+\beta^2}.
 \label{eq:supp-palatini-inverse}
\end{equation}
For the Palatini--Holst action the spinless connection Euler equation is
$D(P_{\alpha,\beta}B)=0$.  For a nondegenerate coframe, the Cartan map
\[
 \mathcal K_e(T)^{IJ}=T^I\wedge e^J-e^I\wedge T^J
\]
is injective.  Consequently connection stationarity implies $T=0$; at finite
cutoff, a uniform lower singular-value bound on $\mathcal K_{e,X}$ controls the
torsion norm by the connection Euler residual.
\end{theorem}

\begin{proof}
The inverse formula follows from $\star^2=-I$ on Lorentzian bivectors.  Varying
the connection and integrating by parts gives the Euler source.  To prove
injectivity, expand $T^I=\frac12T^I{}_{KL}e^K\wedge e^L$ in a coframe basis.
Choosing a fourth index in
$T^I\wedge e^J-e^I\wedge T^J=0$ first eliminates all components with
$K,L\ne I$; the remaining coefficients satisfy pairwise sign relations that,
in four dimensions, force them all to vanish.  Inverting the two maps gives the
finite estimate.
\end{proof}


\section{Noncollapse, curvature compactness and the Einstein limit}
\label{supp:continuum}

This appendix proves the estimates and limiting argument of
Section~\ref{sec:einstein-limit}. It uses the finite geometry and action
from Sections~\ref{sec:adm-main} and~\ref{sec:einstein-dynamics}; temporal
compactness, curvature propagation and variational consistency remain
separate hypotheses on one common physical sequence. The alternative
connection routes are kept in distinct subsections.
\subsection{Cut estimates and spatial compactness}

For a nonempty half-volume cut $A$, write $A^c=V_X\setminus A$ and define
the mean-zero demand
\begin{equation}
 b_{A,X}(v)=\mu_X(A)^{2/3}
 \left[\frac{m_{v,X}\one_A(v)}{\mu_X(A)}
       -\frac{m_{v,X}\one_{A^c}(v)}{\mu_X(A^c)}\right].
 \label{eq:main-cut-demand}
\end{equation}
This fixes the normalization of the dual cut witness below.
\begin{equation}
 C_{\rm S}=\frac{128D_*}{(I_*^{\rm sp})^2},
 \qquad
 C_{\rm P}=V_*^{2/3}C_{\rm S}.
 \label{eq:main-sobolev-constants}
\end{equation}

Let a finite physical endpoint graph have vertex masses $m_{v,X}$, measure
$\mu_X$, symmetric conductances $c_{uv,X}$, and efficient edge capacities
$u_{e,X}$.  We use the two cut constants $I_X^{\rm eff}$ and
$I_X^{\rm sp}$ of \eqref{eq:main-cut-constant}; on the positive spatial branch
they are related by the finite capacity--conductance comparison
\eqref{eq:main-capacity-conductance}.

\begin{definition}[Spatial Dirichlet normalization]
\label{def:supp-spatial-dirichlet}
On the finite weighted graph $G_X=(V_X,E_X)$ set
\begin{equation}
 \mathcal E_X^{\rm sp}(f)
 =\sum_{\{u,v\}\in E_X}c_{uv,X}|f(u)-f(v)|^2,
 \qquad
 \langle f,\Delta_X^{\rm sp}f\rangle_{L^2(\mu_X)}
 =\mathcal E_X^{\rm sp}(f).
 \label{eq:supp-spatial-dirichlet}
\end{equation}
All scalar Sobolev estimates below are first stated for real-valued functions;
the real weighted Laplacian has the same spectrum after complexification.  Let
$P_{0,X}$ be the orthogonal projection onto the mean-zero subspace and put
$T_X(t)=e^{-t\Delta_X^{\rm sp}}P_{0,X}$.  Since the conductances are symmetric
and nonnegative, $e^{-t\Delta_X^{\rm sp}}$ is a finite Markov semigroup and is
contractive on every $L^q(\mu_X)$, $1\le q\le\infty$.
\end{definition}

\begin{lemma}[Cut inequality implies the half-volume $L^1$ Sobolev bound]
\label{lem:supp-half-volume-sobolev}
Assume
\begin{equation}
 h_Xc_X(\partial A)\ge I_*^{\rm sp}\,\mu_X(A)^{2/3}
 \label{eq:supp-cut-isoperimetry}
\end{equation}
whenever $0<\mu_X(A)\le\mu_X(V_X)/2$.  If $g\ge0$ and
$\mu_X(\supp g)\le\mu_X(V_X)/2$, then
\begin{equation}
 \|g\|_{L^{3/2}(\mu_X)}
 \le \frac{h_X}{I_*^{\rm sp}}
 \sum_{\{u,v\}\in E_X}c_{uv,X}|g(u)-g(v)|.
 \label{eq:supp-l1-sobolev}
\end{equation}
\end{lemma}

\begin{proof}
For $t\ge0$ put $A_t=\{v:g(v)>t\}$.  Since $A_t\subseteq\supp g$, every
nonempty $A_t$ lies on the half-volume side of
\eqref{eq:supp-cut-isoperimetry}.  The layer-cake identity gives
$g=\int_0^\infty\one_{A_t}\,\dd t$.  Minkowski's integral inequality in
$L^{3/2}(\mu_X)$ therefore yields
\begin{equation}
 \|g\|_{3/2}
 \le\int_0^\infty\|\one_{A_t}\|_{3/2}\,\dd t
 =\int_0^\infty\mu_X(A_t)^{2/3}\,\dd t.
 \label{eq:supp-layer-minkowski}
\end{equation}
Applying the cut inequality to every level set and then using the finite
coarea identity gives
\begin{align}
 \|g\|_{3/2}
 &\le \frac{h_X}{I_*^{\rm sp}}
       \int_0^\infty c_X(\partial A_t)\,\dd t \\
 &=\frac{h_X}{I_*^{\rm sp}}
   \sum_{\{u,v\}\in E_X}c_{uv,X}|g(u)-g(v)|,
\end{align}
which is \eqref{eq:supp-l1-sobolev}.
\end{proof}

\begin{lemma}[Discrete $L^6$ Sobolev and Poincar\'e inequalities]
\label{lem:supp-l6-poincare}
Assume \eqref{eq:supp-cut-isoperimetry},
$\mu_X(V_X)\le V_*$, and
\begin{equation}
 \sum_uc_{uv,X}\le D_*h_X^{-2}m_{v,X}
 \qquad(v\in V_X).
 \label{eq:supp-weighted-degree}
\end{equation}
Then every mean-zero $f$ satisfies
\begin{equation}
 \|f\|_6^2\le C_{\rm S}\mathcal E_X^{\rm sp}(f),
 \qquad
 C_{\rm S}=\frac{128D_*}{(I_*^{\rm sp})^2},
 \label{eq:supp-sobolev}
\end{equation}
and
\begin{equation}
 \|f\|_2^2\le C_{\rm P}\mathcal E_X^{\rm sp}(f),
 \qquad C_{\rm P}=V_*^{2/3}C_{\rm S}.
 \label{eq:supp-poincare}
\end{equation}
In particular
$\lambda_1(\Delta_X^{\rm sp}|_{\one^\perp})\ge C_{\rm P}^{-1}$.
\end{lemma}

\begin{proof}
First suppose $f\ge0$ and $\mu_X(\supp f)\le\mu_X(V_X)/2$.  Apply
\Cref{lem:supp-half-volume-sobolev} to $g=f^4$.  Since
$\|f^4\|_{3/2}=\|f\|_6^4$ and
$|a^4-b^4|\le4(a^3+b^3)|a-b|$ for $a,b\ge0$,
\begin{align}
 \|f\|_6^4
 &\le \frac{4h_X}{I_*^{\rm sp}}
 \sum_{\{u,v\}}c_{uv,X}(f(u)^3+f(v)^3)|f(u)-f(v)| \\
 &\le \frac{4h_X}{I_*^{\rm sp}}
 \left[\sum_{\{u,v\}}c_{uv,X}(f(u)^3+f(v)^3)^2\right]^{1/2}
 \mathcal E_X^{\rm sp}(f)^{1/2}.
 \label{eq:supp-sobolev-cs}
\end{align}
Using $(a^3+b^3)^2\le2(a^6+b^6)$ and
\eqref{eq:supp-weighted-degree},
\begin{align}
 \sum_{\{u,v\}}c_{uv,X}(f(u)^3+f(v)^3)^2
 &\le2\sum_v f(v)^6\sum_uc_{uv,X} \\
 &\le2D_*h_X^{-2}\|f\|_6^6.
 \label{eq:supp-degree-six}
\end{align}
Substitution into \eqref{eq:supp-sobolev-cs} and division by
$\|f\|_6^3$ gives
\begin{equation}
 \|f\|_6^2
 \le \frac{32D_*}{(I_*^{\rm sp})^2}\mathcal E_X^{\rm sp}(f)
 \label{eq:supp-half-sobolev-six}
\end{equation}
for every half-volume-supported nonnegative $f$.

For a general real $f$, choose a weighted median $m$ so that
$\mu_X(f>m)$ and $\mu_X(f<m)$ are each at most half the total volume.  Put
$f_+=(f-m)_+$ and $f_-=(m-f)_+$.  Truncation is $1$-Lipschitz and, edgewise,
\begin{equation}
 |f_+(u)-f_+(v)|^2+|f_-(u)-f_-(v)|^2
 \le |f(u)-f(v)|^2,
 \label{eq:supp-truncation-energy}
\end{equation}
so
$\mathcal E(f_+)+\mathcal E(f_-)\le\mathcal E(f)$.  Since the two functions
have half-volume support, \eqref{eq:supp-half-sobolev-six} gives
\begin{align}
 \|f-m\|_6^2
 &=\bigl(\|f_+\|_6^6+\|f_-\|_6^6\bigr)^{1/3} \\
 &\le \|f_+\|_6^2+\|f_-\|_6^2
 \le \frac{32D_*}{(I_*^{\rm sp})^2}\mathcal E(f).
 \label{eq:supp-median-six}
\end{align}
If $f$ has mean zero, then
\begin{equation}
 |m|\,\mu_X(V_X)^{1/6}
 \le \|f-m\|_6,
 \label{eq:supp-median-mean}
\end{equation}
because
$|m|=|\mu_X(V_X)^{-1}\int(f-m)\,\dd\mu_X|$ and H\"older gives
$\|f-m\|_1\le\mu_X(V_X)^{5/6}\|f-m\|_6$.  Therefore
$\|f\|_6\le2\|f-m\|_6$, which yields \eqref{eq:supp-sobolev} with the
stated factor four.  Finally,
\begin{equation}
 \|f\|_2^2\le\mu_X(V_X)^{2/3}\|f\|_6^2
 \le V_*^{2/3}C_{\rm S}\mathcal E(f),
\end{equation}
which is \eqref{eq:supp-poincare}.  Rayleigh--Ritz gives the spectral-gap
bound.
\end{proof}

\begin{lemma}[Nash inequality, ultracontractivity, and spectral counting]
\label{lem:supp-nash-counting}
Under the hypotheses of \Cref{lem:supp-l6-poincare}, every mean-zero $f$ obeys
\begin{equation}
 \|f\|_2^{10/3}
 \le C_{\rm S}\mathcal E_X^{\rm sp}(f)\|f\|_1^{4/3}.
 \label{eq:supp-nash}
\end{equation}
Moreover, for every $t>0$,
\begin{equation}
 \|T_X(t)\|_{L^1\to L^\infty}
 \le4\left(\frac{3C_{\rm S}}{2t}\right)^{3/2}.
 \label{eq:supp-ultracontractive}
\end{equation}
Consequently, if
$N_X(R)=\rank\one_{[0,R]}(\Delta_X^{\rm sp})$, then for $R>0$
\begin{equation}
 N_X(R)\le1+4e^{3/2}V_*(C_{\rm S}R)^{3/2}.
 \label{eq:supp-global-weyl-count}
\end{equation}
\end{lemma}

\begin{proof}
Interpolation between $L^1$ and $L^6$ gives
$\|f\|_2\le\|f\|_1^{2/5}\|f\|_6^{3/5}$.  Raising this inequality to the
power $10/3$ and using \eqref{eq:supp-sobolev} yields
\eqref{eq:supp-nash}.

Let $F_t=T_X(t)f$ and $u(t)=\|F_t\|_2^2$.  Since $F_t$ has mean zero and the
finite Markov semigroup is $L^1$-contractive,
\begin{equation}
 \|F_t\|_1
 \le\|P_{0,X}f\|_1
 \le2\|f\|_1.
 \label{eq:supp-p0-l1}
\end{equation}
Self-adjointness and finite dimensionality permit differentiation under the
spectral decomposition:
\begin{equation}
 u'(t)=-2\mathcal E_X^{\rm sp}(F_t).
 \label{eq:supp-energy-decay}
\end{equation}
Combining \eqref{eq:supp-nash}, \eqref{eq:supp-p0-l1}, and
\eqref{eq:supp-energy-decay} gives
\begin{equation}
 u'(t)
 \le-\frac{2}{C_{\rm S}2^{4/3}\|f\|_1^{4/3}}u(t)^{5/3}.
 \label{eq:supp-nash-ode}
\end{equation}
Integrating the differential inequality for $u^{-2/3}$ yields
\begin{equation}
 \|T_X(t)f\|_2^2
 \le \frac{3^{3/2}}{2}C_{\rm S}^{3/2}t^{-3/2}\|f\|_1^2.
 \label{eq:supp-l1-l2}
\end{equation}
Because $T_X(t/2)$ is self-adjoint and
$T_X(t)=T_X(t/2)T_X(t/2)$ on the mean-zero subspace, duality and the semigroup
property give
\begin{align}
 \|T_X(t)\|_{1\to\infty}
 &\le\|T_X(t/2)\|_{1\to2}\|T_X(t/2)\|_{2\to\infty} \\
 &=\|T_X(t/2)\|_{1\to2}^2
 \le4\left(\frac{3C_{\rm S}}{2t}\right)^{3/2},
\end{align}
which is \eqref{eq:supp-ultracontractive}.  This is the standard
Sobolev--Nash ultracontractivity route in finite weighted form; compare
\cite[Chs.~10--11]{Chung1997} and \cite{SaloffCoste2002}.

Let $K_t^0(v,w)$ be the integral kernel of $T_X(t)$ with respect to $\mu_X$.
The spectral expansion gives
$K_t^0(v,v)=\sum_{j\ge1}e^{-t\lambda_j}|\varphi_j(v)|^2\ge0$, and therefore
\begin{align}
 \Tr(T_X(t))
 &=\sum_v m_{v,X}K_t^0(v,v) \\
 &\le \mu_X(V_X)\|T_X(t)\|_{1\to\infty} \\
 &\le4V_*\left(\frac{3C_{\rm S}}{2t}\right)^{3/2}.
 \label{eq:supp-heat-trace}
\end{align}
Since the positive cut margin implies connectedness, the zero eigenvalue is
simple.  Hence
\begin{equation}
 (N_X(R)-1)e^{-tR}
 \le\sum_{0<\lambda_j\le R}e^{-t\lambda_j}
 \le\Tr(T_X(t)).
 \label{eq:supp-counting-trace}
\end{equation}
Choosing $t=3/(2R)$ gives \eqref{eq:supp-global-weyl-count}.
\end{proof}

\begin{theorem}[Operational Sobolev, Poincar\'e, and compact-screen bounds]
\label{thm:supp-spatial-screen}
Assume uniformly that $\mu_X(V_X)\le V_*$,
\eqref{eq:supp-weighted-degree} holds, the capacity--conductance comparison
\eqref{eq:main-capacity-conductance} holds, and
$I_X^{\rm eff}\ge I_*^{\rm eff}>0$.  Then
$I_X^{\rm sp}\ge I_*^{\rm eff}/a_+=I_*^{\rm sp}>0$ and the estimates
\eqref{eq:supp-sobolev}, \eqref{eq:supp-poincare}, and
\eqref{eq:supp-global-weyl-count} hold.  In addition,
\begin{equation}
 \|(I-P_{X,R})f\|_2^2
 \le R^{-1}\mathcal E_X^{\rm sp}(f),
 \qquad P_{X,R}=\one_{[0,R]}(\Delta_X^{\rm sp}).
 \label{eq:supp-screen-count-tail}
\end{equation}
Thus $P_{X,R}$ is a uniformly exhausted finite-rank compact-screen family.
A sequence of positive cut constants tending to zero gives the normalized
cut/potential obstruction \eqref{eq:main-neck-obstruction}; a zero-capacity
cut is retained directly as disconnected support.
\end{theorem}

\begin{proof}
The comparison \eqref{eq:main-cut-comparison} converts the efficient-capacity
margin into \eqref{eq:supp-cut-isoperimetry}.  The Sobolev and Poincar\'e
estimates are \Cref{lem:supp-l6-poincare}; the counting law is
\Cref{lem:supp-nash-counting}.  By the spectral theorem,
$\Delta_X^{\rm sp}\succeq R(I-P_{X,R})$, so
\eqref{eq:supp-screen-count-tail} follows immediately.  For a positive minimizing cut $A$, set
$\phi=\one_A/u_X(\partial A)$.  Then
$\sum_eu_{e,X}|\dd\phi(e)|=1$, and \eqref{eq:main-cut-demand} gives
$|\sum_vb_{A,X}(v)\phi(v)|=\mu_X(A)^{2/3}/u_X(\partial A)
=(I_X^{\rm eff})^{-1}$.  This is an explicit feasible dual witness.
A zero-capacity cut is recorded directly as disconnected support.
\end{proof}

\begin{proposition}[Compact positive-screen upgrade]
\label{prop:supp-compact-upgrade}
Let $Y_X\rightharpoonup Y$ in a Hilbert space and let positive operators
$\mathbb A_X$ satisfy
$mI\preceq\mathbb A_X\preceq MI$ and
$\mathbb A_X^{-1/2}\to\mathbb A^{-1/2}$ strongly.  Put
$Z_X=\mathbb A_X^{1/2}Y_X$ and suppose $Z_X\rightharpoonup Z$.  If for every
$R$ there are finite-rank screens $S_{X,R}\to S_R$ in operator norm, with
$S_R$ compact, and
\[
 \lim_{R\to\infty}\sup_X\|(I-S_{X,R})Z_X\|=0,
 \qquad
 \lim_{R\to\infty}\|(I-S_R)Z\|=0,
\]
then $Y_X\to Y$ strongly.  Consequently every bounded quadratic insertion in
$Y_X$ converges in $L^1$ on the corresponding compact spacetime cylinder.
\end{proposition}

\begin{proof}
For fixed $R$, compactness of $S_R$ converts weak convergence into
$S_RZ_X\to S_RZ$ strongly, and operator-norm convergence gives
$S_{X,R}Z_X-S_RZ_X\to0$.  Decompose
\[
 Z_X-Z=(I-S_{X,R})Z_X+(S_{X,R}Z_X-S_RZ)+(S_R-I)Z.
\]
Taking the limsup and then $R\to\infty$ gives $Z_X\to Z$ strongly.  Uniform
coercivity and strong convergence of the inverse square roots then give
$Y_X\to Y$.  The quadratic statement follows from
$\|B(Y_X,Y_X)-B(Y,Y)\|_1\le C\|Y_X-Y\|_2(\|Y_X\|_2+\|Y\|_2)$.
\end{proof}

\subsection{Strong coframe interpolation}

\begin{lemma}[Strong interpolation for the reconstructed coframes]
\label{lem:supp-coframe-interpolation}
Let $K$ be a compact spacetime cylinder with finite reconstructed volume measure.
Suppose the common-cylinder interpolants of the finite coframes satisfy
\begin{equation}
 e_X\to e\quad\hbox{strongly in }L^2(K),
 \qquad
 \sup_X\|e_X\|_{L^6(K)}\le M_6<\infty.
 \label{eq:supp-coframe-l2-l6}
\end{equation}
Then $e\in L^6(K)$ and, for every $2\le q<6$,
\begin{equation}
 e_X\to e\quad\hbox{strongly in }L^q(K).
 \label{eq:supp-strong-lq}
\end{equation}
More precisely, if
$\theta_q=3/q-1/2\in(0,1]$, then
\begin{equation}
 \|e_X-e\|_{L^q}
 \le \|e_X-e\|_{L^2}^{\theta_q}
      \|e_X-e\|_{L^6}^{1-\theta_q}.
 \label{eq:supp-interpolation-explicit}
\end{equation}
If $p>3/2$ and $p'=p/(p-1)$, then $2p'<6$ and
\begin{equation}
 e_X\wedge e_X\longrightarrow e\wedge e
 \quad\hbox{strongly in }L^{p'}(K).
 \label{eq:supp-bivector-strong}
\end{equation}
\end{lemma}

\begin{proof}
The uniform $L^6$ bound gives a weakly convergent subsequence in $L^6(K)$.
Because $K$ has finite measure, weak $L^6$ convergence implies distributional
convergence, while the strong $L^2$ convergence has the unique distributional
limit $e$.  Hence every weak $L^6$ cluster point is $e$, so $e\in L^6(K)$ and
$\|e\|_6\le M_6$.  The difference $e_X-e$ is therefore uniformly bounded in
$L^6$.  The Riesz--Thorin/H\"older interpolation identity
$1/q=\theta_q/2+(1-\theta_q)/6$ gives
\eqref{eq:supp-interpolation-explicit}; its right-hand side tends to zero for
$q<6$, proving \eqref{eq:supp-strong-lq}.

For $p>3/2$, $p'<3$ and hence $q=2p'<6$.  The exterior product is a bounded
bilinear map on the finite-dimensional coframe fibres, so pointwise
$|a\wedge b|\le C|a||b|$.  Therefore H\"older's inequality gives
\begin{align}
 \|e_X\wedge e_X-e\wedge e\|_{L^{p'}}
 &\le C\|(e_X-e)\wedge e_X+e\wedge(e_X-e)\|_{L^{p'}} \\
 &\le C(\|e_X\|_{L^{2p'}}+\|e\|_{L^{2p'}})
       \|e_X-e\|_{L^{2p'}}\longrightarrow0,
\end{align}
which proves \eqref{eq:supp-bivector-strong}.
\end{proof}

\subsection{Curvature-energy propagation}

The differential estimate underlying
\Cref{thm:main-curvature-propagation} is
\begin{equation}
 \frac12\frac{\dd}{\dd t}z_X^2
 \le a_Xz_X^2+z_X\|f_X\|_{M_X},
 \label{eq:supp-curvature-energy-differential}
\end{equation}

\begin{proof}[Proof of \Cref{thm:main-curvature-propagation}]
Differentiate $z_X^2=y_X^TM_Xy_X$ and substitute
\eqref{eq:main-curvature-propagation-writer}.  The principal term cancels
exactly because $M_XK_X+K_X^TM_X=0$.  By
\eqref{eq:main-curvature-growth}, the remaining quadratic term is at most
$2a_Xz_X^2$, while Cauchy--Schwarz in the positive $M_X$-Gram bounds the source
term by $2z_X\|f_X\|_{M_X}$.  Applying the integrating-factor argument to
$(z_X^2+\varepsilon^2)^{1/2}$ and sending $\varepsilon\downarrow0$ gives
\eqref{eq:main-curvature-energy-propagation}.  The interpolation estimate in
\eqref{eq:main-curvature-incidence}, integration against the physical lapse and
slice measure, and the uniform budgets give
\eqref{eq:main-derived-l2-curvature}.
\end{proof}

\subsection{Fully discrete propagation with physical transfers}

The preceding estimate has a fully discrete counterpart, so a differentiable
curvature path is not needed merely to propagate its size.  At one spatial
cutoff, take a partition of the selected time interval with lengths $s_j>0$.
Let $Y_j=M_j^{1/2}y_j$ be the normalized actual curvature records and
let $T_j$ be their recorded physical transfer in these coordinates.  Suppose
\begin{equation}
 Y_{j+1}-T_jY_j
 =s_j(K_j+L_j)\frac{Y_{j+1}+T_jY_j}{2}+\varepsilon_j,
 \label{eq:main-discrete-curvature-update}
\end{equation}
where $K_j^T=-K_j$, $(L_j+L_j^T)/2\preceq a_jI$,
$\|T_j\|_{\op}\le e^{d_j}$, and $a_j,d_j\ge0$.  The residual
$\varepsilon_j$ includes the complete source, solver and coordinate-transfer
errors of the actual step.  In particular, nonlinear changes of field
coordinates are not assumed to commute with midpoint discretization.

\begin{theorem}[Fully discrete initial-energy curvature certificate]
\label{thm:main-discrete-curvature}
Assume $s_ja_j/2\le b<1$, and set
\begin{equation}
 D_X=\sum_jd_j,\quad A_X^{\rm d}=\sum_js_ja_j,\quad
 \mathcal D_X=\sum_j\frac{\|\varepsilon_j\|^2}{s_j},\quad
 T_X=\sum_js_j.
 \label{eq:main-discrete-curvature-budgets}
\end{equation}
Each recorded linear step is uniquely solvable, and
\begin{equation}
 \max_j\|Y_j\|
 \le e^{D_X+A_X^{\rm d}/(1-b)}
       \left(\|Y_0\|+\frac{\sqrt{T_X\mathcal D_X}}{1-b}\right).
 \label{eq:main-discrete-curvature-bound}
\end{equation}
There is no restriction involving $s_j\|K_j\|_{\op}$.  With a uniformly
stable interval interpolation of these actual curvature records, a uniform
lapse bound, and a bounded reconstruction error as in
\eqref{eq:main-curvature-incidence}, uniform bounds on the quantities in
\eqref{eq:main-discrete-curvature-bound} imply
$\sup_X\|R_X\|_{L^2(K)}<\infty$.
\end{theorem}

The statement concerns the represented linear step with its coefficients
fixed.  For an implicit nonlinear evolution, existence of that evolution is
not inferred from this matrix inverse; its actual coefficient and residual
records must satisfy \eqref{eq:main-discrete-curvature-update}.
The proof in this appendix uses the skew cancellation before
taking norms and retains the physical transfer between successive Grams.


\begin{proof}
Write $A_j=K_j+L_j$.  For any real vector $v$,
\[
 \langle v,(I-s_jA_j/2)v\rangle
       \ge(1-s_ja_j/2)\|v\|^2.
\]
Cauchy--Schwarz gives
$\|(I-s_jA_j/2)v\|\ge(1-s_ja_j/2)\|v\|$.  Thus the finite matrix is
invertible and its inverse norm is at most $(1-s_ja_j/2)^{-1}$.

Let $v=T_jY_j$ and first let $w$ solve the step with zero residual.  Pair
$w-v=s_jA_j(w+v)/2$ with $(w+v)/2$ to obtain
\[
 \|w\|^2-\|v\|^2
 \le\frac{s_ja_j}{2}(\|w\|+\|v\|)^2.
\]
Dividing by $\|w\|+\|v\|$ when this number is nonzero, and treating the
zero case directly, gives
\[
 \|w\|\le\frac{1+s_ja_j/2}{1-s_ja_j/2}\|v\|.
\]
The inhomogeneous solution adds
$(I-s_jA_j/2)^{-1}\varepsilon_j$.  Consequently
\[
 \|Y_{j+1}\|\le e^{d_j}
       \frac{1+s_ja_j/2}{1-s_ja_j/2}\|Y_j\|
       +\frac{\|\varepsilon_j\|}{1-b}.
\]
For $0\le x\le b$,
$\log((1+x)/(1-x))\le2x/(1-b)$.  Every product of step factors is
therefore bounded by $\exp(D_X+A_X^{\rm d}/(1-b))$.
Iteration and weighted Cauchy--Schwarz,
\[
 \sum_j\|\varepsilon_j\|
       \le\sqrt{T_X\sum_j\|\varepsilon_j\|^2/s_j},
\]
prove \eqref{eq:main-discrete-curvature-bound}.  The size of the skew matrix
$K_j$ has entered neither estimate.

For raw curvature records $y_j$ with Gram $M_j$, the transfer is
$T_j=M_{j+1}^{1/2}\mathcal T_jM_j^{-1/2}$, where $\mathcal T_j$ is the
actual recorded frame/transport map.  Its norm excess is retained in $d_j$.
Stable interval interpolation bounds the reconstructed slice norm by a fixed
multiple of the largest endpoint norm.  Integrating with the physical lapse
and adding the bounded reconstruction error yields, for an appropriate
interpolation constant $C_I$,
\[
 \|R_X\|_{L^2(K)}
 \le C_I\sqrt{N_*T_X}\,e^{D_X+A_X^{\rm d}/(1-b)}
       \left(\|Y_0\|+\frac{\sqrt{T_X\mathcal D_X}}{1-b}\right)+C_\zeta.
\]
All interface and incoming boundary contributions must already be included
in the actual transfer, residual or reconstruction comparison.  This proof
does not assert that arbitrary nonlinear coordinate changes preserve the
midpoint recurrence: any such discrepancy is part of $\varepsilon_j$.
\end{proof}

\subsection{Differentiated action-current control}

The forcing budget can itself be checked against the reconstructed action,
provided the actual differentiated current has the required incidence.  Write
$\mathfrak a_X=D\mathcal A_X$ for its action-gradient covector and suppose
\begin{equation}
 f_X=\mathsf T_X\mathfrak a_X+r_X^{\rm src}.
 \label{eq:main-current-incidence}
\end{equation}
The transfer $\mathsf T_X$ includes the physical metric-test lift, the nonzero
Palatini normalization, trace reversal, the differentiated Bianchi current,
and its frame comparison.  Torsion, Cartan, boundary and discretization
contributions not in this map remain in $r_X^{\rm src}$.

\begin{proposition}[Sharp differentiated action-current certificate]
\label{prop:main-action-current}
Let $H_X\succ0$ be the actual source Gram, transferred to the curvature-energy
space.  The least squared amplification from the action dual norm is
\begin{equation}
 \Gamma_X
   =\|H_X^{1/2}\mathsf T_XG_X^{1/2}\|_{\op}^2,\qquad
 \|\mathsf T_X\mathfrak a\|_{H_X}^2
       \le\Gamma_X\|\mathfrak a\|_{G_X^{-1}}^2.
 \label{eq:main-sharp-current-transfer}
\end{equation}
Consequently a pointwise certificate
$\|\mathfrak a_X(t)\|_{G_X^{-1}}^2\le C_X(t)\Delta_X(t)$ gives
\begin{equation}
 \int_0^{T_*}\|f_X\|_{H_X}\,\dd t
 \le\sqrt{2T_*\int_0^{T_*}
       \bigl(\Gamma_XC_X\Delta_X+
                         \|r_X^{\rm src}\|_{H_X}^2\bigr)\,\dd t}.
 \label{eq:main-action-source-budget}
\end{equation}
For the normalized discrete source
$\varepsilon_j/s_j=\mathsf T_j\mathfrak a_j+r_j^{\rm src}$, the
corresponding squared-energy estimate is
$\mathcal D_X\le2\sum_js_j(\Gamma_jC_j\Delta_j+
\|r_j^{\rm src}\|_{H_j}^2)$.
\end{proposition}

This replaces a separately supplied source-size bound by an explicitly scaled
action-side bound on the verified transfer branch.  It does not turn a small
undifferentiated Euler residual into a small Bianchi current: a single centered
grid derivative can have $\Gamma_X=h_X^{-2}$.  The sharpness example and the
expected-energy version of the certificate are proved in
this appendix.  An estimate only in expectation yields a bound
in probability; almost-sure uniform curvature control additionally uses the
pathwise or summable source conditions stated there.

\paragraph{Differentiated gravitational source.}
The source norm is not the norm of an undifferentiated Einstein residual.
On the actual shared-link branch the forcing contains the differentiated,
volume-normalized Palatini metric row, together with the differentiated
Cartan/torsion defect and the differential Bianchi defect.  Lower-order
transport terms are retained in $a_X$.  To state a quantitative criterion,
use the physical incidence \eqref{eq:main-current-incidence}, with all
untransferred contributions collected in $r_X^{\rm src}$.  Its verification
is an assertion about the actual action, connection and observation records;
it is not obtained simply by introducing an auxiliary source variable.

The pointwise source estimate used in
\Cref{prop:main-action-current} is
\begin{equation}
 \|f_X\|_{H_X}^2
       \le2\Gamma_XC_X\Delta_X+2\|r_X^{\rm src}\|_{H_X}^2.
 \label{eq:supp-curvature-source-budget}
\end{equation}
With deterministic supported upper bounds on $\Gamma_X$, it also holds in
expectation when the action-gradient estimate is only in expectation.
The proof below retains the expected product when no such upper bound is
available.

\begin{proof}[Proof of \Cref{prop:main-action-current}]
Write $\mathfrak a=G_X^{1/2}v$.  Then
$\|\mathfrak a\|_{G_X^{-1}}=\|v\|$ and
\[
 \|\mathsf T_X\mathfrak a\|_{H_X}
       =\|H_X^{1/2}\mathsf T_XG_X^{1/2}v\|.
\]
The maximum squared ratio is the largest squared singular value.  A unit
maximizing right singular vector attains it, so the constant is sharp and
also supplies a finite witness to large source amplification.
Apply $\|u+v\|^2\le2\|u\|^2+2\|v\|^2$ to
\eqref{eq:main-current-incidence}.  Time Cauchy--Schwarz proves
\eqref{eq:main-action-source-budget}; summing the same pointwise inequality
with weights $s_j$ proves its discrete counterpart.  For varying field
values the argument is pointwise before averaging.  Without a deterministic
supported bound, the expected estimate must retain
$\mathbb E[\Gamma_X\|\mathfrak a_X\|_{G_X^{-1}}^2]$, not a product of
separately averaged factors.
\end{proof}

\begin{proposition}[Sharp grid-scale source amplification]
\label{prop:supp-current-derivative-loss}
On a periodic grid of $m$ sites, with $m$ divisible by four, spacing $h$ and
inner product $h\sum_i u_iv_i$, let
$(D_hu)_i=(u_{i+1}-u_{i-1})/(2h)$.  With the same positive source and target
norms,
\begin{equation}
 \|D_h\|_{2\to2}^2=h^{-2}.
 \label{eq:supp-current-derivative-loss}
\end{equation}
In particular, a vanishing unscaled action-gradient norm does not imply a
bounded differentiated current.
\end{proposition}
\begin{proof}
The Fourier multiplier is $i\sin(2\pi k/m)/h$, whose maximum modulus is
$h^{-1}$ and is attained at $k=m/4$.  The same singular value occurs on the
real sine/cosine subspace.  Choose a unit vector $v_h$ in that subspace and
set $\mathfrak a_h=h^{1/2}v_h$.  Then
$\|\mathfrak a_h\|_2^2=h\to0$, whereas
$\|D_h\mathfrak a_h\|_2^2=h^{-1}\to\infty$.
\end{proof}

Thus, for this derivative transfer and uniformly controlled remaining
factors, $C_X\Delta_X=O(h_X^2)$ suffices for bounded source and
$o(h_X^2)$ suffices for vanishing source.  These rates are not inferred from
$\Delta_X\to0$.  Cancellations in the actual Bianchi transfer may improve
$\Gamma_X$, but they must be retained in that map and its physical Gram.

\begin{corollary}[Expected curvature energy and a sufficient pathwise upgrade]
\label{cor:supp-stochastic-curvature}
Suppose the finite propagation equation holds pathwise on $[0,T_*]$, with
$\int a_X\le A_*$ deterministically.  Suppose its forcing satisfies
\eqref{eq:main-current-incidence}, with the target norm already transferred
to the curvature-energy norm, and
$\mathbb E\|\mathfrak a_X(t)\|_{G_X^{-1}}^2\le C_X(t)\Delta_X(t)$.
Using deterministic supported upper bounds on $\Gamma_X$, one obtains
\begin{equation}
 \begin{split}
 \mathbb E\sup_{t\le T_*}z_X(t)^2
 \le 2e^{2A_*}\bigg[&\mathbb Ez_X(0)^2
       +2T_*\int_0^{T_*}\Gamma_XC_X\Delta_X\,\dd t\\
       &+2T_*\mathbb E\int_0^{T_*}\|r_X^{\rm src}\|_{H_X}^2\,\dd t\bigg].
 \end{split}
 \label{eq:supp-expected-curvature}
\end{equation}
Together with the stable interpolation and lapse bounds, uniformity of these
budgets gives a uniform expected squared spacetime curvature norm.  For a
countable cofinal sequence, a sufficient pathwise strengthening is
\begin{equation}
 \sup_Xz_X(0)<\infty\ \hbox{almost surely},\qquad
 \sum_X\mathbb E\int_0^{T_*}\|f_X(t)\|_{H_X}^2\,\dd t<\infty,
 \label{eq:supp-pathwise-source-budget}
\end{equation}
with pathwise uniform growth, interpolation, lapse and boundary budgets.
Under these conditions $\sup_X\|R_X\|_{L^2(K)}<\infty$ almost surely.
\end{corollary}
\begin{proof}
The pathwise estimate gives
$\sup_tz_X(t)\le e^{A_*}(z_X(0)+\int\|f_X\|)$.
Square and apply Cauchy--Schwarz in time to obtain
\[
 \mathbb E\sup_tz_X(t)^2
 \le2e^{2A_*}\left(\mathbb Ez_X(0)^2
       +T_*\mathbb E\int_0^{T_*}\|f_X\|_{H_X}^2\,\dd t\right).
\]
Apply \Cref{prop:main-action-current} and the interpolation estimate.
For the last assertion, Tonelli makes the sum of the nonnegative integrated
source energies finite almost surely.  Their supremum is then finite, and
the deterministic propagation estimate applies to the whole realized
sequence.  A uniform expected norm alone gives boundedness in probability
by Markov's inequality; it is not substituted for this pathwise argument.
\end{proof}


\subsection{Literal-link anisotropic compactness}
\label{subsubsec:supp-literal-link-compactness}

For the periodic route of \Cref{thm:main-literal-link-compactness}, let
$S_i$ denote the positive spatial shift and write
$U_i=\exp(hA_i)=I+hZ_i$.  Define the literal electric coefficient by the
exact link identity
\begin{equation}
 E_iU_i=\partial_tU_i/h-A_0U_i+U_iS_iA_0,
 \label{eq:supp-literal-electric-definition}
\end{equation}
with the same sign convention as the Cartan writer.  The spatial magnetic
record $F_{ij}$ is the $h^{-2}$ logarithmic plaquette coefficient on the
common identity chart.  These are actual finite link observables rather than
curvatures assigned after interpolation.

\begin{lemma}[Exact link evolution and negative-norm time control]
\label{lem:supp-literal-link-evolution}
The coordinate $Z_i=(U_i-I)/h$ satisfies
\begin{equation}
 \partial_tZ_i
 =E_iU_i+D_i^+A_0+Z_iS_iA_0-A_0Z_i.
 \label{eq:supp-literal-Z-equation}
\end{equation}
On the chart $h|A_i|\le\delta_*$,
\begin{align}
 |Z_i|&\le C|A_i|,
 &|Z_i-A_i|&\le Ch|A_i|^2,
 &|D_j^+Z_i|&\le C|D_j^+A_i|,
 \label{eq:supp-literal-Z-comparison}\\
 \|\partial_tZ\|_{H_h^{-2}}
 &\le C\left(\|E\|_{2,h}
 +(1+\|Z\|_{2,h})\|A_0\|_{2,h}\right).
 \label{eq:supp-literal-Z-time}
\end{align}
The constants are independent of the cutoff.
\end{lemma}
\begin{proof}
Multiply \eqref{eq:supp-literal-electric-definition} on the right by
$U_i^{-1}$ and use $U_i=I+hZ_i$; rearrangement gives
\eqref{eq:supp-literal-Z-equation}.  The first two pointwise estimates are
the exponential integral formula and its quadratic remainder.  The
differential of $Q_h(A)=(e^{hA}-I)/h$ has norm at most $e^{h|A|}$ on the
fixed chart, so integration along the segment between neighboring
coefficients gives the difference estimate.

Pair \eqref{eq:supp-literal-Z-equation} with a matrix test $\phi$ of
$H_h^2$ norm at most one.  The uniform discrete embedding bounds
$\|\phi\|_{\infty,h}$ and $\|D^-\phi\|_{2,h}$.  Summation by parts moves
$D_i^+$ from $A_0$ to $\phi$, while the two product terms have $L_h^1$ norm
at most $C\|Z\|_{2,h}\|A_0\|_{2,h}$.  The electric term is bounded by
$C\|E\|_{2,h}$ because $U_i$ is uniformly bounded.  Taking the supremum over
tests proves \eqref{eq:supp-literal-Z-time}.
\end{proof}

\begin{proof}[Proof of \Cref{thm:main-literal-link-compactness}]
The hypotheses and \Cref{lem:supp-literal-link-evolution} give a uniform
$L_t^\infty L_h^2\cap L_t^2H_h^1$ bound for $Z_h$ and an
$L_t^2H_h^{-2}$ bound for $\partial_tZ_h$.  For a fixed Fourier radius
$R$, every coefficient of the projection onto $|k|\le R$ is bounded in
$H^1(0,T)$; hence the projected family is precompact in $L^2$.  The spatial
$H^1$ estimate gives the uniform tail bound
\begin{equation}
 \|(I-P_R)\mathcal I_hZ_h\|_{L^2_{t,x}}
 \le CR^{-1}\|Z_h\|_{L_t^2H_h^1}.
 \label{eq:supp-literal-Fourier-tail}
\end{equation}
A diagonal extraction followed by $R\to\infty$ therefore yields strong
$L^2$ compactness of $\mathcal I_hZ_h$.

The spatial Sobolev inequality and interpolation give
\begin{equation}
 \int_0^T\|A_h\|_{10/3,h}^{10/3}\,\dd t
 \le \|A_h\|_{L_t^\infty L_h^2}^{4/3}
      \|A_h\|_{L_t^2L_h^6}^{2}\le C.
 \label{eq:supp-literal-integrability}
\end{equation}
Using $h|A_i|\le\delta_*$ and
\eqref{eq:supp-literal-Z-comparison},
\begin{equation}
 \|Z_h-A_h\|_{L^2}^2
 \le Ch^2\int|A_h|^4
 \le C\delta_*^{2/3}h^{4/3}\int|A_h|^{10/3}
 \longrightarrow0.
 \label{eq:supp-literal-log-error}
\end{equation}
Thus $A_h$ has the same strong spatial limit, while the assumed bound gives
$A_{0,h}\rightharpoonup A_0$ after extraction.

Pass to the exact temporal identity.  Since $U_i-I=hZ_i$,
$\|E_i(U_i-I)\|_1\le h\|E_i\|_2\|Z_i\|_2\to0$.
The shifted temporal arrays $S_iA_{0,h}$ have the same weak limit as
$A_{0,h}$, and the remaining products pass by weak--strong duality.
Summation by parts gives
\begin{equation}
 \partial_tA_i
 =E_i+\partial_iA_0+A_iA_0-A_0A_i
 \label{eq:supp-literal-electric-limit}
\end{equation}
in distributions.

For the magnetic record, the four-link plaquette expansion on the common
identity chart has linear part equal to the discrete curl and quadratic part
$[A_i,A_j]$.  Translation errors are bounded in $L^1$ by
$Ch\|D^+A\|_2\|A\|_2$, while the cubic remainder is bounded by
$Ch\|A\|_3^3$ and tends to zero by
\eqref{eq:supp-literal-integrability}.  Hence
\begin{equation}
 F_{ij}=\partial_iA_j-\partial_jA_i+[A_i,A_j]
 \label{eq:supp-literal-magnetic-limit}
\end{equation}
in the weak limit.  Equations
\eqref{eq:supp-literal-electric-limit}--\eqref{eq:supp-literal-magnetic-limit}
are exactly the components of
$R(\omega)=\dd\omega+\omega\wedge\omega$.  No strong compactness of
$A_0$ has been used.
\end{proof}

\subsection{Full-connection compactness: a Hodge--temporal criterion}

\paragraph{A separate connection compactness certificate.}
\begin{assumption}[Connection compactness on a common physical cylinder]
\label{ass:main-connection-compactness}
After the declared physical frame identifications, the actual interpolated
connections form a relatively compact family in $L^2(K)$ on each compact
cylinder $K$. The identifications and positive norms are fixed or uniformly
controlled on overlaps. No rapidly varying unrecorded gauge transformation
is used to assert convergence. A common subsequence is selected on a
countable compact exhaustion.
\end{assumption}

The assumption is distinct from a curvature-size bound and from smallness
of the connection Euler residual. The following gives a sufficient
realization in finite quantities. It also identifies the additional records
needed to replace the assumption by estimates.

On the periodic odd Cartesian regulator let $A_h=(A_{1,h},A_{2,h},A_{3,h})$
be the spatial connection coefficients and $A_{0,h}$ the temporal
coefficient, all in one common trivialization. With the forward and backward
differences of \eqref{eq:main-open-differences}, set
\begin{equation}
 (\dd_h A_h)_{ij}=D_i^+A_{j,h}-D_j^+A_{i,h},\qquad
 \delta_h A_h=-\sum_iD_i^-A_{i,h}.
 \label{eq:main-connection-hodge}
\end{equation}
The internal coefficient norm is positive; it is not the indefinite
Lorentz contraction. All norms in the following theorem are over the common
finite time interval and spatial torus, with uniformly comparable physical
measures. Let $\mathcal I_h$ be the trigonometric spatial interpolation.

\begin{theorem}[Finite Hodge--temporal sufficient criterion]
\label{thm:main-hodge-connection}
Suppose the actual spatial curvature writer obeys
\begin{equation}
 \dd_hA_h=R_h^{\rm sp}-\mathcal Q_h(A_h)+\epsilon_h,
 \qquad
 \|\mathcal Q_h(A_h)\|_{L^2_{t,h}}
       \le C\|A_h\|_{L^4_{t,h}}^2,
 \label{eq:main-connection-hodge-incidence}
\end{equation}
and that the following budget is uniform in $h$:
\begin{equation}
 \begin{split}
 &\|(A_{0,h},A_h)\|_{L^2_{t,h}}
 +\|R_h^{\rm sp}\|_{L^2_{t,h}}+\|A_h\|_{L^4_{t,h}}^2
 +\|\epsilon_h\|_{L^2_{t,h}}+\|\delta_hA_h\|_{L^2_{t,h}}\\
 &\qquad
 +\|D^+A_{0,h}\|_{L^2_{t,h}}
 +\|\partial_t\mathcal I_h(A_{0,h},A_h)\|_{L^2_tH^{-1}_x}
 \le C_K.
 \end{split}
 \label{eq:main-connection-hodge-budget}
\end{equation}
Then the interpolated full connections are relatively compact in
$L^2(K)$ and satisfy \Cref{ass:main-connection-compactness}.
\end{theorem}

The proof in this appendix uses an exact discrete Hodge
identity and a finite-mode temporal compactness argument. The codifferential,
$L^4$ product and temporal budgets are additional physical conditions; they
are not inferred from $\|R_h\|_2$ alone. For piecewise reconstructions, jumps
and frame transfers must be included in the actual temporal derivative and
incidence residual. The theorem does not establish existence of a gauge
satisfying these estimates for an arbitrary Lorentz connection.


The compactness result here is a quantitative sufficient criterion in a
common gauge, not a gauge-selection theorem. It is an instance of the
spatial compactness and temporal regularity mechanism underlying
\cite{SimonCompact}; the particular finite normalization is proved directly.

\begin{lemma}[Exact periodic discrete Hodge identity]
\label{lem:supp-connection-hodge}
For a vector-valued periodic spatial one-cochain $A=(A_1,A_2,A_3)$,
\begin{equation}
 \sum_{i,j}\|D_i^+A_j\|_h^2
   =\sum_{i<j}\|D_i^+A_j-D_j^+A_i\|_h^2
      +\left\|\sum_iD_i^-A_i\right\|_h^2.
 \label{eq:supp-discrete-hodge-identity}
\end{equation}
All coefficient norms use the same positive inner product.
\end{lemma}
\begin{proof}
For a Fourier frequency $k$ let $d_i=(e^{2\pi ihk_i}-1)/h$.
The symbol of $D_i^-$ is $-\overline d_i$. On the coefficient vector
$z=(\widehat A_1,\widehat A_2,\widehat A_3)$ the identity is
\[
 \sum_{i<j}|d_i z_j-d_j z_i|^2
       +\left|\sum_i\overline d_i z_i\right|^2
       =\left(\sum_i|d_i|^2\right)\left(\sum_j|z_j|^2\right).
\]
Expand the two squares; their cross terms cancel. Apply this identity to
each internal coordinate and sum using Parseval. No continuum product rule
or curvature identity is used in this step.
\end{proof}

\begin{proof}[Proof of \Cref{thm:main-hodge-connection}]
The incidence \eqref{eq:main-connection-hodge-incidence} and the stated
product bound give a uniform $L^2_{t,h}$ bound for $\dd_h A_h$.
Equation \eqref{eq:supp-discrete-hodge-identity} then bounds every spatial
forward derivative of $A_h$ in that norm. The temporal coefficient has its
own spatial derivative bound in
\eqref{eq:main-connection-hodge-budget}. Uniform equivalence of the
forward-difference Sobolev norm and the trigonometric $H^1$ norm on the odd
frequency cube therefore gives
\[
 \|\Omega_h\|_{L^2_tH^1_x}
       +\|\partial_t\Omega_h\|_{L^2_tH^{-1}_x}\le C_K,
 \qquad \Omega_h=\mathcal I_h(A_{0,h},A_h).
\]
For completeness fix a Fourier radius $R$ and let $P_R$ be the common
orthogonal projection onto $|k|\le R$. The spatial estimate gives
\[
 \|(I-P_R)\Omega_h\|_{L^2_{t,x}}^2
       \le C_K(1+R^2)^{-1}.
\]
Each of the finitely many coefficients of $P_R\Omega_h$ is bounded in
$L^2(0,T)$, and its distributional time derivative is bounded in $L^2(0,T)$:
pair the $H^{-1}$ derivative with that fixed smooth Fourier vector. Thus
these coefficients are bounded in $H^1(0,T)$ and relatively compact in
$L^2(0,T)$. Choose a subsequence successively for $R=1,2,\ldots$.
The uniform tail estimate makes the diagonal subsequence Cauchy in
$L^2_{t,x}$. Uniformly comparable physical norms give the asserted strong
compactness of the actual connections. Restriction and a second diagonal
argument give a common subsequence on a countable cylinder exhaustion.
\end{proof}
\subsection{Transported screens and general curvature identification}

\begin{proposition}[Transported-screen version of connection compactness]
\label{prop:supp-connection-screens}
Let $\Omega_X\in L^2(0,T;\mathcal H)$ be uniformly bounded after common
physical interpolation into a Hilbert space $\mathcal H$. For every $R$
suppose there are time-independent finite-rank operators $P_{X,R}$ and
$P_R$ such that $\|P_{X,R}-P_R\|_{\mathrm{op}}\to0$, and
\begin{equation}
 \lim_{R\to\infty}\sup_X
       \|(I-P_{X,R})\Omega_X\|_{L^2_t\mathcal H}=0.
 \label{eq:supp-connection-screen-tail}
\end{equation}
Assume the coefficients of $P_R\Omega_X$ in a fixed finite basis of its
range form a relatively compact family in $L^2(0,T)$ for each $R$.
Then $\Omega_X$ is relatively compact in $L^2(0,T;\mathcal H)$.
\end{proposition}
\begin{proof}
For fixed $R$, operator-norm convergence and the uniform $L^2$ bound give
$(P_{X,R}-P_R)\Omega_X\to0$ in $L^2_t\mathcal H$. Extract successive
subsequences on which the finitely many coefficients of $P_R\Omega_X$
converge. For the diagonal subsequence, split a difference into its two
screen tails and its finite-rank part. First take the cutoff limit at fixed
$R$ and then let $R\to\infty$, using
\eqref{eq:supp-connection-screen-tail}. This gives a Cauchy sequence in the
complete space $L^2_t\mathcal H$.
\end{proof}

A bounded $H^1$ time norm for each transported coefficient is one sufficient
condition in the proposition. If screens or frames depend on time, their
actual derivatives contribute to those coefficient bounds. Spatial rank
counts by themselves do not provide this temporal compactness. This
formulation allows non-Cartesian regulators without identifying arbitrary
finite graph connections with smooth-gauge connections in advance.


\paragraph{Identification with the actual connection curvature.}
Curvature size and curvature identification are separate questions.  Let
$\mathcal T_{X,K}$ be the finite determining curvature-test bank, and let
$\mathcal R_X(\omega_X)$ denote the finite Cartan writer on the common-cylinder
interpolation.  Define
\begin{equation}
 \mathfrak C_X(K):=
 \sup_{\substack{\Phi\in\mathcal T_{X,K}\\\|\Phi\|_{L^2(K)}\le1}}
 \left|\int_K\!\langle R_X-\mathcal R_X(\omega_X),\Phi\rangle\right|.
 \label{eq:main-cartan-curvature-residual}
\end{equation}
The banks exhaust a fixed countable determining smooth core, with compatible
test restrictions and vanishing test-interpolation errors.  Their geometric
consistency is recorded separately by
\begin{equation}
 \mathfrak I_X(\Phi):=
 \int_K\!\langle\mathcal R_X(\omega_X),\Phi\rangle
 -\langle\dd_{\rm dist}\omega_X+\omega_X\wedge\omega_X,\Phi\rangle.
 \label{eq:main-cartan-incidence}
\end{equation}
The distributional derivative includes every interface contribution of a
piecewise connection.  Thus matching two finite writers alone is not used to
identify the nonlinear curvature of a weakly convergent connection.

For the propagated $p=2$ route the operational hypotheses are
\begin{equation}
 \omega_X\to\omega\ \hbox{strongly in }L^2_{\rm loc},\qquad
 \mathfrak C_X(K)\to0,\qquad
 \mathfrak I_X(\Phi)\to0
 \label{eq:main-identified-curvature-hypotheses}
\end{equation}
The theorem below gives the more general $L^p$ statement.


\begin{theorem}[Weak curvature compactness with an identified connection limit]
\label{thm:supp-curvature-compactness}
Let $K$ be a compact cylinder and $1<p<\infty$.  A bound
\begin{equation}
 \sup_X\|R_X\|_{L^p(K)}\le C_{K,p}<\infty
 \label{eq:supp-curvature-lp-bound}
\end{equation}
gives a weakly convergent subsequence
\begin{equation}
 R_X\rightharpoonup R\quad\hbox{in }L^p(K).
 \label{eq:supp-curvature-reflexive-limit}
\end{equation}
If, in addition, $\omega_X\to\omega$ strongly in $L^2_{\rm loc}$ and
\begin{equation}
 R_X-[\dd_{\rm dist}\omega_X+\omega_X\wedge\omega_X]
       \longrightarrow0\quad\hbox{in distributions},
 \label{eq:supp-curvature-identification}
\end{equation}
then $R=\dd\omega+\omega\wedge\omega$ distributionally.
The finite-writer conditions
\eqref{eq:main-cartan-curvature-residual}--\eqref{eq:main-cartan-incidence},
with exhaustive determining tests and their interpolation control, give
\eqref{eq:supp-curvature-identification}.  In particular, the propagated $p=2$ bound gives the stated connection-curvature conclusion.
\end{theorem}
\begin{proof}
For $1<p<\infty$, reflexivity gives weak sequential compactness of bounded
balls of $L^p(K)$.  Strong $L^2$ convergence of the connection gives
convergence of its distributional derivative, and
\[
 \|\omega_X\wedge\omega_X-\omega\wedge\omega\|_{L^1(K)}
 \le C(\|\omega_X\|_2+\|\omega\|_2)\|\omega_X-\omega\|_2\to0.
\]
Thus \eqref{eq:supp-curvature-identification} identifies the distributional
limit of $R_X$.  Weak $L^p$ convergence has the same distributional limit,
proving the assertion.  For the finite-writer version, split the tested
residual into $R_X-\mathcal R_X(\omega_X)$ and the discrepancy
\eqref{eq:main-cartan-incidence}.  The first tends to zero on the exhaustive
core by \eqref{eq:main-cartan-curvature-residual}, and the second by the
stated consistency.  The declared test-interpolation control extends the
identity to the determining test topology.  A diagonal subsequence works on
a countable compact exhaustion.

If the curvature norms are instead unbounded, choose a subsequence on which
they diverge and retain the unit fields $R_X/\|R_X\|_{L^p(K)}$ together
with their diverging scales.  This witnesses failure of the specified
curvature-size certificate, not nonexistence of every possible weaker limit.
\end{proof}

For a piecewise smooth connection, its full distributional derivative includes
face measures.  A bound on cell interiors alone therefore does not control the
curvature entering the continuum variation.  The following construction
specifies a sufficient volume realization without discarding these measures.

Work on a periodic rectangular four-dimensional domain with a uniform cubical
mesh of side $h$, or on a compact interior subdomain together with its
neighbouring cells.  Polynomial degrees are fixed, positive Grams and physical
measures are uniformly comparable to the reference ones, and all face
coefficients are compared in a common local trivialization.  For an oriented
interior face $f$, set $J_f=n_f^\flat\wedge[\omega_h]_f$, with the sign
chosen so that
\subsection{Interfaces and distributional curvature}

For discontinuous connection reconstructions, the identification has a concrete
sufficient realization.  On a uniformly regular cubical chart put
$J_f=n_f^\flat\wedge[\omega_h]_f$ at each interior face and
\begin{equation}
 \mathcal J_h^2=\sum_fh^{-1}\|J_f\|_{L^2(f)}^2.
 \label{eq:main-interface-budget}
\end{equation}
A bounded-overlap volume lifting $\mathcal L_hJ$ of the face measures gives
\begin{equation}
 \begin{aligned}
 R_h^{\rm lift}
   &=\dd_b\omega_h+\omega_h\wedge\omega_h+\mathcal L_hJ,\\
 \|R_h^{\rm lift}\|_{L^2(K)}
   &\le\|\dd_b\omega_h\|_2+C\|\omega_h\|_4^2+C\mathcal J_h.
 \end{aligned}
 \label{eq:main-lifted-curvature}
\end{equation}
Its discrepancy from the full distributional curvature is at most
$C_Kh\mathcal J_h\|\nabla\Phi\|_\infty$ on a fixed smooth test.  The
construction and proof are given in \Cref{thm:supp-interface-curvature}.
These are explicit additional interface budgets, not consequences of a small
connection Euler residual.  The lifting is used only when the action's actual
curvature writer agrees with it on the determining tests, up to a vanishing
recorded discrepancy.

\begin{equation}
 \dd_{\rm dist}\omega_h=\dd_b\omega_h+\sum_fJ_f\delta_f.
 \label{eq:supp-connection-jumps}
\end{equation}
Let $\eta\ge0$ be a fixed smooth function supported in $(-1/4,1/4)$ with
integral one.  On the normal prism of $f$, define
\begin{equation}
 (\mathcal L_fJ_f)(y+rn_f)=h^{-1}\eta(r/h)J_f(y),\qquad
 \mathcal L_hJ=\sum_f\mathcal L_fJ_f,
 \label{eq:supp-interface-lifting}
\end{equation}
and extend each term by zero outside its prism.

\begin{lemma}[Interface lifting and its distributional error]
\label{lem:supp-interface-lifting}
With $\mathcal J_h$ defined in \eqref{eq:main-interface-budget},
\begin{equation}
 \|\mathcal L_hJ\|_2\le C\mathcal J_h,\qquad
 \left|\left\langle\mathcal L_hJ-\sum_fJ_f\delta_f,\Phi\right\rangle\right|
       \le C_Kh\mathcal J_h\|\nabla\Phi\|_\infty
 \label{eq:supp-interface-lifting-bounds}
\end{equation}
for each fixed smooth compactly supported or periodic test $\Phi$.
\end{lemma}
\begin{proof}
Integration in the normal coordinate gives
\[
 \|\mathcal L_fJ_f\|_2^2
       =h^{-1}\|\eta\|_{L^2(\R)}^2\|J_f\|_{L^2(f)}^2.
\]
The prisms have bounded overlap depending only on dimension.  Pointwise
Cauchy--Schwarz over the overlapping prisms proves the first bound.
For the second, subtract $\Phi(y)$ from $\Phi(y+rn_f)$ and use
$|r|\le h/4$.  The total error is at most
$Ch\|\nabla\Phi\|_\infty\sum_f\|J_f\|_{L^1(f)}$.  Moreover
\[
 \sum_fh\|J_f\|_{L^1(f)}
 \le\left(\sum_fh^{-1}\|J_f\|_{L^2(f)}^2\right)^{1/2}
       \left(\sum_fh^3|f|\right)^{1/2}
 \le C_Kh\mathcal J_h.
\]
There are $O_K(h^{-4})$ faces of area $h^3$, which proves the last
inequality and the result.
\end{proof}

\begin{theorem}[Constructive connection-and-interface curvature criterion]
\label{thm:supp-interface-curvature}
Suppose
\begin{equation}
 \|\dd_b\omega_h\|_{L^2(K)}+\|\omega_h\|_{L^4(K)}^2+
       \mathcal J_h\le C_K.
 \label{eq:supp-connection-interface-budget}
\end{equation}
Then the volume curvature $R_h^{\rm lift}$ in
\eqref{eq:main-lifted-curvature} is uniformly bounded in $L^2(K)$ and has
the displayed distributional error.  If also
$\omega_h\to\omega$ strongly in $L^2_{\rm loc}$, every weak $L^2$ cluster
point equals $\dd\omega+\omega\wedge\omega$ distributionally.
\end{theorem}
\begin{proof}
The matrix exterior product satisfies
$|\omega_h\wedge\omega_h|\le C|\omega_h|^2$.  H\"older's inequality and
\Cref{lem:supp-interface-lifting} give
\eqref{eq:main-lifted-curvature}.  The discrepancy from
$\dd_{\rm dist}\omega_h+\omega_h\wedge\omega_h$ is exactly the lifting
error in \eqref{eq:supp-interface-lifting-bounds}, which tends to zero on
each fixed test.  Apply \Cref{thm:supp-curvature-compactness} with $p=2$.
\end{proof}

The theorem is a sufficient identification and size criterion with its own
explicit connection and jump budgets.  A stationary connection estimate in
$L^2$ does not imply \eqref{eq:supp-connection-interface-budget}.
Nor is $R_h^{\rm lift}$ silently substituted for the action's curvature
writer: that writer must agree with the lift on the determining tests up to
a vanishing recorded discrepancy.  On a conforming connection reconstruction
$J_f=0$, and the interface part is absent.

\subsubsection{The weak--strong Palatini passage}

\subsection{Weak--strong variational passage}

\begin{lemma}[Weak--strong curvature pairing and volume convergence]
\label{lem:supp-weak-strong-palatini}
Assume \eqref{eq:supp-coframe-l2-l6}, let $p>3/2$, and suppose
\begin{equation}
 R_X\rightharpoonup R\quad\hbox{weakly in }L^p(K).
 \label{eq:supp-curvature-weak}
\end{equation}
Then, for every bounded compactly supported coefficient tensor $\Phi$,
\begin{equation}
 \int_K\!\langle\Phi(e_X\wedge e_X),R_X\rangle
 \longrightarrow
 \int_K\!\langle\Phi(e\wedge e),R\rangle.
 \label{eq:supp-weak-strong-pairing}
\end{equation}
The same conclusion holds with an internal Hodge star applied to either
bivector factor.  Moreover
\begin{equation}
 e_X\wedge e_X\wedge e_X\wedge e_X
 \longrightarrow
 e\wedge e\wedge e\wedge e
 \quad\hbox{strongly in }L^1(K).
 \label{eq:supp-volume-strong}
\end{equation}
\end{lemma}

\begin{proof}
Weak convergence in $L^p$ implies $\sup_X\|R_X\|_{L^p}<\infty$.  Put
$B_X=\Phi(e_X\wedge e_X)$ and $B=\Phi(e\wedge e)$.  By
\Cref{lem:supp-coframe-interpolation}, $B_X\to B$ strongly in $L^{p'}$.
Split the difference as
\begin{align}
 \int_K\langle B_X,R_X\rangle-\int_K\langle B,R\rangle
 &=\int_K\langle B_X-B,R_X\rangle
   +\int_K\langle B,R_X-R\rangle.
 \label{eq:supp-pairing-split}
\end{align}
The first term tends to zero by H\"older and the uniform $L^p$ bound on
$R_X$; the second tends to zero by the definition of weak $L^p$ convergence,
because $B\in L^{p'}$.  The internal Hodge star is a fixed bounded linear map
on the bivector fibre and changes none of these estimates.

For the volume term, \Cref{lem:supp-coframe-interpolation} with $q=4$ gives
$e_X\to e$ strongly in $L^4(K)$ and uniform $L^4$ boundedness.  Expanding the
fourfold difference one factor at a time and applying H\"older with four
$L^4$ factors gives
\begin{align}
 &\|e_X^{\wedge4}-e^{\wedge4}\|_{L^1} \\
 &\qquad\le C\|e_X-e\|_{L^4}
 (\|e_X\|_{L^4}+\|e\|_{L^4})^3\longrightarrow0,
\end{align}
which is \eqref{eq:supp-volume-strong}.
\end{proof}

\begin{remark}[Weak--strong passage without compensated compactness]
\label{rem:supp-no-divcurl}
The product passage above does not use a div--curl or compensated-compactness
lemma.  Such a device would be needed if both the coframe bivector and curvature
were only weakly convergent.  Here the compact-screen hypothesis gives strong
$L^2$ convergence of the coframe together with a uniform $L^6$ bound; the
explicit interpolation in \Cref{lem:supp-coframe-interpolation} upgrades this
to strong $L^{2p'}$, so ordinary weak--strong duality suffices.  At this stage
all fields are already represented by their canonical common-cylinder
interpolants; equivalently, the same estimates may be read cellwise with the
physical cell measures because the displayed $L^q$ hypotheses are precisely
the cutoff-uniform interpolation bounds used in the handoff.
\end{remark}

\begin{theorem}[Localized renewal--Palatini handoff]
\label{thm:supp-renewal-palatini}
Assume the face-linearization and complete memory-short branches, the
polynomial local class of \Cref{def:main-gravitational-class}, real Lorentz
naturality, scalar residual typed multiplicity commutant, and represented
determinant volume. Assume in addition the physical first-variation
remainder condition \eqref{eq:main-first-variation-reduction}.  Let the constant calibrated coefficients
converge to $(\alpha,\beta,\lambda)$ with $\beta\ne0$, and let the spinless
connection Euler residual vanish with the uniform Cartan inverse of
\Cref{thm:supp-palatini-torsion}.  On every compact $K$, assume
\eqref{eq:supp-coframe-l2-l6}, nondegeneracy of the limiting coframe, and either
the continuous or the discrete initial-energy curvature certificate.
Assume also that the connection curvature is identified either by
\Cref{thm:supp-curvature-compactness} or by the literal-link route of
\Cref{thm:main-literal-link-compactness}, and assume the bounded, convergent
physical test lifts of \textup{(E1)} in \Cref{thm:main-renewal-einstein}.
Then the limiting score has the classified form
\begin{equation}
 \alpha L_{\rm Holst}+\beta L_{\rm Palatini}+\lambda L_{\rm vol},
 \label{eq:supp-phv-classification}
\end{equation}
limiting torsion vanishes, and the metric first variations converge to
\begin{equation}
 \chi\int_K\sqrt{-g}\,
       (G_{\mu\nu}+\Lambda g_{\mu\nu})k^{\mu\nu},\qquad\chi\ne0.
 \label{eq:supp-einstein-insertion}
\end{equation}
Here $\chi$ is the nonzero Palatini normalization and $\Lambda$ is the
corresponding calibrated volume-to-Palatini coefficient ratio, in the density
conventions of \eqref{eq:main-phv-densities}.
\end{theorem}
\begin{proof}
Both coframe bivectors and Lorentz curvature lie in
$\Lambda^2\R^{1,3}$.  After complexification the bivector module splits into
nonisomorphic self-dual and anti-self-dual irreducibles.  The real commutant is
therefore
\begin{equation}
 \operatorname{End}_{SO^+(1,3)}(\Lambda^2\R^{1,3})
       =\Span_\R\{I,\star\}.
 \label{eq:supp-lorentz-commutant}
\end{equation}
The residual typed commutant supplies no additional internal tensor, and the
only alternating four-coframe scalar at this order is determinant volume.
Thus \eqref{eq:supp-phv-classification} follows in the stated first-order
class.  The calibrated coefficient limits ensure that these are the
coefficients of the limiting variational functional, rather than unrelated
values chosen after taking the limit. The first-variation remainder condition
removes the unclassified part on the actual physical tests; no interchange
of differentiation and merely pointwise convergence is being assumed.

By \Cref{thm:supp-palatini-torsion}, the vanishing connection Euler residual
and the uniform inverse bounds give vanishing torsion.  On the full-connection
route, strong local $L^2$ convergence of $\omega_X$ passes the Cartan product
directly.  On the literal-link route the spatial coefficients converge
strongly in $L^2$ while the temporal coefficient converges weakly in $L^2$;
strong convergence of the coframe is sufficient to pass every torsion product
by weak--strong duality.  The distributional derivative of $e_X$ passes
against every fixed smooth test, so the limit is a metric-compatible
torsion-free represented connection.  Its curvature is identified by the
selected connection route, not merely by matching finite curvature labels.

The continuous or discrete propagation theorem supplies
$\sup_X\|R_X\|_{L^2(K)}<\infty$.  After passage to the identified weak
subsequence, \Cref{lem:supp-coframe-interpolation} gives
$e_X\wedge e_X\to e\wedge e$ strongly in $L^2$.  The curvature terms
therefore converge by \Cref{lem:supp-weak-strong-palatini} with $p=p'=2$.
The same argument applies to their physical coframe variations: their
coefficient tensors are uniformly bounded and converge, so multiplying the
strong $L^2$ bivectors by these tensors preserves strong convergence.
The four-coframe volume term and its test insertions converge in $L^1$ by
the same lemma.  No product of two uncontrolled weak limits is used.

For the torsion-free connection, the first Bianchi identity removes the
Holst coframe variation.  The Palatini and volume variations give the
Einstein and cosmological insertions in
\eqref{eq:supp-einstein-insertion}.  The coefficient of the first is nonzero
because $\beta\ne0$.  A nonzero coefficient pair alone would only justify
the torsion step, and would not justify division by this Einstein coefficient.
\end{proof}

\begin{corollary}[Vacuum equation from certified accepted-action variations]
\label{cor:supp-renewal-einstein}
Under the preceding hypotheses, vanishing realized common-action variations
on the determining metric-test core imply
\[
 G_{\mu\nu}+\Lambda g_{\mu\nu}=0
\]
distributionally on every compact cylinder.  The scaled hypothesis
\eqref{eq:supp-scaled-test-gap} is a sufficient stochastic route on the
almost-sure subsequence of \Cref{prop:supp-physical-test-stationarity},
provided the remaining compactness and identification budgets hold on the
same realized sequence.
\end{corollary}
\begin{proof}
The handoff identifies the limit of each certified first variation with
\eqref{eq:supp-einstein-insertion}.  Its value is zero on the determining
core, and $\chi\ne0$ permits division.  Continuity in the declared test
topology extends the identity to the closure of that core.  In the stochastic
case, first take the realized vanishing subsequence supplied by
\Cref{prop:supp-physical-test-stationarity}; the assumed pathwise compactness
budgets then permit the curvature subsequence used in the handoff.
A mean-square bound is not interpreted as stationarity of every finite
sample.
\end{proof}

\subsection{The failed-certificate witnesses}
\label{subsec:supp-failed-certificates}

For reference, the complete failure bank consists of a common-action curl
or period; a nonvanishing realized metric variation or failure of its scaled
action-gap bound; a vanishing Palatini normalization or failed coupling
calibration; a cut/potential witness; a face-linearity, memory or torsion
defect; an unbounded initial, growth, transfer or differentiated-source
budget; an unstable curvature interpolation; a Cartan or interface
mismatch; failure of the literal-link chart/energy criterion on that route;
or failure of the common-cylinder compactness conditions.


\begin{proof}[Proof of Corollary~\ref{thm:main-einstein-alternative}]
For a complete cylinder, \Cref{thm:main-determining-kernel} reconstructs the
finite determining field and accepted kernel.  The reward-pressure and
common-action theorems either reconstruct the common action or expose its
curl/period failure.  The accepted-action theorem and the physical test-lift
estimate give a quantitative stationarity certificate with every cutoff
constant retained.  The spatial cut theorem gives the stated spatial screen
or a cut witness, while the common-cylinder comparison packet separately
records the temporal and transported-screen conditions needed for strong
coframe convergence.

Face linearization and memory subtraction give the classified first-order
score on their positive branches.  The calibrated coupling and Cartan tests
supply the nonzero Einstein coefficient and torsion control.  Continuous or
fully discrete propagation bounds the actual curvature from the finite
initial energy and growth/transfer/source records; the sharp current transfer
may supply the last budget from scaled action errors.  Finally the connection
and interface comparisons identify that curvature.  If all certificates pass,
\Cref{thm:supp-renewal-palatini,cor:supp-renewal-einstein} prove the positive branch, using a diagonal subsequence on the compact exhaustion.

For completeness, the cofinal witnesses in the complementary branch are obtained
from the quantifiers of the failed certificates.  A nonnegative residual
which does not tend to zero has a subsequence bounded below by some
$\epsilon>0$.  An unbounded budget has a subsequence tending to infinity.
A failed strictly positive lower margin has a subsequence tending to zero;
for the cut margin the finite graph supplies the corresponding potential
witness.  Failure of a uniform compactness tail supplies an $\epsilon>0$
and normalized components outside arbitrarily large declared screens.
Failure of a required calibrated Cauchy comparison retains the separated
pairs of finite records themselves.  A finite algebraic failure retains its
actual curl, period or operator residual.  Discarding a finite prefix is
allowed when checking the cofinal tail.  These records identify which
sufficient certificate failed; they do not establish that every alternative
compactness argument or every other subsequence must fail.
\end{proof}

\begin{remark}[What is discharged and what remains dynamical]
Neither a continuum curvature bound nor an unidentified nonlinear connection
product is inserted into the conclusion.  The former follows from the
continuous or discrete propagation certificate; the latter is handled by
strong connection convergence and the full distributional Cartan comparison,
including interfaces.  The differentiated source budget can follow from
\eqref{eq:main-action-source-budget}, whereas bounded symmetric growth,
physical occurrence and noncollapse remain genuine conditions of the
microscopic evolution.  A bounded mean curvature norm does not by itself
supply the pathwise version.  The fixed-amplitude nonsymmetric Cauchy theorem
is supplied independently by Appendix~\ref{supp:open-3plus1} for its declared
local writer; the Gowdy sector below remains a complementary nonlinear
symmetry-reduced class with pathwise full-curvature control.
\end{remark}


\section{Constructive nonsymmetric evolution and finite renewal realization}
\label{supp:open-3plus1}

This section proves \Cref{thm:main-open-3plus1} directly from the declared
local writer.  The purpose is to supply a fixed-radius, genuinely
three-dimensional nonlinear population theorem on the same reconstructed
Lorentzian coefficient chart used by the general certification theorem,
without invoking \textup{(E1)--(E5)} as independent assumptions.  The
evolution law is declared explicitly and is kept logically distinct from the
exact finite Euler flow studied in Appendix~\ref{supp:exact-action-audit}.
No symmetry reduction or analytic Fourier tail is used.

\subsection{Execution graph, exact difference calculus, and the local writer}

Let $N=2m+1$, $h=N^{-1}$ and
$\Lambda_h=(h\mathbb Z/\mathbb Z)^3$.  Let $\mathcal I_h$ denote the real
trigonometric interpolation on the signed frequency cube
$\{-m,\ldots,m\}^3$, and let $\mathcal S_h$ denote grid sampling.  The
tetrahedral $A_3\cong D_3$ root
execution directions may be represented as
\begin{equation}
 \mathcal R_F=\{\pm(e_i+e_j),\ \pm(e_i-e_j):1\le i<j\le3\}.
 \label{eq:supp-open-fcc-roots}
\end{equation}
All these integer vectors have even coordinate sum.  This harmless parity at
fixed smooth frequency becomes relevant for a cutoff-uniform Cartesian
Sobolev realization.

\begin{lemma}[Root-parity obstruction and a minimal execution connector]
\label{lem:supp-open-parity}
On the odd periodic grid, the root Dirichlet form
\begin{equation}
 \mathcal D_F(u)=\sum_{r\in\mathcal R_F/\{\pm1\}}
       \|(S_r-I)u/h\|_h^2
 \label{eq:supp-open-root-dirichlet}
\end{equation}
is not uniformly coercive for
$\|\nabla\mathcal I_hu\|_{L^2}^2$.  Moreover the root-graph distance from
$0$ to $e_1$ is $N-1$.  If the oriented execution edges $\pm e_1$ are
adjoined, then with
\begin{equation}
 \mathcal D_+(u)=\mathcal D_F(u)+\|D_{e_1}^+u\|_h^2,
 \qquad \mathcal D_C(u)=\sum_i\|D_i^+u\|_h^2,
 \label{eq:supp-open-connector-energy}
\end{equation}
one has
\begin{equation}
 \frac15\mathcal D_C(u)\le\mathcal D_+(u)\le9\mathcal D_C(u).
 \label{eq:supp-open-connector-coercivity}
\end{equation}
At least one added fixed direction of odd coordinate sum is necessary within
this fixed-direction class.  The connector is an execution incidence only;
it does not enter the $A_3$ moment reconstruction.
\end{lemma}
\begin{proof}
Take
$u_h(x)=\exp(2\pi i m(x_1+x_2+x_3))$.  Every difference in a minus root
vanishes and every plus-root multiplier remains bounded, while
$\|\nabla\mathcal I_hu_h\|_2\asymp m$ diverges.  The same argument applies to
every fixed-radius word in even-parity directions.  A root path ending at
$e_1$ has integer displacement $e_1+Nz$ with even coordinate sum; since $N$
is odd, $\sum z_i$ is odd and one coordinate has modulus at least $N-1$.
The explicit alternating pair $-e_1+e_2,-e_1-e_2$, repeated $m$ times,
attains that distance.

For the connector, the identity
\[
 S_{e_2}-I=S_{e_1}^{-1}\bigl[(S_{e_1+e_2}-I)-(S_{e_1}-I)\bigr]
\]
and its $e_3$ analogue give
$\mathcal D_C\le5\mathcal D_+$.  Conversely each root difference is bounded
by the sum of the two participating coordinate differences, so
$\mathcal D_F\le8\mathcal D_C$.  Adding the connector gives the upper bound.
If every added direction had even coordinate sum, the high-frequency test
above would remain unchanged.
\end{proof}

Define $S_i,D_i^\pm,D_i^0$ by \eqref{eq:main-open-differences}, use
$\langle u,v\rangle_h=h^3\sum_xu(x)\cdot v(x)$, and for a multi-index put
$D^\alpha=\prod_i(D_i^+)^{\alpha_i}$ and
$S^\alpha=\prod_iS_i^{\alpha_i}$.  Set
\begin{equation}
 \|u\|_{r,h}^2=\sum_{|\alpha|\le r}\|D^\alpha u\|_h^2,
 \qquad
 \|X\|_{\mathcal X_h^r}^2=\|q\|_{r+1,h}^2+\|v\|_{r,h}^2.
 \label{eq:supp-open-norms}
\end{equation}

\begin{lemma}[Exact shifted product rule and uniform Sobolev calculus]
\label{lem:supp-open-calculus}
On the periodic array space,
\begin{equation}
 (D_i^+)^*=-D_i^-,\qquad (D_i^0)^*=-D_i^0,
 \label{eq:supp-open-adjoints}
\end{equation}
and
\begin{equation}
 D^\alpha(fw)=\sum_{\beta\le\alpha}\binom\alpha\beta
 (S^{\alpha-\beta}D^\beta f)D^{\alpha-\beta}w.
 \label{eq:supp-open-shifted-product}
\end{equation}
For $r\ge3$, the norms \eqref{eq:supp-open-norms} are uniformly equivalent
to the trigonometric $H^r$ norms, form a uniform product algebra, and the
shifted commutator
\begin{equation}
 \mathcal C_\alpha(f,w)=D^\alpha(fw)-(S^\alpha f)D^\alpha w
 \label{eq:supp-open-commutator}
\end{equation}
satisfies, for a constant coefficient $f_*$,
\begin{align}
 \|\mathcal C_\alpha(f,w)\|_h
 &\le C_r\|f-f_*\|_{r,h}\|w\|_{r-1,h},\\
 \|D_i^-\mathcal C_\alpha(f,w)\|_h
 +\|D_i^0\mathcal C_\alpha(f,w)\|_h
 &\le C_r\|f-f_*\|_{r+1,h}\|w\|_{r,h}.
 \label{eq:supp-open-commutator-bounds}
\end{align}
The constants are independent of $h$.
\end{lemma}
\begin{proof}
Translation of the finite sum gives the adjoints, and
$D_i^+(fw)=(S_if)D_i^+w+(D_i^+f)w$ gives
\eqref{eq:supp-open-shifted-product} by iteration.  The forward symbol
$d_i(k)=(e^{2\pi ihk_i}-1)/h$ satisfies
$4|k_i|\le|d_i(k)|\le2\pi|k_i|$ on $|k_i|\le m$.  Parseval therefore gives
uniform norm equivalence.  Multiplication is convolution on the finite
frequency group and the exact identity
$d_i([k+p])=d_i(k)+e^{2\pi ihk_i}d_i(p)$ gives the usual uniform weighted
convolution estimate.  In every commutator term at least one difference falls
on $f$; one additional output difference gives the second line.  Since
$r>3/2$, the required Fourier $\ell^1$ weights are summable uniformly in the
finite cube.  Analytic coefficient compositions follow from their convergent
power series on a smaller ball.
\end{proof}

For the harmonic reduced vacuum equation use the coefficient functions in
\eqref{eq:main-harmonic-normal-row}.  The skew transport is
\eqref{eq:main-open-skew-transport} and the first-jet compensator is
\eqref{eq:main-open-compensator}.  Pointwise multiplication operators are
denoted $M_f$.

\begin{proposition}[Locality and continuum consistency of the writer]
\label{prop:supp-open-locality}
The acceleration in \eqref{eq:main-open-writer} at a site depends only on the
site and the twelve coordinate-stencil neighbours
$\{\pm e_i,\pm(e_i-e_j):i<j\}$.  These lie within two hops of the execution
graph obtained in \Cref{lem:supp-open-parity}.  On every smooth metric jet,
the continuum limit of the second row is exactly
\eqref{eq:main-harmonic-normal-row}.
\end{proposition}
\begin{proof}
The flux at $x$ is
\[
 h^{-2}\sum_{i,j}\{c^{ij}(x)[q(x+he_j)-q(x)]
 -c^{ij}(x-he_i)[q(x-he_i+he_j)-q(x-he_i)]\}.
\]
The skew transport uses only $x\pm he_i$, and the compensator only first
forward differences.  In the smooth limit the divergence contributes
$(\partial_ic^{ij})\partial_jq$ and the skew transport contributes
$(\partial_ib^i)v$ in addition to the desired second- and mixed-derivative
terms.  The two coefficient-differential terms in
\eqref{eq:main-open-compensator} cancel exactly these additions by the
ordinary chain rule.  No discrete Leibniz rule is invoked.
\end{proof}

\subsection{Shifted nonlinear energy and a cutoff-independent lifespan}

On a sufficiently small fixed $\mathcal X_h^s$ ball, $a\ge a_*>0$ and
$c\succeq c_*I$.  Put
$a_\alpha=S^\alpha a$, $c_\alpha=S^\alpha c$,
$q_\alpha=D^\alpha q$, and $v_\alpha=D^\alpha v$, and define
\begin{equation}
 \mathscr E_{s,h}(X)=\frac12\sum_{|\alpha|\le s}
 \left[\langle v_\alpha,a_\alpha v_\alpha\rangle_h
 +\sum_{i,j}\langle D_i^+q_\alpha,c_\alpha^{ij}D_j^+q_\alpha\rangle_h
 +\|q_\alpha\|_h^2\right].
 \label{eq:supp-open-energy}
\end{equation}
It is uniformly equivalent to $\|X\|_{\mathcal X_h^s}^2$.

\begin{theorem}[All-component shifted nonlinear energy]
\label{thm:supp-open-energy}
For $s\ge11$, every solution of \eqref{eq:main-open-writer} remaining in the
fixed chart satisfies
\begin{equation}
 \frac{\dd}{\dd t}\mathscr E_{s,h}
 \le C_s\mathscr E_{s,h}+C_s\mathscr E_{s,h}^{3/2},
 \label{eq:supp-open-energy-inequality}
\end{equation}
with $C_s$ independent of $h$.
\end{theorem}
\begin{proof}
Apply $D^\alpha$ to the mass row.  The exact shifted product rule gives
\begin{equation}
 a_\alpha\dot v_\alpha
 =\sum_{i,j}D_i^-(c_\alpha^{ij}D_j^+q_\alpha)
 -\mathsf K_{b_\alpha}v_\alpha+\mathsf R_\alpha,
 \label{eq:supp-open-commuted-row}
\end{equation}
where every term in $\mathsf R_\alpha$ contains either a shifted commutator
or the quadratic first-jet source.  In particular the outer-divergence
commutator is retained, rather than treated as an undifferentiated error.
\Cref{lem:supp-open-calculus} and the analytic coefficient bounds imply
\begin{equation}
 \left(\sum_{|\alpha|\le s}\|\mathsf R_\alpha\|_h^2\right)^{1/2}
 \le C_s\|X\|_{\mathcal X_h^s}^2.
 \label{eq:supp-open-remainder}
\end{equation}
Pair \eqref{eq:supp-open-commuted-row} with $v_\alpha$.
Summation by parts cancels the divergence against the derivative of the
spatial energy, while
$\langle v_\alpha,\mathsf K_{b_\alpha}v_\alpha\rangle_h=0$.  Thus
\begin{equation}
 \begin{split}
 \dot{\mathscr E}_{s,h}=\sum_{|\alpha|\le s}\bigg[&
 \tfrac12\langle v_\alpha,\dot a_\alpha v_\alpha\rangle_h
 +\tfrac12\sum_{i,j}\langle D_i^+q_\alpha,
          \dot c_\alpha^{ij}D_j^+q_\alpha\rangle_h\\
 &+\langle q_\alpha,v_\alpha\rangle_h
 +\langle v_\alpha,\mathsf R_\alpha\rangle_h\bigg].
 \end{split}
 \label{eq:supp-open-energy-identity}
\end{equation}
The coefficient time derivatives are analytic multiples of $v$ and have
$\ell^\infty$ size $O(\|X\|_{\mathcal X_h^s})$.  Energy equivalence,
Cauchy--Schwarz and \eqref{eq:supp-open-remainder} prove the result.  The
translated coefficient in each differentiated row is essential: replacing
it by an unshifted continuum Leibniz rule would leave uncancelled top-order
terms.
\end{proof}

\begin{corollary}[Common nonlinear time and forced stability]
\label{cor:supp-open-lifespan}
There are $\varepsilon_s,T_s,C_s>0$, independent of $h$, such that every
$\|X_h(0)\|_{\mathcal X_h^s}\le\varepsilon_s$ has a unique solution through
$T_s$ satisfying \eqref{eq:main-open-uniform-time}.  If an acceleration defect
$f_h$ is added, then on the same smaller chart
\begin{equation}
 (\mathscr E_{s,h}^{1/2})'
 \le C_s\mathscr E_{s,h}^{1/2}+C_s\mathscr E_{s,h}
       +C_s\|f_h\|_{s,h}.
 \label{eq:supp-open-forced-energy}
\end{equation}
Two top-bounded records satisfy the corresponding lower-order difference
estimate
\begin{equation}
 \sup_{t\le T_s}\|X_h-Y_h\|_{\mathcal X_h^r}
 \le C_s\left(\|X_h(0)-Y_h(0)\|_{\mathcal X_h^r}
 +\int_0^{T_s}\|f_h-\widetilde f_h\|_{r,h}\,\dd t\right)
 \label{eq:supp-open-difference}
\end{equation}
for $3\le r\le s-1$.
\end{corollary}
\begin{proof}
The finite ODE is analytic on the coefficient chart.  A scalar comparison for
\eqref{eq:supp-open-energy-inequality} keeps its energy strictly inside a
smaller fixed ball on a time chosen before the cutoff; finite-dimensional
continuation then gives the common lifespan.  Adding the differentiated force
to the remainder proves \eqref{eq:supp-open-forced-energy}.
For the difference $X-Y=(q-p,v-w)=(\delta q,\delta v)$, write the principal
row with the coefficients of $q$.  Apart from $a(q)(f-\widetilde f)$,
the additional source is exactly
\begin{equation}
 \begin{split}
 \mathcal R_{X,Y}={}&
 \sum_{i,j}D_i^-\bigl((c^{ij}(q)-c^{ij}(p))D_j^+p\bigr)
 -\mathsf K_{b(q)-b(p)}w\\
 &+\mathsf G_h(q,v)-\mathsf G_h(p,w)
       -(a(q)-a(p))V_{0,h}(p,w),
 \end{split}
 \label{eq:supp-open-difference-source}
\end{equation}
where $V_{0,h}$ is the unforced acceleration.  Writing the mass difference
this way cancels the contribution of $\widetilde f$ inside $w_t$; it does
not require a pointwise bound on that force.  For $3\le r\le s-1$,
\[
 \|\mathcal R_{X,Y}\|_{r,h}
 \le C_s\|X-Y\|_{\mathcal X_h^r}.
\]
Indeed the flux term uses $p\in H_h^{r+2}$, the transport term uses
$w\in H_h^{r+1}$, and the mass term uses
$V_{0,h}(p,w)\in H_h^r$.  These are all supplied by the top ball and
$r\le s-1$.  Analytic coefficient differences are bounded by
$\|\delta q\|_{r+1,h}$; the first-jet source uses this norm and
$\|\delta v\|_{r,h}$.

Use the quadratic energy \eqref{eq:supp-open-energy} for
$(\delta q,\delta v)$ with the shifted coefficients of $q$.
The same principal cancellations and commutator bounds give
\[
 (\mathscr E_{{\rm diff},r}^{1/2})'
 \le C_s\mathscr E_{{\rm diff},r}^{1/2}
                         +C_s\|f-\widetilde f\|_{r,h}.
\]
Gronwall proves \eqref{eq:supp-open-difference}.  No top-order force has
been estimated by a bare inverse-mesh operator norm, so the constant has
no $e^{CT/h}$ factor.
\end{proof}

\subsection{Compatible initial data, consistency, and subsidiary propagation}

All Sobolev norms in this subsection are spatial norms on the unit torus;
$C_tH^r$ denotes the supremum over the common time interval.  Constants may
depend on $s$ and the fixed coefficient chart, but not on $h$.  In
particular, interpolation is not assumed to commute with multiplication or
with the nonlinear coefficient maps.

\begin{lemma}[Quantitative sampling and nonlinear interpolation]
\label{lem:supp-open-interpolation}
Let $r\ge2$ be an integer, let $\mathcal P_h=\mathcal I_h\mathcal S_h$,
and let $u_h,w_h$ be grid arrays.  On the odd frequency cube,
\begin{align}
 \|\mathcal P_h u-u\|_{H^r}
   &\le C_rh\|u\|_{H^{r+1}},
   &\|\mathcal P_hu\|_{H^r}&\le C_r\|u\|_{H^r},
   \label{eq:supp-open-sampling-estimate}\\
 \|\mathcal I_hD_i^\pm u_h-\partial_i\mathcal I_hu_h\|_{H^r}
   &\le C_rh\|u_h\|_{r+2,h},
   \label{eq:supp-open-difference-consistency}\\
 \|\mathcal I_h(u_hw_h)-(\mathcal I_hu_h)(\mathcal I_hw_h)\|_{H^r}
   &\le C_rh\|u_h\|_{r+1,h}\|w_h\|_{r+1,h}.
   \label{eq:supp-open-product-consistency}
\end{align}
The first estimate also holds with $h$ replaced by $h^j$ and $H^{r+1}$
by $H^{r+j}$ for every fixed positive integer $j$.  If $A$ is an analytic
coefficient map on the common small chart, then
\begin{equation}
 \|\mathcal I_h A(u_h)-A(\mathcal I_hu_h)\|_{H^r}
 \le C_{r,A}h\|u_h-u_*\|_{r+1,h}
 \label{eq:supp-open-composition-consistency}
\end{equation}
for a constant base point $u_*$.  These estimates admit time-differentiated
versions by the ordinary chain rule, with one spatial order reserved for
each controlled time derivative.
\end{lemma}
\begin{proof}
For $u$ with Fourier coefficients $\widehat u(k)$, the coefficient of its
interpolant at $k\in\{-m,\ldots,m\}^3$ is
$\sum_{\ell\in\mathbb Z^3}\widehat u(k+N\ell)$.  The omitted high-frequency
tail has $H^r$ norm at most $C N^{-j}\|u\|_{H^{r+j}}$.  For the folded
part, apply Cauchy--Schwarz in $\ell\ne0$ with weights
$\langle k+N\ell\rangle^{r+j}$.  Uniformly in the retained $k$,
\[
 \sum_{\ell\ne0}\langle k+N\ell\rangle^{-2(r+j)}
 \le C_r N^{-2(r+j)}.
\]
Multiplication by $\langle k\rangle^{2r}\le C_rN^{2r}$ and summation in
$k$ prove the error estimate.  The same argument with $j=0$, together with
the retained coefficients, proves stability since $r>3/2$.
The symbol estimate
$|(e^{\pm ih\xi}-1)/(\pm h)-i\xi|\le Ch|\xi|^2$
proves \eqref{eq:supp-open-difference-consistency}; averaging gives the
corresponding estimate for $D_i^0$.

Since $\mathcal I_h(u_hw_h)=\mathcal P_h[(\mathcal I_hu_h)
(\mathcal I_hw_h)]$, the product estimate follows from sampling consistency
and the $H^{r+1}$ product inequality.  For composition use
$\mathcal I_h A(u_h)=\mathcal P_h[A(\mathcal I_hu_h)]$ and the analytic
composition bound on the smaller common chart.  Constant terms are
interpolated exactly.  Differentiating these identities, rather than
postulating a discrete chain rule, gives the last assertion.
\end{proof}

\begin{lemma}[Uniform physical time jets]
\label{lem:supp-open-time-jets}
An unforced solution on the common top ball satisfies
\begin{equation}
 \|q_h\|_{C_tH_h^{s+1}}+\|q_{h,t}\|_{C_tH_h^s}
 +\|q_{h,tt}\|_{C_tH_h^{s-1}}
 +\|q_{h,ttt}\|_{C_tH_h^{s-2}}\le C_s\varepsilon.
 \label{eq:supp-open-time-jet-bound}
\end{equation}
The same bound holds uniformly for the coefficient family
\eqref{eq:main-open-law-family} with $B$ constant during the history.
\end{lemma}
\begin{proof}
The first two terms follow from the nonlinear energy.  The divergence row
maps $H_h^{s+1}$ to $H_h^{s-1}$, the skew transport maps $H_h^s$ to
$H_h^{s-1}$, and the first-jet source is controlled by the product calculus.
Multiplication by $a(g_h)^{-1}$ preserves that space.  Differentiating this
same equation once replaces $q_h$ by $q_{h,t}$ in the principal spatial
term and $v_h$ by $q_{h,tt}$ in the transport term.  Every coefficient time
derivative is an analytic multiple of $v_h$.  The resulting terms are
bounded in $H_h^{s-2}$.  For the enlarged family,
$h^2\Lambda_h^2:H_h^{r+2}\to H_h^r$ is uniformly bounded, by
$h^2\ell_h\le12$ in Fourier variables; apply this first to $q_h$ and then
to $v_h$.  No time derivative of $B$ occurs.
\end{proof}

For $(\gamma,K)\in\mathcal D_s(\varepsilon)$ choose
\begin{equation}
 \begin{gathered}
 g_{00}^0=-1,\qquad g_{0i}^0=0,\qquad g_{ij}^0=\gamma_{ij},\\
 v_{ij}^0=-2K_{ij},\qquad v_{00}^0=2\tr_\gamma K,\qquad
 v_{0i}^0=\gamma^{jk}(\partial_j\gamma_{ki}
                         -\tfrac12\partial_i\gamma_{jk}).
 \end{gathered}
 \label{eq:supp-open-harmonic-initial}
\end{equation}
Initialize $q_h(0)=\mathcal S_h(g^0-\eta)$ and
$v_h(0)=\mathcal S_hv^0$.  These are samples of initial data only, not of a
future Einstein solution.  Finite rounding of this preparation is addressed
in \Cref{thm:supp-open-finite-words}.

Use the lower-index gauge covector
$c_b(g)=g_{b\mu}g^{\alpha\beta}\Gamma^\mu_{\alpha\beta}$ and set
\begin{equation}
 \mathfrak H_{ab}(g,c)=\nabla_{(a}c_{b)}
                   -\tfrac12g_{ab}\nabla^dc_d,
 \qquad r_{h,ab}=G_{ab}(g_h)-\mathfrak H_{ab}(g_h,c(g_h)),
 \quad g_h=\eta+\mathcal I_hq_h.
 \label{eq:supp-open-reduced-residual}
\end{equation}
Thus $r_h$ is the actual continuum reduced residual of the interpolant.
The conversion from the normal row
\eqref{eq:main-harmonic-normal-row} to $r_h$ is a bounded analytic
pointwise trace-reversal map on the common chart.  This is the harmonic
reduction convention used here; compare \cite[Sec.~2.1]{Lindblom2006}.

\begin{lemma}[Consistency of the complete reduced residual]
\label{lem:supp-open-reduced-consistency}
For the unforced central writer,
\begin{equation}
 \|r_h\|_{C_tH^{s-2}}+
 \|\partial_tr_h\|_{C_tH^{s-3}}\le C_sh\varepsilon.
 \label{eq:supp-open-reduced-rate}
\end{equation}
For the $B$-writer the same estimate holds with a uniform constant on
$\overline{\mathcal B}$.  In particular the residual includes all
interpolation, mass, flux, transport and source errors.
\end{lemma}
\begin{proof}
Write $\widetilde q=\mathcal I_hq_h$, $\widetilde v=\mathcal I_hv_h$.
For a representative spatial flux insert the two intermediate quantities
$\partial_i\mathcal I_h(c_h^{ij}D_j^+q_h)$ and
$\partial_i[(\mathcal I_hc_h^{ij})(\mathcal I_hD_j^+q_h)]$.
Their differences from the discrete flux and the continuum flux are bounded
in $H^{s-2}$, respectively, by
\[
 Ch\|c_h^{ij}D_j^+q_h\|_{H_h^s},\qquad
 Ch\|c_h^{ij}\|_{H_h^s}\|D_j^+q_h\|_{H_h^s},\qquad
 Ch\|q_h\|_{H_h^{s+1}},
\]
with constant coefficient factors separated when necessary.  The last term
uses composition consistency and
$\|\mathcal I_hD_j^+q_h-\partial_j\widetilde q\|_{H^{s-1}}
\le Ch\|q_h\|_{H_h^{s+1}}$.
The continuum expansion of the flux contains
$(\partial_ic^{ij})\partial_j\widetilde q$.  The expansion of the skew
transport contains $(\partial_ib^i)\widetilde v$.  The two compensator
terms cancel precisely these contributions, with the signs in
\eqref{eq:main-open-compensator}.  Their remaining discrepancies are
estimated by \Cref{lem:supp-open-interpolation}.

For the mass term one must also retain
\[
 \mathcal I_h(a_hv_{h,t})-a(g_h)\mathcal I_hv_{h,t}.
\]
Its $H^{s-2}$ norm is $O(h\varepsilon)$ because
$v_{h,t}$ is bounded in $H_h^{s-1}$.  The two skew-transport terms use
$v_h\in H_h^s$; the analytic first-jet source is controlled in the same
way.  This proves the undifferentiated normal residual bound.

Differentiate these decompositions once in physical time.  The highest
spatially differentiated unknown is now $v_h\in H_h^s$, and the highest
mass jet is $v_{h,tt}\in H_h^{s-2}$.  The identical estimates therefore
apply in $H^{s-3}$.  Trace reversal and its time derivative are uniformly
bounded multipliers, proving \eqref{eq:supp-open-reduced-rate}.  Finally,
for the enlarged writer, the extra residual and its time derivative satisfy
\[
 h^2\|B\Lambda_h^2q_h\|_{H_h^{s-2}}
 +h^2\|B\Lambda_h^2v_h\|_{H_h^{s-3}}
 \le C_sh\|B\|\varepsilon,
\]
by the multiplier bounds.  This proves uniformity for the family.
\end{proof}

\begin{lemma}[Initial harmonic and normal-constraint errors]
\label{lem:supp-open-initial-gauge}
For \eqref{eq:supp-open-harmonic-initial} and its sampled preparation,
\begin{equation}
 \|c(g_h)(0)\|_{H^{s-2}}
 +\|\partial_tc(g_h)(0)\|_{H^{s-3}}
 \le C_sh\varepsilon.
 \label{eq:supp-open-initial-gauge-bound}
\end{equation}
The estimate uses the actual finite acceleration, not an assigned
continuum acceleration.
\end{lemma}
\begin{proof}
At a unit-lapse, zero-shift initial slice, direct contraction gives
\[
 c^0=\tfrac12(v_{00}+\gamma^{ij}v_{ij}),\qquad
 c^i=-\gamma^{ij}v_{0j}
       +\gamma^{jk}\Gamma^i_{jk}(\gamma).
\]
The unrounded preparation therefore has $c=0$ exactly.  Sampling stability,
product consistency and the one-derivative definition of $c$ give
$\|c(g_h)(0)\|_{H^{s-2}}\le Ch\varepsilon$.

The Gauss--Codazzi normal components $G_{0b}$ at this slice depend only on
the spatial metric, its first two spatial derivatives and its extrinsic
curvature; they contain no metric acceleration.  The vacuum constraint
identities in \eqref{eq:main-open-data-class} and sampling consistency give
$\|G_{0b}(g_h)(0)\|_{H^{s-3}}\le Ch\varepsilon$.
In \eqref{eq:supp-open-reduced-residual}, the coefficient of
$\partial_tc_b$ in $\mathfrak H_{0b}$ is $1/2$ on this slice.  All its
other terms involve $c$, tangential derivatives of $c$, and bounded
first-jet coefficients.  Consequently
\[
 \|\partial_tc(0)\|_{H^{s-3}}
 \le C\bigl(\|G_{0\cdot}(0)\|_{H^{s-3}}
       +\|r_{0\cdot}(0)\|_{H^{s-3}}
       +\|c(0)\|_{H^{s-2}}\bigr).
\]
The residual lemma proves the assertion.  This calculation explains exactly
where the four initial constraints, in addition to $c(0)=0$, are used.
\end{proof}

\begin{lemma}[Forced subsidiary energy on the common interval]
\label{lem:supp-open-subsidiary}
For the harmonic initialization,
\begin{equation}
 \sup_{t\le T_s}\bigl(
 \|c(g_h)\|_{H^{s-2}}+\|\partial_tc(g_h)\|_{H^{s-3}}
 +\|G(g_h)\|_{H^{s-3}}\bigr)\le C_sh\varepsilon.
 \label{eq:supp-open-subsidiary}
\end{equation}
\end{lemma}
\begin{proof}
The contracted Bianchi identity applied to the actual smooth interpolant
and \eqref{eq:supp-open-reduced-residual} gives the exact equation
\begin{equation}
 \Box_{g_h}c_b+\Ric_b{}^d(g_h)c_d=-2\nabla^ar_{h,ab}.
 \label{eq:supp-open-exact-subsidiary}
\end{equation}
There is no additional unspecified discretization term: all such terms are
already in $r_h$.  The identity follows by commuting the two covariant
derivatives in $\nabla^a\mathfrak H_{ab}$; it is the forced form of the
usual harmonic subsidiary equation \cite[Sec.~2.1]{Lindblom2006}.

For $k=s-3$ use the positive coordinate wave energy
\[
 \mathcal E_{c,k}=\frac12\sum_{|\alpha|\le k}\int_{\mathbb T^3}
 \left(a|\partial_t\partial^\alpha c|^2
 +c^{ij}\partial_i\partial^\alpha c\cdot
             \partial_j\partial^\alpha c
 +|\partial^\alpha c|^2\right)\dd x.
\]
Here the index contractions in the energy are Euclidean positive
contractions on the four components, and $a,c^{ij}$ are the hyperbolic
principal coefficients.  The top metric and time-jet bounds give uniform
coercivity and uniform coefficient, commutator and Ricci-product bounds.
Pairing each commuted equation with
$\partial_t\partial^\alpha c$, integrating the spatial principal term by
parts, and estimating the mixed transport and lower-order terms gives
\begin{equation}
 (\mathcal E_{c,k}^{1/2})'
 \le C_s\mathcal E_{c,k}^{1/2}
       +C_s\bigl(\|r_h\|_{H^{k+1}}
                         +\|\partial_tr_h\|_{H^k}\bigr).
 \label{eq:supp-open-subsidiary-energy}
\end{equation}
In particular the differentiated source is not bounded by the
undifferentiated Euler residual alone.  The initial-gauge lemma and
\eqref{eq:supp-open-reduced-rate} give the first two terms of
\eqref{eq:supp-open-subsidiary} by Gronwall.  Its last term follows from
$G=\mathfrak H+r$.  The common nonlinear existence interval was established
before this subsidiary estimate and is not a constraint-preserving
bootstrap assumption.
\end{proof}

\begin{lemma}[Lower-order Cauchy comparison and acceleration recovery]
\label{lem:supp-open-continuum-comparison}
Let $g$ be a solution of the continuum harmonic normal row on the common
top ball and let $X_h$ solve the central writer.  Set
$Y_h=\mathcal S_h(g-\eta,\partial_tg)$ and
$e_{0,h}=\|X_h(0)-Y_h(0)\|_{\mathcal X_h^{s-2}}$.
Then
\begin{equation}
 \begin{split}
 &\|q_h-\mathcal S_h(g-\eta)\|_{C_tH_h^{s-1}}
 +\|v_h-\mathcal S_h\partial_tg\|_{C_tH_h^{s-2}}\\
 &\qquad
 +\|q_{h,tt}-\mathcal S_h\partial_t^2g\|_{C_tH_h^{s-3}}
 \le C_s(e_{0,h}+h\varepsilon).
 \end{split}
 \label{eq:supp-open-sampled-comparison}
\end{equation}
\end{lemma}
\begin{proof}
Sampling the continuum normal row gives the finite row with an acceleration
force of $C_tH_h^{s-2}$ norm at most $Ch\varepsilon$.
This follows from the same flux, transport and composition decompositions
as in \Cref{lem:supp-open-reduced-consistency}, now applied to a sampled
continuum field.  The three derivatives used by the second-order spatial
consistency error are available from $g-\eta\in H^{s+1}$.
Apply \eqref{eq:supp-open-difference} with $r=s-2$.  This gives the first
two terms without an inverse-mesh exponential.  For the last term subtract
the mass rows and divide by $a(g_h)$: two spatial derivatives of the
$H_h^{s-1}$ position difference and one of the $H_h^{s-2}$ velocity
difference lie in $H_h^{s-3}$; coefficient differences multiply uniformly
bounded continuum jets.  The sampled residual is smaller than the same
right-hand side.  This proves the acceleration estimate independently of
any differentiation of a convergence statement.
\end{proof}

\begin{theorem}[Whole-sequence nonsymmetric vacuum limit]
\label{thm:supp-open-einstein}
For sufficiently small $(\gamma,K)\in\mathcal D_s(\varepsilon)$, the
complete sequence of interpolated solutions converges on
$[0,T_s]\times\mathbb T^3$ to the unique harmonic vacuum development with
those initial data, and
\begin{equation}
 G(g)=0.
 \label{eq:supp-open-vacuum}
\end{equation}
Moreover,
\begin{equation}
 \sup_{t\le T_s}\left(
 \|g_h-g\|_{H^{s-1}}+\|\partial_tg_h-\partial_tg\|_{H^{s-2}}
 +\|\partial_t^2g_h-\partial_t^2g\|_{H^{s-3}}\right)
 \le C_sh\varepsilon,
 \label{eq:supp-open-metric-rate}
\end{equation}
and
\begin{equation}
 \|\operatorname{Riem}(g_h)-\operatorname{Riem}(g)\|_{C_tH^{s-3}}
 \le C_sh\varepsilon.
 \label{eq:supp-open-curvature-rate}
\end{equation}
\end{theorem}
\begin{proof}
The top energy and \Cref{lem:supp-open-time-jets} give compactness in
$C_tH^{s+1-\delta}\times C_tH^{s-\delta}$ for $0<\delta<1$.
Indeed, a fixed finite spatial Fourier screen has uniformly Lipschitz
coefficients, while its complementary tails are bounded by the top norms;
interpolation between adjacent Sobolev orders gives uniform time
continuity at the displayed orders.  Arzel\`a--Ascoli on each screen and a
diagonal extraction produce a cluster point, retaining the top norms by
weak lower semicontinuity.  The time-jet bounds identify its velocity and
acceleration distributionally.

The complete residual lemma passes the finite row to the continuum
harmonic reduced equation.  The subsidiary lemma gives $c(g)=0$, hence
$G(g)=0$.  Sampling consistency identifies its initial data.  To see
uniqueness, subtract two continuum reduced equations with these initial
data.  The coordinate wave energy at order $s-2$ controls the position in
$H^{s-1}$ and velocity in $H^{s-2}$; coefficient differences multiply the
bounded top jets of the other solution.  The product estimate and
integration by parts give $E_{\rm diff}^{1/2}(t)\le
C\int_0^t E_{\rm diff}^{1/2}$, so the difference vanishes.  Therefore all
cluster points coincide and the whole cutoff sequence converges.

Only after constructing this limit do we sample it for an error estimate.
Its initial samples agree with the initialization, so $e_{0,h}=0$ in
\Cref{lem:supp-open-continuum-comparison}.  Add
\eqref{eq:supp-open-sampling-estimate} for $g$, $\partial_tg$ and
$\partial_t^2g$ to obtain \eqref{eq:supp-open-metric-rate}.
On a nondegenerate chart, curvature is a smooth expression in the inverse
metric, its first derivatives and its full second jet.  The three slots in
\eqref{eq:supp-open-metric-rate} control all second space--time derivatives
in $H^{s-3}$.  The product algebra at this index proves
\eqref{eq:supp-open-curvature-rate}.  No future continuum metric enters the
finite construction; it is used here only for the error comparison.
\end{proof}

\begin{corollary}[Direct constructive Einstein limit]
\label{cor:supp-open-direct-route}
The Einstein conclusion of \Cref{thm:supp-open-einstein} follows from the
cutoff-uniform nonlinear energy, consistency of the declared local writer,
harmonic subsidiary estimate, and uniqueness of the limiting reduced Cauchy
problem.  No item \textup{(E1)--(E5)} of
\Cref{thm:main-renewal-einstein} is invoked as an independent hypothesis.
Accordingly those five items function as sufficient certification conditions
for the broader class of complete renewal cylinders, while
\Cref{thm:main-open-3plus1} supplies an independently constructed
non-symmetry-reduced positive population.
\end{corollary}

\subsection{Positive rate lift and finite renewal words}

The small metric chart has an explicit positive $A_3$ rate lift.  In the
$D_3$ realization
\begin{equation}
 r_a\in\{(1,1,0),(1,-1,0),(1,0,1),(1,0,-1),(0,1,1),(0,1,-1)\},
 \label{eq:supp-open-d3-roots}
\end{equation}
put $B=N_{\rm ADM}^2\gamma^{-1}$, $\varrho=\sqrt{\det\gamma}$ and, for a
coordinate pair $ij$ with remaining index $k$,
\begin{equation}
 c_{ij,\pm}=\frac{B_{ii}+B_{jj}-B_{kk}}4\pm\frac{B_{ij}}2,
 \qquad d_a=\frac14r_a\cdot\beta,\qquad
 k_a^\pm=\frac{c_a}{2h^2}\pm\frac{d_a}{2h}.
 \label{eq:supp-open-rate-lift}
\end{equation}
At the flat point all $c_a=1/4$, so a fixed small metric chart makes all
rates positive and comparable to $h^{-2}$.  Direct calculation gives
\begin{equation}
 h^2\sum_a(k_a^++k_a^-)r_ar_a^T=B,
 \qquad h\sum_a(k_a^+-k_a^-)r_a=\beta,
 \label{eq:supp-open-rate-inversion}
\end{equation}
so the ADM reconstruction returns the same metric.  The resulting weighted
periodic graph has a cutoff-uniform cut margin by perturbation of the flat
$D_3$ regulator and \Cref{thm:main-spatial-screen}.

To obtain finite histories, use a fixed-degree local Taylor--Hermite
construction.  The construction below makes both the spatial derivative
loss and the precision convention explicit.  It applies to the central
writer and to every fixed $B\in\overline{\mathcal B}$; when the coefficient
is selected operationally it is retained as a finite initial mark.

\begin{lemma}[Inverse-mesh time jets for a local finite writer]
\label{lem:supp-open-local-jets}
Let $\mathcal F_{B,h}(q,v)=(v,V_{B,h}(q,v))$ denote the differential writer
on the common $\mathcal X_h^s$ ball.  For each fixed $1\le j\le9$,
\begin{equation}
 \|D^j\mathcal F_{B,h}(X)\|_{(\mathcal X_h^s)^j\to\mathcal X_h^s}
 \le C_{s,j}h^{-1},\qquad
 \|\mathcal F_{B,h}(X)\|_{\mathcal X_h^s}
 \le C_sh^{-1}\|X\|_{\mathcal X_h^s}.
 \label{eq:supp-open-local-vector-bounds}
\end{equation}
Consequently an exact local solution from a source of size at most
$\varepsilon$ has
\begin{equation}
 \|\partial_t^j q\|_{H_h^s}\le C_{s,j}\varepsilon h^{1-j}
 \quad(1\le j\le8).
 \label{eq:supp-open-inverse-time-jets}
\end{equation}
Each initial jet through order seven is a function of a fixed-radius
neighbourhood of the source, independent of $h$.
\end{lemma}
\begin{proof}
The inverse inequality
$\|u\|_{r+1,h}\le C_rh^{-1}\|u\|_{r,h}$ follows from the bounded
frequency cube.  The acceleration has two spatial derivatives of $q$ and
one of $v$, whereas the phase norm contains one extra derivative of $q$.
Thus it loses at most one order in that norm.  The product calculus bounds
all analytic coefficient derivatives with the same single loss.  The added
term satisfies
$h^2\|\Lambda_h^2q\|_{H_h^s}\le Ch^{-1}\|q\|_{H_h^{s+1}}$,
so the bounds are uniform in $B$.  Define successive vector jets by
$J_1=\mathcal F_{B,h}$ and $J_{j+1}=DJ_j\mathcal F_{B,h}$.
Repeated differentiation gives
$\|J_j(X)\|_{\mathcal X_h^s}\le C_{s,j}\varepsilon h^{-j}$.
For $j=1$ the assertion is $q_t=v$; for $j\ge2$ the $q$-jet is
the $v$-component of $J_{j-1}$. This proves
\eqref{eq:supp-open-inverse-time-jets}.  Each differentiation composes a
fixed local stencil a finite number of times; it does not enlarge its
radius with the cutoff.  These estimates deliberately retain inverse-mesh
powers and do not assume additional continuum regularity.
\end{proof}

At a source $X=(q,v)$, let $T_X$ be the degree-seven Taylor polynomial of
the position, using the jets of the actual differential writer at that
source.  For $0\le t\le\tau$, set
\begin{equation}
 T_X(t)=q+\sum_{j=1}^7\frac{t^j}{j!}q^{(j)}(0),\qquad
 X_1=(T_X(\tau),T_X'(\tau)),\qquad
 J_{B,h}(X)=D V_{B,h}(X)[\mathcal F_{B,h}(X)].
 \label{eq:supp-open-local-taylor}
\end{equation}
The endpoint-corrected polynomial $Q_X$ is the unique polynomial of degree
at most seven whose value, velocity, acceleration and jerk at zero are
$(q,v,V_{B,h}(X),J_{B,h}(X))$, and whose corresponding jets at $\tau$
are $(X_1,V_{B,h}(X_1),J_{B,h}(X_1))$.

\begin{lemma}[Complete Taylor--Hermite residual]
\label{lem:supp-open-hermite-residual}
For $c_-h^2\le\tau\le c_+h^2$ and sufficiently small $h$, the polynomial
$Q_X$ remains in the common chart and its actual defect
$f_X=Q_X''-V_{B,h}(Q_X,Q_X')$ satisfies
\begin{equation}
 \|f_X\|_{C(I;H_h^s)}\le C_s\varepsilon h^5,
 \qquad
 \|\partial_tf_X\|_{C(I;H_h^{s-3})}\le C_s\varepsilon h^3.
 \label{eq:supp-open-hermite-local-defect}
\end{equation}
These bounds concern the complete interval, not only collocation points.
\end{lemma}
\begin{proof}
Compare $T_X$ with the exact local solution from $X$, used only in this
proof.  Taylor's integral remainder and
\eqref{eq:supp-open-inverse-time-jets} give
\begin{equation}
 \|T_X^{(j)}-q^{(j)}\|_{C(I;H_h^s)}
 \le C_s\varepsilon h^{-7}\tau^{8-j}\quad(0\le j\le3).
 \label{eq:supp-open-taylor-remainder}
\end{equation}
The position error in $H_h^{s+1}$ costs one inverse power of $h$.
Together with the velocity error, this gives a phase error
$O(\varepsilon h^7)$.  The acceleration Lipschitz bound in
\eqref{eq:supp-open-local-vector-bounds} then yields
$\|T_X''-V_{B,h}(T_X,T_X')\|_{H_h^s}\le C\varepsilon h^5$.
Differentiating this defect once gives a bound $C\varepsilon h^3$ in
$H_h^s$: the third Taylor remainder is $O(\varepsilon h^3)$,
$DJ_{B,h}$ has norm at most $Ch^{-2}$ from $\mathcal X_h^s$ to $H_h^s$,
and the remaining term is $D_vV_{B,h}$ applied to the acceleration defect.

Put $\delta_2=V_{B,h}(X_1)-T_X''(\tau)$ and
$\delta_3=J_{B,h}(X_1)-T_X'''(\tau)$.  The preceding bounds give
$\|\delta_2\|_{H_h^s}\le C\varepsilon h^5$ and
$\|\delta_3\|_{H_h^s}\le C\varepsilon h^3$.
The correction is explicitly
\begin{equation}
 Q_X(t)-T_X(t)=\tau^2\delta_2 H_2(t/\tau)
                         +\tau^3\delta_3 H_3(t/\tau),
 \quad
 \begin{cases}
 H_2(z)=-\tfrac12z^4(z-1)^2(4z-5),\\
 H_3(z)=\tfrac16z^4(z-1)^3.
 \end{cases}
 \label{eq:supp-open-hermite-correction}
\end{equation}
Their jets through order three vanish at zero, and their endpoint jets
are the unit acceleration and unit jerk jets, respectively.  Existence
and uniqueness also follow because a degree-seven polynomial with zero
jets through order three at both endpoints is divisible by
$t^4(t-\tau)^4$ and must vanish.
The correction has second derivative $O(\varepsilon h^5)$, third
derivative $O(\varepsilon h^3)$ in $H_h^s$, and phase norm
$O(\varepsilon h^7)$.  The same local Lipschitz estimates therefore prove
\eqref{eq:supp-open-hermite-local-defect}, even at the stronger spatial
order for its differentiated local term.  Only $H_h^{s-3}$ is retained in
the asserted physical residual topology.
\end{proof}

\begin{lemma}[Sobolev-scaled finite precision]
\label{lem:supp-open-hermite-precision}
Let $Q_X$ be the local polynomial of
\Cref{lem:supp-open-hermite-residual}, and let $\widetilde Q_X$ be a
degree-seven endpoint Hermite perturbation. Set $\delta Q=\widetilde Q_X-Q_X$ and let
\begin{equation}
 b_h=\max_{\substack{e\in\{0,1\}\\0\le j\le3}}
       \tau^j\|\delta Q^{(j)}(e\tau)\|_{H_h^s}.
 \label{eq:supp-open-hermite-precision}
\end{equation}
Then $\|\delta Q^{(j)}\|_{C(I;H_h^s)}\le C\tau^{-j}b_h$
for $0\le j\le3$. If $b_h\le c\varepsilon\tau^2h^5$ and the
source lies in a smaller common chart, the perturbation remains in the
local phase tube. The changes of the complete residual and its time
derivative are bounded by $C\tau^{-2}b_h$ and $C\tau^{-3}b_h$,
respectively. Hence the bounds \eqref{eq:supp-open-hermite-local-defect}
are preserved.
\end{lemma}
\begin{proof}
The Hermite map on $[0,1]$ is a fixed finite-dimensional linear inverse.
Rescaling its input jets and its derivatives proves the first estimate.
For the nonlinear terms use
$\|\delta(Q,Q')\|_{\mathcal X_h^s}
\le C(h^{-1}+\tau^{-1})b_h$ and
\eqref{eq:supp-open-local-vector-bounds}.  Since $\tau/h=O(h)$,
these terms are bounded by $C\tau^{-2}b_h$ in the residual.  Differentiating
once, write $\partial_tf=Q^{(3)}-J_{B,h}(Q,Q')
-D_vV_{B,h}(Q,Q')f$.
The local bounds for $DJ_{B,h}$, $D_vV_{B,h}$ and its differential,
together with $\|f_X\|_{H_h^s}\le C\varepsilon h^5$, give
$C\tau^{-3}b_h$. The quadratic perturbation term is absorbed using
the stated smallness of $b_h$. These estimates are applied to a
perturbation of the controlled local polynomial, not to arbitrary
polynomials having only bounded position and velocity.
\end{proof}

\begin{theorem}[Finite local successor words]
\label{thm:supp-open-finite-words}
Fix $T<T_s$.  There is a finite chronological Taylor--Hermite realization
with $\tau_j\asymp h^2$, fixed execution radius and shared endpoint jets
through order three, such that
\begin{equation}
 \sup_{t\le T}\|(Q_h,Q_h')\|_{\mathcal X_h^s}\le C_s\varepsilon,
 \quad \|f_h^{\rm num}\|_{C_tH_h^s}\le C_s\varepsilon h^5,
 \quad \|\partial_tf_h^{\rm num}\|_{C_tH_h^{s-3}}
                         \le C_s\varepsilon h^3.
 \label{eq:supp-open-finite-defects}
\end{equation}
In particular the weaker $O(\varepsilon h^4)$ and
$O(\varepsilon h^2)$ residual bounds in the main text hold.
A sufficient precision convention is to approximate each scaled endpoint
jet $\tau_j^\ell Q^{(\ell)}$ pointwise to error at most
$c_s\varepsilon h^{s+10}$ relative to the local construction, using finer
intermediate accuracy as needed.  The starting finite record may differ
from the compatible preparation by $O(\varepsilon h^4)$ in
$\mathcal X_h^s$.  Every resulting history has the same Einstein limit.
\end{theorem}
\begin{proof}
At each source compute the finite jet bank
\eqref{eq:supp-open-local-taylor}, round its successor value and velocity,
and attach finite approximations to the acceleration and jerk of that
successor.  Reuse these latter caches at the next source.  Interpolate the
two endpoint banks by the Hermite polynomial.  The field and all the jet
maps are finite local analytic maps on bounded source sets, so their finite
arithmetic tolerances can be chosen to meet the stated endpoint error.
At fixed $h$ all registers may be chosen from bounded rational grids; with
a finite horizon this gives a genuinely finite alphabet.

A pointwise error $\eta$ on the unit-volume grid obeys
$\|\delta u\|_{H_h^s}\le C_sh^{-s}\eta$.  The stated scaled precision
therefore makes $b_h\le C_s\varepsilon h^{10}$, which is smaller than
$\varepsilon\tau_j^2h^5\asymp\varepsilon h^9$ on all sufficiently fine
grids.  Source cache and successor rounding errors are both included in
this endpoint comparison.  The two preceding lemmas give the residual
bounds on every interval.  Shared jets make $Q_h$ globally $C^3$ in time;
hence $f_h^{\rm num}$ is $C^1$, with no uncounted delta function at a cut.
The radius is bounded by the fixed number of local jet compositions and
endpoint evaluations, not by the number of sites or time steps.

Until a putative first guard exit, the forced energy bounds the cumulative
forcing by $C\varepsilon Th^5$.  After decreasing the initial radius,
this cannot close the strict guard margin, so the bound holds on the whole
interval.  The lower-order difference estimate gives position and velocity
shadowing, and the actual acceleration equation gives the second time-jet
estimate.  Combining with \Cref{thm:supp-open-einstein,thm:supp-law-einstein}
proves the claimed limit.  Moreover the physical stage cost is
$c_{h,j}\le C\tau_j(h^{10}+h^6)$, strictly below $M\tau_jh^4$ for any
fixed $M>0$ on a sufficiently fine tail.  This supplies the strict margin
used for nonempty probability-row neighborhoods.  The construction concerns
finite information range and precision; it makes no hardware-latency claim.
\end{proof}

The results above prove the population statement required in the main text.
They do not identify the local writer with the exact finite Euler equations of
the reconstructed action.  That stronger question has a different finite
constraint and response structure, to which we now turn.


\section{Law-space robustness, finite preparations and chronological compatibility}
\label{supp:open-law-basin}

This appendix proves \Cref{thm:main-open-law-basin} and
\Cref{thm:main-same-cylinder-einstein}. The cutoff-uniform estimates of
Appendix~\ref{supp:open-3plus1} are the input: first the law family is
enlarged, then finite initial constraints are prepared, and finally the
chronological coupling is compared with the independently reconstructed
action without identifying the two finite dynamics.

\subsection{A fixed open family of local differential writers}

Put
\begin{equation}
 \Lambda_h=-\sum_{i=1}^3D_i^-D_i^+,
 \qquad
 \ell_h(k)=4h^{-2}\sum_i\sin^2(\pi hk_i),
 \qquad0\le h^2\ell_h(k)\le12.
 \label{eq:supp-law-laplacian}
\end{equation}
On the metric chart of Appendix~\ref{supp:open-3plus1}, let
$c\succeq c_*I$.  Choose $0<b_0<c_*/24$ and set
\begin{equation}
 \mathcal B=\{B=B^T\in\mathbb R^{10\times10}:\|B\|_{\op}<b_0\}.
 \label{eq:supp-law-ball}
\end{equation}
This is an open ball of dimension $55$.  For $B\in\mathcal B$ define
\begin{equation}
 q_t=v,\qquad
 a(g)v_t=\sum_{i,j}D_i^-\!\bigl(c^{ij}(g)D_j^+q\bigr)
       -\mathsf K_bv+\mathsf G_h(q,v)-h^2B\Lambda_h^2q.
 \label{eq:supp-law-family}
\end{equation}
The coefficient $B$ is constant during one physical history.  It may vary
with the cutoff or be sampled once as a finite initial mark.

\begin{lemma}[Uniform fourth-order multiplier bounds]
\label{lem:supp-law-multiplier}
For every array $w$,
\begin{equation}
 h^2\|\Lambda_hw\|_h^2\le12\sum_i\|D_i^+w\|_h^2.
 \label{eq:supp-law-domination}
\end{equation}
At every fixed index used below,
\begin{align}
 h^2\|\Lambda_h^2q\|_{s-1,h}&\le C_s\|q\|_{s+1,h},&
 h^2\|\Lambda_h^2q\|_{s-2,h}&\le C_sh\|q\|_{s+1,h},\\
 h^2\|\Lambda_h^2q\|_{s-3,h}&\le C_sh^2\|q\|_{s+1,h},&
 h^2\|\Lambda_h^2v\|_{s-3,h}&\le C_sh\|v\|_{s,h}.
 \label{eq:supp-law-orders}
\end{align}
\end{lemma}
\begin{proof}
Parseval reduces the first claim to
$h^2\ell_h(k)^2\le12\ell_h(k)$.  Uniform equivalence of the forward-difference
Sobolev norm with the Fourier weight $(1+\ell_h)^{r/2}$ reduces the four
remaining claims to the bounded ratios
\[
 \frac{h^2\ell_h^2}{1+\ell_h},\quad
 \frac{h^2\ell_h^2}{(1+\ell_h)^{3/2}},\quad
 \frac{h^2\ell_h^2}{(1+\ell_h)^2},
\]
which are bounded by $12$, $\sqrt{12}h$, and $h^2$, respectively.
\end{proof}

The stencil $\Lambda_h^2$ has fixed coordinate radius two.  Its contribution
is the negative $q$-gradient of
$h^2\langle\Lambda_hq,B\Lambda_hq\rangle_h/2$; this identity concerns the
added term only and does not identify the complete local writer with the raw
finite Euler flow of the reconstructed signed action.

For $q_\alpha=D^\alpha q$ define
\begin{equation}
 \mathscr E^B_{s,h}=\mathscr E_{s,h}
   +\frac{h^2}{2}\sum_{|\alpha|\le s}
       \langle\Lambda_hq_\alpha,B\Lambda_hq_\alpha\rangle_h,
 \label{eq:supp-law-energy}
\end{equation}
where $\mathscr E_{s,h}$ is \eqref{eq:supp-open-energy}.

\begin{theorem}[Fixed-radius propagation on an open writer ball]
\label{thm:supp-law-stability}
For $s\ge11$, uniformly for $B\in\overline{\mathcal B}$ on a sufficiently
small common Cauchy chart,
\begin{equation}
 c_s\|X\|_{\mathcal X_h^s}^2\le\mathscr E^B_{s,h}
 \le C_s\|X\|_{\mathcal X_h^s}^2,
 \qquad
 (\mathscr E^B_{s,h})'\le C_s\mathscr E^B_{s,h}
                     +C_s(\mathscr E^B_{s,h})^{3/2}.
 \label{eq:supp-law-energy-bound}
\end{equation}
Consequently there are $\varepsilon_s,T_s>0$, independent of $h$ and $B$,
such that every initial record with
$\|X_h(0)\|_{\mathcal X_h^s}\le\varepsilon_s$ has a unique solution through
$T_s$ and
\begin{equation}
 \sup_{t\le T_s}\|X_h(t)\|_{\mathcal X_h^s}\le C_s\|X_h(0)\|_{\mathcal X_h^s}.
 \label{eq:supp-law-lifespan}
\end{equation}
With an additional acceleration defect $f$,
\begin{equation}
 (\sqrt{\mathscr E^B_{s,h}})'
 \le C_s\sqrt{\mathscr E^B_{s,h}}+C_s\mathscr E^B_{s,h}+C_s\|f\|_{s,h}.
 \label{eq:supp-law-forced}
\end{equation}
\end{theorem}
\begin{proof}
By \Cref{lem:supp-law-multiplier}, the absolute value of the added quadratic
energy at each commuted level is at most
$6b_0\sum_i\|D_i^+q_\alpha\|_h^2$.  The pre-existing spatial energy is at
least $c_*\sum_i\|D_i^+q_\alpha\|_h^2/2$, so
$b_0<c_*/24$ leaves a cutoff-independent positive margin.

Commute \eqref{eq:supp-law-family} with $D^\alpha$.  Since $B$ and
$\Lambda_h$ commute with every lattice shift, the new principal term is
$-h^2B\Lambda_h^2q_\alpha$.  Pairing with $v_\alpha$ gives exactly
\[
 -h^2\langle v_\alpha,B\Lambda_h^2q_\alpha\rangle_h
 =-\frac{\dd}{\dd t}\frac{h^2}{2}
       \langle\Lambda_hq_\alpha,B\Lambda_hq_\alpha\rangle_h.
\]
Thus the additional high-frequency term cancels at principal order rather
than being estimated as a small top-order force.  Every remaining commutator
is one of those controlled in \Cref{thm:supp-open-energy}; the only change is
that the actual acceleration appears in analytic mass-coefficient terms, and
\Cref{lem:supp-law-multiplier} bounds it uniformly in $H_h^{s-1}$.  This gives
\eqref{eq:supp-law-energy-bound}.  The finite-dimensional vector field is
analytic on the common chart, and the same first-exit argument used for
\eqref{eq:supp-open-forced-energy} proves the lifespan and the forced estimate.
\end{proof}

\begin{theorem}[Uniform Einstein limit for the whole coefficient family]
\label{thm:supp-law-einstein}
Let $X_{B,h}$ and $X_{0,h}$ have the same sufficiently small compatible
initial data.  On their common interval,
\begin{equation}
 \sum_{j=0}^3
 \|q_{B,h}^{(j)}-q_{0,h}^{(j)}\|_{C_tH_h^{s-1-j}}
       \le C_sh\|B\|_{\op}\varepsilon.
 \label{eq:supp-law-jet-comparison}
\end{equation}
For any sequence $B_h\in\overline{\mathcal B}$, with no convergence of $B_h$
required, the whole interpolated sequence converges to the same harmonic
vacuum development as the central writer.  In particular the metric and its
first two time derivatives satisfy the $O(h\varepsilon)$ bounds of
\Cref{thm:supp-open-einstein}, and the subsidiary field, rate positivity,
cut margin and noncollapse estimates remain uniform.
\end{theorem}
\begin{proof}
For comparison only, regard the added term on the $B$-trajectory as an
external force
\[
 f_{B,h}=-h^2a(g_{B,h})^{-1}B\Lambda_h^2q_{B,h}.
\]
Both trajectories already exist on the common top chart by
\Cref{thm:supp-law-stability}.  The lower-order bounds in
\Cref{lem:supp-law-multiplier} give
\[
 \|f_{B,h}\|_{C_tH_h^{s-2}}
 +\|\partial_tf_{B,h}\|_{C_tH_h^{s-3}}
 \le C_sh\|B\|\varepsilon,
\]
while the top $H_h^{s-1}$ norm need not be small.  Apply the forced difference
estimate \eqref{eq:supp-open-difference} at the reserved lower orders to get
the first two jet slots.  Subtracting the acceleration rows gives the second
time derivative; differentiating once and using the velocity differential of
the local vector field gives the third.  This proves
\eqref{eq:supp-law-jet-comparison}.  Add the central-writer convergence of
\Cref{thm:supp-open-einstein}.  All constants are uniform on
$\overline{\mathcal B}$, so no convergence of $B_h$ is needed.
\end{proof}

\begin{proposition}[Rapid switching is not included]
\label{prop:supp-law-switching}
For two distinct scalar marks $B_i=\sigma_iI$ in the admissible interval, one
resolved corner mode has frequency
\begin{equation}
 \Omega_{\sigma,h}(k)^2=\ell_h(k)(1+\sigma h^2\ell_h(k)).
 \label{eq:supp-law-dispersion}
\end{equation}
Holding each law for a quarter period gives transfer matrices whose two-step
product has eigenvalues
$-\Omega_{\sigma_1,h}/\Omega_{\sigma_2,h}$ and its reciprocal.  Their modulus
ratio stays separated from one at the ultraviolet corner when
$\sigma_1\ne\sigma_2$.  Repeated $O(h)$-time switching can therefore amplify
that mode like $\exp(cT/h)$ on a fixed slab.
\end{proposition}
\begin{proof}
The flat scalar mode is an oscillator.  A quarter-period transfer is
$\bigl(\begin{smallmatrix}0&\Omega^{-1}\\-\Omega&0\end{smallmatrix}\bigr)$.
Multiplying the two transfers gives the stated diagonal matrix.  At the
corner $h^2\ell_h(k)\to12$, so the frequency ratio approaches
$\sqrt{(1+12\sigma_1)/(1+12\sigma_2)}\ne1$.
\end{proof}

\subsection{Relatively open neighborhoods of chronological transition rows}

Fix $B\in\overline{\mathcal B}$ and $T<T_s$.  Partition $[0,T]$ into cells
$I_j$ with $c_-h^2\le\tau_j\le c_+h^2$.  A chronological letter is a finite
physical polynomial metric segment $Q_e$ with shared endpoint value,
velocity, acceleration and jerk records.  Before a possible guard exit define
\begin{equation}
 f_e(t)=Q_e''(t)-V_{B,h}(Q_e(t),Q_e'(t))
 \label{eq:supp-law-residual}
\end{equation}
and the cost \eqref{eq:main-law-cost}; a grammar or execution failure has cost
one and terminates the word.  The common-energy margin is chosen so that any
word with total cost below one is evaluated entirely inside the coefficient
chart.

\begin{lemma}[A stopped physical cost controls the whole word]
\label{lem:supp-law-cost-control}
If $C_h\le d_h^2<1$ and $d_h$ is sufficiently small, then no failure or guard
exit occurs and
\begin{equation}
 \|f_h\|_{L_t^2H_h^s}+\|\partial_tf_h\|_{L_t^2H_h^{s-3}}
       \le2\varepsilon d_h.
 \label{eq:supp-law-cost-force}
\end{equation}
Relative to the exact unforced $B$-history with the same initial record,
\begin{equation}
 \|Q_h-q_{B,h}\|_{C_tH_h^s}
 +\|Q_h'-q_{B,h}'\|_{C_tH_h^{s-1}}
 +\|Q_h''-q_{B,h}''\|_{C_tH_h^{s-3}}
 \le C_s\varepsilon d_h.
 \label{eq:supp-law-cost-shadow}
\end{equation}
\end{lemma}
\begin{proof}
Until a first guard hit, Cauchy--Schwarz gives
$\int_0^{t_*}\|f_h\|_{s,h}\dd t\le\sqrt T\,\varepsilon d_h$.
For small $d_h$ the forced energy \eqref{eq:supp-law-forced} contradicts a
first exit through the strict guard margin; a failure is excluded because it
would already contribute unit cost.  Hence the integrals are over the whole
interval and \eqref{eq:supp-law-cost-force} follows. Shared endpoint
value, velocity, acceleration and jerk make $f_h$ globally $H^1$ in time
with values in $H_h^{s-3}$, rather than merely piecewise $H^1$. The
Hilbert-valued one-dimensional Sobolev estimate therefore gives
\[
 \|f_h\|_{C_tH_h^{s-3}}
 \le C_T(\|f_h\|_{L_t^2H_h^{s-3}}+
                         \|\partial_tf_h\|_{L_t^2H_h^{s-3}})
 \le C_T\varepsilon d_h.
\]
Subtract the exact and forced mass equations at the reserved lower order,
use \eqref{eq:supp-open-difference} and Gronwall, and then use this
pointwise force bound in the actual acceleration difference. This proves
the final $C_t$ slot without omitting a temporal jump contribution.
\end{proof}

\begin{theorem}[A chronological law neighborhood with the same Einstein limit]
\label{thm:supp-law-row-neighbourhood}
Let $K_{h,j}(e\mid\pi)$ be an arbitrary finite transition row at each live
prefix, possibly depending on the entire past.  If
\begin{equation}
 \sum_eK_{h,j}(e\mid\pi)c_{h,j}(\pi,e)<M\tau_jh^4
 \label{eq:supp-law-row-moment}
\end{equation}
for every live prefix, then
\begin{equation}
 \mathbb EC_h\le MTh^4,
 \qquad\mathbb P\{C_h>h^2\}\le MTh^2.
 \label{eq:supp-law-row-probability}
\end{equation}
On $\{C_h\le h^2\}$ the history obeys
\Cref{lem:supp-law-cost-control} with $d_h=h$ and therefore has the same
Einstein limit as \Cref{thm:supp-law-einstein}.  At fixed cutoff and finite
grammar the strict inequalities \eqref{eq:supp-law-row-moment} define
nonempty relatively open convex subsets of the product of row simplexes when
a uniformly accurate completion is available at every live source.  They
contain strictly positive rows whenever all candidate letters have finite
cost.  Along the odd sequence $h=(2m+1)^{-1}$, the exceptional probabilities
are summable, so the conclusion holds eventually almost surely under any
common coupling.
\end{theorem}
\begin{proof}
Assign zero future cost to stopped prefixes.  Conditional on any live prefix,
\eqref{eq:supp-law-row-moment} bounds the expected stage cost by
$M\tau_jh^4$, independently of the probability of reaching that prefix.
Summation gives the expectation bound; Markov at threshold $h^2$ gives the
probability bound.  The deterministic cost lemma applies on the complement.
The series $\sum_m(2m+1)^{-2}$ converges, so the first Borel--Cantelli lemma
gives the eventual almost-sure assertion without independence between
cutoffs.  At fixed finite grammar each strict conditional mean bound is an
open affine half-space in one row simplex; finite products preserve relative
openness and convexity.  Small positive mixing preserves a strict margin.
\end{proof}

For an observable neighborhood description, attach to each ordinary letter
the profile
\begin{equation}
 \Phi_e=(f_e/\varepsilon,\partial_tf_e/\varepsilon,0),
 \label{eq:supp-law-profile}
\end{equation}
and to a failure the vector $(0,0,1)$ in
$L^2(I_j;H_h^s)\oplus L^2(I_j;H_h^{s-3})\oplus\mathbb R$.  Let
$d(e,e')=\|\Phi_e-\Phi_{e'}\|$ and let $W_{2,d}$ be the finite optimal
transport distance induced by $d$.

\begin{proposition}[Physical profile transport certifies a row neighborhood]
\label{prop:supp-law-transport}
For finite rows $K,P$,
\begin{equation}
 \left(\sum_eK(e)c_e\right)^{1/2}
 \le\left(\sum_eP(e)c_e\right)^{1/2}+W_{2,d}(K,P).
 \label{eq:supp-law-transport-bound}
\end{equation}
If $\sum_eP(e)c_e\le M_0\tau_jh^4$, then every row satisfying
\begin{equation}
 W_{2,d}(K,P)<\gamma\sqrt{\tau_j}\,h^2
 \label{eq:supp-law-transport-ball}
\end{equation}
obeys \eqref{eq:supp-law-row-moment} for any
$M>(\sqrt{M_0}+\gamma)^2$.  These strict balls are relatively open in the
finite probability simplex.
\end{proposition}
\begin{proof}
For any coupling of $K$ and $P$, Minkowski in the profile Hilbert space gives
$\|\Phi_e\|_{L^2}\le\|\Phi_{e'}\|_{L^2}+\|\Phi_e-\Phi_{e'}\|_{L^2}$.
Average under the coupling and minimize.  The finite coupling polytope is
compact, so a minimizer exists.  Continuity in the row probabilities follows
from finite-dimensional optimal transport.
\end{proof}

The scale in \eqref{eq:supp-law-transport-ball} is intentional.  A fixed
unscaled total-variation ball is not stable under refinement: if there are
$J_h\asymp h^{-2}$ stages and an absorbing failure of probability $p_h$ is
added independently at every live stage, then the failed-word probability is
$1-(1-p_h)^{J_h}$, which tends to zero exactly when $J_hp_h\to0$.  Thus even
$p_h\to0$ can be too large.  The theorem controls the accumulated physical
defect rather than hiding this refined-word effect.

\subsection{Exact finite initial constraints from free records}

We now remove the need to start the constructive population from a continuum
constraint solution.  This subsection uses the original reconstructed finite
action only at the initial cut.  Let the odd phase derivative be
\begin{equation}
 \widehat{\delta_i u}(k)=i\kappa_i(k)\widehat u(k),
 \qquad \kappa_i(k)=2h^{-1}\sin(\pi hk_i),
 \qquad \Omega_h^2=-\sum_i\delta_i^2.
 \label{eq:supp-initial-phase-derivative}
\end{equation}
For $r\ge3$ put
\begin{equation}
 \mathcal P_h^r=H_h^{r+2}(\operatorname{Sym}_3)
                 \times H_h^{r+1}(\operatorname{Sym}_3),
 \qquad \mathcal Y_h^r=H_h^r(\mathbb R\oplus\mathbb R^3),
 \label{eq:supp-initial-spaces}
\end{equation}
and write $X=(u,\pi)$, $\gamma=I+u$.  Let
$\mathscr C_h:\mathcal P_h^r\to\mathcal Y_h^r$ be
\eqref{eq:main-action-initial-constraints}.

\begin{lemma}[Uniform original-action constraint calculus]
\label{lem:supp-initial-calculus}
On a common ball $\mathscr C_h$ is analytic, $\mathscr C_h(0)=0$, and
\begin{equation}
 D\mathscr C_h(0)(u,\pi)=
 \left(\delta_i\delta_j u_{ij}-\sum_i\delta_i^2\tr u,
                         -2\delta_j\pi^{ij}\right)=:L_h(u,\pi).
 \label{eq:supp-initial-linear-map}
\end{equation}
Writing $\mathscr C_h=L_h+\mathscr N_h$,
\begin{equation}
 \|\mathscr N_h(X)\|_{\mathcal Y_h^r}\le C_r\|X\|^2,
 \qquad \|D\mathscr N_h(X)\|\le C_r\|X\|,
 \label{eq:supp-initial-nonlinear-bounds}
\end{equation}
with cutoff-independent constants.  No lapse or shift time derivative and no
metric acceleration appears in $\mathscr C_h$.
\end{lemma}
\begin{proof}
Stationary connection and Legendre elimination are analytic on the retained
small chart.  In a lapse or shift derivative, every spatial load has one phase
derivative; exact skew-adjointness transfers it to the stationary connection,
so the output loses at most the declared Sobolev orders.  At flat data the
quadratic phase-compatible action has the Fierz--Pauli rows displayed in
\eqref{eq:supp-initial-linear-map}; the full link remainder starts at cubic
order in the flat connection and contributes no linear initial row.  Taylor's
theorem gives \eqref{eq:supp-initial-nonlinear-bounds}.
\end{proof}

Let $P_0$ be spatial averaging in $\mathcal Y_h^r$ and
$P_\perp=I-P_0$.  On mean-zero targets define
\begin{equation}
 R_h(f,m)=\left(\tfrac12(\Omega_h^2)^{-1}f\,I,
  \tfrac12\mathcal L_h^{\rm vec}\mathcal D_h^{-1}m\right),
 \label{eq:supp-initial-right-inverse}
\end{equation}
where
$(\mathcal L_h^{\rm vec}Y)_{ij}=\delta_iY_j+\delta_jY_i
 -\frac23\delta_{ij}\delta_kY_k$ and
$\mathcal D_hY=-\delta_j(\mathcal L_h^{\rm vec}Y)_{ij}$.

\begin{lemma}[Range inverse and four missing means]
\label{lem:supp-initial-range}
One has $L_hR_h=I$ on the mean-zero target and
$\|R_h\|\le C_r$ uniformly.  For $d=N^3$,
\begin{equation}
 \operatorname{rank}L_h=4d-4,
 \qquad\dim\ker L_h=8d+4.
 \label{eq:supp-initial-rank}
\end{equation}
For each sufficiently small $z\in\ker L_h$ there is a unique mean-zero
$y_h(z)=O(\|z\|^2)$ such that
\begin{equation}
 P_\perp\mathscr C_h(z+R_hy_h(z))=0.
 \label{eq:supp-initial-range-zero}
\end{equation}
The remaining mean map
$\Theta_h(z)=P_0\mathscr C_h(z+R_hy_h(z))$ is analytic and has zero constant
and linear jets.
\end{lemma}
\begin{proof}
For $u=\phi I$ the scalar row equals $2\Omega_h^2\phi$.  For
$\pi=\mathcal L_h^{\rm vec}Y$ the vector row equals $2\mathcal D_hY$, whose
nonzero-mode symbol is $|\kappa|^2I+\kappa\kappa^T/3$ and is uniformly
invertible on the unit odd torus.  This proves the right inverse and rank.
The nonlinear range equation is
$y=-P_\perp\mathscr N_h(z+R_hy)$; the bounds of
\Cref{lem:supp-initial-calculus} make it a uniform contraction on a ball of
radius $C\|z\|^2$.  The expansion of $\Theta_h$ follows because
$P_0L_h=0$.
\end{proof}

To resolve the four means, let $T_i\in\operatorname{Sym}_3$ have its two
transverse off-diagonal entries equal to $1/\sqrt2$ and all other entries
zero, and define
\begin{equation}
 u_*=\sum_{i=1}^3T_i\cos(2\pi x_i),\qquad
 p(\vartheta)=-\tfrac23\kappa I
       +\sum_{i=1}^3b_iT_i\sin(2\pi x_i),
 \qquad \vartheta=(\kappa,b_1,b_2,b_3).
 \label{eq:supp-initial-seed}
\end{equation}
These are balancing directions only; arbitrary complementary modes remain
free.

\begin{lemma}[Exact four-mode mean calibration]
\label{lem:supp-initial-balance}
With $\omega_h=2h^{-1}\sin(\pi h)$, the quadratic mean jet of the full
retained action on the seed \eqref{eq:supp-initial-seed} is
\begin{equation}
 Q_h(\vartheta)=
 \begin{pmatrix}
 \frac23\kappa^2-\frac12\sum_i b_i^2-\frac38\omega_h^2\\
 -\frac12\omega_hb_1\\-\frac12\omega_hb_2\\-\frac12\omega_hb_3
 \end{pmatrix}.
 \label{eq:supp-initial-quadratic-mean}
\end{equation}
It has the exact root
\begin{equation}
 \vartheta_h^*=(3\omega_h/4,0,0,0),
 \qquad
 DQ_h(\vartheta_h^*)=
 \operatorname{diag}(\omega_h,-\omega_h/2,-\omega_h/2,-\omega_h/2),
 \label{eq:supp-initial-balance-root}
\end{equation}
whose inverse is uniformly bounded.  The opposite sign of the first component
gives the second branch.
\end{lemma}
\begin{proof}
Keep spatially constant lapse and shift independent and expand the canonical
Hamiltonian to quadratic order about flat data.  For the transverse trace-free
configuration seed and momentum equal to a homogeneous trace plus transverse
trace-free modes, the coefficient is
\[
 \mathscr H_h^{[2]}(u,\pi;N,\beta)=
 N\left(\|\pi\|_h^2-\tfrac12\|\tr\pi\|_h^2
          +\tfrac14\sum_i\|\delta_i u\|_h^2\right)
 +\sum_i\beta^i\langle\pi,\delta_i u\rangle_h.
\]
The complete link remainder starts at cubic order in the stationary connection
at flat data, so it does not modify this quadratic mean coefficient.  Taking
minus the lapse derivative and the three shift derivatives, then using
orthogonality of the constant and unit Fourier modes on every odd grid, gives
\eqref{eq:supp-initial-quadratic-mean}.  Direct differentiation yields
\eqref{eq:supp-initial-balance-root}; $\omega_h$ is uniformly separated from
zero on the declared grids.
\end{proof}

Let $\mathcal Z_h=\ker L_h$ and remove the four balancing seed directions
from an arbitrary small free record, leaving $W$ in a fixed complementary
subspace.  Write $Z_h(\vartheta)$ for the seed and, for $t\ne0$, define
\begin{equation}
 \mathcal Q_h(t,\vartheta,W)=t^{-2}\Theta_h(t[Z_h(\vartheta)+W]).
 \label{eq:supp-initial-normalized-mean}
\end{equation}
Since $\Theta_h$ has zero constant and linear jets, the integral Taylor formula
extends this map analytically through $t=0$.

\begin{theorem}[Exact finite preparation from free initial records]
\label{thm:supp-action-prepared-chart}
There are $t_0,\rho,d_0,h_0>0$, independent of $h$, such that for
$0<|t|<t_0$, $h<h_0$, and every complementary free record $W$ with
$\|W\|<\rho$, there is a unique $\vartheta_h(t,W)$ near
$\vartheta_h^*$ for which \eqref{eq:supp-initial-normalized-mean} vanishes.
The resulting record
\begin{equation}
 \mathfrak I_h(t,W)
 =t[Z_h(\vartheta_h(t,W))+W]
   +R_hy_h(t[Z_h(\vartheta_h(t,W))+W])
 \label{eq:supp-initial-prepared-record}
\end{equation}
satisfies
\begin{equation}
 \mathscr C_h(\mathfrak I_h(t,W))=0,
 \qquad \|\mathfrak I_h(t,W)\|_{\mathcal P_h^r}\le C_r|t|.
 \label{eq:supp-initial-exact-zero}
\end{equation}
At every such record $D\mathscr C_h$ is onto and admits a right inverse with
\begin{equation}
 \|\mathscr R_{h,t}f\|_{\mathcal P_h^r}
 \le C_r\bigl(\|P_\perp f\|_{\mathcal Y_h^r}+|t|^{-1}|P_0f|\bigr).
 \label{eq:supp-initial-full-inverse}
\end{equation}
For fixed $t\ne0$, the image contains a relatively open neighborhood in the
complete finite zero set; that zero set has codimension $4N^3$ in the
$12N^3$-dimensional canonical phase.  No symmetry, CMC, conformal-flatness,
frequency-support restriction, or continuum constraint solution is imposed
on $W$.
\end{theorem}
\begin{proof}
The analytic extension of \eqref{eq:supp-initial-normalized-mean} converges in
$C^1$ uniformly to the quadratic map of
\Cref{lem:supp-initial-balance} as $|t|+\|W\|\to0$.  Its root has a uniform
inverse margin.  Therefore
$\vartheta\mapsto\vartheta-[DQ_h(\vartheta_h^*)]^{-1}
\mathcal Q_h(t,\vartheta,W)$ is a contraction on one cutoff-independent ball
after choosing its radius first and then $t_0,\rho$.  The range equation was
already solved in \Cref{lem:supp-initial-range}, giving the complete zero.

For the derivative, eliminate the mean-zero block with the uniformly bounded
range inverse.  Along a balancing direction the remaining mean derivative is
$t^2D_\vartheta\mathcal Q_h$, while a range direction contributes only
$O(t)$ to the mean.  Solving for the balancing coefficient and multiplying by
the physical factor $t$ gives exactly the $|t|^{-1}$ term in
\eqref{eq:supp-initial-full-inverse}.  The inverse function theorem on the
zero set then gives relative openness and the codimension count.  The same
estimate shows why uniformity through the symmetric tip $t=0$ is neither
claimed nor true in general.
\end{proof}

The harmonic completion of a prepared record sets
$g_{00}=-1$, $g_{0i}=0$, $g_{ij}=\gamma_{ij}$ and uses its actual canonical
metric velocity $w$ for the spatial slot, with
\begin{equation}
 v_{ij}=w_{ij},\qquad
 v_{00}=-\gamma^{ij}w_{ij},\qquad
 v_{0i}=\gamma^{jk}\left(\delta_j\gamma_{ki}
             -\tfrac12\delta_i\gamma_{jk}\right).
 \label{eq:supp-initial-harmonic-completion}
\end{equation}
These lapse--shift velocities do not enter the four initial rows
\eqref{eq:main-action-initial-constraints}.

\begin{assumption}[Cofinal consistency of initial-action elimination]
\label{ass:supp-initial-consistency}
Let $\mathscr W_h$ denote the actual canonical spatial metric velocity
used in \eqref{eq:supp-initial-harmonic-completion}.  The continuum maps
$\mathscr C$ and $\mathscr W$ use the same action normalization and momentum
convention; $\mathscr C=0$ is equivalent to the vacuum constraints for
$(\gamma,K)$ with $K=-\mathscr W/2$ at unit lapse and zero shift.
On the bounded preparation chart in $\mathcal P_h^{s+2}$, $s\ge11$, the
actual eliminated maps satisfy
\begin{equation}
 \begin{split}
 &\|\mathcal I_h\mathscr C_h(X_h)
                 -\mathscr C(\mathcal I_hX_h)\|_{H^s}
 +\|\mathcal I_h\mathscr W_h(X_h)
                 -\mathscr W(\mathcal I_hX_h)\|_{H^s}\\
 &\hspace{4em}\le\chi_h\|X_h\|_{\mathcal P_h^{s+2}},
 \qquad \chi_h\longrightarrow0.
 \end{split}
 \label{eq:supp-initial-consistency-budget}
\end{equation}
The same comparisons hold for the first derivatives used in the normalized
range and four-mean equations, with their declared Sobolev buffer.
The phase derivative, right inverse, complementary free spaces and balancing
seeds use the common Fourier identifications.  This condition concerns the
original initial-action map, not the subsequently chosen harmonic writer.
\end{assumption}

\begin{lemma}[Consistency survives a uniformly regular elimination]
\label{lem:supp-initial-elimination-consistency}
Suppose an eliminated variable $z_h$ is defined by an actual equation
$F_h(x_h,z_h)=0$.  On a common comparison ball assume a preconditioner $A_h$
has uniformly bounded norm and
$\|I-A_hD_zF_h\|\le\kappa<1$.  If a comparison value $z_h^{\rm cmp}$
satisfies $\|F_h(x_h,z_h^{\rm cmp})\|\le\epsilon_h$, then
\begin{equation}
 \|z_h-z_h^{\rm cmp}\|
 \le\frac{\|A_h\|}{1-\kappa}\epsilon_h.
 \label{eq:supp-initial-elimination-error}
\end{equation}
With corresponding $C^1$ consistency of $F_h$ and its derivatives, the
eliminated map is $C^1$ consistent.  Thus finite compositions of the phase
differences, pointwise analytic coefficient maps and uniformly regular
stationary/Legendre solves satisfy
\Cref{ass:supp-initial-consistency} when their raw equations have the stated
Sobolev consistency.  Bounds on the inverse blocks must be retained; fixed
cutoff analyticity alone is not sufficient.
\end{lemma}
\begin{proof}
The map $z\mapsto z-A_hF_h(x_h,z)$ is a contraction on the comparison ball.
Subtract its fixed-point equation at $z_h$ from its value at
$z_h^{\rm cmp}$ and absorb $\kappa\|z_h-z_h^{\rm cmp}\|$.  This gives
\eqref{eq:supp-initial-elimination-error}.  For derivatives differentiate
$F_h(x_h,z_h(x_h))=0$ and use
$D_xz_h=-(D_zF_h)^{-1}D_xF_h$.  The inverse difference identity
$U^{-1}-V^{-1}=U^{-1}(V-U)V^{-1}$ and the common inverse bounds transfer
$C^1$ consistency.  The symbol
$2h^{-1}\sin(\pi hk_i)$ differs from $2\pi k_i$ by a controlled
higher-order symbol; the sampling and product estimates of
\Cref{lem:supp-open-interpolation} therefore verify the raw comparisons
for fixed compositions with the stated derivative buffer.  Iterating the
elimination argument proves the sufficient criterion.
\end{proof}

\begin{proposition}[Convergent preparation charts and subsequential preparation]
\label{prop:supp-initial-cofinal-preparation}
Fix $0<|t|<t_0$ and assume \Cref{ass:supp-initial-consistency}.
Let $W_h$ lie in the fixed complementary free ball of
\Cref{thm:supp-action-prepared-chart}, with its bound taken in
$\mathcal P_h^{s+2}$.  If $\mathcal I_hW_h$ converges in
$\mathcal P^{s+1}$ under the declared comparisons, then the prepared
records and their harmonic completions converge in the Cauchy topology
$H^{s+1}\times H^s$ to compatible vacuum initial data.
For an arbitrary bounded sequence $W_h$, the same conclusion holds after
subsequence extraction, but the limiting data need not be unique.
\end{proposition}
\begin{proof}
The prepared phase records have a uniform $\mathcal P_h^{s+2}$ bound.
Trigonometric interpolation and Rellich compactness give subsequential
convergence in $\mathcal P^{s+1}$.  Equation
\eqref{eq:supp-initial-consistency-budget} and the exact finite equality
$\mathscr C_h=0$ make every such limit solve $\mathscr C=0$.
The velocity comparison and the phase-difference consistency transfer the
harmonic completion to its continuum counterpart.

For a convergent free record, retain the normalized range fixed point and
the four-mean contraction in the proof of
\Cref{thm:supp-action-prepared-chart}.  Their comparison maps converge in
$C^1$ by \Cref{ass:supp-initial-consistency}; the explicit Fourier right
inverse and the four balancing modes have the same limits.  Their
contraction constants remain strictly below one on the chosen ball.
Every cluster point is consequently a fixed point of the same continuum
range/mean maps with the same free record.  Fixed-point uniqueness identifies
all cluster points.  Quantitatively, for either preconditioned contraction
$\mathcal T_h$, a sampled limiting fixed point $z^{\rm cmp}_h$ obeys
\[
 \|z_h-z^{\rm cmp}_h\|
 \le (1-\kappa)^{-1}
           \|\mathcal T_h(z^{\rm cmp}_h)-z^{\rm cmp}_h\|.
\]
The right side contains the free-record comparison and the actual map
consistency defect.  Its constants may depend on the fixed nonzero $t$;
no uniform estimate through the singular tip $t=0$ is inferred.
For merely bounded free records, first extract a convergent free subsequence
and use the same argument.  Alternating between two different limiting free
records is not a claim of whole-sequence convergence.
\end{proof}

\begin{corollary}[Action-prepared inputs populate the open law family]
\label{cor:supp-action-prepared-einstein}
Fix sufficiently small nonzero $t$ and the cofinally consistent branch of
\Cref{ass:supp-initial-consistency}.  For a convergent sequence of free
records as in \Cref{prop:supp-initial-cofinal-preparation}, and any sequence
$B_h\in\overline{\mathcal B}$, the harmonically completed histories exist
on a common interval and converge to the vacuum Einstein development of
the limiting prepared data.  No convergence of $B_h$ is required.  With
only the bounded free-record hypothesis, this is a subsequential statement.
The four original lapse--shift rows vanish exactly at the initial cut,
independently of $B_h$.
\end{corollary}
\begin{proof}
The preceding proposition supplies convergent compatible Cauchy data and
keeps them in the fixed top ball.  Apply the lower-order stability estimate
to their initial discrepancy and the $O(h)$ consistency force.  The proof
of \Cref{thm:supp-open-einstein} then gives the continuum development, with
an error bounded by the initial comparison error plus $C_sh\varepsilon$;
no $O(h)$ rate is assigned to an unquantified convergent input sequence.
The uniform comparison in \Cref{thm:supp-law-einstein} removes any need for
convergence of $B_h$.  The chronological row theorem applies with these same
initial records and its stated residual conditions.  Exact vanishing of
the initial constraints is a fixed-cutoff fact already established before
the acceleration law is selected.
\end{proof}

The three results of this section give an open population in a precise sense,
but not universal attraction.  The coefficient ball is fixed and genuinely
open; the transition-row neighborhoods are open in the scaled complete-profile
topology appropriate to $O(h^{-2})$ stages; and the initial records are
relatively open in the exact finite constraint surface.  The raw exact-action
evolution itself remains the separate question addressed next.


\subsection{Same-cylinder action occurrence and geometric compactness}
\label{supp:same-cylinder}

This section proves \Cref{thm:main-same-cylinder-einstein}.  The finite
correction below changes only the chronological occurrence weights on the
already reconstructed action support; it does not fit a new action to the
corrected transition law.

Let $X,Y$ be finite stage spaces, let $\alpha$ be a faithful source law and
let $Q(y\mid x)$ be a stochastic comparator on an admitted support
$\mathcal E$.  Set
$R(x,y)=\alpha(x)Q(y\mid x)$,
$\rho(y)=\sum_xR(x,y)$ and $d=\beta-\rho$ for a prescribed target law
$\beta$.  On edge arrays use
\begin{equation}
 \|h\|_{R^{-1}}^2=\sum_{(x,y)\in\mathcal E}\frac{h(x,y)^2}{R(x,y)},
 \qquad
 (Ah)(x)=\sum_yh(x,y),\qquad(Bh)(y)=\sum_xh(x,y).
 \label{eq:supp-same-cylinder-edge-geometry}
\end{equation}

For the main-text notation, the source-preserving correction is
When $d_{h,j}\in\Ran\mathcal L_{h,j}$, put
$u_{h,j}=\mathcal L_{h,j}^{\dagger}d_{h,j}$ and
\begin{equation}
 h_{h,j}^*(x,y)=R_{h,j}(x,y)
 \{u_{h,j}(y)-(Q_{h,j}u_{h,j})(x)\}.
 \label{eq:main-thomson-correction}
\end{equation}
If the leverage is at most $1/2$, then
$\pi_{h,j}^{\rm T}=R_{h,j}+h_{h,j}^*$ is nonnegative and
$K_{h,j}^{\rm T}(y\mid x)=\pi_{h,j}^{\rm T}(x,y)/\nu_{h,j}(x)$ is a
stochastic row with source marginal $\nu_{h,j}$ and target marginal
$\nu_{h,j+1}$.  For an edge observable $c$, write
\begin{theorem}[Chronological source-shorted Thomson correction]
\label{thm:supp-chronological-thomson}
Put
\begin{equation}
 P_s=I-A^*\diag(\alpha)^{-1}A,
 \qquad
 \mathcal L_{\alpha,Q}=BP_sB^*.
 \label{eq:supp-thomson-Ps}
\end{equation}
Then
\begin{equation}
 \mathcal L_{\alpha,Q}
 =\diag(\rho)-Q^{\mathsf T}\diag(\alpha)Q\succeq0,
 \label{eq:supp-thomson-L}
\end{equation}
and
\begin{equation}
 \langle u,\mathcal L_{\alpha,Q}u\rangle
 =\sum_x\alpha(x)\operatorname{Var}_{Q(x,\cdot)}(u).
 \label{eq:supp-thomson-variance}
\end{equation}
The constraints $Ah=0$, $Bh=d$ are solvable if and only if
$d\in\Ran\mathcal L_{\alpha,Q}$.  On that branch the unique
minimum-$R^{-1}$-norm solution is
\begin{equation}
 h_*(x,y)=R(x,y)\{u(y)-(Qu)(x)\},
 \qquad u=\mathcal L_{\alpha,Q}^{\dagger}d,
 \label{eq:supp-thomson-hstar}
\end{equation}
with
\begin{equation}
 \|h_*\|_{R^{-1}}^2
 =\langle d,\mathcal L_{\alpha,Q}^{\dagger}d\rangle.
 \label{eq:supp-thomson-energy}
\end{equation}
If $\max_{(x,y)\in\mathcal E}|u(y)-(Qu)(x)|\le1/2$, then
$R+h_*$ is nonnegative and has source marginal $\alpha$ and target marginal
$\beta$.
\end{theorem}
\begin{proof}
The weighted adjoints are
$(A^*f)(x,y)=R(x,y)f(x)$ and
$(B^*g)(x,y)=R(x,y)g(y)$.  Since $AA^*=\diag(\alpha)$,
$P_s$ is the orthogonal projection onto $\ker A$.  Direct multiplication
gives \eqref{eq:supp-thomson-L}, and expanding the quadratic form gives
\eqref{eq:supp-thomson-variance}.  The Moore--Penrose normal equations for
$B|_{\ker A}$ give \eqref{eq:supp-thomson-hstar} and
\eqref{eq:supp-thomson-energy}.  Finally
$R+h_*=R\{1+u(y)-(Qu)(x)\}$, so the leverage bound gives positivity while
the two linear constraints give the required marginals.
\end{proof}

For a real edge observable $c$, write
$\bar c(x)=(Qc)(x)$ and
\begin{equation}
 \mathcal V_R(c)=\sum_{x,y}R(x,y)|c(x,y)-\bar c(x)|^2.
 \label{eq:supp-same-cylinder-variance}
\end{equation}

\begin{lemma}[Variance-sensitive source-preserving transfer]
\label{lem:supp-same-cylinder-transfer}
If $Ah=0$, then
\begin{equation}
 \left|\sum_{x,y}h(x,y)c(x,y)\right|
 \le\|h\|_{R^{-1}}\mathcal V_R(c)^{1/2}.
 \label{eq:supp-same-cylinder-transfer}
\end{equation}
Hence the Thomson coupling $\pi^{\rm T}=R+h_*$ satisfies
\begin{equation}
 |\mathbb E_{\pi^{\rm T}}c-\mathbb E_Rc|
 \le\sqrt{\mathcal E_{\rm occ}\mathcal V_R(c)}.
 \label{eq:supp-same-cylinder-transfer-thomson}
\end{equation}
If $0\le c\le B$ and $m=\mathbb E_Rc$, then
$\mathcal V_R(c)\le Bm$.
\end{lemma}
\begin{proof}
Because $Ah=0$, the source-only function $\bar c(x)$ is orthogonal to $h$.
Thus $\sum hc=\sum h(c-\bar c)$, and weighted Cauchy--Schwarz proves the
first inequality.  Insert $h=h_*$ for the second.  The final statement
follows rowwise from $\operatorname{Var}(c)\le\mathbb E(c^2)\le B\mathbb Ec$.
\end{proof}

Consider now chronological stages $j=0,\ldots,J-1$ with faithful marginals
$\nu_j$ and corrected couplings $\pi_j^{\rm T}$.  Starting with
$X_0\sim\nu_0$ and using the corrected transition rows gives
$X_j\sim\nu_j$ for every $j$ by induction.  Therefore, for nonnegative costs
$c_j$,
\begin{equation}
 \mathbb E\sum_jc_j(X_j,X_{j+1})
 \le\sum_j\left(m_j+
 \sqrt{\mathcal E_{{\rm occ},j}\mathcal V_j}\right),
 \label{eq:supp-same-cylinder-chain-cost}
\end{equation}
where $m_j=\mathbb E_{R_j}c_j$ and
$\mathcal V_j=\mathcal V_{R_j}(c_j)$.  Under
\eqref{eq:main-same-cylinder-writer-scale},
\begin{equation}
 \mathbb EC_h\le(M_0+\sqrt{M_1V_0})Th^4,
 \qquad
 \Pr\{C_h>h^2\}\le(M_0+\sqrt{M_1V_0})Th^2.
 \label{eq:supp-same-cylinder-badprob}
\end{equation}
Markov's inequality proves the probability estimate, and the first
Borel--Cantelli lemma gives an almost-sure successful tail along any
refinement with $\sum_nh_n^2<\infty$; no independence is used.

The same centered estimate applies to the field-action gap.  If
$\Delta_j\ge0$ is bounded by $B_j$ on the finite support and
$a_j=\mathbb E_{R_j}\Delta_j$, then
\begin{equation}
 \mathbb E_{\pi_j^{\rm T}}\Delta_j
 \le a_j+\sqrt{\mathcal E_{{\rm occ},j}
                    \mathcal V_{R_j}(\Delta_j)}
 \le a_j+\sqrt{B_j\mathcal E_{{\rm occ},j}a_j}.
 \label{eq:supp-same-cylinder-action-gap}
\end{equation}
For moving chronological marginals, this gap transfer alone is not a
stationarity theorem. Apply \Cref{prop:supp-nonstationary-action-gap} to the
corrected coupling, using the same common-action chart at both ends.
With $C_{1,h,j},C_{2,h,j}$ as in that proposition and
$\bar\Delta^{\rm T}_{h,j}=\mathbb E_{\pi^{\rm T}_{h,j}}\Delta_{h,j}$,
put
\begin{equation}
 U_{h,j}(k)=B_{h,j}(k)
 \left[(C_{1,h,j}+C_{2,h,j})\bar\Delta^{\rm T}_{h,j}
                      +C_{2,h,j}d^+_{\mathcal A,h,j}\right]^{1/2}.
 \label{eq:supp-same-cylinder-stage-test}
\end{equation}
If the actual determining first variation is represented as a finite sum
of stage insertions with coefficients $w_{h,j}(k)$, Minkowski gives
\begin{equation}
 \bigl(\mathbb E|\delta\mathcal A_h[k]|^2\bigr)^{1/2}
 \le\sum_j|w_{h,j}(k)|U_{h,j}(k)=:\mathfrak U_h(k).
 \label{eq:supp-same-cylinder-physical-scale}
\end{equation}
The required scaling is $\mathfrak U_h(k)\to0$ for each fixed test.
For a single represented field insertion the sum has one term.  The
weights, chart pullbacks and lift bounds are those of the physical action;
they are not normalized after the estimate.  An independent direct
mean-square first-variation certificate may be used instead.  Small gap
alone cannot suppress a nonzero chronological action descent.

\begin{proposition}[Constructive metric jets imply the literal-link budget]
\label{prop:supp-same-cylinder-link-budget}
On the successful event $C_h\le h^2$ of the nonsymmetric constructive
writer, the canonical metric interpolation remains in the fixed small chart
and obeys the cutoff-uniform metric-jet bounds of
\Cref{thm:main-open-3plus1,thm:main-open-law-basin}.  On the spinless regular
Cartan branch these bounds imply
\begin{equation}
 \sup_h\left(
 \|A_h\|_{L_t^\infty H_h^s}
 +\|\partial_tA_h\|_{L_t^\infty H_h^{s-1}}
 +\|A_{0,h}\|_{L^2}+\|E_h\|_{L^2}+\|F_h\|_{L^2}\right)<\infty
 \label{eq:supp-same-cylinder-link-budget}
\end{equation}
for fixed $s\ge3$.  In particular
\eqref{eq:main-link-chart}--\eqref{eq:main-link-energy} hold on every
sufficiently fine successful cutoff.
\end{proposition}
\begin{proof}
The canonical ADM coframe is an analytic finite-dimensional function of the
metric on the fixed nondegenerate chart.  Uniform discrete Moser estimates
therefore transfer the metric bounds and their first two time derivatives to
$e_h$.  The stationary Cartan equation has a uniformly invertible algebraic
principal block and an $h$-weighted quadratic connection remainder; the same
Neumann estimate used in the spinless Cartan branch gives a uniform
$H_h^s$ bound for $A_h$.  Differentiating that algebraic equation once in
time gives the $H_h^{s-1}$ bound for $\partial_tA_h$.  The normalized
plaquette is an analytic four-link function whose apparent $h^{-2}$
singularity is removable, and the electric link is the exact temporal link
identity; these give the final two $L^2$ bounds.  Sobolev embedding gives
$h\|A_h\|_\infty\to0$, hence one common identity chart on a cofinal tail.
\end{proof}

\begin{proof}[Proof of \Cref{thm:main-same-cylinder-einstein}]
The Thomson theorem gives the corrected chronological marginals and support.
Equation \eqref{eq:supp-same-cylinder-chain-cost} with the assumed scaling
gives \eqref{eq:main-same-cylinder-probability}.  On its successful event,
\Cref{prop:supp-same-cylinder-link-budget} and
\Cref{thm:main-literal-link-compactness} identify the connection curvature.
Equation \eqref{eq:supp-same-cylinder-action-gap} and the nonstationary
estimate \eqref{eq:supp-same-cylinder-physical-scale}, with its explicitly
assumed scaling, give mean-square vanishing of the represented first variations.  The diagonal/Tonelli argument of
\Cref{prop:supp-physical-test-stationarity} supplies a deterministic
subsequence on which the countable determining core vanishes almost surely.
The classified Palatini--Holst--volume representative, nonzero Palatini
normalization, noncollapse and strong coframe compactness then permit the
weak--strong Palatini passage.  On the literal-link route the curvature is
already identified by \Cref{thm:main-literal-link-compactness}; hence the
limiting variation is
$\chi\int\sqrt{-g}(G_{\mu\nu}+\Lambda g_{\mu\nu})k^{\mu\nu}$ and vanishes on
the determining core.  Borel--Cantelli intersects the successful residual
tail with the full-measure stationarity subsequence without any independence
assumption.
\end{proof}


\section{Exact-action response, preparation and boundary realization}
\label{supp:exact-action-audit}
This appendix proves the response and fixed-cutoff realization assertions
summarized in \Cref{thm:main-exact-action-summary}. It uses the original
phase-compatible action, not the harmonic evolution law. Static all-cutoff
results, the three-site graded boundary chart, and the separate invariant
normal-factor chart keep their individual hypotheses and constants;
Appendix~\ref{supp:exact-action-slow} treats the subsequent slow limit.

\subsection{Canonical constraints and the finite characteristic quotient}

Let $X=(u,\pi)\in\mathbb R^{12V}$ and let
$\lambda=(N,\beta)\in\mathbb R^{4V}$.  For the canonical Hamiltonian
$H_h(X,\lambda)$ put
\begin{equation}
 P_h=\nabla_\lambda H_h,\qquad
 F_h^{\rm can}=\mathbb J\nabla_XH_h,\qquad
 B_h=D_XP_hF_h^{\rm can},\qquad M_h=D_\lambda P_h.
 \label{eq:supp-exact-canonical-objects}
\end{equation}
The exact Euler system is
\begin{equation}
 X_t=F_h^{\rm can}(X,\lambda),\qquad P_h(X,\lambda)=0,
 \label{eq:supp-exact-euler}
\end{equation}
and differentiating the constraint gives
\begin{equation}
 B_h+M_h\dot\lambda=0.
 \label{eq:supp-exact-first-consistency}
\end{equation}
Phase-compatible multiplier homogeneity gives
\begin{equation}
 H_h(X,c\lambda)=cH_h(X,\lambda),\qquad
 M_h\lambda=0,
 \qquad \langle\lambda,B_h\rangle_h=0.
 \label{eq:supp-exact-homogeneity}
\end{equation}
These are global scaling identities, not local discrete diffeomorphism
identities.

Introduce multiplier momenta and the extended constraint bracket
\begin{equation}
 \mathcal A_h=D_XP_h\,\mathbb J(D_XP_h)^*,\qquad
 \mathfrak D_h=\begin{pmatrix}0&-M_h\\M_h&\mathcal A_h\end{pmatrix}.
 \label{eq:supp-exact-dirac-matrix}
\end{equation}

\begin{theorem}[Finite characteristic quotient on a regular one-null branch]
\label{thm:supp-exact-constraint-rank}
Suppose $D_XP_h$ is onto, $\ker M_h=\operatorname{span}\{\lambda\}$ and
\eqref{eq:supp-exact-first-consistency} is solvable.  Then
\begin{equation}
 \ker\mathfrak D_h=\operatorname{span}\left\{
 \binom{\lambda}{0},\binom{-M_h^\dagger B_h}{\lambda}\right\},
 \qquad \rank\mathfrak D_h=8V-2,
 \label{eq:supp-exact-dirac-rank}
\end{equation}
and the local characteristic quotient has dimension $12V-2$.
\end{theorem}
\begin{proof}
If $(r,s)\in\ker\mathfrak D_h$, the first row gives $s=c\lambda$.
Differentiating the homogeneity identity in $X$ and pairing with the
Hamiltonian vector field gives $\mathcal A_h\lambda=B_h$.  The second row
therefore gives $r=-cM_h^\dagger B_h+d\lambda$.  The extended phase has
dimension $20V$ and the $8V$ independent primary and secondary constraints
leave dimension $12V$; quotienting the two characteristic directions gives
$12V-2$.
\end{proof}

For the flat Cartan tangent, let $\mathsf P_{\rm tr}$ be the orthogonal
projection onto scalar spatial matrices. Its signed operator is
\begin{equation}
 K_0=2I-6\mathsf P_{\rm tr},\qquad
 K_0^{-1}=\tfrac12I-\tfrac34\mathsf P_{\rm tr}.
 \label{eq:main-exact-action-K0}
\end{equation}

\subsection{An exact regular transverse block and lossless source elimination}

Use the odd-grid phase derivatives
\begin{equation}
 \widehat{\delta_i u}(k)=i\kappa_i(k)\widehat u(k),
 \qquad \kappa_i(k)=2h^{-1}\sin(\pi hk_i),
 \qquad \Omega_h^2=-\sum_i\delta_i^2.
 \label{eq:supp-exact-phase-derivative}
\end{equation}
Let
\begin{equation}
 \mathcal T_h=\{\beta:\widehat\beta(0)=0,
 \ \kappa(k)^T\widehat\beta(k)=0\ (k\ne0)\}
 \label{eq:supp-exact-transverse-space}
\end{equation}
be the transverse-shift space, of dimension $2(V-1)$.  On the scalar-momentum
family the first source restricted to $\mathcal T_h$ has the form
\begin{equation}
 A_h\beta=-\frac{h\varkappa_h}{2}\,\mathcal L_h
          \left(-\tfrac12\operatorname{curl}_\delta\beta\right),
 \label{eq:supp-exact-transverse-source}
\end{equation}
where $\varkappa_h$ is uniformly bounded above and below and
$\mathcal L_hw$ is the symmetric matrix with zero diagonal and
$(\mathcal L_hw)_{ij}=S_iw_j+S_jw_i$.

\begin{theorem}[Exact all-cutoff transverse rank and natural-order bounds]
\label{thm:supp-exact-transverse}
At every odd cutoff,
\begin{equation}
 \rank A_h=2(V-1),\qquad \mathsf P_{\rm tr}A_h=0,
 \label{eq:supp-exact-transverse-rank}
\end{equation}
and
\begin{equation}
 c h^2\|\Omega_h\beta\|_h^2
 \le\|A_h\beta\|_h^2
 \le C h^2\|\Omega_h\beta\|_h^2,
 \qquad \beta\in\mathcal T_h.
 \label{eq:supp-exact-transverse-bounds}
\end{equation}
The same estimates hold after multiplication by a scalar Fourier Sobolev
weight, and
\begin{equation}
 A_h^*K_0^{-1}A_h=\tfrac12A_h^*A_h\succ0.
 \label{eq:supp-exact-transverse-positive}
\end{equation}
\end{theorem}
\begin{proof}
At one Fourier mode the off-diagonal map is unitarily equivalent to
$w\mapsto(w_1+w_2,w_1+w_3,w_2+w_3)$, whose squared singular values are
$4,1,1$.  For transverse $\beta$,
$\|\operatorname{curl}_\delta\beta\|_h=\|\Omega_h\beta\|_h$.
Substitution in \eqref{eq:supp-exact-transverse-source} proves the two-sided
estimate and injectivity.  The output diagonal vanishes, so the trace
projection vanishes and $K_0^{-1}$ acts by $1/2$ on this range.
\end{proof}

Write the complete leading graded source as $F_h=(A_h,G_h)$ with target
$(0,d_h)$, and put $\overline G_h=(I-\Pi_{A,h})G_h$.

\begin{corollary}[Lossless elimination of the transverse source columns]
\label{cor:supp-exact-transverse-elimination}
One has
\begin{equation}
 \rank F_h=(2V-2)+\rank\overline G_h,
 \label{eq:supp-exact-rank-split}
\end{equation}
and
\begin{equation}
 F_h^*K_0^{-1}F_h\sim
 \operatorname{diag}\left(\tfrac12A_h^*A_h,
          \overline G_h^*K_0^{-1}\overline G_h\right).
 \label{eq:supp-exact-Gram-split}
\end{equation}
The triangular column transformation producing this congruence leaves the
actual target unchanged and preserves every remaining source equation.
\end{corollary}
\begin{proof}
Use the coefficient
$C_h=(A_h^*A_h)^{-1}A_h^*G_h$ and right-multiply by
$\left(\begin{smallmatrix}I&-C_h\\0&I\end{smallmatrix}\right)$.
The transformed source is $(A_h,\overline G_h)$ with orthogonal column
ranges.  Since $K_0^{-1}A_h=A_h/2$, the signed cross term vanishes as well.
\end{proof}

\subsection{Exact scalar--trace reconstruction of the signed response}

Work after any regular positive block elimination and let $F$ denote the
remaining full-rank source with actual target $b$.  Define
\begin{equation}
 \Pi=F(F^*F)^{-1}F^*,\qquad
 w=F(F^*F)^{-1}b,
 \qquad E\tau=\tau I/\sqrt3,
 \label{eq:supp-exact-Pi-w}
\end{equation}
then
\begin{equation}
 C=E^*\Pi E,\qquad H=2I-3C,\qquad \zeta=E^*w.
 \label{eq:supp-exact-trace-data}
\end{equation}

\begin{theorem}[Exact scalar--trace response reconstruction]
\label{thm:supp-exact-scalar-reconstruction}
The signed Gram $F^*K_0^{-1}F$ is invertible if and only if $H$ is
invertible.  If $\chi=H^{-1}\zeta$, then the original physical response is
\begin{equation}
 v=w-3(I-\Pi)E\chi,\qquad E^*v=-\chi,
 \label{eq:supp-exact-scalar-v}
\end{equation}
and
\begin{equation}
 \|v\|^2=\|w\|^2+9\langle\chi,(I-C)\chi\rangle.
 \label{eq:supp-exact-scalar-norm}
\end{equation}
Consequently
\begin{equation}
 v_h\to0\quad\Longleftrightarrow\quad
 w_h\to0\ \hbox{and}\ H_h^{-1}E^*w_h\to0
 \label{eq:supp-exact-response-equivalence}
\end{equation}
on any regular cofinal family.
\end{theorem}
\begin{proof}
Write $F=QR$ with $Q^*Q=I$.  The signed Gram is congruent to
$(2I-3Q^*EE^*Q)/4$.  The nonzero eigenvalues of
$Q^*EE^*Q$ and $E^*QQ^*E=C$ agree, proving the invertibility equivalence.
The response is characterized by $F^*v=b$ and $K_0v\in\Ran F$.  The first
identity gives $\Pi v=w$ and the second gives
$(I-\Pi)v=3(I-\Pi)EE^*v$.  Applying $E^*$ yields
$H(E^*v)=-\zeta$, hence $E^*v=-\chi$ and
\eqref{eq:supp-exact-scalar-v}.  Orthogonality of the two displayed terms
gives \eqref{eq:supp-exact-scalar-norm}.
\end{proof}

Diagonalizing the compression $C$ gives the trace-angle form.  Equivalently,
with the positive polar decomposition $F=QR$ and
$T=Q^*\mathsf P_{\rm tr}Q$, $\eta=R^{-1}b$, one obtains
\begin{equation}
 v=(2I-3\mathsf P_{\rm tr})Q(2I-3T)^{-1}\eta.
 \label{eq:supp-exact-response-polar}
\end{equation}
If $Tu_j=\mu_ju_j$ and $\eta_j=\langle u_j,\eta\rangle$, then
\begin{equation}
 \|\mathcal E^0\|^2
 =\sum_j\frac{4-3\mu_j}{(2-3\mu_j)^2}|\eta_j|^2,
 \label{eq:supp-exact-physical-norm}
\end{equation}
so $\mu=2/3$ is the unique signed trace resonance.  This is the spectral
form of \Cref{thm:supp-exact-scalar-reconstruction}, not a different response
law.

For a cofinal regular family define
$\nu_h=\sum_j|\eta_{h,j}|^2\delta_{\mu_{h,j}}$.  Then
$\|\mathcal E_h^0\|\to0$ if and only if
$\nu_h([0,1])\to0$ and
\begin{equation}
 \lim_{\delta\downarrow0}\limsup_{h\to0}
 \int_{|2-3\mu|<\delta}\frac{\dd\nu_h(\mu)}{|2-3\mu|^2}=0.
 \label{eq:supp-exact-resonance-ui}
\end{equation}
The proof is the direct split of \eqref{eq:supp-exact-physical-norm} into a
region separated from $2/3$ and a shrinking neighborhood of the resonance.

\subsection{Time-dependent response and the necessary clock test}

Static response control is insufficient for a spacetime comparison.  Let
$F(t),b(t)$ be $C^r$ with full column rank and regular signed Gram in a fixed
physical tangent space, and let $\Pi,w,C,H,\zeta,\chi,v$ be the preceding
objects.

\begin{theorem}[Derivative hierarchy in scalar--trace coordinates]
\label{thm:supp-exact-time-hierarchy}
Put $\chi_0=\chi$ and, for $1\le j\le r$,
\begin{equation}
 \chi_j=H^{-1}\left[\zeta^{(j)}+
 3\sum_{\ell=1}^j\binom j\ell C^{(\ell)}\chi_{j-\ell}\right].
 \label{eq:supp-exact-time-sources}
\end{equation}
Then $\chi_j=\chi^{(j)}=-E^*v^{(j)}$.  In particular
\begin{equation}
 \chi'=H^{-1}(\zeta'+3C'\chi),\qquad
 v'=w'+3\Pi'E\chi-3(I-\Pi)E\chi'.
 \label{eq:supp-exact-first-time-source}
\end{equation}
Thus a small static response does not control the differentiated response
without bounds on the transported source geometry.
\end{theorem}
\begin{proof}
Differentiate $H\chi=\zeta$ repeatedly and use $H^{(\ell)}=-3C^{(\ell)}$.
The identity $E^*v=-\chi$ holds at every time; differentiating
\eqref{eq:supp-exact-scalar-v} gives the second formula.
\end{proof}

The normalized exact continuation has an active/weak amplitude grading
$S_a=\operatorname{diag}(aI_A,a^2I_W)$.  If $u_a$ is the normalized source
coefficient entering the multiplier solve, then
\begin{equation}
 \|\lambda_t\|_h^2=a^{-2}\|(u_a)_A\|_h^2
                    +a^{-4}\|(u_a)_W\|_h^2.
 \label{eq:supp-exact-rate-norm}
\end{equation}

\begin{proposition}[Necessary lapse--shift clock condition for harmonic shadowing]
\label{prop:supp-exact-clock-test}
Suppose $a(h)\to0$ and an exact-action metric is to shadow a small harmonic
metric in a topology controlling the initial first time derivative, with the
same lapse, shift and spatial metric at the initial cut.  Then necessarily
\begin{equation}
 a(h)^{-1}\|(u_{a(h)})_A\|_h\to0,
 \qquad a(h)^{-2}\|(u_{a(h)})_W\|_h\to0.
 \label{eq:supp-exact-rate-necessary}
\end{equation}
This requirement is independent of smallness of the instantaneous electric
coefficient.
\end{proposition}
\begin{proof}
At $N=1$, $\beta=0$ in ADM coordinates,
$\partial_tg_{00}=-2\partial_tN$ and
$\partial_tg_{0i}=\gamma_{ij}\partial_t\beta^j$.  First-jet shadowing
therefore forces $\|\lambda_t\|\to0$.  Equation
\eqref{eq:supp-exact-rate-norm} gives the two conditions.
\end{proof}

\subsection{Prepared exact-action families}

The generic probe-source calculation remains useful for comparison.  For the
smooth first jet used in the original response audit, the actual weak source
obeys
\begin{equation}
 \|\mathcal I_hd_{h,\rho}-h^2d_\rho^{(2)}\|_{H^r}\le C_rh^4,
 \qquad c_rh^2\le\|d_{h,\rho}\|_{H_h^r}\le C_rh^2.
 \label{eq:supp-exact-source-expansion}
\end{equation}
Thus the generic weak Ward defect is genuinely second order, but
\Cref{thm:supp-exact-scalar-reconstruction} shows why its norm alone does not
control the physical response.

We next record two exact preparations with different purposes.

\paragraph{A rationally detuned all-cutoff active family.}
Let $N=2m+1$, $h=N^{-1}$ and define the one-dimensional product-defect operator
\begin{equation}
 D_hu=\delta(fu)-f\delta u-(\delta f)u,
 \qquad f(x)=\cos(2\pi x).
 \label{eq:supp-exact-cosine-defect}
\end{equation}
It is self-adjoint, has kernel the constants, and its nonconstant part is an
irreducible zero-diagonal Jacobi matrix.  For positive amplitudes
$\alpha=(\alpha_1,\alpha_2,\alpha_3)$ put
\begin{equation}
 U_\alpha=\sum_{i=1}^3\alpha_iT_i\cos(2\pi x_i),\qquad
 p_{h,\alpha}=-s_{h,\alpha}I,
 \qquad
 s_{h,\alpha}=\omega_h\sqrt{\frac{\alpha_1^2+\alpha_2^2+\alpha_3^2}{6}},
 \label{eq:supp-exact-detuned-seed}
\end{equation}
where $T_i$ is the transverse off-diagonal tensor used in
\eqref{eq:supp-initial-seed}.  The diagonal lapse part of the complete first
source is
\begin{equation}
 \mathscr D_{h,\alpha}n=
 \begin{pmatrix}
 (\alpha_2D_{2,h}-\alpha_3D_{3,h})n\\
 (\alpha_3D_{3,h}-\alpha_1D_{1,h})n\\
 (\alpha_1D_{1,h}-\alpha_2D_{2,h})n
 \end{pmatrix}.
 \label{eq:supp-exact-diagonal-source}
\end{equation}

\begin{lemma}[Finite rational detuning]
\label{lem:supp-exact-rational-detuning}
For every sufficiently fine odd cutoff there is a rational
$\theta_h$ with $1<\theta_h<1+h$ such that, for
$\alpha_h=(1,\theta_h,1)$,
\begin{equation}
 \|\mathscr D_{h,\alpha_h}n\|_h\ge ch^6\|n\|_h
 \qquad(\overline n=0).
 \label{eq:supp-exact-diagonal-margin}
\end{equation}
The selection can be made by a finite algebraic search at each cutoff.
\end{lemma}
\begin{proof}
The Jacobi representation of $D_h$ has simple algebraic spectrum
$\Sigma_h$.  Its reciprocal-eigenvalue sum is $O(h^{-3})$.  The union of
ratios $\theta$ for which
$|\lambda_a-\theta\lambda_b|<4h^6$ has measure $O(h^2)$, smaller than a
fixed fraction of the interval $(1,1+h)$ for small $h$.  Choose a dyadic
rational outside that union and sufficiently close to the selected point.
In the tensor product eigenbasis, the three squared components of
\eqref{eq:supp-exact-diagonal-source} vanish simultaneously only for the
constant mode, and the spectral separation gives the stated margin.
\end{proof}

\begin{theorem}[All-cutoff regular active source]
\label{thm:supp-exact-active-regular}
The exact original-action initial constraints have analytic preparations
\begin{equation}
 \mathscr C_h(X_{h,\alpha_h}(a))=0,
 \qquad X_{h,\alpha_h}(a)=a(U_{\alpha_h},p_{h,\alpha_h})+O(a^2),
 \label{eq:supp-exact-detuned-preparation}
\end{equation}
uniformly for all sufficiently fine odd cutoffs.  On their first jets the
complete leading active source has
\begin{equation}
 \ker\mathcal A_h^{\rm act}=\mathbb R(1,0)
 \oplus\{(0,\beta):\operatorname{curl}_\delta\beta=0\},
 \qquad \rank\mathcal A_h^{\rm act}=3V-3,
 \label{eq:supp-exact-active-rank}
\end{equation}
its range is trace-free, and
\begin{equation}
 (\mathcal A_h^{\rm act})^*K_0^{-1}\mathcal A_h^{\rm act}
 =\tfrac12(\mathcal A_h^{\rm act})^*\mathcal A_h^{\rm act}
 \succeq ch^{16}I.
 \label{eq:supp-exact-active-Gram}
\end{equation}
Hence the complete remaining signed rank problem is confined to the weak
columns.
\end{theorem}
\begin{proof}
The exact preparation follows from the finite initial-constraint chart and
the nondegenerate four-mean calibration; the $O(h)$ detuning changes the
prepared data by $O(|a|h)$ and not their continuum first jet.  The diagonal
lapse margin is \Cref{lem:supp-exact-rational-detuning}.  The transverse
shift block is controlled by \Cref{thm:supp-exact-transverse}.  Together they
leave only constant lapse and curl-free shifts in the first-source kernel.
The entire first source is trace-free on this three-amplitude family, so
$K_0^{-1}$ acts by $1/2$ on its range.  The lapse estimate costs $h^{-6}$ and
the transverse reconstruction an additional two powers, giving a sufficient
smallest singular value $ch^8$ and therefore the displayed Gram bound.
\end{proof}

\paragraph{A source-neutral nontrivial family.}
For each axis let $T_i$ be the transverse off-diagonal tensor and let $P_i$
be the transverse diagonal trace-free tensor with entries $+1,-1$ on the two
coordinates orthogonal to axis $i$.  Put
\begin{equation}
 U=\sum_{i=1}^3r_i\left[T_i\cos(2\pi x_i+\phi_i)
          +\sigma_iP_i\sin(2\pi x_i+\phi_i)\right],
 \qquad r_i>0,\quad\sigma_i\in\{-1,1\}.
 \label{eq:supp-exact-balanced-family}
\end{equation}
The two polarizations on each axis are orthogonal and have equal Frobenius
norm.

\begin{theorem}[Exact source-neutral finite preparation]
\label{thm:supp-exact-neutral-preparation}
For the family \eqref{eq:supp-exact-balanced-family}, the complete weak
quadratic target vanishes at every odd cutoff.  On compact nonzero parameter
sets the original finite constraint map admits exact preparations
\begin{equation}
 \mathscr C_h(X_h(a))=0,
 \qquad X_h(a)=a(U,p_h)+O(a^2)
 \label{eq:supp-exact-neutral-prepared}
\end{equation}
uniformly for sufficiently fine odd cutoffs.  At a concrete three-site seed
the complete leading signed Gram is nonsingular; hence a relatively open
balanced neighborhood has an exact local original-action continuation for
small nonzero $a$.  On that branch the literal initial temporal and spatial
link curvatures obey
\begin{equation}
 \|E_h(a,0)\|_h+\|F_h^{\rm sp}(a,0)\|_h\le C_h|a|.
 \label{eq:supp-exact-initial-curvature}
\end{equation}
The constant in this last estimate is fixed-cutoff; no common physical time
is inferred.
\end{theorem}
\begin{proof}
The weak quadratic drift on one axis is a sum of two polarization invariants:
the norm difference and the Frobenius cross product.  Equal norms and
orthogonality make both vanish, independently of $h$.  The linear constraints
hold because the configurations and oscillatory momenta are transverse and
trace-free.  The nonconstant rows are solved by the uniform range inverse of
the initial-constraint chart.  The four means are solved by a homogeneous
trace momentum and three oscillatory balancing parameters; the derivative of
the quadratic mean map is uniformly invertible on compact nonzero
polarization sets.  This gives \eqref{eq:supp-exact-neutral-prepared}.  At the
explicit three-site seed, finite exact elimination gives a nonzero determinant
for the complete signed Gram, and analyticity preserves it on a neighborhood.
Because the normalized source target starts one amplitude order later on the
balanced branch, its true finite-amplitude response is $O_h(a)$.  Substitution
in the literal temporal-link derivative and spatial plaquette formulas gives
\eqref{eq:supp-exact-initial-curvature}.
\end{proof}

\subsection{Direct transverse response and the positive scalar flux equation}

A further weak-neutral preparation can be chosen so that the leading active
range is trace-free at every sufficiently fine odd cutoff.  On this family
the transverse source admits an exact least-squares factorization.  Let
$T_h$ be the transverse source, $B_h=T_h^*T_h$, and let the normalized
positive normal matrix $\mathsf H_h$ be the finite-band operator obtained by
factoring the phase curl and coefficient multiplication.  It satisfies
\begin{equation}
 \frac9{16}I\preceq\mathsf H_h\preceq324I.
 \label{eq:supp-exact-H-bounds}
\end{equation}
The actual pseudoinverse has the factorization
\begin{equation}
 T_h^\dagger=\frac{2\sqrt2}{h\omega_h}
 \Omega_h^{-1}\mathsf H_h^{-1}R_h\widehat{\mathcal L}_h^*\iota_o^*,
 \label{eq:supp-exact-T-dagger}
\end{equation}
with the operators defined by the physical Hodge and coefficient
factorization of the source; no inverse-then-project substitution is made.

\begin{theorem}[Actual direct transverse geometric suppression]
\label{thm:supp-exact-direct-transverse}
For the actual prepared transverse target $b_{\beta,h}$ and direct solution
$\beta_{0,h}=-2B_h^{-1}b_{\beta,h}$,
\begin{equation}
 \|\beta_{0,h}\|_{H_h^{s+3}}\le C_sh^{-1},
 \qquad \|\mathfrak E_h(\beta_{0,h})\|_{H_h^s}\le C_sh,
 \label{eq:supp-exact-direct-result}
\end{equation}
while the folded-face tail is exponentially small.  These are estimates of
the actual direct source response, not of a prescribed low-frequency
projection.
\end{theorem}
\begin{proof}
The factorization \eqref{eq:supp-exact-T-dagger} and
\eqref{eq:supp-exact-H-bounds} give two derivatives of Sobolev gain for
$B_h^{-1}$ at cost $h^{-2}$.  The prepared transverse target is $O(h)$ on a
fixed Fourier bank.  Applying the geometric observation operator, whose
symbol contributes the compensating $h^2$ factor at the required order,
gives the second estimate.  Exponential off-diagonal decay of
$\mathsf H_h^{-1}$ in cyclic Fourier distance gives the folded-face tail.
\end{proof}

After eliminating the transverse variable, the actual leading lapse solves
\begin{equation}
 \mathscr S_hn_h=g_h.
 \label{eq:supp-exact-scalar-equation}
\end{equation}
Let $Q_h$ project onto the physical off-diagonal complement and let
$K_{Q,h}$ be the induced positive coefficient metric.

\begin{theorem}[Exact positive scalar energy]
\label{thm:supp-exact-scalar-flux}
At every odd cutoff,
\begin{equation}
 \frac1{324}I\preceq K_{Q,h}\preceq\frac{16}{9}I,
 \label{eq:supp-exact-KQ-bounds}
\end{equation}
and the unchanged scalar operator is
\begin{equation}
 \mathscr S_h
 =\tfrac12\mathscr D_h^*\mathscr D_h
  +\tfrac{h^2}{2}\sigma_h^*K_{Q,h}^{-1}\sigma_h.
 \label{eq:supp-exact-S-flux}
\end{equation}
Hence
\begin{equation}
 2\langle n,\mathscr S_hn\rangle_h
 =\|\mathscr D_hn\|_h^2
  +h^2\langle\sigma_h(n),K_{Q,h}^{-1}\sigma_h(n)\rangle_h.
 \label{eq:supp-exact-scalar-energy}
\end{equation}
On the bare diagonal soft subspace the normalized full energy diverges under
refinement, so a pole predicted by the diagonal compression alone is not a
pole of the full physical scalar energy.  This compression result does not by
itself prove a uniform inverse for \eqref{eq:supp-exact-scalar-equation}.
\end{theorem}
\begin{proof}
The coefficient map giving the off-diagonal source is uniformly invertible,
so compression of its positive inverse metric gives
\eqref{eq:supp-exact-KQ-bounds}.  Orthogonal projection of the physical
source yields the exact residual factor
$hM_h^{-*}Q_hK_{Q,h}^{-1}\sigma_h(n)$; squaring and adding the diagonal
source gives \eqref{eq:supp-exact-S-flux}.  On the finite-dimensional soft
subspace the limiting flux is nonzero on every nonzero vector, while the
bare diagonal contribution is higher order; compactness therefore forces the
normalized full energy to diverge.
\end{proof}

\subsection{Dynamical tangency: an exact obstruction and an upstream escape direction}

Exact preparation is not the same as invariance of the prepared set under the
original action vector field.  Let $\mathcal W$ be the weak curl-free
multiplier space and let $\mathbb K_1(Y)$ be the first Poisson matrix at a
canonical first jet $Y$.  For a leading multiplier tilt $\ell$ define the
weak null space
\begin{equation}
 \mathcal N_W(Y)=\{\nu\in\mathcal W:
 \mathbb K_1(Y)[\nu,\mu]=0\ \hbox{for every multiplier }\mu\},
 \label{eq:supp-exact-weak-null}
\end{equation}
and the first tangency functional
\begin{equation}
 \Delta_\nu(Y,\ell)=
 Dd_{2,W}(Y)[F_1Y+G\ell,\nu]
 +\mathbb K_1(F_1Y+G\ell)[\nu,\ell].
 \label{eq:supp-exact-tangency-functional}
\end{equation}
It is determined entirely by the canonical first jet and leading multiplier
tilt; higher preparation coefficients do not enter this leading null test.

\begin{theorem}[Initial-layer obstruction on an exactly prepared branch]
\label{thm:supp-exact-initial-layer}
There is a nontrivial exactly prepared three-site original-action family
$X_a$, $\lambda_a$ with
\begin{equation}
 \mathsf B_W(X_a,\lambda_a)=0,
 \qquad (\lambda_t)_W(a,0)=0,
 \qquad \|\operatorname{Riem}(g_h)(a,0)\|=O_h(a),
 \label{eq:supp-exact-prepared-zero-rate}
\end{equation}
but with a weak null test $\nu_*$ satisfying
$\Delta_{\nu_*}(Y,\ell)\ne0$.  Along the exact continuation,
\begin{equation}
 \left.\frac{\dd}{\dd t}\mathsf B_W[\nu_*]\right|_{t=0}
 =a^2\Delta_{\nu_*}(Y,\ell)+O_h(a^3),
 \label{eq:supp-exact-tangency-growth}
\end{equation}
which forces a weak multiplier acceleration of order $a^{-2}$ in the
normalized grading.  For every fixed $S$ the exact histories exist on
$0\le t\le a^3S$ for sufficiently small $a$, and the unscaled rates and
curvature remain $O_h(a)$ there.  Nevertheless no constants $T,C,a_0>0$ can
make
\begin{equation}
 \sup_{0\le t\le T}\|R_{0i0j}(g_h)(a,t)\|\le Ca
 \label{eq:supp-exact-no-linear-fixed-time}
\end{equation}
hold for every $0<a<a_0$ on this family.
\end{theorem}
\begin{proof}
Taylor expansion of the exact weak drift gives
$D_X\mathsf B_W=a\mathcal C_\ell+O_h(a^2)$ and
$D_\lambda\mathsf B_W=a\mathcal K_W+O_h(a^2)$.  Pairing with a Poisson-null
$\nu_*$ removes the unknown active multiplier acceleration at leading order
and yields \eqref{eq:supp-exact-tangency-growth}.  Rescale time by
$t=a^3s$; the finite-dimensional analytic vector field converges on every
fixed $s$-interval to a linear displacement system, giving the shrinking
layer and the nonzero limiting curvature slope.  If a uniform $Ca$ bound
held on a fixed physical interval, choose a fixed $s$ beyond the corresponding
normalized bound and then take $a$ sufficiently small so that $a^3s<T$,
contradicting that limiting slope.
\end{proof}

The leading multiplier tilt has a finite affine freedom after the lower
quadratic, normalization and cubic-mean conditions are imposed.  On the
preceding canonical first jet the tangency functional is constant and nonzero
on that entire admissible tilt family; hence changing only the leading
multiplier tilt cannot repair the obstruction.

The canonical first jet itself has additional admissible directions.  Write
$\mathcal L(Y,\ell)$ for the complete lower incidence consisting of the
quadratic drift, multiplier normalizations and cubic mean rows.  At the
selected base point its multiplier derivative has constant rank, and one
canonical direction $Z_2$ admits an exact multiplier tangent $m_2$ with
\begin{equation}
 D\mathcal L(Y,\ell)[Z_2,m_2]=0.
 \label{eq:supp-exact-lower-tangent}
\end{equation}

\begin{theorem}[The fixed-first-jet tangency obstruction is deformable]
\label{thm:supp-exact-upstream-deformation}
For the exact lower-incidence tangent $(Z_2,m_2)$ above, one weak null test
$\nu_\diamond$ remains Poisson-null to first order and
\begin{equation}
 D\Delta_{\nu_\diamond}(Y,\ell)[Z_2,m_2]\ne0.
 \label{eq:supp-exact-upstream-nonzero}
\end{equation}
Thus the nonzero tangency obstruction on the preceding prepared first jet is
not invariant under all admissible canonical first-jet deformations of the
unchanged finite action.
\end{theorem}
\begin{proof}
Differentiate every row of the finite lower-incidence map in the selected
canonical direction and solve the resulting characteristic-zero linear system
for $m_2$.  The same weak test has zero first variation of its Poisson-null
condition in each coordinate direction of the local first-jet chart.  The
tangency functional is quadratic in the simultaneous first-jet/multiplier
increment, so its derivative is obtained exactly by polarization.  Direct
finite algebra gives a nonzero rational coefficient; hence the derivative is
nonzero over characteristic zero.  The calculation establishes a tangent,
not an exact nonlinear root.
\end{proof}


\subsection{Certified same-event weak Poisson and rapid spectral packet}
\label{subsec:supp-exact-certified-rapid}

The spectral input for the conditional fixed-cutoff boundary theorem is
recorded at one prepared event of the unchanged $N=3$, $h=1/3$ action.
The rank and spectral certificates in this subsection do not by themselves
verify the complete graded-chart assumption used by that theorem.  The
preparation equations define a regular finite root box $\mathcal B_*$ around
$(Y_*,\ell_*)$; all statements in this subsection hold uniformly on that
box.  The weak multiplier space splits as
\begin{equation}
 \mathbb R^{29}=\mathcal H\oplus\mathcal G,
 \qquad \dim\mathcal H=3,
 \qquad \dim\mathcal G=26,
 \label{eq:supp-certified-weak-split}
\end{equation}
where $\mathcal H$ is the exact constant-shift space.  Translation invariance
gives
\begin{equation}
 K_{WW}=\begin{pmatrix}0&0\\0&K_g\end{pmatrix},
 \qquad K_g^{\mathsf T}=-K_g.
 \label{eq:supp-certified-Kblock}
\end{equation}
After eliminating the regular active block, let $D_W$ be the symmetric weak
mass and let $D_{26}$ be its Schur complement on $\mathcal G$.  The finite
mass certificate gives
\begin{equation}
 \operatorname{inertia}(D_W)=(18,11,0),
 \qquad
 \operatorname{inertia}(D_{26})=(15,11,0).
 \label{eq:supp-certified-mass-inertia}
\end{equation}

\begin{theorem}[Exact rank of the weak first-Poisson block]
\label{thm:supp-certified-weak-rank}
On $\mathcal B_*$,
\begin{equation}
 \rank K_g=24,\qquad\rank K_{WW}=24.
 \label{eq:supp-certified-rank24}
\end{equation}
Hence
\begin{equation}
 \ker K_{WW}=\mathcal Z_g\oplus\mathcal H,
 \qquad \dim\mathcal Z_g=2,
 \qquad \dim\mathcal H=3.
 \label{eq:supp-certified-kernel}
\end{equation}
\end{theorem}
\begin{proof}
For the lower bound, delete gradient indices $2$ and $24$ in the fixed
zero-based ordering.  The resulting principal $24\times24$ interval block
$K_{24}$ admits a verified preconditioner $R_{24}$ with
\begin{equation}
 \|I-R_{24}K_{24}\|_\infty<8.80\times10^{-14},
 \qquad
 \|K_{24}^{-1}\|_\infty<2.56\times10^{-2}.
 \label{eq:supp-certified-K24}
\end{equation}
The Neumann series therefore proves $\rank K_g\ge24$ throughout the box.

For the upper bound, on the ten-profile carrier the exact gradient block has
the form
\begin{equation}
 K_g(Y)=1458\sum_{r=1}^4\xi_rM_r,
 \qquad M_r^{\mathsf T}=-M_r,
 \qquad \max_{i,j,r}|(M_r)_{ij}|\le6,
 \label{eq:supp-certified-pencil}
\end{equation}
with integral $26\times26$ matrices $M_r$.  Let
$P(\xi)=\operatorname{Pf}(\sum_r\xi_rM_r)$, a homogeneous degree-thirteen
integer polynomial.  After setting $\xi_4=1$, its degree in each of the
remaining three variables is at most thirteen.  Exact evaluation on the
full grid $\{0,\ldots,13\}^3$ gives zero for each of the primes
\begin{equation}
 1000000007,\quad1000000009,\quad1000000033,\quad1000000087.
 \label{eq:supp-certified-primes}
\end{equation}
Repeated univariate interpolation therefore makes every coefficient
divisible by their product.  On the other hand each coefficient obeys
\begin{equation}
 |c|\le25!!\,24^{13}
 =6929388469032406070392258560000,
 \label{eq:supp-certified-pf-bound}
\end{equation}
and the prime product is
$1000000136000004854000053496000180873>2|c|$.  Thus every coefficient is
zero in characteristic zero.  The Pfaffian vanishes identically, so every
$K_g$ has rank at most twenty-four.  Combining both bounds proves
\eqref{eq:supp-certified-rank24}; the three exact constant shifts give
\eqref{eq:supp-certified-kernel}.
\end{proof}

The constant shifts factor from the generalized weak pencil:
\begin{equation}
 \det(K_{WW}+\lambda D_W)
 =\lambda^3\det(D_{HH})\det(K_g+\lambda D_{26}).
 \label{eq:supp-certified-pencil-factor}
\end{equation}
Thus the nontrivial rapid spectrum is carried by the $26\times26$ generator
$A_g=-D_{26}^{-1}K_g$.

\begin{theorem}[Certified weak spectral trichotomy]
\label{thm:supp-certified-rapid-spectrum}
On $\mathcal B_*$ the reduced generator $A_g$ has six simple positive real
eigenvalues and their six negative partners, one simple complex quartet with
two eigenvalues in each open half-plane, four simple purely imaginary
conjugate pairs, and a semisimple double zero.  The nonzero hyperbolic
spectrum obeys
\begin{equation}
 |\operatorname{Re}\lambda|>0.0015069.
 \label{eq:supp-certified-hyperbolic-gap}
\end{equation}
Consequently the full $29\times29$ weak generator $B_W$ has an invariant
real splitting
\begin{equation}
 \mathbb R^{29}=E_-\oplus E_0\oplus E_+,
 \qquad(\dim E_+,\dim E_-,\dim E_0)=(8,8,13),
 \label{eq:supp-certified-8813}
\end{equation}
with constants $M\ge1$, $\kappa>0$ such that
\begin{equation}
 \|e^{-sB_W}|_{E_+}\|\le Me^{-\kappa s},\qquad
 \|e^{sB_W}|_{E_-}\|\le Me^{-\kappa s},\qquad
 \sup_{s\in\mathbb R}\|e^{sB_W}|_{E_0}\|\le M.
 \label{eq:supp-certified-semigroup}
\end{equation}
\end{theorem}
\begin{proof}
A verified rational similarity transform for $A_g$ puts the interval family
inside disjoint Gershgorin components.  The similarity error on the complete
box is less than $1.558\times10^{-25}$.  The separated components consist of
six positive and six negative real singletons, one quartet separated from the
imaginary axis, four imaginary conjugate singleton pairs, and one
two-eigenvalue component at the origin.  Reality and the skew/symmetric pencil
symmetry $\lambda\mapsto-\lambda$ imply that the imaginary singleton pairs
are exactly imaginary.  The only possible final two-eigenvalue cluster is a
real pair, an imaginary pair, or a double zero.  By
\Cref{thm:supp-certified-weak-rank}, $\ker A_g=\ker K_g$ has dimension two;
therefore the final cluster is exactly a semisimple double zero.  Counting
modes gives eight stable, eight unstable and ten reduced centre directions;
restoring the three exact constant shifts gives
\eqref{eq:supp-certified-8813}.  Finite-dimensional spectral calculus gives
\eqref{eq:supp-certified-semigroup}.
\end{proof}

The preceding theorems prove \Cref{thm:supp-certified-rapid-spectrum}.  Their
role is local and fixed-cutoff: no spatial-refinement statement is extracted
from the numerical constants.

\subsection{A complete amplitude-graded finite chart}
\label{subsec:supp-exact-chart}

We now fix $N=3$, $h=1/3$, $V=27$ and the mean lapse equal to one. The
parameter $a>0$ is amplitude, not mesh size. We give the exact hypotheses
needed for the boundary and slow-time results, rather than suppressing the
finite coefficient problem in the phrase ``regular branch.'' All sources,
normalizations, constraints and readouts below are restrictions of the same
original Hamiltonian $H_h$.

\begin{assumption}[Amplitude-graded exact-action chart]
\label{ass:supp-graded-exact-chart}
The original finite system has an analytic prepared record
$Q_a^0=(X_a^0,\lambda_a^0)$ with $P_h(Q_a^0)=0$,
$X_a^0=aY_*+O(a^2)$ and $\lambda_a^0=e+a\ell_*+O(a^2)$.
The normalized multiplier rate $r$ solves
$B_h+M_hr=0$ with zero mean-lapse derivative. The following data and
inequalities refer to a common neighborhood and fixed-cutoff constants.
\begin{enumerate}[label=\textup{(C\arabic*)},leftmargin=3em]
\item All $108$ constraint rows are retained, in $104$ nonconstant and four
mean rows. Their divided canonical normal derivative is an onto map
$\mathcal C_0:\mathbb R^{324}\to\mathbb R^{108}$. Choose
$J_c:\mathbb R^{108}\to\mathbb R^{324}$ with $\mathcal C_0J_c$ invertible,
and $E:\mathbb R^{216}\to\ker\mathcal C_0$, with free projection
$\Lambda E=I$, $\Lambda J_c=0$. The normalized multipliers split into
$78$ active and $29$ weak coordinates, with a fixed shear $H$.
\item For $w=(q,p_A,p_W)$ the substitution
\begin{equation}
 \begin{split}
 X&=X_a^0+a^2(Eq+J_c n),\\
 \lambda_A&=\lambda_{a,A}^0+a^2(p_A-Hp_W),\qquad
 \lambda_W=\lambda_{a,W}^0+ap_W
 \end{split}
 \label{eq:supp-exact-normal-chart}
\end{equation}
into $P_h$, followed by division of the nonconstant and mean rows by $a^2$
and $a^3$, defines an analytic map $\mathcal P(a,n,w)$. It satisfies
$\mathcal P(a,0,0)=0$, $\mathcal P(0,n,w)=\mathcal C_0J_cn$,
and $D_q\mathcal P=O(a)$, $D_p\mathcal P=O(a^2)$, with these factors
preserved by the fixed-order derivatives used below.
\item In unbalanced free coordinates $v=(v_s,v_f)\in\mathbb R^{294}
\oplus\mathbb R^{29}$, and time $s=t/a^3$, the original constrained field is
\begin{equation}
 \begin{split}
 v_s'&=a^3b_s+R_s(a,v),\\
 v_f'&=C_{\rm red}v_s+Bv_f+a^3b_w+R_f(a,v),
 \end{split}
 \label{eq:supp-exact-unbalanced-field}
\end{equation}
where $R_s,R_f$ are analytic and
\begin{equation}
 \begin{aligned}
 |R_s|&\le C(a^4+a^2|v|+a|v|^2),&
 |D_vR_s|&\le C(a^2+a|v|),\\
 |R_f|&\le C(a^4+a|v|+|v|^2),&
 |D_vR_f|&\le C(a+|v|).
 \end{aligned}
 \label{eq:supp-exact-graded-remainders}
\end{equation}
The transport identity is $C_{\rm red}b_s+Bb_w=0$.
For $D=\operatorname{diag}(0_{294},B)$ there are commuting projections
$\Pi_++\Pi_-=I$ and constants $M_D,\kappa>0$ such that
\begin{equation}
 \|e^{sD}\Pi_-\|\le M_D,\qquad
 \|e^{-sD}\Pi_+\|\le M_De^{-\kappa s}\quad(s\ge0).
 \label{eq:supp-exact-dichotomy}
\end{equation}
No rank is prescribed in the boundary argument.  At the certified event of
\Cref{thm:supp-certified-rapid-spectrum}, the weak block $B$ is the generator
$B_W$ of \eqref{eq:supp-certified-8813}, $\Pi_+$ is the eight-dimensional
unstable projector and $\Pi_-$ contains the stable and bounded-centre
subspaces.
\item The original canonical vector field $F=\mathbb J\nabla_XH_h$ has
bounded derivatives on the physical chart. The metric and ordered-link
readouts are their actual smooth fixed-cutoff jet functions. In particular,
metric curvature uses $(X,\lambda,X_t,\lambda_t,X_{tt})$, with no second
time derivative of lapse or shift. The stationary spatial connection $z$
has derivative orders $z_X=O(1)$, $z_{\lambda_A}=O(a)$ and
$z_{\lambda_W}=O(a^2)$ on the relevant tubes, with their differentiated
bounds. The temporal connection has bounded derivatives. At zero amplitude
the reference curvature vanishes.
\end{enumerate}
\end{assumption}

These are analytic divisibilities and actual source/readout comparisons,
not merely formal powers along one curve. For example, the divided mean
rows in item \textup{(C2)} retain the four means rather than remove them.
The certified event of \Cref{thm:supp-certified-weak-rank,thm:supp-certified-rapid-spectrum}
supplies the weak rank and spectral part of this chart at one concrete
prepared shape.  The analytic divisibility, source and readout bounds in the
present assumption are the remaining finite chart conditions used by the
boundary argument.  No constant in this fixed-$N=3$ certificate is asserted
to be uniform at other lattice sizes.

\begin{lemma}[Exact nonlinear constraint coordinates]
\label{lem:supp-exact-normal-coordinates}
Under \Cref{ass:supp-graded-exact-chart}, there is an analytic
$n=n(a,w)$ such that $P_h(\Theta_a(w))=0$, where $\Theta_a$ is
\eqref{eq:supp-exact-normal-chart} with this value of $n$. Moreover,
\begin{equation}
 \begin{gathered}
 n(a,0)=0,\qquad n=O(a|q|+a|q|^2+a^2|p|),\\
 D_qn=O(a),\qquad D_pn=O(a^2),\\
 D_q\mathcal X_a=O(a^2),\quad D_p\mathcal X_a=O(a^4),\quad
 D_q^2\mathcal X_a=O(a^3),\quad
 D_qD_p\mathcal X_a,D_p^2\mathcal X_a=O(a^4),
 \end{gathered}
 \label{eq:supp-exact-normal-derivatives}
\end{equation}
where $\mathcal X_a$ is the canonical component of $\Theta_a$.
For positive small $a$ this is a coordinate chart on the full
$323$-dimensional normalized constraint leaf.
\end{lemma}
\begin{proof}
The derivative $D_n\mathcal P(0,0,0)=\mathcal C_0J_c$ is invertible.
The analytic implicit-function theorem therefore solves
$\mathcal P(a,n,w)=0$ locally. At $a=0$ the unique solution is $n=0$,
and $\mathcal P(a,0,0)=0$ gives $n(a,0)=0$. Differentiating the implicit
equation gives
$D_qn=-(D_n\mathcal P)^{-1}D_q\mathcal P$ and the analogous formula for
$p$. The assumed analytic factors give the stated derivative orders,
including their further derivatives. Integration from $w=0$ yields the
bound on $n$. Differentiating \eqref{eq:supp-exact-normal-chart} yields the
remaining bounds. The free projection recovers $q$, and the affine
multiplier formulas recover $p_A,p_W$; invertibility of the normal
derivative completes the coordinate inverse. All original constraint rows
vanish because their divided equations vanish, with $a>0$.
\end{proof}

Evaluation of $\Theta_a$ is a coordinate solve, not a force projecting an
unconstrained physical history back onto the constraint set.

\subsection{Nonlinear two-point construction on an amplitude-squared interval}
\label{subsec:supp-exact-boundary}

In the balanced coordinate $v=(ax,y)$, $w=(x,y)$, $x=(q,p_A)$, the exact
reduced vector field $\mathcal H_a$ is
\begin{equation}
 (\mathcal H_a)_q=a\Lambda F(\Theta_a(w)),\qquad
 (\mathcal H_a)_A=a(r_A+aHr_W)(\Theta_a(w)),\qquad
 (\mathcal H_a)_W=a^2r_W(\Theta_a(w)).
 \label{eq:supp-exact-reduced-field}
\end{equation}
The inverse multiplier solve in $r$ is the actual signed solve, not a
positive surrogate. Set
\begin{equation}
 w_{\rm ref}(s)=(a^2s b_s,a^3s b_w),\qquad
 f_a=\mathcal H_a(w_{\rm ref})-w_{\rm ref}'.
 \label{eq:supp-exact-reference}
\end{equation}

\begin{lemma}[Balanced derivative and reference residual estimates]
\label{lem:supp-exact-balanced-field}
On $|w|\le M_0a$ and for $1\le S\le c/a$,
\begin{equation}
 \begin{gathered}
 \|D_w\mathcal H_a-D\|\le C_{M_0}a,\qquad
 \|D_w^j\mathcal H_a\|\le C_j\quad(2\le j\le4),\\
 \|D_w^j(\mathcal H_a)_q\|\le C_j a^2\quad(1\le j\le3),\\
 \|f_a\|_\infty\le Ca^3,\qquad
 \|f_a'\|_\infty\le Ca^4,\qquad
 \|w_{\rm ref}\|_\infty\le Cca,\quad |w_{\rm ref}'|\le Ca^2.
 \end{gathered}
 \label{eq:supp-exact-balanced-estimates}
\end{equation}
\end{lemma}
\begin{proof}
Conjugate \eqref{eq:supp-exact-unbalanced-field} by
$\operatorname{diag}(aI_{294},I_{29})$. The perturbation of $D$ has blocks
$D_sR_s$, $a^{-1}D_fR_s$, $aC_{\rm red}+aD_sR_f$, and $D_fR_f$.
The bounds \eqref{eq:supp-exact-graded-remainders} make each $O(a)$ on
this tube. Analyticity and the factor $a$ in the entire slow remainder
remove the apparent singular division and bound higher balanced derivatives.
The canonical row in \eqref{eq:supp-exact-reduced-field}, together with
\eqref{eq:supp-exact-normal-derivatives} and bounded derivatives of $F$,
gives its extra $a^2$ factor.
The unbalanced reference is $v_{\rm ref}=a^3s(b_s,b_w)$.
Its linear terms cancel by $C_{\rm red}b_s+Bb_w=0$. The balanced residual
is consequently bounded by sums of $a^3$, $a^4s$ and $a^6s^2$ in the
slow row and $a^4$, $a^4s$, $a^6s^2$ in the weak row. These are $O(a^3)$
for $s\le c/a$. Differentiating along $v_{\rm ref}'=a^3(b_s,b_w)$ and
using the derivative bounds gives $O(a^4)$. The other two estimates follow
directly from the reference formula.
\end{proof}

Define the forward--backward Green operator
\begin{equation}
 (\mathcal G_Sf)(s)=\int_0^s e^{(s-r)D}\Pi_-f(r)\,\dd r
                  -\int_s^S e^{(s-r)D}\Pi_+f(r)\,\dd r.
 \label{eq:supp-exact-green}
\end{equation}
It satisfies $\|\mathcal G_S\|\le C_G(1+S)$. The factor $S$ comes from
the bounded center, not the exponentially decaying kernels. Put
$N_a(s,e)=\mathcal H_a(w_{\rm ref}+e)-\mathcal H_a(w_{\rm ref})-De$
and iterate
\begin{equation}
 e^{[0]}=0,\qquad
 e^{[m+1]}=\mathcal G_S\{f_a+N_a(\cdot,e^{[m]})\},\qquad
 w^{[m]}=w_{\rm ref}+e^{[m]}.
 \label{eq:supp-exact-nonlinear-sweeps}
\end{equation}
The forward--backward placement is the standard decoupling mechanism for
unstable boundary problems \cite{Mattheij1985}; the amplitude estimates
below pertain to this particular finite field.

\begin{theorem}[Exact constrained boundary history]
\label{thm:supp-exact-boundary}
For sufficiently small fixed $c,a_0>0$, $0<a<a_0$ and $1\le S\le c/a$,
\eqref{eq:supp-exact-nonlinear-sweeps} is a contraction on
$\|e\|_\infty\le K a^2$. Its unique fixed point gives
\begin{equation}
 \begin{gathered}
 w_*'=\mathcal H_a(w_*),\qquad
 \Pi_-(w_*(0)-w_{\rm ref}(0))=0,\qquad
 \Pi_+(w_*(S)-w_{\rm ref}(S))=0,\\
 \|w_*-w_{\rm ref}\|_\infty\le Ca^2,\qquad
 \|(w_*-w_{\rm ref})'\|_\infty\le Ca^3.
 \end{gathered}
 \label{eq:supp-exact-boundary-estimate}
\end{equation}
The reconstructed curve $\Theta_a(w_*(t/a^3))$ solves the original full
normalized action equations. Every ideal iterate also satisfies all
constraints, without in general satisfying the evolution equation.
\end{theorem}
\begin{proof}
On the indicated ball, \Cref{lem:supp-exact-balanced-field} gives
$\operatorname{Lip}_eN_a\le L_0a$. Choose $c$ and then $a_0$ so that
$q=C_G(1+S)L_0a\le1/4$. Also
$\|\mathcal G_S f_a\|\le C(1+S)a^3\le C_0a^2$.
Taking $K=2C_0$ makes the ball invariant, and the contraction theorem gives
the unique fixed point. The integral equation is equivalent to the
specified differential and boundary conditions.

For the derivative, write $F_e=f_a+N_a(s,e_*)$. Integration by parts in
\eqref{eq:supp-exact-green} gives
\[
 e_*'(s)=e^{sD}\Pi_-F_e(0)+e^{(s-S)D}\Pi_+F_e(S)
                    +\mathcal G_S F_e'(s).
\]
Here $\|F_e\|\le Ca^3$ and
$\|F_e'\|\le Ca^4+L_0a\|e_*'\|$, because the derivative at fixed $e$
is a difference of $D\mathcal H_a$ applied to $w_{\rm ref}'$.
The endpoint factors are bounded in their prescribed time directions.
Absorption using the same $q<1$ proves $\|e_*'\|\le Ca^3$.

The reconstructed velocity and the original vector field agree in every
free coordinate by \eqref{eq:supp-exact-reduced-field}. Both are tangent to
$P_h=0$. Their difference is therefore a canonical normal vector annihilated
by the invertible normal constraint derivative, and is zero. This recovers
the complete canonical evolution, not just its free projection. Constraint
satisfaction of all iterates follows separately from the definition of
$\Theta_a$.
\end{proof}

\begin{theorem}[Exact-action boundary realization on a graded chart]
\label{thm:main-exact-boundary}
Fix $N=3$ and $h=1/3$, and assume the complete amplitude-graded
exact-action chart of \Cref{ass:supp-graded-exact-chart}, including its
analytic divisibilities, source estimates, transport identity and physical
readout bounds on one common neighborhood. The certified weak spectrum
supplies the spectral part of that assumption but not its other rows.
Then there are $c,a_0,C>0$ such that, for every $0<a<a_0$, the unchanged original finite
action admits an exact boundary-selected history
$Q_{*,a}=(X_{*,a},\lambda_{*,a})$ of
\eqref{eq:supp-exact-euler}, with all $108$ original constraint rows
identically zero, on $0\le t\le ca^2$.  Its initial record preserves the
certified first field and first multiplier tilt,
\begin{equation}
 X_{*,a}(0)=aY_*+O(a^2),\qquad
 \lambda_{*,a}(0)=e+a\ell_*+O(a^2).
 \label{eq:main-exact-boundary-first-jet}
\end{equation}
Relative to the exact prepared record through the same first jet, the
boundary selector spends only higher preparation coefficients:
\begin{equation}
 \Delta X(0)=O(a^4),\qquad
 \Delta\lambda_A(0)=O(a^4),\qquad
 \Delta\lambda_W(0)=O(a^3).
 \label{eq:main-exact-boundary-selector-cost}
\end{equation}
In the actual metric and ordered-link readouts,
\begin{equation}
 \sup_{0\le t\le ca^2}
 \bigl(\|\operatorname{Riem}(g_{*,a})\|
   +\|E_{*,a}^{\rm link}\|+\|F_{*,a}^{\rm sp,link}\|\bigr)\le Ca.
 \label{eq:main-exact-boundary-curvature}
\end{equation}
The boundary construction is rank-free: the nonexpanding coordinates are
fixed at the initial cut and the $E_+$ coordinates of
\eqref{eq:supp-certified-8813} are fixed at the terminal cut.  Thus no
prescribed value of the unstable dimension is used in the Green argument.

The ideal nonlinear iterates remain exactly constrained and converge to this
history.  After $m\ge2$ iterations their readout errors are bounded by
$C2^{-m}$, $Ca^2 2^{-m}$ and $Ca^4 2^{-m}$ for the metric, temporal-link
and spatial-link curvatures, respectively.  Hence
$m=3+\lceil p\log_2(1/a)\rceil$ gives errors
$O(a^p)$, $O(a^{p+2})$, and $O(a^{p+4})$ on the entire interval.
\end{theorem}
The same history also satisfies $X_t,\lambda_t,X_{tt}=O(a)$ and
$Q_*(t)-Q_a^0=O(a^3)$ on its full interval.

\begin{proof}[Proof of the anchoring and curvature assertions]
At $s=0$ the forward--backward correction is $O(a^2)$ in the balanced
coordinates.  Inserting this estimate in
\eqref{eq:supp-exact-normal-chart}, together with
$D_qn=O(a)$ and $D_pn=O(a^2)$, gives the three orders in
\eqref{eq:main-exact-boundary-selector-cost}; in particular the first field and
first multiplier tilt are unchanged.  Throughout the interval the slow
coordinates of $w_*$ are $O(a)$, its weak coordinates are $O(a^2)$, and
their scaled rates are $O(a^2)$ and $O(a^3)$, respectively. The physical
coordinate formulas and $t=a^3s$ give displacement $O(a^3)$ and
$X_t,\lambda_t=O(a)$. Differentiating the unchanged canonical equation
gives $X_{tt}=D_XF X_t+D_\lambda F\lambda_t=O(a)$.
Spatial differentiation is bounded at this fixed cutoff. The smooth
metric readout therefore has curvature $O(a)$; equivalently its
Gauss--Codazzi expressions require these jets and first, but not second,
multiplier derivatives. The stationary connection orders in the certificate
give the same bound for its actual ordered-link curvatures. No equality of
metric and link curvature at a fixed cutoff is used.
\end{proof}


\subsection{Approximate curves, time jets, and readout accuracy}
\label{subsec:supp-exact-accuracy}

\begin{proposition}[Accuracy of the constrained nonlinear iterates]
\label{prop:supp-exact-sweep-accuracy}
For $m\ge2$ write $\rho_m=2^{-m}$ and
$\delta_m=w^{[m]}-w_*$. On the entire boundary interval,
\begin{equation}
 \max_{0\le j\le2}\|\partial_s^j\delta_m\|_\infty\le Ca^2\rho_m,
 \qquad
 \|\partial_s^j(\delta_m)_q\|_\infty\le Ca^4\rho_m\quad(j=1,2).
 \label{eq:supp-exact-sweep-jets}
\end{equation}
After nonlinear reconstruction the physical errors are
\begin{equation}
 \begin{array}{c|ccccccc}
 \text{quantity}&X&\lambda_A&\lambda_W&X_t&(\lambda_A)_t&(\lambda_W)_t&X_{tt}\\\hline
 \text{error}/(C\rho_m)&a^4&a^4&a^3&a^3&a&1&1.
 \end{array}
 \label{eq:supp-exact-physical-error-table}
\end{equation}
The actual readout errors satisfy
\begin{equation}
 \begin{split}
 \|\operatorname{Riem}(g_m)-\operatorname{Riem}(g_*)\|_\infty&\le C\rho_m,\\
 \|E_m^{\rm link}-E_*^{\rm link}\|_\infty&\le Ca^2\rho_m,\\
 \|F_m^{\rm sp,link}-F_*^{\rm sp,link}\|_\infty&\le Ca^4\rho_m.
 \end{split}
 \label{eq:supp-exact-reader-accuracy}
\end{equation}
\end{proposition}
\begin{proof}
The contraction gives $\|\delta_m\|\le Ca^2 4^{-m}$. Differentiate the
Green equation:
$\delta_m'=D\delta_m+N_a(s,e^{[m-1]})-N_a(s,e_*)$.
This bounds its first derivative with a fixed index loss. A second
differentiation uses bounded second derivatives of $\mathcal H_a$,
$w_{\rm ref}'=O(a^2)$ and $e_*'=O(a^3)$, with another fixed index loss.
These losses are absorbed by $C2^{-m}$. The free canonical row of $D$ is
zero, and $(w^{[m]})_q'=(\mathcal H_a)_q(w^{[m-1]})$; its derivative
bounds gain the extra $a^2$ in \eqref{eq:supp-exact-sweep-jets}.

A physical derivative costs $a^{-3}$. The multiplier coordinates are affine,
so their errors follow immediately. For $X_t$ and $X_{tt}$ use the chain
rule with \eqref{eq:supp-exact-normal-derivatives}. The leading scaled
canonical derivative errors are $a^2(a^4\rho_m)$ and
$a^4(a^2\rho_m)$, both $O(a^6\rho_m)$. The second scaled derivative has
the same leading order; products involving second chart derivatives are
smaller. Division by $a^3$ and $a^6$ yields the table.
Smoothness of the metric jet readout gives its stated error. For the
stationary spatial connection, the three source derivative orders give
$\Delta z=O(a^4\rho_m)$ and
$\Delta z_t=O(a^2\rho_m)$: the direct contributions to the latter are
$a^3\rho_m$, $a\cdot a\rho_m$ and $a^2\cdot\rho_m$.
Changes of coefficients multiplied by the exact $O(a)$ rates are no larger.
The temporal connection has error $O(a^3\rho_m)$. The actual temporal-link
formula uses $z_t$ and the shifted temporal connection; the spatial-link
formula uses only spatial connection values. Their smooth fixed-cutoff
expressions yield the last two estimates.
\end{proof}

Taking $m=3+\lceil p\log_2(1/a)\rceil$ proves the accuracy assertion in
\Cref{thm:main-exact-boundary}. The ideal sweeps contain exact integrals and
exact constraint-coordinate solves. A numerical implementation needs a
separate residual certificate.

\begin{proposition}[A posteriori boundary residual estimate]
\label{prop:supp-exact-boundary-residual}
Let $\widetilde w$ remain in the same contraction tube and set
$r=\widetilde w'-\mathcal H_a(\widetilde w)$,
$d_-=\Pi_-(\widetilde w(0)-w_{\rm ref}(0))$,
$d_+=\Pi_+(\widetilde w(S)-w_{\rm ref}(S))$. Then
\begin{equation}
 \|\widetilde w-w_*\|_\infty
 \le\frac{M_D(|d_-|+|d_+|)+C_G(1+S)\|r\|_\infty}{1-q},
 \qquad q\le\tfrac14.
 \label{eq:supp-exact-residual-bound}
\end{equation}
To obtain curvature accuracy one additionally controls the actual
first two scaled derivatives of this error and the differentiated normal
solve; a position residual alone is not a curvature certificate.
\end{proposition}
\begin{proof}
Variation of constants gives the same Green fixed-point equation with
additional terms $e^{sD}d_-+e^{(s-S)D}d_++\mathcal G_Sr$.
Subtract the exact equation and absorb its $q$-Lipschitz term. The final
statement follows from the physical derivative losses in
\eqref{eq:supp-exact-physical-error-table}. For an inexact normal solve,
constraint residuals are its actual divided residuals multiplied back by
$a^2$ and $a^3$, not exact zeros.
\end{proof}

\subsection{Invariant graphs and exact normal expansion}
\label{subsec:supp-exact-normals}

The boundary construction selects an exact history but does not make its
initial-value problem stable. The following formulation separates the
nonlinear invariance required for a residence estimate from a frozen
spectral computation. We use the invariant-graph method in finite dimension;
see \cite{BatesLuZeng} for the broader invariant-foliation setting.

\begin{lemma}[A quantitative forward invariant-graph criterion]
\label{lem:supp-exact-graph-criterion}
Consider a localized $C^3$ field
\begin{equation}
 \dot u=A_u u+b_u+N_u(u,n),\qquad
 \dot n=A_n n+b_n+N_n(u,n)
 \label{eq:supp-exact-graph-field}
\end{equation}
on finite-dimensional spaces, with bounded nonlinear derivatives through
order three. Suppose for $t\ge0$
$\|e^{-tA_u}\|\le Me^{-\beta t}$,
$\|e^{tA_n}\|\le Me^{\alpha t}$, and choose
$\max\{0,\alpha\}<\nu$ with $3\nu<\beta$. If the global nonlinear Lipschitz constant
$L$ satisfies
\begin{equation}
 \max_{j=1,2,3}ML\left((j\nu-\alpha)^{-1}
                             +(\beta-j\nu)^{-1}\right)<1,
 \label{eq:supp-exact-graph-contraction}
\end{equation}
then the forward weighted trajectories define a $C^3$ invariant graph
$u=h(n)$. On any open set where localization agrees with the original
field, the graph is locally invariant for that original field.
\end{lemma}
\begin{proof}
For fixed $n(0)=n_0$ solve
\begin{align*}
 n(t)&=e^{tA_n}n_0+\int_0^t e^{(t-r)A_n}
                          [b_n+N_n(u(r),n(r))]\,\dd r,\\
 u(t)&=-\int_t^\infty e^{(t-r)A_u}
                          [b_u+N_u(u(r),n(r))]\,\dd r.
\end{align*}
In the norm $\sup_{t\ge0}e^{-\nu t}|(u(t),n(t))|$, the two kernels have
bounds $M/(\nu-\alpha)$ and $M/(\beta-\nu)$. The map is a contraction
and defines a unique weighted trajectory. Set $h(n_0)=u(0)$.
Time translation preserves the equations and the weighted class, so
uniqueness proves graph invariance. Differentiating in $n_0$ gives a linear
integral equation with the same contractive principal term. Second and
third derivatives have sources composed of products of lower derivatives;
these belong to the weights $2\nu$ and $3\nu$. The corresponding three
inequalities in \eqref{eq:supp-exact-graph-contraction} solve these equations
and the strict spectral margins justify their continuous dependence.
This proves $C^3$ regularity. Restricting a graph trajectory to an interval
on which the localization is inactive yields the original equation.
\end{proof}

Affine forcing can displace this graph from the coordinate origin.
Consequently neither passage through a selected reference point nor equality
of its tangent with the frozen spectral subspace is a consequence of this
lemma alone. Also, for a nonstationary invariant manifold its tangent need
not be invariant under the Jacobian at each point. These distinctions matter
when extracting the certificate below from an actual action chart.

\begin{assumption}[Two-scale invariant-graph and normal-factor certificate]
\label{ass:supp-exact-normal}
On the original normalized $323$-dimensional constraint leaf there is a
nonempty local rapid-selected graph $\mathcal S_a$ of dimension $312$.
Inside it there is a nonempty local graph $\mathcal D_a$ of dimension $311$.
Both are locally invariant under the original field. They may be obtained
from \Cref{lem:supp-exact-graph-criterion} in successive rapid and slow
coordinates, provided its quantitative extension and spectral-gap
hypotheses hold on the respective charts. The following additional
physical comparisons hold throughout those charts.
\begin{enumerate}[label=\textup{(N\arabic*)},leftmargin=3em]
\item In rapid time $s=t/a^3$, the eleven-component normal
$v=u-h_r(n)$ satisfies
$v'=N_a(u,n)v$, with $\|N_a-J_u\|$ sufficiently small uniformly in the
chart, where $J_u$ has spectrum in the open right half-plane. Its physical
normal displacement is uniformly comparable to $a|v|$.
\item Within $\mathcal S_a$, use slow time $\tau=t/a$. The scalar normal
$\omega=w-h_s(n)$ to $\mathcal D_a$ satisfies
\begin{equation}
 \omega_\tau=[\lambda_+(a)+r_a(w,n)]\omega,\quad
 \lambda_+(a)=\gamma+O(a),\quad |r_a|\le\gamma/4,\quad\gamma>0.
 \label{eq:supp-exact-slow-normal-factor}
\end{equation}
Its physical normal displacement is uniformly comparable to $a|\omega|$.
The certified normal radii are uniformly bounded above; no uniform lower
bound for the slow chart radius is assumed.
\end{enumerate}
\end{assumption}

The normal factorizations are exact consequences of graph invariance and
the corresponding derivative estimates, when those estimates are verified.
Indeed, subtract the graph equation from
$u'-Dh_r(n)n'$ at $u=h_r(n)+v$ and integrate the derivative in $u$
along the segment. The constant term vanishes, leaving a matrix times $v$.
For the slow scalar graph the correction contains
$D_wN_+-Dh_s(n)D_wN_0$ integrated along that segment. Thus its bound
requires control of $Dh_s$ as well as of the nonlinear field derivatives;
a simple unstable eigenvalue by itself is not sufficient.

\begin{proposition}[Rapid Lyapunov growth and slow scalar growth]
\label{prop:supp-exact-normal-growth}
Under \Cref{ass:supp-exact-normal}, there are $c,C,\mu_r,\mu_s>0$
such that, until exit from the corresponding full certified chart,
\begin{equation}
 c e^{\mu_rt/a^3}|v(0)|\le |v(t)|\le C e^{Ct/a^3}|v(0)|,
 \qquad
 e^{\mu_st/a}|\omega(0)|\le |\omega(t)|\le e^{2\gamma t/a}|\omega(0)|,
 \label{eq:supp-exact-normal-growth}
\end{equation}
where the scalar estimate is for histories in $\mathcal S_a$ and one can
take $\mu_s=\gamma/2$ after reducing $a_0$.
\end{proposition}
\begin{proof}
The matrix
$W=\int_0^\infty e^{-rJ_u^*}e^{-rJ_u}\,\dd r$ is positive definite and
satisfies $J_u^*W+WJ_u=I$. For $N_a$ sufficiently close to $J_u$,
\[
 \frac{\dd}{\dd s}(v^*Wv)
       =v^*(N_a^*W+WN_a)v\ge\tfrac12|v|^2
       \ge (2\|W\|)^{-1}v^*Wv.
\]
Integrate and use the extremal eigenvalues of $W$ to obtain the lower
rapid bound. A uniform upper bound on $N_a$ gives the upper estimate by
Gronwall. For the slow scalar factor,
$\gamma/2\le\lambda_+(a)+r_a\le2\gamma$ at small $a$; its exact integral
solution gives the other two inequalities. Converting each time variable
separately yields \eqref{eq:supp-exact-normal-growth}.
\end{proof}

\begin{theorem}[Two necessary normal precision scales for local residence]
\label{thm:main-exact-residence}
On the separate invariant normal-factor chart of \Cref{ass:supp-exact-normal}, suppose the same original finite action satisfies the stated graph and factorization certificate.  This chart is not identified with the rapid spectral packet of \Cref{thm:supp-certified-rapid-spectrum}.  There are nested
locally invariant constrained graphs
\[
 \mathcal D_a\subset\mathcal S_a\subset\mathcal P_a,
 \qquad \dim\mathcal D_a=311,\quad
 \dim\mathcal S_a=312,\quad\dim\mathcal P_a=323.
\]
For a prescribed physical time $T>0$, local residence in the certified
charts requires the rapid physical normal displacement to satisfy
\begin{equation}
 \delta_r(0)\le Ca\exp(-\mu_rT/a^3).
 \label{eq:main-exact-rapid-width}
\end{equation}
Inside $\mathcal S_a$, the additional slow physical normal must satisfy
\begin{equation}
 \delta_s(0)\le Ca\exp(-\mu_sT/a).
 \label{eq:main-exact-slow-width}
\end{equation}
Consequently a tube whose every exact history remains in these charts
through $T$ cannot have algebraic thickness in either indicated normal
sector. The constants are fixed-cutoff constants. These are residence
obstructions, not statements of curvature blowup or nonexistence after
chart exit.
\end{theorem}

\begin{proof}[Proof of \Cref{thm:main-exact-residence}]
Residence through $T$ bounds the terminal rapid normal by its fixed chart
radius. The lower growth bound gives
$|v(0)|\le C\exp(-\mu_rT/a^3)$. Multiply by the physical normal scale
$a$ to obtain \eqref{eq:main-exact-rapid-width}. Within $\mathcal S_a$ the
same argument for $\omega$ gives \eqref{eq:main-exact-slow-width}.
A smaller available slow radius can only decrease that bound. Apply the
rapid inequality to the rapid radius of any putative resident tube, and
apply the slow inequality to its intersection with $\mathcal S_a$.
Since $ae^{-cT/a^k}=o(a^m)$ for every fixed $m$ and $k>0$, neither radius
can be bounded below by a fixed algebraic power of $a$.
The dimension counts subtract eleven and then one normal coordinate from
$323$; the eleven rapidly contracting directions are retained.
\end{proof}

For example, if $|v(0)|\ge c a^m$, the residence time is at most
$(a^3/\mu_r)[m\log(1/a)+C_m]$, as long as the initial point lies in the
chart and the history exists until its exit. The analogous slow estimate is
$(a/\mu_s)[m\log(1/a)+C_m]$. In an anchored control with
$|v(0)|\asymp a^3|e|$ the rapid expression instead contains $m+3$ when
$|e|\ge ca^m$. A trajectory may leave the full chart through a tangential
boundary sooner; none of these estimates asserts growth after that exit.
Exact zero normal coordinates remain zero while local invariance applies,
so the theorem does not exclude a longer trajectory on the selected graph.


\section{Slow exact-action compatibility and obstruction theory}
\label{supp:exact-action-slow}

This appendix proves the slow compatibility and obstruction assertions of
\Cref{thm:main-exact-action-summary}. It uses the amplitude-graded
original-action chart and source projections of
Appendix~\ref{supp:exact-action-audit}, at fixed spatial cutoff.
The first step identifies necessary bounded-rate limits, the second
measures the initial acceleration defect, and the final steps reduce the
quadratic slope condition. Necessary algebraic compatibility is not
identified with nonlinear positive-time propagation.
\subsection{State-dependent principal terms on the slow physical scale}
\label{subsec:supp-exact-slow}

Set $t=a\tau$, $s=\tau/a^2$, $w=(x,az)$ and $\xi=(x,z)$.
The unbalanced coordinate is then $v=a\xi$. On bounded $\xi$-sets the
physical chart gives
\begin{equation}
 X-X_a^0=O(a^2),\qquad
 \lambda_A-\lambda_{a,A}^0=O(a^2),\qquad
 \lambda_W-\lambda_{a,W}^0=O(a^2).
 \label{eq:supp-exact-slow-physical-tube}
\end{equation}

In these coordinates the exact slow system is
\begin{align}
 x_\tau&=b_s+\Sigma_0(\xi)+a\Sigma_1(a,\xi),\notag\\
 z_\tau&=a^{-2}\Phi_0(\xi)+a^{-1}\Phi_1(\xi)
             +b_w+\Phi_2(a,\xi),
 \label{eq:main-exact-slow-system}
\end{align}
\begin{lemma}[Exact anisotropic slow expansion]
\label{lem:supp-exact-slow-expansion}
The graded analytic remainders in
\eqref{eq:supp-exact-graded-remainders} give exactly
\eqref{eq:main-exact-slow-system}, where the remainders $\Sigma_1,\Phi_2$
are uniformly bounded on bounded $\xi$-sets and
\begin{equation}
 \begin{split}
 \Phi_0(x,z)&=C_{\rm red}x+Bz,\\
 \Phi_1(\xi)&=D_aD_vR_f(0,0)[\xi]
                   +\tfrac12D_v^2R_f(0,0)[\xi,\xi],\\
 \Sigma_0(\xi)&=\tfrac12D_a^2D_vR_s(0,0)[\xi]
                   +\tfrac12D_aD_v^2R_s(0,0)[\xi,\xi].
 \end{split}
 \label{eq:supp-exact-slow-coefficients}
\end{equation}
In particular $\Sigma_0(0)=\Phi_1(0)=0$ and both maps have degree at most
two.
\end{lemma}
\begin{proof}
Since $v'=a^3\xi_\tau$, substitution in the unbalanced equations gives
$x_\tau=b_s+a^{-3}R_s(a,a\xi)$ and
$z_\tau=a^{-2}(C_{\rm red}x+Bz)+b_w+a^{-3}R_f(a,a\xi)$.
The analytic order bounds force
$R_s(a,a\xi)=a^3\Sigma_0(\xi)+a^4\Sigma_1(a,\xi)$ and
$R_f(a,a\xi)=a^2\Phi_1(\xi)+a^3\Phi_2(a,\xi)$.
Taylor expansion shows that exactly the displayed mixed derivatives survive
in the leading coefficients; the constant remainders are $O(a^4)$.
Analyticity on the original chart gives uniform boundedness on any fixed
$\xi$-ball for sufficiently small $a$. This proves the expansion before any
trajectory estimate is used.
\end{proof}

The rank assertions needed for the next step are additional finite tests
on the same coefficient matrices. Let $\mathcal H$ be the three-dimensional
constant-shift space and $\mathcal Z_g$ a two-dimensional weak-zero space.
Choose a direct decomposition
\begin{equation}
 \mathbb R^{29}=\Ran B\oplus\mathcal Z_g\oplus\mathcal H,
 \qquad \rank B=24,
 \label{eq:supp-exact-weak-decomposition}
\end{equation}
with $p_g,p_H$ its two indicated projections, annihilating $\Ran B$.
Assume $p_HC_{\rm red}=0$ and $p_gC_{\rm red}$ is onto
$\mathcal Z_g$. In applications the decomposition and covector
identifications must use the retained positive source pairing, not an
unjustified orthogonal projection for an indefinite matrix.

\begin{proposition}[Rank and cokernel of the leading slow constraint]
\label{prop:supp-exact-slow-rank}
Under \eqref{eq:supp-exact-weak-decomposition} and the stated coupling
tests,
\begin{equation}
 \rank\Phi_0=26,\qquad \dim\ker\Phi_0=297,\qquad
 \operatorname{coker}\Phi_0\cong\mathcal H.
 \label{eq:supp-exact-critical-dimensions}
\end{equation}
\end{proposition}
\begin{proof}
The image of $B$ contributes $24$ directions. Modulo that image,
$C_{\rm red}$ contributes exactly the two directions in $\mathcal Z_g$
and none in $\mathcal H$. Therefore $\rank[C_{\rm red}\ B]=26$.
Rank--nullity in the $323$-dimensional domain gives $297$, and the quotient
of the $29$-dimensional codomain is the remaining three-dimensional space.
\end{proof}

Define $\mathfrak m(\xi)=p_H\Phi_1(\xi)$ and
$\mathfrak g(\xi)=p_gC_{\rm red}(b_s+\Sigma_0(\xi))$.
Their ordered pair is
\begin{equation}
 \mathfrak C_{\rm sl}(\xi)=
 \left(p_gC_{\rm red}[b_s+\Sigma_0(\xi)],\ p_H\Phi_1(\xi)\right)
 \in\mathbb R^2\oplus\mathbb R^3.
 \label{eq:main-exact-five-map}
\end{equation} This is a finite
polynomial map extracted from the unchanged action, not a replacement
constraint added to its finite dynamics.

\begin{theorem}[Necessary slow compatibility and the bounded-rate readout]
\label{thm:supp-exact-five-compatibility}
Let $a_j\downarrow0$ and suppose original exact histories exist on
$0\le t\le Ta_j$ in the indicated chart, with
\begin{equation}
 \sup_j\sup_{0\le\tau\le T}
       (|\xi_{a_j}(\tau)|+|\partial_\tau\xi_{a_j}(\tau)|)<\infty.
 \label{eq:supp-exact-bounded-slow-rate}
\end{equation}
Every uniformly convergent subsequence has a Lipschitz limit $\xi_0$ with
\begin{equation}
 \Phi_0(\xi_0)=0,\qquad
 (x_0)_\tau=b_s+\Sigma_0(\xi_0),\qquad
 \mathfrak C_{\rm sl}(\xi_0)=0.
 \label{eq:supp-exact-five-identities}
\end{equation}
The differential identity holds almost everywhere. The physical histories
satisfy $X_t,\lambda_t,X_{tt}=O(a_j)$ and the corresponding metric and link
curvature is $O(a_j)$ on their full intervals, subject to the readout
conditions in \Cref{ass:supp-graded-exact-chart}.
\end{theorem}
\begin{proof}
The bound gives equicontinuity, hence uniform subsequential compactness.
Multiply the weak equation by $a^2$:
\begin{equation}
 \Phi_0(\xi_a)+a\Phi_1(\xi_a)
    =a^2[z_{a,\tau}-b_w-\Phi_2(a,\xi_a)].
 \label{eq:supp-exact-slow-principal-balance}
\end{equation}
Its right-hand side is $O(a^2)$, so $\Phi_0(\xi_0)=0$. The regular $x$
equation passes in integral form to the second identity. Applying $p_H$
to \eqref{eq:supp-exact-slow-principal-balance}, using $p_H\Phi_0=0$ and
dividing by $a$, gives $\mathfrak m(\xi_a)=O(a)$ and hence
$\mathfrak m(\xi_0)=0$. Differentiate
$C_{\rm red}x_0+Bz_0=0$ almost everywhere and apply $p_g$. The $B$ term
vanishes, leaving $\mathfrak g(\xi_0)=0$ after substitution of the regular
$x$ drift. Continuity extends the polynomial identities to all times.

In the physical chart, $p_W=az$ and $t=a\tau$.
The bounded slow rates and \eqref{eq:supp-exact-normal-derivatives} imply
$X_t,\lambda_t=O(a)$. The exact canonical equation with bounded
derivatives then gives $X_{tt}=O(a)$. Apply the original readout estimate
as in \Cref{thm:main-exact-boundary}. These bounds rely on
\eqref{eq:supp-exact-bounded-slow-rate}; the polynomial equalities alone
do not provide them.
\end{proof}

\begin{proposition}[First correction and initial slow tangency]
\label{prop:supp-exact-slow-correction}
For $\xi_0\in\ker\Phi_0$, the equation
$\Phi_0\xi_1+\Phi_1(\xi_0)=0$ is solvable exactly when
$\mathfrak m(\xi_0)=0$. If $R_0$ is a right inverse of $\Phi_0$ on its
range, one solution is $\xi_1=-R_0\Phi_1(\xi_0)$, modulo
$\ker\Phi_0$. Moreover, if a limiting path in
\Cref{thm:supp-exact-five-compatibility} is $C^1$ at zero with
$\xi_0(0)=0$ and $\dot\xi_0(0)=b=(b_s,b_w)$, then
\begin{equation}
 D\mathfrak C_{\rm sl}(0)[b]
       =\left(p_gC_{\rm red}D\Sigma_0(0)[b],\
                    p_HD\Phi_1(0)[b]\right)=0.
 \label{eq:supp-exact-slow-first-tangency}
\end{equation}
\end{proposition}
\begin{proof}
The range of $\Phi_0$ is the kernel of $p_H$ by
\Cref{prop:supp-exact-slow-rank}. Thus the first-correction equation is
solvable if and only if $p_H\Phi_1(\xi_0)=0$, and a right inverse gives
the stated solution. Taylor expansion then gives
$\Phi_0(\xi_0+a\xi_1)+a\Phi_1(\xi_0+a\xi_1)=O(a^2)$.
For the final assertion differentiate
$\mathfrak C_{\rm sl}(\xi_0(\tau))=0$ at zero. The actual derivative of
the limiting path is part of the hypothesis; convergence of derivatives at
one initial point alone is not substituted for it.
\end{proof}

The transport identity gives $b\in\ker\Phi_0$ and
$\mathfrak C_{\rm sl}(0)=0$. It does not evaluate
$D\mathfrak C_{\rm sl}(0)[b]$. Along the particular affine ansatz
$\xi=\tau b$, the three mean conditions reduce to
$D\mathfrak m(0)[b]=0$ and $D^2\mathfrak m(0)[b,b]=0$, since
$\mathfrak m$ is quadratic and vanishes at zero. These six scalar
contractions test that ansatz; they are not imposed on every curved slow
path. A vanishing first tangency vector likewise does not prove invariance
of the nonlinear five-equation zero set.


\subsection{Initial acceleration defect and first Lyapunov--Schmidt reduction}
\label{subsec:supp-exact-initial-slow-gate}

For the initial-cut audit it is useful to retain the regular and singular
source terms separately.  Rewrite \eqref{eq:main-exact-slow-system} as
\begin{align}
 x_\tau&=f(\xi)+a s(a,\xi,u),\label{eq:supp-exact-source-x}\\
 z_\tau&=a^{-2}\Phi_0(\xi)+a^{-1}g(\xi)+h(a,\xi,u),
 \qquad \xi=(x,z),
 \label{eq:supp-exact-source-z}
\end{align}
where
\begin{equation}
 f=b_s+\Sigma_0,
 \qquad g=\Phi_1,
 \qquad h=b_w+\widehat\Phi_2,
 \label{eq:supp-exact-source-identification}
\end{equation}
and $s=\widehat\Sigma_1$ denotes the corresponding regular order-$a$
source.  The companion coordinate $u$ contains all remaining finite
initializer/source variables.  Define
\begin{equation}
 \mathfrak c_g(\xi)=p_gC_{\rm red}f(\xi),
 \qquad
 \mathfrak c_H(\xi)=p_Hg(\xi).
 \label{eq:supp-exact-c-functions}
\end{equation}
The leading bounded-rate differential--algebraic system is therefore
\begin{equation}
 \Phi_0(\xi)=0,\qquad x_\tau=f(\xi),
 \qquad \mathfrak c_g(\xi)=0,\qquad \mathfrak c_H(\xi)=0.
 \label{eq:supp-exact-leading-dae}
\end{equation}
This is a necessary slow limit of exact histories; it is not imposed as an
additional finite equation.

Write $s_0=s(0,0,u_0)$ and $h_0=h(0,0,u_0)$.  Suppose an exact family,
defined on an amplitude-dependent interval if necessary, satisfies
\begin{equation}
 \xi_a(0)=a\kappa+O(a^2),\qquad
 \partial_\tau\xi_a(0)=b+o(1),\qquad b=(b_s,b_w),
 \qquad u_a(0)\to u_0,
 \qquad p_gC_{\rm red}b_s=0.
 \label{eq:supp-exact-initial-jets-corrected}
\end{equation}
No uniform second-time-derivative bound is assumed.  The regular equation
gives
\begin{equation}
 \partial_\tau x_a(0)
 =b_s+a\{Df(0)\kappa+s_0\}+o(a).
 \label{eq:supp-exact-initial-xprime}
\end{equation}
Define
\begin{align}
 \mathcal L_0\kappa
 &=\binom{p_gC_{\rm red}Df(0)\kappa}{p_HDg(0)\kappa},
 \label{eq:supp-L0-definition}\\
 f_0
 &=\binom{p_g\{C_{\rm red}s_0+Dg(0)b\}}
               {p_H(h_0-b_w)}.
 \label{eq:supp-f0-definition}
\end{align}

\begin{theorem}[Initial acceleration-defect identity]
\label{thm:supp-exact-initial-gate}
Under \eqref{eq:supp-exact-initial-jets-corrected},
$\kappa\in\mathcal K_0$ and
\begin{equation}
 (\mathcal L_0\kappa+f_0)_H=0.
 \label{eq:supp-exact-initial-H-automatic}
\end{equation}
If the exact histories and the composed source $h(a,\xi_a,u_a)$ are twice,
respectively once, differentiable at the initial cut, then the following
limit exists and
\begin{equation}
 (\mathcal L_0\kappa+f_0)_g
 =\lim_{a\downarrow0}a\,p_g
 \left\{z_{a,\tau\tau}(0)-
 \frac{\dd}{\dd\tau}h(a,\xi_a,u_a)\bigg|_{\tau=0}\right\}.
 \label{eq:supp-exact-initial-acceleration-defect}
\end{equation}
Hence the complete five-row equation
$\mathcal L_0\kappa=-f_0$ holds if and only if the two-row acceleration
defect in \eqref{eq:supp-exact-initial-acceleration-defect} vanishes.
\end{theorem}

\begin{proof}
Multiplying \eqref{eq:supp-exact-source-z} by $a$ at $\tau=0$ gives
\[
 a\partial_\tau z_a(0)
 =\Phi_0\kappa+O(a)+g(\xi_a(0))+a h(a,\xi_a(0),u_a(0)).
\]
All terms except $\Phi_0\kappa$ tend to zero, so
$\Phi_0\kappa=0$.  Applying $p_H$ to
\eqref{eq:supp-exact-source-z} and using $p_H\Phi_0=0$ gives in the
initial limit
\[
 p_Hb_w=p_HDg(0)\kappa+p_Hh_0,
\]
which is \eqref{eq:supp-exact-initial-H-automatic}.

At positive amplitude the exact identity
\begin{equation}
 \Phi_0\xi_a+a g(\xi_a)
 =a^2\{\partial_\tau z_a-h(a,\xi_a,u_a)\}
 \label{eq:supp-exact-initial-pole-identity}
\end{equation}
holds.  Differentiate it at $\tau=0$ and apply $p_g$:
\begin{equation}
 p_gC_{\rm red}\partial_\tau x_a(0)
 +a p_gDg(\xi_a(0))\partial_\tau\xi_a(0)
 =a^2p_g\left\{z_{a,\tau\tau}(0)-
        \frac{\dd h}{\dd\tau}\bigg|_0\right\}.
 \label{eq:supp-exact-initial-differentiated-pole}
\end{equation}
Subtract $p_gC_{\rm red}b_s=0$, divide by $a$, and use
\eqref{eq:supp-exact-initial-xprime} and
\eqref{eq:supp-exact-initial-jets-corrected}.  The left-hand side converges
to
\[
 p_g\{C_{\rm red}Df(0)\kappa+C_{\rm red}s_0+Dg(0)b\},
\]
which proves \eqref{eq:supp-exact-initial-acceleration-defect}.
\end{proof}

\begin{corollary}[Sufficient acceleration certificate]
\label{cor:supp-exact-initial-repair}
If the projected difference in
\eqref{eq:supp-exact-initial-acceleration-defect} is $o(a^{-1})$, then
$\mathcal L_0\kappa=-f_0$.  In particular, uniform bounds for
$z_{a,\tau\tau}(0)$ and
$\frac{\dd}{\dd\tau}h(a,\xi_a,u_a)|_0$ are sufficient.
\end{corollary}

The qualification is not removable by analyticity of the initial record
alone.

\begin{proposition}[Analytic preparation does not imply the two-row cancellation]
\label{prop:supp-exact-initial-counterexample}
The autonomous polynomial Lagrangian
\begin{equation}
 L(\alpha,z,\lambda;\alpha',z',\lambda')
 =\tfrac12\alpha^2(z')^2-\tfrac12(z-\alpha)^2+\lambda\alpha'
 \label{eq:supp-exact-counterexample-L}
\end{equation}
has exact initial records analytic in $a>0$, with $\alpha(0)=a$ and bounded
analytic initial velocities, for which the analogue of the upper homological
row is nonzero.  The exact solutions nevertheless have uniformly bounded
first derivatives on every fixed finite interval.
\end{proposition}

\begin{proof}
The Euler equations are
\[
 \alpha'=0,\qquad (\alpha^2z')'=\alpha-z,\qquad
 \lambda'=\alpha(z')^2+(z-\alpha).
\]
Set $\alpha=a$, $x=a^2z'$ and
$x(0)=z(0)=\lambda(0)=0$.  Then
\begin{equation}
 x'=-z+a,\qquad z'=a^{-2}x,
 \qquad x=a^2\sin(\tau/a),\quad
 z=a\{1-\cos(\tau/a)\}.
 \label{eq:supp-exact-counterexample-solution}
\end{equation}
The multiplier equation is solved by integration and has derivative bounded
by $2a$.  Thus all Euler equations hold and the first derivatives remain
bounded, but $az''(0)=1$.  The right-hand side of
\eqref{eq:supp-exact-initial-acceleration-defect} is therefore nonzero in
the corresponding one-row reduction.  The example uses one unchanged
polynomial action for every $a$; amplitude labels the initial record rather
than changing the action.
\end{proof}

The action itself supplies the required acceleration test.  On the normalized
multiplier slice let
\begin{equation}
 \mathcal B(Q)+\mathcal M(Q)r(Q)=0,
 \qquad Q=(X,\lambda),
 \qquad V=(F^{\rm can},r).
 \label{eq:supp-exact-original-consistency}
\end{equation}

\begin{proposition}[Original-action multiplier acceleration]
\label{prop:supp-exact-action-acceleration}
Where $\mathcal M$ is invertible on the normalized slice,
\begin{equation}
 \mathcal M\dot r
 =-D_Q\mathcal B[V]-D_Q\mathcal M[V]r.
 \label{eq:supp-exact-action-acceleration}
\end{equation}
Thus the two-row quantity in
\eqref{eq:supp-exact-initial-acceleration-defect} is computable from a
directional derivative of the unchanged finite action together with the
actual amplitude-graded decoder; it is not obtained by differentiating an
independently prescribed target rate.
\end{proposition}
\begin{proof}
Differentiate \eqref{eq:supp-exact-original-consistency} along the exact
physical solution and solve for $\dot r$.
\end{proof}

\subsection{Finite-jet evaluation and preparation dependence of the acceleration defect}
\label{subsec:supp-finite-acceleration}

All coefficients in this subsection are evaluated at the initial cut, at
fixed spatial regulator and in the fixed mean-lapse normalization.  The
symbol $[a^j]$ denotes a Taylor coefficient, including its factorial
normalization.  It does not denote a finite difference of sampled values.
Write $r_a=\lambda_t(0)$ and set
\begin{equation}
 \mathcal S_a=D_Q\mathcal B(Q_a)[V_a]
       +D_Q\mathcal M(Q_a)[V_a]r_a,
 \qquad V_a=(F^{\rm can}(Q_a),r_a).
 \label{eq:supp-finite-acc-source}
\end{equation}
This expression uses the actual initialized rate and the derivative of the
original action.  A prescribed rate at one cut is not differentiated as
though it were an identity along the history.

\begin{proposition}[Exact elimination of the active acceleration]
\label{prop:supp-finite-acc-schur}
Suppose the normalized multiplier mass on the analytic initial family has
\begin{equation}
 \mathcal M_a=\begin{pmatrix}
 a^2 A(a)&a^3 L(a)\\a^3L(a)^{\mathsf T}&a^4 C(a)
 \end{pmatrix},\qquad
 D(a)=C(a)-L(a)^{\mathsf T}A(a)^{-1}L(a),
 \label{eq:supp-finite-acc-mass}
\end{equation}
where $A(0)$ and $D(0)$ are invertible.  For
$\mathcal S_a=(\mathcal S_A,\mathcal S_W)$ put
\begin{equation}
 \chi(a)=\mathcal S_W-aL^{\mathsf T}A^{-1}\mathcal S_A,
 \qquad v(a)=D(a)^{-1}\chi(a).
 \label{eq:supp-finite-acc-effective}
\end{equation}
Then the physical accelerations satisfy exactly
\begin{equation}
 \lambda_{tt,W}=-a^{-4}v(a),\qquad
 \lambda_{tt,A}=-a^{-2}A^{-1}\mathcal S_A
                   +a^{-3}A^{-1}Lv(a).
 \label{eq:supp-finite-acc-elimination}
\end{equation}
\end{proposition}
\begin{proof}
The acceleration equation is $\mathcal M_a\lambda_{tt}=-\mathcal S_a$.
Its active row gives
$\lambda_{tt,A}=-a^{-2}A^{-1}\mathcal S_A-aA^{-1}L\lambda_{tt,W}$.
Substitution into the weak row gives $a^4D\lambda_{tt,W}=-\chi$;
back-substitution proves the other formula.
\end{proof}

In the balanced chart $\lambda_W=\lambda_W^0(a)+a^2z$ and $t=a\tau$,
with chart centre independent of time, one has
$z_{\tau\tau}=\lambda_{tt,W}$.  This decoder identity is part of the
coefficient test; a different time-dependent chart must retain its chain-rule
terms.  Define
\begin{equation}
 \theta(a)=a p_g\frac{\dd}{\dd\tau}h(a,\xi_a,u_a)\bigg|_0.
 \label{eq:supp-finite-acc-theta}
\end{equation}

\begin{theorem}[Coefficient-complete acceleration certificate]
\label{thm:supp-finite-acc-coefficients}
Assume the analytic mass and source hypotheses of
\Cref{prop:supp-finite-acc-schur}.  Write
$D(a)=\sum_{j\ge0}a^jD_j$, $\chi(a)=\sum_{j\ge0}a^j\chi_j$.
The four coefficients $v_0,\ldots,v_3$ are obtained using only the same
invertible matrix $D_0$:
\begin{equation}
 D_0v_j=\chi_j-\sum_{i=1}^jD_i v_{j-i},\qquad 0\le j\le3.
 \label{eq:supp-finite-acc-recursion}
\end{equation}
If $\theta(a)=\sum_{k=-3}^0 a^k\theta_k+O(a)$, the two-row acceleration
defect has a finite limit precisely when
\begin{equation}
 p_gv_j+\theta_{j-3}=0\qquad(0\le j<3).
 \label{eq:supp-finite-acc-poles}
\end{equation}
In that case its value is
\begin{equation}
 \Delta_g=-p_gv_3-\theta_0.
 \label{eq:supp-finite-acc-verdict}
\end{equation}
For a family satisfying \Cref{thm:supp-exact-initial-gate}, this value
also equals the entirely initial-data expression
\begin{equation}
 \Delta_g=p_g\{C_{\rm red}Df(0)\kappa+
                  C_{\rm red}s_0+Dg(0)b\}.
 \label{eq:supp-finite-acc-initial-verdict}
\end{equation}
Thus the action-side coefficient calculation and the initialized slow-source
calculation must give the same two-vector.  When
$\chi_0=\chi_1=\chi_2=0$ and $\theta(a)=o(1)$, this simplifies to
$\Delta_g=-p_gD_0^{-1}\chi_3$.  Neither simplification is inferred merely
from an $O(a)$ initial rate.
\end{theorem}
\begin{proof}
Compare coefficients in $Dv=\chi$.  The exact projected expression is
$-a^{-3}p_gv(a)-\theta(a)$, by
\eqref{eq:supp-finite-acc-elimination}.  Its three possible negative powers
vanish exactly under \eqref{eq:supp-finite-acc-poles}; its constant
coefficient is \eqref{eq:supp-finite-acc-verdict}.  The initial-data formula
is the acceleration-defect identity.  If $\chi_j=0$ for $j<3$, the
recursion gives $v_j=0$ for $j<3$ and $v_3=D_0^{-1}\chi_3$.
\end{proof}

If $h$ is uniformly $C^1$ and both $\partial_\tau\xi_a(0)$ and
$\partial_\tau u_a(0)$ are bounded, the chain rule gives $\theta(a)=O(a)$.
If those companion-velocity bounds are not available, the corresponding
$\theta_k$ must be included.  Lower coefficients can vanish after projection
without $\chi$ itself being divisible by $a^3$.  The recurrence, rather than
a presumed leading order, covers this case.  Inverses of $A(a)$ needed for
$D_j$ and $\chi_j$ are obtained similarly: for $U(a)=A(a)^{-1}$,
$A_0U_0=I$ and
$A_0U_j=-\sum_{i=1}^jA_iU_{j-i}$.  Thus only finitely many action jets and
solves with $A_0,D_0$ enter the certificate.

\begin{proposition}[Dependence on preparation jets]
\label{prop:supp-finite-acc-invariance}
Consider two exact initialized families in the same slow-source chart.
If they have the same $\kappa,b,s_0,Df(0),Dg(0),C_{\rm red},p_g$, then
they have the same $\Delta_g$.  For a differentiable family of initial
preparations with the last four maps fixed,
\begin{equation}
 D_d\Delta_g[\dot d]=p_g\{C_{\rm red}Df(0)D_d\kappa[\dot d]
 +C_{\rm red}D_ds_0[\dot d]+Dg(0)D_db[\dot d]\}.
 \label{eq:supp-finite-acc-total-response}
\end{equation}
Consequently a higher target jet preserving this full determining tuple
cannot remove a nonzero defect.  A nonzero derivative of the acceleration
source alone is not a certificate for the derivative of the complete defect.
\end{proposition}
\begin{proof}
Apply \eqref{eq:supp-finite-acc-initial-verdict} to the two families and
then differentiate that formula.  If the source maps themselves vary, their
product-rule terms must also be retained.  The formula concerns exact
initialized records, so implicit changes of those records are included.
\end{proof}

For example, the same polynomial Lagrangian in
\Cref{prop:supp-exact-initial-counterexample} admits the exact solutions
\[
 x=a^2\sin(\tau/a)+a^4d\cos(\tau/a),\qquad
 z=a\{1-\cos(\tau/a)\}+a^3d\sin(\tau/a).
\]
Their initial weak rate changes by $a^2d$, while $\kappa$, the limiting
first rate and the source coefficients are unchanged.  Nevertheless
$az''(0)=1$ for every fixed $d$.  This example tests the logical distinction;
it does not assign a coefficient to the gravitational event.

\subsection{Defect-aware Lyapunov--Schmidt reduction}
It remains useful to retain the algebraic reduction for arbitrary forcing.
Let $J=\{0,\ldots,25\}\setminus\{2,24\}$ and $F=\{2,24\}$.  The certified
principal block $K_{JJ}$ is invertible, so
\begin{equation}
 Z_g=\begin{pmatrix}-K_{JJ}^{-1}K_{JF}\\ I_2\end{pmatrix}
 \label{eq:supp-exact-Zg-basis}
\end{equation}
spans $\mathcal Z_g$.  Let $K_{Ag}$ be the active/gradient first-Poisson
block and $A_{78}\succ0$ the certified active mass, and define
\begin{equation}
 G_A=K_{Ag}Z_g,
 \qquad
 Q_g=G_A^{\mathsf T}A_{78}^{-1}G_A.
 \label{eq:supp-exact-Qg}
\end{equation}
The verified $LDL^{\mathsf T}$ pivots obey
\begin{equation}
 d_1>1.2143\times10^7,
 \qquad
 d_2>8.00\times10^5,
 \label{eq:supp-exact-Qg-pivots}
\end{equation}
so $Q_g$ is positive definite.  Therefore the induced upper two-row vertical
map $\mathcal A_g:\mathcal Z_g\to\mathcal Z_g$ is an isomorphism; for
$z\in\mathcal Z_g$,
\begin{equation}
 \langle z,\mathcal A_gz\rangle
 =-\langle K_{Ag}z,A_{78}^{-1}K_{Ag}z\rangle.
 \label{eq:supp-exact-Ag-form}
\end{equation}
Choose a complement $\mathcal X_h$ of the two vertical directions in
$\mathcal K_0$, with inclusion $E_x$, and write the complete block map as
\begin{equation}
 \mathcal L_0(J_gz+E_xx)=
 \binom{\mathcal A_gz+\mathcal B_xx}
       {\mathcal C_gz+\mathcal D_xx},\qquad
 \mathcal H_x=\mathcal D_x-\mathcal C_g\mathcal A_g^{-1}\mathcal B_x.
 \label{eq:supp-exact-complete-LS-blocks}
\end{equation}
No assumption that the upper block vanishes on an arbitrary complement is
made.

\begin{theorem}[Lyapunov--Schmidt reduction of the initial gate]
\label{thm:supp-exact-three-row}
For arbitrary $f=(f_g,f_H)\in\mathcal Z_g\oplus\mathcal H$,
$\mathcal L_0\kappa=-f$ is solvable if and only if
\begin{equation}
 f_H-\mathcal C_g\mathcal A_g^{-1}f_g\in\Ran\mathcal H_x,
 \label{eq:supp-exact-hard-forcing}
\end{equation}
and
\begin{equation}
 \rank\mathcal L_0=2+\rank\mathcal H_x.
 \label{eq:supp-exact-slow-rank-split}
\end{equation}
Thus full five-row surjectivity is equivalent to
$\rank\mathcal H_x=3$.  If the actual initialized coefficient is
$\kappa_*=J_gz_*+E_xx_*$ and
$\mathcal L_0\kappa_*+f_0=(\Delta_g,0)$, then
\begin{equation}
 \mathcal H_xx_*+f_{0,H}-\mathcal C_g\mathcal A_g^{-1}f_{0,g}
            =-\mathcal C_g\mathcal A_g^{-1}\Delta_g.
 \label{eq:supp-exact-defect-feedback}
\end{equation}
There exists some coefficient $\widetilde\kappa$ with
$\mathcal L_0\widetilde\kappa=-f_0$ if and only if
$\mathcal C_g\mathcal A_g^{-1}\Delta_g\in\Ran\mathcal H_x$.
This algebraic existence is distinct from the assertion that the given
initializer has zero defect.
\end{theorem}
\begin{proof}
For arbitrary $x$, solve the upper rows by
$z=-\mathcal A_g^{-1}(\mathcal B_xx+f_g)$.  Substitution gives
$\mathcal H_xx+f_H-\mathcal C_g\mathcal A_g^{-1}f_g=0$ and the rank
formula follows by invertible block operations.  For the actual initializer,
the upper equation instead gives
$z_*=-\mathcal A_g^{-1}(\mathcal B_xx_*+f_{0,g})
      +\mathcal A_g^{-1}\Delta_g$.
Its original lower rows vanish, so substitution proves
\eqref{eq:supp-exact-defect-feedback}.  Finally,
$\mathcal L_0(\widetilde\kappa-\kappa_*)=-(\Delta_g,0)$ has a solution
exactly when $\mathcal C_g\mathcal A_g^{-1}\Delta_g$ lies in
$\Ran\mathcal H_x$, by the same block elimination.
\end{proof}

\begin{proof}[Proof of
\Cref{thm:supp-exact-initial-gate,thm:supp-exact-three-row}]
The initial identity and its qualification are
\Cref{thm:supp-exact-initial-gate,cor:supp-exact-initial-repair};
\eqref{eq:supp-exact-action-acceleration} is
\Cref{prop:supp-exact-action-acceleration}.  Positivity of
\eqref{eq:supp-exact-Qg} gives invertibility of $\mathcal A_g$, and
\Cref{thm:supp-exact-three-row} gives the finite block reduction.
\end{proof}

\subsection{Quadratic slow-slope gate}
\label{subsec:supp-exact-slope-gate}

We next eliminate the two generalized-zero equations before differentiating
in slow time.  Let
$J_g:\mathcal Z_g\hookrightarrow\mathbb R^{29}$ and
$J_H:\mathcal H\hookrightarrow\mathbb R^{29}$ be the weak insertions.  When composed with a map on $\mathcal K_0$,
these symbols mean the insertions $(0,J_g\cdot)$ and $(0,J_H\cdot)$.
On the certified prepared event the literal source identities are
\begin{equation}
 \begin{gathered}
 \mathfrak c_g(0)=\mathfrak c_H(0)=0,\qquad Df(0)J_H=0,\\
 D\mathfrak c_g(0)J_H=D\mathfrak c_H(0)J_H=0,\qquad
 \mathcal A_g=D\mathfrak c_g(0)J_g\in GL_2(\mathbb R).
 \end{gathered}
 \label{eq:supp-exact-slope-entrance}
\end{equation}
The harmonic identities in
\eqref{eq:supp-exact-slope-entrance} are proved below from the finite
constant-shift Ward structure; they are pointwise identities at the prepared
event and are not extrapolated to an arbitrary neighbouring rate.

Let
\begin{equation}
 \mathcal X=\ker(p_gC_{\rm red})\subset\mathbb R^{294},
 \qquad \dim\mathcal X=292,
 \label{eq:supp-exact-X-space}
\end{equation}
choose an isomorphism $X:\mathbb R^{292}\to\mathcal X$ and inverse $L_X$,
and let $B^\#$ denote the inverse of $B$ on $\Ran B$, extended by zero on
$\ker B$.  Put
\begin{equation}
 \mathsf E r=(Xr,-B^\#C_{\rm red}Xr),
 \qquad
 \Xi_0(r,z_g,y)=\mathsf E r+J_gz_g+J_Hy.
 \label{eq:supp-exact-linear-critical-coordinates}
\end{equation}

\begin{lemma}[Exact two-row elimination]
\label{lem:supp-exact-two-row-elimination}
The map $\Xi_0$ is a linear isomorphism onto $\mathcal K_0$.  There is a
unique local analytic map $\gamma(r,y)\in\mathcal Z_g$, with
$\gamma(0,0)=0$, such that
\begin{equation}
 \Xi(r,y)=\mathsf E r+J_g\gamma(r,y)+J_Hy,
 \qquad
 \mathfrak c_g(\Xi(r,y))=0.
 \label{eq:supp-exact-Xi-reduced}
\end{equation}
Moreover $\gamma_y(0,0)=0$.
\end{lemma}
\begin{proof}
If $\Phi_0(x,z)=0$, applying $p_g$ gives $x\in\mathcal X$.  Conversely,
for $x\in\mathcal X$, $C_{\rm red}x$ lies in $\Ran B$ because both its
$p_g$ and $p_H$ projections vanish.  Every solution of
$Bz=-C_{\rm red}x$ is $-B^\#C_{\rm red}x$ plus one unique element of
$\ker B=\mathcal Z_g\oplus\mathcal H$, proving the linear statement.
The derivative of $\mathfrak c_g(\Xi_0)$ in $z_g$ at zero is
$\mathcal A_g$, so the analytic implicit-function theorem gives
\eqref{eq:supp-exact-Xi-reduced}.  Differentiating in $y$ and using
\eqref{eq:supp-exact-slope-entrance} gives $\gamma_y(0,0)=0$.
\end{proof}

Define
\begin{equation}
 \mathcal F(r,y)=L_Xf(\Xi(r,y)),
 \qquad
 \mathcal H(r,y)=\mathfrak c_H(\Xi(r,y)).
 \label{eq:supp-exact-reduced-functions}
\end{equation}

\begin{theorem}[Reduced three-constraint leading system]
\label{thm:supp-exact-reduced-leading-system}
A local absolutely continuous path in the coordinate neighbourhood solves
\eqref{eq:supp-exact-leading-dae} if and only if it is
$\xi=\Xi(r,y)$ with
\begin{equation}
 r_\tau=\mathcal F(r,y),\qquad \mathcal H(r,y)=0.
 \label{eq:supp-exact-reduced-dae}
\end{equation}
At zero,
\begin{equation}
 \mathcal H(0,0)=0,
 \qquad \mathcal F_y(0,0)=\mathcal H_y(0,0)=0.
 \label{eq:supp-exact-paired-zero}
\end{equation}
\end{theorem}
\begin{proof}
The coordinate map enforces $\Phi_0=0$ and the implicit graph enforces
$\mathfrak c_g=0$; the remaining algebraic equation is
$\mathfrak c_H=0$.  The regular equation has $x=Xr$, so it is equivalent
to $r_\tau=L_Xf$.  The identities at zero follow from
\eqref{eq:supp-exact-slope-entrance} and
$\gamma_y(0,0)=0$.
\end{proof}

Set
\begin{equation}
 v=\mathcal F(0,0),\qquad
 u_0=\tfrac12\mathcal F_r(0,0)v,
 \qquad
 \rho_0=\mathcal H_r(0,0)v,
 \label{eq:supp-exact-slope-data}
\end{equation}
and define
\begin{equation}
 \mathcal Q(c)=
 \mathcal H_r\mathcal F_rv+\mathcal H_{rr}[v,v]
 +2\mathcal H_{ry}[v,c]+\mathcal H_{yy}[c,c].
 \label{eq:supp-exact-slope-polynomial}
\end{equation}

\begin{proposition}[Slope coefficients from the original compatibility rows]
\label{prop:supp-exact-original-slope-jets}
Use the physical harmonic row convention of
\eqref{eq:supp-exact-reduced-functions}, so that
$\mathcal Q_{\rm can}=\mathcal Q$ in these coordinates.  Define
\begin{equation}
 \widehat E=\mathsf E-J_g\mathcal A_g^{-1}D\mathfrak c_g(0)\mathsf E,
 \quad b_*=\widehat Ev,\quad b(c)=b_*+J_Hc,
 \quad a_*=\widehat E L_XDf(0)b_*.
 \label{eq:supp-exact-original-slope-directions}
\end{equation}
Then the complete constant coefficient and the slope polynomial are
\begin{equation}
 \begin{split}
 \mathcal Q(c)={}&D\mathfrak c_H(0)a_*
   +D^2\mathfrak c_H(0)[b(c),b(c)]\\
   &-\mathcal C_g\mathcal A_g^{-1}
                   D^2\mathfrak c_g(0)[b(c),b(c)].
 \end{split}
 \label{eq:supp-exact-original-slope-jets}
\end{equation}
In another fixed invertible harmonic row convention, the entire formula,
including $q_{0,\rm can}$ and $\mathcal C_g$, is transformed together.
\end{proposition}
\begin{proof}
Implicit differentiation of $\mathfrak c_g(\Xi)=0$ gives
$\Xi_r(0)=\widehat E$, $\Xi_y(0)=J_H$, and
$\mathcal A_g\gamma_{\alpha\beta}
 =-D^2\mathfrak c_g(0)[\Xi_\alpha,\Xi_\beta]$.
The chain rule for $\mathcal H=\mathfrak c_H\circ\Xi$ therefore gives the
last two terms.  Since $\mathcal F_r=L_XDf(0)\widehat E$, the regular
acceleration contribution is $D\mathfrak c_H(0)a_*$.  Insert these
identities into \eqref{eq:supp-exact-slope-polynomial}.
\end{proof}

\begin{theorem}[Necessary quadratic slope gate]
\label{thm:supp-exact-slope-gate}
Let $(r,y)$ be a Lipschitz solution of
\eqref{eq:supp-exact-reduced-dae} on $[0,T]$ with
$r(0)=y(0)=0$.  Then $\rho_0=0$ and
\begin{equation}
 r(\tau)=\tau v+\tau^2u_0+O(\tau^3),
 \qquad
 \mathcal Q\!\left(y(\tau)/\tau\right)=O(\tau).
 \label{eq:supp-exact-slope-estimate}
\end{equation}
Every cluster point of $y(\tau)/\tau$ is a real zero of $\mathcal Q$.
If the admissible slope ball contains no real zero, no such Lipschitz
trajectory exists.
\end{theorem}
\begin{proof}
Lipschitz continuity gives $r,y=O(\tau)$.  The integral regular equation
first yields $r=\tau v+O(\tau^2)$.  Taylor expansion of
$\mathcal H(r,y)=0$, using $\mathcal H_y(0,0)=0$, gives
$\rho_0=0$.  Since $\mathcal F_y(0,0)=0$,
\[
 r_\tau=v+\mathcal F_r r+O(\tau^2)
       =v+\tau\mathcal F_rv+O(\tau^2),
\]
so integration gives the displayed expansion for $r$.  With
$k(\tau)=y(\tau)/\tau$, the second-order Taylor expansion of
$\mathcal H$ gives
\[
 0=\frac{\tau^2}{2}\mathcal Q(k(\tau))+O(\tau^3),
\]
which proves the estimate and the cluster-point statement.
\end{proof}


\subsection{Source-specific rank-two factorization of the slope gate}
\label{subsec:supp-exact-rank-two-slope}

We now use the constant-shift structure of the same finite Hamiltonian.  For
a multiplier test $\nu$ write
\begin{equation}
 P_\nu=\partial_\lambda H_h[\nu],\qquad
 \mathcal B_\nu=D_XP_\nu F^{\rm can},\qquad
 \mathcal M=H_{\lambda\lambda},
 \label{eq:supp-exact-ward-objects}
\end{equation}
and let $H_c:\mathbb R^3\to\mathbb R^{108}$ insert the literal constant
shifts.  On the prepared first-field chart the finite coefficient identities
used below are
\begin{equation}
 F^{\rm can}(0,e)=0,
 \qquad \mathcal M(X,\lambda)=O(\|X\|^2),
 \label{eq:supp-exact-ward-mass-order}
\end{equation}
For $c\in\Ran H_c$, the coefficient branch used below retains
\begin{equation}
 \begin{gathered}
 P_{c,1}=0,\qquad P_{c,2}(U,p)=\langle p,\delta_cU\rangle,\\
 \mathcal M_2(X,e)c=0\quad\hbox{for every }X,\qquad
 [\delta_c,\delta_d]=0,\qquad \mathbb K_1(Y)c=0.
 \end{gathered}
 \label{eq:supp-exact-harmonic-mass-column}
\end{equation}
Here subscripts denote homogeneous degree in $X$ at $\lambda=e$ and
$\mathcal M_2(X,e)=\tfrac12D_X^2\mathcal M(0,e)[X,X]$.
The quadratic mass-column identity is required on the coefficient space,
not just at one value of $X$.  It is an explicit finite-action identity of
this branch; neither a weak-rank count nor $\mathcal M=O(\|X\|^2)$ alone
implies it.  The joint constraint initializer is analytic in the first field and first
multiplier tilt and has
\begin{equation}
 I(a,Y,\ell)=aY+a^2X_2(Y)+O(a^3),
 \label{eq:supp-exact-joint-initializer}
\end{equation}
with $X_2$ independent of $\ell$.  Its harmonic drift has the exact graded
form
\begin{equation}
 H_c^{\mathsf T}\mathcal B(I(a,Y,\ell),e+a\ell)
 =a^3\{t(Y)+\mathcal R_Y\ell\}+a^4\mathcal U(a,Y,\ell),
 \qquad \mathcal R_YH_c=0,
 \label{eq:supp-exact-harmonic-cubic-law}
\end{equation}
with analytic $\mathcal U$.  These are coefficient identities of the finite
action on the declared chart, not continuum symmetry assumptions.

\begin{lemma}[Constant-shift derivative orders]
\label{lem:supp-exact-ward-orders}
Along $X_a=aY+O(a^2)$ and $\lambda_a=e+O(a)$,
\begin{equation}
 \|D_X\mathcal B_c\|=O(a),\qquad
 \|D_\lambda\mathcal B_c\|=O(a^2),
 \qquad
 D_\lambda\mathcal B_c[H_cd]=O(a^3)
 \label{eq:supp-exact-ward-orders}
\end{equation}
for literal constant shifts $c,d$, uniformly on compact sets of the finite
parameters.
\end{lemma}
\begin{proof}
The vanishing linear constant-shift constraint coefficient gives
$D_XP_c=O(a)$, while $F^{\rm can}=O(a)$ and their required derivatives
are bounded.  Differentiation of $\mathcal B_c$ gives the first estimate.
The identity $\mathcal M_2(X,e)c=0$ for every $X$, symmetry of the mass,
and analyticity give
$D_X\mathcal M(X_a,\lambda_a)[c,\nu]=O(a^2)$:
the cubic remainder at $e$ contributes $O(a^2)$, and changing
$\lambda$ by $O(a)$ changes the quadratic derivative by $O(a^2)$.
For the multiplier derivative use
\[
 D_\lambda\mathcal B_c[\nu]
 =D_X\mathcal M[c,\nu]F^{\rm can}+\{P_c,P_\nu\}.
\]
The first term is $O(a^3)$.  The linear Poisson coefficient vanishes by the
first-Poisson constant-shift identity, giving $O(a^2)$.  When
$\nu=H_cd$, the quadratic bracket is the bracket of two commuting phase
translations and vanishes, leaving $O(a^3)$.
\end{proof}

Fix the actual initial rate $ar_*(a)=O(a)$ and put
\begin{equation}
 \mathcal F_a(Q)=\mathcal B(Q)+\mathcal M(Q)ar_*(a).
 \label{eq:supp-exact-fixed-target-residual}
\end{equation}
A pure harmonic second-record direction $J_Hc$ is realized by an
exact-constrained static tangent satisfying
\begin{equation}
 \delta X=O(a^5),\qquad
 \delta\lambda=a^2H_cc+O(a^3).
 \label{eq:supp-exact-pure-harmonic-ray}
\end{equation}
Let $q_4(\xi;J_Hc)$ denote the $a^4$ coefficient of the literal
constant-shift row of the directional residual at a fixed-first-field
second-record base $\xi$.

\begin{theorem}[Base-uniform harmonic quartic annihilation]
\label{thm:supp-exact-harmonic-q4-zero}
For every sufficiently small second-record base with the prepared first field
and first tilt fixed,
\begin{equation}
 q_4(\xi;J_Hc)=0\qquad(c\in\mathbb R^3).
 \label{eq:supp-exact-harmonic-q4-zero}
\end{equation}
Consequently its derivative with respect to any admissible second-record base
direction also vanishes.
\end{theorem}
\begin{proof}
At the initialized base, differentiate
\eqref{eq:supp-exact-harmonic-cubic-law} with respect to the first multiplier
tilt.  The physical multiplier tangent is $a^2H_cc$ and the associated
canonical tangent starts at order $a^4$; since $\mathcal R_YH_c=0$, the
$a^4$ constant-row coefficient vanishes.  The higher canonical and
multiplier parts of the tangent do not change that coefficient by
\Cref{lem:supp-exact-ward-orders} and
\eqref{eq:supp-exact-ward-mass-order}.

At a nearby second-record base the canonical contribution is
$O(a)O(a^5)=O(a^6)$, the leading harmonic multiplier contribution is
$O(a^3)O(a^2)=O(a^5)$, and every nonharmonic $O(a^3)$ multiplier correction
contributes only $O(a^5)$.  Finally
$D_Q\mathcal M[\delta Q]ar_*(a)=O(a^5)$.  Thus no $a^4$ coefficient is
created.  Differentiating the resulting identity with respect to the base
gives the final assertion.
\end{proof}

This coefficient audit also gives the harmonic identities in
\eqref{eq:supp-exact-slope-entrance}: the pure harmonic weak direction has no
regular first derivative and no first slow-compatibility column at the
prepared event.

Let
\begin{equation}
 \mathfrak g_2(u,w)
 :=D^2\mathfrak c_g(0)[u,w]
 =p_gC_{\rm red}D^2\Sigma_0(0)[u,w]
 \in\mathcal Z_g,
 \label{eq:supp-exact-g2}
\end{equation}
let $b_*$ be the zero-harmonic initial tangent after the two-row vertical
correction, and put
\begin{equation}
 \mathscr C=\mathcal C_g\mathcal A_g^{-1},
 \qquad
 \mathscr R(c)=2\mathfrak g_2(b_*,J_Hc)
                  +\mathfrak g_2(J_Hc,J_Hc).
 \label{eq:supp-exact-R-two-row}
\end{equation}
Let $\mathcal Q_{\rm can}$ denote the physical three-row normalization of
\eqref{eq:supp-exact-slope-polynomial} and
$q_{0,\rm can}=\mathcal Q_{\rm can}(0)$.

\begin{theorem}[Canonical rank-two factorization]
\label{thm:supp-exact-rank-two-factorization}
On the record-preserving native-lift branch,
\begin{equation}
 \mathcal Q_{\rm can}(c)
 =q_{0,\rm can}-\mathscr C\mathscr R(c).
 \label{eq:supp-exact-rank-two-factorization}
\end{equation}
Consequently
\begin{equation}
 D\mathcal Q_{\rm can}(c)
 =-\mathscr C D\mathscr R(c),
 \qquad
 \rank D\mathcal Q_{\rm can}(c)
 \le\rank\mathcal C_g\le2.
 \label{eq:supp-exact-rank-two-bound}
\end{equation}
\end{theorem}
\begin{proof}
Before applying the Ward identity, the linear and quadratic harmonic-slope
coefficients split into a lower regular-row Hessian contribution and a
quartic constant-row contribution.  Theorem
\ref{thm:supp-exact-harmonic-q4-zero} removes the latter whenever one
argument is a pure harmonic weak direction.  The surviving coefficients are
\[
 L_{\rm can}e_i
 =-2\mathcal C_g\mathcal A_g^{-1}
     \mathfrak g_2(b_*,J_He_i),
\]
and
\[
 B_{\rm can}[e_i,e_j]
 =-\mathcal C_g\mathcal A_g^{-1}
     \mathfrak g_2(J_He_i,J_He_j).
\]
Reassembling the quadratic polynomial gives
\eqref{eq:supp-exact-rank-two-factorization}; differentiation gives
\eqref{eq:supp-exact-rank-two-bound}.
\end{proof}

\begin{corollary}[Cokernel obstruction]
\label{cor:supp-exact-slope-cokernel}
For every $\lambda\in\ker\mathcal C_g^{\mathsf T}$,
\begin{equation}
 \lambda^{\mathsf T}\mathcal Q_{\rm can}(c)
 =\lambda^{\mathsf T}q_{0,\rm can}\quad\text{for every }c.
 \label{eq:supp-exact-cokernel-constant}
\end{equation}
Hence
\begin{equation}
 q_{0,\rm can}\notin\Ran\mathcal C_g
 \Longrightarrow
 \mathcal Q_{\rm can}(c)\ne0\quad\text{for every }c.
 \label{eq:supp-exact-cokernel-no-root}
\end{equation}
\end{corollary}
\begin{proof}
The image of $\mathscr C$ equals the image of $\mathcal C_g$.  Pair
\eqref{eq:supp-exact-rank-two-factorization} with a left-null covector.
\end{proof}

Assume now
\begin{equation}
 \rank\mathcal C_g=2,
 \qquad
 q_{0,\rm can}\in\Ran\mathcal C_g.
 \label{eq:supp-exact-rank-two-membership}
\end{equation}
Then $\mathscr C$ is an isomorphism from $\mathcal Z_g$ onto its range; let
$y_0\in\mathcal Z_g$ be defined by $\mathscr C y_0=q_{0,\rm can}$.

\begin{theorem}[Two-equation physical slope reduction]
\label{thm:supp-exact-two-equation-slope}
Under \eqref{eq:supp-exact-rank-two-membership},
\begin{equation}
 \mathcal Q_{\rm can}(c)=0
 \quad\Longleftrightarrow\quad
 \mathscr R(c)=y_0.
 \label{eq:supp-exact-two-equation-slope}
\end{equation}
At every root where $D\mathscr R$ has rank two, the zero set is locally a
one-dimensional real-analytic submanifold of $\mathbb R^3$.
\end{theorem}
\begin{proof}
Equation \eqref{eq:supp-exact-rank-two-factorization} becomes
$\mathcal Q_{\rm can}=\mathscr C(y_0-\mathscr R)$.  Injectivity of
$\mathscr C$ proves the equivalence, and the regular-level-set theorem gives
the local dimension.
\end{proof}

For a rigorous numerical certificate one should work on a transverse
$2$-plane, rather than invert the singular $3\times3$ slope Jacobian.  Choose
an affine splitting
\begin{equation}
 c=c_*+Ey+sn,
 \qquad E:\mathbb R^2\to\mathbb R^3,
 \qquad [E\ n]\in GL_3(\mathbb R),
 \label{eq:supp-exact-slope-slice}
\end{equation}
and put
\begin{equation}
 F_s(y)=\mathscr R(c_*+Ey+sn)-y_0.
 \label{eq:supp-exact-slope-slice-map}
\end{equation}

\begin{proposition}[Two-dimensional contraction certificate]
\label{prop:supp-exact-slope-contraction}
Fix $s=s_0$, a centre $\bar y\in\mathbb R^2$, an invertible
$M\in\mathbb R^{2\times2}$ and $r>0$.  Write
\begin{equation}
 F_{s_0}(\bar y+q)=F_{s_0}(\bar y)+Jq+\mathbb B[q,q]
 \label{eq:supp-exact-slope-quadratic-Taylor}
\end{equation}
with symmetric bilinear $\mathbb B$, and set
\begin{equation}
 \eta=\|MF_{s_0}(\bar y)\|,
 \qquad \theta=\|I-MJ\|,
 \qquad \beta=2\|M\mathbb B\|,
 \qquad q_*=\theta+\beta r.
 \label{eq:supp-exact-slope-contraction-data}
\end{equation}
If
\begin{equation}
 q_*<1,\qquad \eta+q_*r<r,
 \label{eq:supp-exact-slope-contraction-condition}
\end{equation}
then $F_{s_0}$ has a unique zero in the closed radius-$r$ ball about
$\bar y$, lying in its interior.  If the resulting $2\times2$ derivative
remains invertible on a product neighbourhood in $(s,y)$, these slice roots
continue to a local analytic root curve.
\end{proposition}
\begin{proof}
The map $T(q)=q-MF_{s_0}(\bar y+q)$ has derivative norm at most
$q_*<1$ on the ball and satisfies
$\|T(0)\|+q_*r<r$.  Banach's fixed-point theorem gives the unique zero;
the parameterized implicit-function theorem gives the final statement.
\end{proof}

The only quadratic source data required after the range test are
\begin{equation}
 \mathfrak g_2(b_*,J_He_i),\quad 1\le i\le3,
 \qquad
 \mathfrak g_2(J_He_i,J_He_j),\quad1\le i\le j\le3,
 \label{eq:supp-exact-g2-packet}
\end{equation}
namely eighteen scalar entries in the two-dimensional output.  Since
$\mathfrak c_g$ is quadratic on this source class, they can be obtained by
exact polarization; for example
\begin{equation}
 \mathfrak g_2(u,w)
 =\mathfrak c_g(u+w)-\mathfrak c_g(u)-\mathfrak c_g(w)+\mathfrak c_g(0)
 \label{eq:supp-exact-g2-polarization}
\end{equation}
when the arguments are expressed in the normalized coefficient chart.  No
quartic constant-row packet is required for these harmonic slope directions.


\subsection{Finite range certificates and a quantitative slope obstruction}
\label{subsec:supp-finite-range}

The range condition can be tested before any quadratic slope coefficient is
computed.  All inner products in this subsection are fixed positive
coefficient inner products, not the indefinite Cartan form.  Choose an
orthonormal basis of the three harmonic output rows and write
$C_*:=\mathcal C_g\in\mathbb R^{3\times2}$ and
$q_*:=q_{0,\rm can}$.  These must be coefficients at the same prepared
record, in the same row convention as the slope factorization.

\begin{theorem}[All-rank range test and sharp affine obstruction]
\label{thm:supp-finite-range-test}
Put $e_*=(I-C_*C_*^\dagger)q_*$.  Then
\begin{equation}
 q_*\in\Ran C_*\ \Longleftrightarrow\
 \rank(C_*\ q_*)=\rank C_*\ \Longleftrightarrow\ e_*=0.
 \label{eq:supp-finite-range-test}
\end{equation}
For every harmonic slope $c$,
\begin{equation}
 \|\mathcal Q_{\rm can}(c)\|\ge\|e_*\|.
 \label{eq:supp-finite-range-floor}
\end{equation}
This is the exact distance from zero to the affine space
$q_*-\Ran C_*$; it need not be the minimum over the quadratic image of
$\mathscr R$.  If $C_*=(c_1\ c_2)$ has rank two, set
$n=c_1\times c_2$.  Then
\begin{equation}
 \|e_*\|=\frac{|n^{\mathsf T}q_*|}{\|n\|}
 =\frac{|\det(c_1,c_2,q_*)|}{\|c_1\times c_2\|}.
 \label{eq:supp-finite-range-distance}
\end{equation}
For rank one, a nonzero column $u$ gives membership exactly when
$u\times q_*=0$; for rank zero it is exactly $q_*=0$.
\end{theorem}
\begin{proof}
$C_*C_*^\dagger$ is the orthogonal projection onto $\Ran C_*$.  Projection
of the slope factorization onto its orthogonal complement therefore gives
$e_*$, independently of $c$.  Orthogonal projection is contractive and the
nearest point in an affine subspace proves the distance assertion.  In rank
two its orthogonal complement is spanned by $n$; the other cases are the
corresponding line and zero-space tests.
\end{proof}

This bound is unaffected by a change of the two-dimensional domain basis.
Under a change of harmonic row coordinates the physical inner product must
be transported as well.  Membership is invariant without a norm choice;
a numerical Euclidean distance is not invariant under an arbitrary
unaccompanied row rescaling.

\paragraph{Certification with enclosing errors.}
Suppose $\|c_i-\widehat c_i\|\le\epsilon_i$ and
$\|q_*-\widehat q\|\le\epsilon_q$.  Let
$\widehat n=\widehat c_1\times\widehat c_2$ and put
\begin{equation}
 \begin{split}
 \epsilon_n&=\epsilon_1\|\widehat c_2\|
             +\|\widehat c_1\|\epsilon_2+\epsilon_1\epsilon_2,\\
 \epsilon_\Omega&=\epsilon_n\|\widehat q\|
                    +(\|\widehat n\|+\epsilon_n)\epsilon_q.
 \end{split}
 \label{eq:supp-finite-range-enclosures}
\end{equation}
The bilinear cross-product estimate and the triangle inequality give
$\|n-\widehat n\|\le\epsilon_n$ and
$|n^{\mathsf T}q_*-\widehat n^{\mathsf T}\widehat q|
\le\epsilon_\Omega$.  In particular, if
$|\widehat n^{\mathsf T}\widehat q|>\epsilon_\Omega$, rank two and
nonmembership are both certified, and every slope satisfies
\begin{equation}
 \|\mathcal Q_{\rm can}(c)\|\ge
 \frac{|\widehat n^{\mathsf T}\widehat q|-\epsilon_\Omega}
      {\|\widehat n\|+\epsilon_n}>0.
 \label{eq:supp-finite-range-certified-floor}
\end{equation}
An interval containing zero proves neither membership nor rank loss.  Exact
membership instead requires an exact dependence identity, or a validated
root argument when the prepared event is allowed to vary.

\paragraph{Evaluation before constructing a kernel basis.}
A lower-order coefficient realization may supply an invertible harmonic mass
$D_c$ and a map $R_c:\mathbb R^{26}\to\mathbb R^3$ with the exact incidence
\begin{equation}
 C_*=-D_c^{-1}R_cZ_g,\qquad d=D_cq_*,
 \label{eq:supp-finite-range-incidence}
\end{equation}
where $Z_g$ is the kernel basis in \eqref{eq:supp-exact-Zg-basis}.  This
incidence must be checked in the physical source normalization; it is not
inferred from the rank of $K_g$.  Membership is then exactly
$d\in R_c(\ker K_g)$.  The minus sign is immaterial for membership but is
retained when solving for a coefficient.

For an even skew matrix $K$, use the convention
$\operatorname{Pf}\left(\begin{smallmatrix}0&k\\-k&0\end{smallmatrix}\right)=k$
and define its polynomial Pfaffian adjugate $P(K)$ by
$P(K)_{ij}=(-1)^{i+j}\operatorname{Pf}(K_{\widehat i\widehat j})$ for $i<j$
with one-based indices, skew completion, and
$\operatorname{Pf}(0_{0\times0})=1$.  Then
$KP(K)=\operatorname{Pf}(K)I$.  Consider
\begin{equation}
 \mathbb B(K,R,d)=
 \begin{pmatrix}K&R^{\mathsf T}&0\\-R&0_3&d\\0&-d^{\mathsf T}&0\end{pmatrix}.
 \label{eq:supp-finite-range-border}
\end{equation}

\begin{proposition}[One bordered Pfaffian for the range obstruction]
\label{prop:supp-finite-range-pfaffian}
Let $K\in\mathbb R^{2m\times2m}$ have rank $2m-2$ and set
$S=RP(K)R^{\mathsf T}$, $n_P=(S_{23},-S_{13},S_{12})^{\mathsf T}$.
Then
\begin{equation}
 \operatorname{Pf}\mathbb B(K,R,d)=n_P^{\mathsf T}d.
 \label{eq:supp-finite-range-pfaffian}
\end{equation}
One has $n_P\ne0$ exactly when $\rank(R|_{\ker K})=2$; on that branch
the displayed Pfaffian vanishes exactly when $d\in R(\ker K)$.
At membership the border has rank $2m+2$, and away from membership it is
invertible.  Thus for $K=K_g$ and \eqref{eq:supp-finite-range-incidence}
this is a $30\times30$ test of the same range obstruction, not a different
slow-time compatibility condition.
\end{proposition}
\begin{proof}
The image of $P(K)$ is contained in $\ker K$.  A nonzero principal
$(2m-2)$-Pfaffian exists, so $P(K)$ is nonzero and has rank two.  For any
ordered kernel basis $(z_1,z_2)$ it is therefore a nonzero scalar multiple
of $z_1z_2^{\mathsf T}-z_2z_1^{\mathsf T}$.  Its pushforward is nonzero
exactly when $Rz_1,Rz_2$ are independent, and $n_P$ is then a nonzero
multiple of their cross product.  Expanding the border Pfaffian in its final
index, and then in the two remaining added indices, gives
$d_1S_{23}-d_2S_{13}+d_3S_{12}$ with the stated convention.

To check the rank at membership, split off an invertible skew
$(2m-2)$-block of $K$ by congruence.  The remaining six-dimensional skew
matrix couples the two kernel coordinates to the three harmonic coordinates
by a rank-two map, so it has rank at least four.  Membership makes its
Pfaffian zero, hence its even rank is at most four.  Restoring the invertible
block gives rank $2m+2$.  Away from membership the full Pfaffian is nonzero.
\end{proof}

The smaller pivot implementation is sufficient: solve
$K_{JJ}Z_J=-K_{JF}$, set $Z_g=\binom{Z_J}{I_2}$,
and test the triple product of the two columns of $R_cZ_g$ with $d$.
Its zero set agrees with that of the full border.  Exact skew elimination
or a rigorous interval implementation may be used; floating-point Pfaffian
algorithms alone do not certify a zero \cite{Wimmer2012}.

\paragraph{The quadratic calculation after membership.}
On the rank-two membership branch, put
$\varphi(c)=\mathfrak c_g(b_*+J_Hc)$.  Provided the displayed arguments lie
in the coefficient chart, the quadratic source identities imply
\begin{equation}
 \mathscr R(c)=2\{\varphi(c)-\varphi(0)\}.
 \label{eq:supp-finite-range-ten-point}
\end{equation}
Indeed $D\mathfrak c_g(0)J_H=0$ and the Hessian of $\mathfrak c_g$ is
$\mathfrak g_2$.  This identity does not require a real slope root.  For a
nonzero sufficiently small $\delta$, the ten values at
$0$, $\pm\delta e_i$, and $\delta(e_i+e_j)$ for $i<j$ determine the entire
quadratic $\varphi$: if
$\varphi(c)=\varphi(0)+\sum_i\ell_ic_i+\sum_i b_{ii}c_i^2+
2\sum_{i<j}b_{ij}c_ic_j$, then
\begin{align}
 \ell_i&=\frac{\varphi(\delta e_i)-\varphi(-\delta e_i)}{2\delta},\notag\\
 b_{ii}&=\frac{\varphi(\delta e_i)+\varphi(-\delta e_i)-2\varphi(0)}{2\delta^2},\notag\\
 b_{ij}&=\frac{\varphi(\delta(e_i+e_j))-\varphi(\delta e_i)
                    -\varphi(\delta e_j)+\varphi(0)}{2\delta^2}.
 \label{eq:supp-finite-range-interpolation}
\end{align}
These are exact polynomial identities, not asymptotic differencing formulas.
The root test remains \Cref{prop:supp-exact-slope-contraction}.  A root curve
of this necessary polynomial condition is not, by itself, a trajectory of
the nonlinear slow differential--algebraic system.

\begin{remark}[Scope of a finite coefficient verdict]
\label{rem:supp-finite-verdict-scope}
The acceleration recurrence and the range tests use actual coefficient data
at one identified event.  The rank and spectral inequalities in
\Cref{subsec:supp-exact-certified-rapid} do not specify the entries of
$\chi_j$, the differentiated-source terms, $C_*$ or $q_*$.  No value or
sign of either event-specific test is asserted here without those entries
and their exact or enclosing arithmetic.  In particular, a verified algebraic
identity for arbitrary input matrices is distinct from a verified numerical
outcome for the gravitational event.  Neither test alone supplies
positive-time propagation or spatial-refinement uniformity.
\end{remark}

\subsection{Scope of exact-action realization and the remaining propagation problem}

There are now four distinct exact-action statements. The original finite
response identities retain their all-cutoff scope wherever their own rank
hypotheses hold. The boundary realization on the amplitude-graded chart
yields actual constrained short histories, with readout accuracy measured
against that same history.  The exact initialized family pays the three mean rows at the first
finite-amplitude cut, while the remaining two rows are governed by the
acceleration certificate above; their feedback is retained in the reduced
mean rows.  The coefficient-complete acceleration recurrence and all-rank
range test are finite evaluation procedures, while arbitrary forcing admits
the defect-aware Lyapunov--Schmidt reduction.  Bounded-rate positive-time limits additionally
obey the quadratic slope gate and its source-specific rank-two factorization.
The normal factorizations and slow expansion identify separate obstructions
to forward residence and longer bounded-rate evolution. Their fixed-cutoff hypotheses cannot be interchanged with the
cofinal hypotheses of the Einstein handoff.

In particular, the $a^{-1}\Phi_1(\xi)$ contribution in slow time must be
retained in the principal compatibility problem. A frozen spectrum cannot
remove it: the elementary system $x_\tau=1$,
$h_\tau=a^{-1}x^2$ has no growing linear mean term at the origin but gives
$h_\tau=\tau^2/a$ from zero initial data. This example establishes only the
logical need for the nonlinear mean test, not a value or sign for a
coefficient of the gravitational action.

A direct exact-action Einstein limit would require a common family of full
finite coefficient certificates, a regular solution and propagation of the
nonlinear slow compatibility system, differentiated source and readout
control, and a noncollapsed physical interval independent of spatial
cutoff. A physical forward selection mechanism must additionally account
for the rapid and slow precision requirements without using future
endpoint data. None of these conclusions follows from homogeneous
blocking, a successful static determinant test, or an improved numerical
quadrature alone. Conversely, local residence failure is not a no-go theorem
for other exact histories, other gauges with controlled comparisons, or
singular restoration of continuum gauge symmetry.


\section{Homogeneous and Gowdy sectors; explicit geometric regulators}
\label{supp:sectors}

This appendix proves the sector-specific statements of
Section~\ref{sec:vacuum-main}. The homogeneous part retains independent
lapse variation and the complete trial normalizer; the Gowdy part retains
its smooth-slab and finite-offset assumptions. The final part collects the
flat and de Sitter spatial and exact-link estimates. These use the earlier
ADM map and classified action but have their own evolution and access
conditions.
\subsection{Homogeneous dynamics and action matching}
\label{supp:homogeneous-dynamics}

This appendix proves the homogeneous dynamical statements of
Section~\ref{sec:vacuum-main}.  Arbitrary positive homogeneous renewal variables
are first varied with the full lapse constraint retained; the flat/de~Sitter
profiles are solved only after the stationary equations have been obtained.  A second layer constructs a finite positive
trial law whose measured probabilities reconstruct the corresponding finite
action off shell.

\subsubsection{Homogeneous ADM map and reduced vacuum action}

On the root-symmetric branch
$k_{\alpha,h}=\kappa/(8h^2)$ and $m_h=\varrho h^3$.  Since
$\sum_{\alpha\in\Phi_{A_3}}\alpha\alpha^{\mathsf T}=8I$, one has
$\mathcal B=\kappa I$ and $\beta=0$.  Substitution in
\eqref{eq:adm-inversion-main} gives
\eqref{eq:main-isotropic-renewal-map}.

For completeness, let
\[
 \dd s^2=-N(t)^2\dd t^2+a(t)^2\delta_{ij}\dd x^i\dd x^j.
\]
The nonzero Christoffel symbols are
\[
 \Gamma^0_{00}=\dot N/N,\qquad
 \Gamma^0_{ij}=a\dot a\,\delta_{ij}/N^2,\qquad
 \Gamma^i_{0j}=\dot a\,\delta^i_j/a.
\]
Direct contraction gives
\begin{align}
 R_{00}&=-3\left(\ddot a/a-\dot a\dot N/(aN)\right),\\
 R_{ij}&=\left(a\ddot a/N^2+2\dot a^2/N^2
              -a\dot a\dot N/N^3\right)\delta_{ij},\\
 R&=6\left(\ddot a/(aN^2)+\dot a^2/(a^2N^2)
              -\dot a\dot N/(aN^3)\right).
 \label{eq:supp-homogeneous-curvature}
\end{align}
Thus
\begin{equation}
 Na^3R=6\frac{\dd}{\dd t}\left(\frac{a^2\dot a}{N}\right)
       -\frac{6a\dot a^2}{N}.
 \label{eq:supp-homogeneous-total-derivative}
\end{equation}
With the orientation/sign convention used here, the bare torsion-free
Einstein--Hilbert representative on $[t_0,t_1]\times\mathbb T^3$ is therefore
\begin{equation}
 \frac{S_{\rm EH}}{\chi V_0}
 =6\left[\frac{a^2\dot a}{N}\right]_{t_0}^{t_1}
 +\int_{t_0}^{t_1}\left[-\frac{6a\dot a^2}{N}
                         -2\Lambda Na^3\right]\dd t.
 \label{eq:supp-homogeneous-eh-split}
\end{equation}
The Gibbons--Hawking--York Dirichlet boundary term
\cite{York1972,GibbonsHawking1977} contributes
\begin{equation}
 \frac{S_{\rm GHY}}{\chi V_0}
 =-6\left[\frac{a^2\dot a}{N}\right]_{t_0}^{t_1},
 \label{eq:supp-homogeneous-ghy}
\end{equation}
so $S_{\rm EH}+S_{\rm GHY}=\chi V_0S[a,N]$ with $S[a,N]$ exactly
\eqref{eq:main-homogeneous-action-an}.  The corresponding Dirichlet problem
fixes the induced spatial metric at the two endpoint slices,
$\delta a(t_0)=\delta a(t_1)=0$, while the lapse remains independently variable
in the interior.  The finite action \eqref{eq:main-discrete-homogeneous-action}
is a midpoint variational discretization of this already boundary-completed
first-order functional; hence fixing $q_0,q_M$ is precisely the reduced discrete
counterpart of the same induced-metric boundary condition.  Substituting
$a=\varrho^{1/3}$ and $N=\varrho^{1/3}\sqrt\kappa$ gives
\eqref{eq:main-renewal-rate-action}.

\begin{proof}[Detailed proof of \Cref{thm:main-homogeneous-friedmann}]
Varying the lapse and scale factor independently gives
\begin{equation}
 E_N=2a^3(3\mathcal H^2-\Lambda),\qquad
 E_a=6Na^2(2D_\tau\mathcal H+3\mathcal H^2-\Lambda).
 \label{eq:supp-homogeneous-EaEN}
\end{equation}
Because
$a=e^{v/3}$ and $N=e^{u/2+v/3}$,
\[
 E_u=\frac N2E_N,\qquad
 E_v=\frac a3E_a+\frac N3E_N,
\]
which proves \eqref{eq:main-renewal-euler}.  The Jacobian of
$(u,v)\mapsto(\log a,\log N)$ is $-1/6\ne0$, so this change of
variables loses no variation.  The reconstructed Einstein residual is
\begin{equation}
 G_{00}+\Lambda g_{00}=N^2(3\mathcal H^2-\Lambda),
 \qquad
 G_{ij}+\Lambda g_{ij}
 =-a^2(2D_\tau\mathcal H+3\mathcal H^2-\Lambda)\delta_{ij},
 \label{eq:supp-homogeneous-Einstein-residual}
\end{equation}
with $G_{0i}=0$.  Hence the two renewal Euler equations are equivalent to every
Einstein component within this ansatz.
\end{proof}

\begin{proof}[Proof of \Cref{cor:main-derived-desitter-rates}]
For $\Lambda>0$, the lapse constraint gives
$\mathcal H^2=H_0^2$ with $H_0=\sqrt{\Lambda/3}$.  Continuity on a
connected interval fixes one sign $\sigma$, while
$D_\tau\mathcal H=0$ makes it constant.  Integrating
$D_\tau\log a=\sigma H_0$ gives
\eqref{eq:main-stationary-homogeneous-solutions}; the renewal variables then
follow from \eqref{eq:main-isotropic-renewal-map}.  The cases $\Lambda=0$ and
$\Lambda<0$ follow from the same squared constraint.
\end{proof}

\begin{remark}[Why the lapse is varied before gauge fixing]
If $q=a^{3/2}$ and one sets $N=1$ before variation, the remaining equation is
$q''=\omega^2q$ with $\omega=\sqrt{3\Lambda/4}$.  Its general solution
$q=Ae^{\omega t}+Be^{-\omega t}$ satisfies the omitted lapse constraint
$(q')^2=\omega^2q^2$ only when $AB=0$.  Thus a positive cosh-type solution is a
spurious critical point of the prematurely gauge-fixed action.
\end{remark}

\subsubsection{Finite lapse-varied action and stationary recurrence}

This is the midpoint discretization of the already
Gibbons--Hawking--York-completed reduced functional
\eqref{eq:main-homogeneous-action-an}, not of the bare second-derivative
Einstein--Hilbert expression.  The Dirichlet endpoint data are therefore
$q_0,q_M$ fixed (equivalently the induced endpoint three-metrics are fixed),
while every interior $q_j$ and every lapse-weighted interval $s_j$ are varied
independently.  Consequently
\begin{equation}
 \delta S_{\rm d}
 =\sum_{j=1}^{M-1}E_{q_j}\,\delta q_j
  +\sum_{j=0}^{M-1}E_{s_j}\,\delta s_j,
 \qquad \delta q_0=\delta q_M=0,
 \label{eq:main-discrete-dirichlet-variation}
\end{equation}
with no omitted endpoint acceleration term.  If one discretized the uncompleted
Einstein--Hilbert expression instead, a matching discrete gravitational boundary
term would have to be retained explicitly.  Through
\begin{equation}
 \varrho_j=q_j^2,\qquad
 N_{j+1/2}=s_j/\Delta t_j,\qquad
 \kappa_{j+1/2}=\frac{(s_j/\Delta t_j)^2}{\bar q_j^{4/3}},
 \label{eq:main-discrete-renewal-map}
\end{equation}
this is a smooth finite-dimensional action in positive renewal coordinates; no
exponential history appears in its definition.

Put $q_j>0$, $s_j>0$, $\bar q_j=(q_j+q_{j+1})/2$, and define
\eqref{eq:main-discrete-homogeneous-action}.  Because this discretizes the
boundary-completed first-order action, the discrete Dirichlet problem fixes
$q_0,q_M$ and varies $(q_1,\ldots,q_{M-1},s_0,\ldots,s_{M-1})$ independently.
A direct summation-by-parts calculation gives exactly
\eqref{eq:main-discrete-dirichlet-variation}; the endpoint momenta multiply
$\delta q_0$ and $\delta q_M$ and therefore vanish under the fixed induced
endpoint geometry.  No condition is imposed on $\sum_js_j$, since that would
modify the lapse variation and hence the Hamiltonian constraint.  The change of
variables \eqref{eq:main-discrete-renewal-map}, together with
$q_j=\sqrt{\varrho_j}$, is a smooth bijection on the positive finite domain.
Hence critical points in $(q,s)$ and in the corresponding positive renewal
coordinates are equivalent.

\begin{proof}[Proof of the Euler equations and stationary classification in
\Cref{thm:main-finite-homogeneous-stationarity}]
Differentiating an edge term with respect to $s_j$ gives
\[
 A_0\frac{(q_{j+1}-q_j)^2}{s_j^2}-B_0\bar q_j^2,
\]
which is \eqref{eq:main-discrete-lapse-constraint}.  At an interior node the
right derivative of the preceding edge and left derivative of the following
edge give \eqref{eq:main-discrete-scale-euler}.

For $\Lambda>0$, the lapse equation says
$v_j=\sigma_j\omega\bar q_j$ with $\sigma_j\in\{\pm1\}$.  Solving for the
terminal amplitude gives
\[
 q_{j+1}=q_j\frac{1+\sigma_j\omega s_j/2}
                      {1-\sigma_j\omega s_j/2}.
\]
Positivity requires $\omega s_j/2<1$.  Substitution in the interior momentum
matching condition shows
$2A_0\omega q_j(\sigma_j-\sigma_{j-1})=0$, so every positive stationary path
has one common sign.  Conversely a common-sign recurrence satisfies every
finite Euler equation.  For $\Lambda=0$ the lapse equation forces constant
$q$; for $\Lambda<0$ it has no positive real solution.

To prove \eqref{eq:main-finite-profile-convergence}, set
$x_j=\omega s_j/2$.  For $0\le x\le b$,
\[
 0\le2\operatorname{arctanh}x-2x
 =2\int_0^x\frac{t^2}{1-t^2}\dd t
 \le\frac{2x^3}{3(1-b^2)}.
\]
Taking logarithms of the recurrence and summing gives
\[
 \left|\log(q_j/q_0)-\sigma\omega\tau_j\right|
 \le\frac{\omega^3}{12(1-b^2)}\sum_i s_i^3
 \le\frac{\omega^3T}{12(1-b^2)}h_t^2.
\]
The power relations between $q,a,\varrho,\kappa$ then give the stated
convergence of renewal variables.
\end{proof}

For later residual estimates it is useful to define
\begin{equation}
 R_j^\sigma=q_{j+1}-q_j-\sigma\omega s_j\bar q_j.
 \label{eq:supp-homogeneous-residual}
\end{equation}
A direct expansion gives the exact factorization
\begin{equation}
 S_{\rm d}=-A_0\sum_j\frac{(R_j^\sigma)^2}{s_j}
 -A_0\sigma\omega(q_M^2-q_0^2).
 \label{eq:supp-homogeneous-square-factor}
\end{equation}
Thus at fixed endpoints the finite stationary branch is exactly the zero set of
the nonnegative residual energy
$\mathcal E=A_0\sum_j(R_j^\sigma)^2/s_j$.

\begin{lemma}[Full finite variation controlled by the residual energy]
\label{lem:supp-homogeneous-euler-control}
For endpoint-vanishing nodal variations $w_j$ and lapse variations
$\delta s_j=s_jn_j$, put
\begin{equation}
 \mathcal B_s(w,n;q)=\sum_js_j\left(
 \frac{w_{j+1}-w_j}{s_j}-\sigma\omega\bar w_j
 -\sigma\omega\bar q_jn_j\right)^2.
 \label{eq:supp-homogeneous-variation-norm}
\end{equation}
Then
\begin{equation}
 |\delta S_{\rm d}[w,sn]|
 \le2\sqrt{A_0\mathcal E}\sqrt{\mathcal B_s(w,n;q)}
 +\mathcal E\|n\|_\infty.
 \label{eq:supp-homogeneous-euler-bound}
\end{equation}
If $q_j^*$ is the exact common-sign stationary recurrence with the same
$q_0,s_j$, and $\omega\max s_j/2\le b<1$, then on a fixed positive amplitude
range
\begin{equation}
 \max_j|q_j-q_j^*|\le C\sqrt{\mathcal E}.
 \label{eq:supp-homogeneous-path-control}
\end{equation}
\end{lemma}

\begin{proof}
Differentiate \eqref{eq:supp-homogeneous-square-factor}; weighted
Cauchy--Schwarz gives \eqref{eq:supp-homogeneous-euler-bound}.  For the path
estimate, write the forced recurrence as
$q_{j+1}=r_jq_j+R_j/d_j$ with
$d_j=1-\sigma\omega s_j/2$ and
$r_j=(1+\sigma\omega s_j/2)/d_j$.  Iteration, $d_j\ge1-b$, and
$\sum_j|R_j|\le\sqrt{T\sum_jR_j^2/s_j}$ give the result.
\end{proof}

\subsubsection{Measured quadratic costs and exact off-shell action matching}

We now prove that the finite action can be read from an actual normalized trial
law rather than appended as an unrelated deterministic writer.  Let
$E_s(x,y)$ be \eqref{eq:main-measured-quadratic-cost}.

\begin{proposition}[Three-cost identification and rank defect]
\label{prop:supp-three-costs}
Fix $s>0$ and distinct $u,v>0$.  Within the declared quadratic class,
$E_s(u,u)$, $E_s(u,v)$, and $E_s(v,u)$ determine
\begin{align}
 B&=\frac{E_s(u,u)}{su^2},\\
 C&=\frac{E_s(u,v)-E_s(v,u)}{2(v^2-u^2)},\\
 A&=\frac{s}{2(v-u)^2}\left[
 E_s(u,v)+E_s(v,u)
 -2Bs\left(\frac{u+v}{2}\right)^2\right].
 \label{eq:supp-three-cost-identification}
\end{align}
Moreover, with $\mathfrak D=AB-C^2$, $c=-C/A$, and
$\nu=\mathfrak D/A$,
\begin{equation}
 E_s(x,y)=\frac A s\left[y-x-cs\frac{x+y}{2}\right]^2
 +\nu s\left(\frac{x+y}{2}\right)^2.
 \label{eq:supp-cost-completion}
\end{equation}
Thus $\mathfrak D=0$ is exactly the rank-one compatibility condition.
\end{proposition}

\begin{proof}
The diagonal measurement determines $B$.  Reversal exchanges $x$ and $y$,
leaving the first two terms invariant and changing the sign of the $C$ term,
which determines $C$; the reversal sum then determines $A$.  Completing the
square gives \eqref{eq:supp-cost-completion}.
\end{proof}

\begin{proof}[Proof of \Cref{thm:main-operational-action-matching}]
From \eqref{eq:main-accepted-row-normalizer},
$Z_sQ_s=p_sa_s$.  Taking logarithms and using
$a_s=e^{-E_s/\varepsilon}$ proves \eqref{eq:main-observed-cost}.  Summing over
one path gives
\[
 \varepsilon\sum_j\left[\log(Q_j/p_j)+\log Z_j\right]
 =-\sum_jE_{s_j}(q_j,q_{j+1}).
\]
The $C$ term telescopes to $C(q_M^2-q_0^2)$, which is exactly cancelled by the
endpoint term in \eqref{eq:main-observed-homogeneous-action}; this proves the
off-shell identity \eqref{eq:main-offshell-action-match}.  Differentiating that
identity on the smooth positive parametric domain proves equality of all
interior and lapse first variations.

For the mismatch statement, minimizing $y\mapsto E_s(x,y)$ gives, for small
$s$, a ratio with logarithmic drift $c=-C/A$, so the limiting scale factor has
Hubble rate $2c/3$.  Hence its geometric cosmological parameter is
$\Lambda_{\rm geom}=4C^2/(3A^2)$.  Dividing the bulk action by $A/A_0$ shows
that its parameter is $\Lambda_{\rm act}=4B/(3A)$.  Therefore
\[
 G[g]+\Lambda_{\rm act}g
 =(\Lambda_{\rm act}-\Lambda_{\rm geom})g
 =\frac{4(AB-C^2)}{3A^2}g,
\]
which is \eqref{eq:main-measured-einstein-mismatch}.
\end{proof}

\begin{remark}[Why the acceptance normalizer is physical]
Retaining only $p_s$ and $Q_s$ would not determine the action: replacing
$E_s(x,y)$ by $E_s(x,y)+c_s(x)$ leaves $Q_s$ unchanged but multiplies $Z_s(x)$
by $e^{-c_s(x)/\varepsilon}$.  Such a source-dependent term can change the
lapse variation.  The total acceptance frequency is therefore a load-bearing
part of the complete physical trial table.
\end{remark}

\subsubsection{An explicit finite stochastic cylinder}

We give the construction underlying
\Cref{thm:main-explicit-operational-family}.  Set $s=T/M$ and, for large $M$,
\begin{equation}
 d=1-\sigma\omega s/2,\qquad
 r=\frac{1+\sigma\omega s/2}{1-\sigma\omega s/2},\qquad
 \delta_M=M^{-3},\qquad
 \varepsilon_M=\frac{M^{-5}}{1+\log M}.
 \label{eq:supp-explicit-scales}
\end{equation}
Starting from $F_0=\{q_0\}$, define support intervals recursively by
\[
 L_0=q_0/2,\quad U_0=3q_0/2,\qquad
 L_{j+1}=rL_j-\delta_M,\quad U_{j+1}=rU_j+\delta_M,
\]
and let $F_j$ be finite grids of mesh at most $\delta_M$ in these intervals.
Each source proposes uniformly from $F_{j+1}$ and accepts with
\begin{equation}
 a_j(x,y)=\exp\left[-\frac{A_0d^2(y-rx)^2}{s\varepsilon_M}\right].
 \label{eq:supp-explicit-acceptance}
\end{equation}
Because
$d(y-rx)=y-x-\sigma\omega s(x+y)/2$, this is exactly the rank-one cost
\eqref{eq:main-rank-one-cost}.

\begin{lemma}[History-level entropy--energy estimate]
\label{lem:supp-entropy-energy}
At stage $j$, let a finite reference row $p_j$ and nonnegative cost $E_j$ define
$Q_j=p_je^{-E_j/\varepsilon}/Z_j$.  Let $\mathbb Q$ be the product of these
accepted rows and $\mathbb P\ll\mathbb Q$ any complete path law with the same
initial value, possibly with memory.  If each relevant row contains a proposal
of cost at most $e_j$ and every proposal weight is at least $p_{*,j}$, then
\begin{equation}
 \mathbb E_{\mathbb P}\sum_jE_j
 \le\sum_je_j+\varepsilon\left[
 \sum_j\log(1/p_{*,j})+D_{\rm KL}(\mathbb P\Vert\mathbb Q)\right].
 \label{eq:supp-entropy-energy}
\end{equation}
\end{lemma}

\begin{proof}
For a fixed history prefix, direct substitution gives
\[
 \mathbb E_P E_j
 =\varepsilon D_{\rm KL}(P_j\Vert Q_j)
 -\varepsilon D_{\rm KL}(P_j\Vert p_j)-\varepsilon\log Z_j.
\]
The candidate bound gives
$Z_j\ge p_{*,j}e^{-e_j/\varepsilon}$.  Integrating over prefixes, summing, and
using the chain rule for path relative entropy proves
\eqref{eq:supp-entropy-energy}.
\end{proof}

\begin{proof}[Proof of \Cref{thm:main-explicit-operational-family}]
The rational stationary multiplier and its inverse are uniformly bounded on
$0\le j\le M$ when $\omega s/2\le b<1$.  Iterating the support recursion shows
that all $F_j$ lie in one fixed interval $0<q_-\le q\le q_+<\infty$ and have
cardinality $O(M^3)$.  For every source $x$, the exact target $rx$ lies inside
the next support and has a grid point within $\delta_M$.  The corresponding
cost is at most
$C\delta_M^2/s=CM^{-5}$.  Hence the proposal floor is $p_*\ge cM^{-3}$ and the
total acceptance probability has a positive explicit lower bound.

Apply \Cref{lem:supp-entropy-energy} to the ideal accepted product law
$\mathbb Q_M$ itself.  With the scales
\eqref{eq:supp-explicit-scales},
\begin{equation}
 \mathbb E_{\mathbb Q_M}\mathcal E_M\le CM^{-4}.
 \label{eq:supp-explicit-energy}
\end{equation}
The deterministic estimates
\eqref{eq:supp-homogeneous-euler-bound} and
\eqref{eq:supp-homogeneous-path-control}, Jensen's inequality, and the
second-order stationary-profile estimate then give the two bounds in
\eqref{eq:main-stochastic-profile} and an $O(M^{-2})$ expected first variation
against every fixed smooth scale/lapse test.

To make every cutoff history genuinely finite, permit at most
\[
 R_M=\left\lceil7\underline Z_M^{-1}\log M\right\rceil
\]
independent proposals at each stage, where $\underline Z_M$ is the explicit
row-wise acceptance lower bound.  A union bound gives total failure probability
at most $M^{-6}$.  Conditional on complete success, the path-law correction has
relative entropy at most $-\log(1-M^{-6})$, which changes
\eqref{eq:supp-explicit-energy} only by a summable term.  The expected residual
energies are summable in $M$, so Tonelli and the first Borel--Cantelli lemma give
$\mathcal E_M\to0$ almost surely under any common coupling.  The deterministic
path estimates then imply uniform and strong $H^1$ convergence to
$q_0e^{\sigma\omega\tau}$.

Finally, \eqref{eq:main-stochastic-root-rates} is the positive ADM inverse on
the proper-time nodal branch.  The $A_3$ root identity gives exactly
$\mathcal B_j=\kappa_jI$, and ordinary Taylor expansion of the symmetric root
generator yields the spatial continuum coefficient
$\tfrac12q_*^{-4/3}\Delta$.  Together with the homogeneous Einstein identity,
this proves the vacuum limit with $\Lambda=4\omega^2/3$.
\end{proof}

\begin{remark}[Exact scope of the operational construction]
The quadratic duration law, the rank-one condition $\mathfrak D=0$, the
nonzero action normalization, and the value of $\Lambda$ are not deduced from
the bare predictive quotient.  They are finite operational conditions that can
be measured and can fail.  The construction proves that this matched class is
nonempty and robust under vanishing measured action/likelihood errors; it does
not prove generic microscopic attraction to that class.
\end{remark}

\subsubsection{Exact conservative blocking and the relative defect}
\label{subsec:supp-homogeneous-blocking}

Put $u=A/s+Bs/4$, $v=A/s-Bs/4$, $\mu=\sqrt{AB}$ and
$\theta=2\operatorname{artanh}((s/2)\sqrt{B/A})$.  Then
$u^2-v^2=\mu^2$, $u=\mu\coth\theta$ and
$v=\mu\operatorname{csch}\theta$, giving
\eqref{eq:main-hyperbolic-cost}.

\begin{lemma}[Two-edge Schur identity]
\label{lem:supp-two-edge-schur}
Let
$E_i(x,y)=p_ix^2-2q_ixy+r_iy^2$ with $p_i,q_i,r_i>0$ and
$d_i=p_ir_i-q_i^2\ge0$.  Then the minimizer of
$E_1(x,z)+E_2(z,y)$ is positive on positive endpoints and the blocked quadratic
has
\begin{equation}
 p'=p_1-\frac{q_1^2}{r_1+p_2},\qquad
 q'=\frac{q_1q_2}{r_1+p_2},\qquad
 r'=r_2-\frac{q_2^2}{r_1+p_2},
 \label{eq:supp-schur-coefficients}
\end{equation}
with
\begin{equation}
 p'r'-(q')^2=\frac{r_2d_1+p_1d_2}{r_1+p_2}.
 \label{eq:supp-schur-determinant}
\end{equation}
\end{lemma}
\begin{proof}
Complete the square in $z$.  The minimizer is
$z_*=(q_1x+q_2y)/(r_1+p_2)>0$.  The three coefficients in
\eqref{eq:supp-schur-coefficients} are the remaining quadratic coefficients;
direct expansion gives \eqref{eq:supp-schur-determinant}.
\end{proof}

\begin{proof}[Proof of \Cref{thm:main-no-blocking-attractor}]
For $\mathcal E_{\mu,\theta_i,C}$ put
$u_i=\mu\coth\theta_i$ and $v_i=\mu\operatorname{csch}\theta_i$.
The two intermediate $Cz^2$ terms cancel.  Applying
\Cref{lem:supp-two-edge-schur} and $u_i^2-v_i^2=\mu^2$ gives
\begin{equation}
 u'=\frac{u_1u_2+\mu^2}{u_1+u_2},\qquad
 v'=\frac{v_1v_2}{u_1+u_2},\qquad C'=C.
\end{equation}
The hyperbolic addition formulas identify these with
$\mu\coth(\theta_1+\theta_2)$ and
$\mu\operatorname{csch}(\theta_1+\theta_2)$, proving
\eqref{eq:main-hyperbolic-blocking}.  Converting the two-identical-edge result
back to midpoint coefficients gives \eqref{eq:main-homogeneous-block-map}.
Hence $AB$, $C$, $\mathfrak D$, and
$\Delta_{\rm rel}=\mathfrak D/(AB)$ are invariant.

For a change of amplitude and time units, $x=r\widetilde x$ and
$s=\ell\widetilde s$, followed by common action rescaling by $d>0$, one has
$\widetilde A=r^2A/(d\ell)$,
$\widetilde B=r^2\ell B/d$, and $\widetilde C=r^2C/d$; therefore
$C^2/(AB)$ is unchanged.

For the finite positive statement, let $F\subset(0,\infty)$ be finite and
write the subprobability kernel as
\begin{equation}
 L_C(x,y)=e^{Cx^2/\varepsilon}L_0(x,y)e^{-Cy^2/\varepsilon},
 \qquad L_0(x,y)=L_0(y,x)>0.
 \label{eq:supp-finite-twist-similarity}
\end{equation}
Then $L_C^n=D_CL_0^nD_C^{-1}$ and hence, for $x\ne y$,
\begin{equation}
 \frac{\varepsilon}{2(x^2-y^2)}
 \log\frac{L_C^n(x,y)}{L_C^n(y,x)}=C.
 \label{eq:supp-finite-twist-read}
\end{equation}
Thus finite blocking preserves the measured endpoint twist exactly.  If the
intermediate positive amplitudes are integrated over a grid of covering error
$d$ and cardinality $m$, completion of the positive interior quadratic form
around its exact bridge minimizer gives
\begin{equation}
 0\le E_{F,n}^{\varepsilon}(x,y)-
 \mathcal E_{\mu,n\theta,C}(x,y)
 \le4u(n-1)d^2+\varepsilon(n-1)\log m.
 \label{eq:supp-finite-block-error}
\end{equation}
At fixed total duration one may take $d_n=n^{-2}$,
$m_n=O(n^2)$, and $\varepsilon_n=n^{-3}/(1+\log n)$, making the right side
$O(n^{-2})$.  The three-probe coefficient reconstruction is continuous on a
compact calibrated range with $AB>0$, so every initial
$\Delta_{\rm rel}>0$ remains bounded away from zero for sufficiently fine
finite blocking.  This proves the operational statement.
\end{proof}


\subsubsection{The cross-scale optimizer test}

\paragraph{Matching is an exact cross-scale consistency condition.}
The absence of an attracting flow does not leave the matched surface without
an intrinsic test.  The same blocking operation characterizes it by asking
whether local minimization commutes with the elimination of an intermediate
amplitude.

\begin{proposition}[Optimizer composition characterizes action--law matching]
\label{prop:main-optimizer-composition}
Let $m_\theta(x)$ minimize
$y\mapsto\mathcal E_{\mu,\theta,C}(x,y)$ over $y>0$ for $x>0$, and use
$\rho=C/\mu$ and $\Delta_{\rm rel}$ from
\eqref{eq:main-relative-rank-defect}.  Then
\begin{equation}
 m_\theta(x)=r_\theta x,\qquad
 r_\theta=(\cosh\theta+\rho\sinh\theta)^{-1},
 \label{eq:main-optimizer-map}
\end{equation}
and, for every $\theta,\phi>0$,
\begin{equation}
 r_{\theta+\phi}^{-1}-(r_\theta r_\phi)^{-1}
       =\Delta_{\rm rel}\sinh\theta\sinh\phi.
 \label{eq:main-optimizer-defect}
\end{equation}
Thus $AB=C^2$ is equivalent to
$m_{\theta+\phi}=m_\phi\circ m_\theta$ for every positive pair, and already
to that equality for one positive pair at one $x>0$.  On this branch
$r_\theta=e^{-\rho\theta}$ with $\rho=\pm1$.
\end{proposition}

The proof is in Appendix~\ref{supp:homogeneous-dynamics}.
Equation~\eqref{eq:main-optimizer-defect} gives a second finite diagnostic of
the relative defect, independent of its reconstruction from three cost
entries.  On compact calibrated ranges away from zero denominators it is
stable under small measurement errors.  It concerns the minimizing maps of
the reconstructed cost; identifying these with observed low-noise updates is
a further test of the accepted law.  It is a characterization of matching,
not a mechanism which makes an arbitrary law approach it.

\begin{proof}[Proof of \Cref{prop:main-optimizer-composition}]
The coefficient of $y^2$ in
$\mathcal E_{\mu,\theta,C}(x,y)$ is $\mu\coth\theta+C>0$, since
$|C|\le\mu$ and $\theta>0$.  Its unique minimizer is therefore
\[
 y=\frac{\mu\operatorname{csch}\theta}{\mu\coth\theta+C}\,x
       =\frac{x}{\cosh\theta+\rho\sinh\theta}>0.
\]
Set $d_\theta=\cosh\theta+\rho\sinh\theta$.  Hyperbolic addition gives
\[
 d_{\theta+\phi}-d_\theta d_\phi
       =(1-\rho^2)\sinh\theta\sinh\phi,
\]
which proves \eqref{eq:main-optimizer-defect}.  All denominators and both
hyperbolic sines are positive.  Equality of the maps for even one positive
pair and one nonzero amplitude therefore forces $1-\rho^2=0$, equivalently
$AB=C^2$.  Conversely, $\rho=\pm1$ gives $d_\theta=e^{\rho\theta}$ and
exact composition for every pair.  For $\Delta_{\rm rel}>0$ the formula
also gives $r_{\theta+\phi}<r_\theta r_\phi$.

Solving the same identity for the defect gives
\begin{equation}
 \Delta_{\rm rel}
 =\frac{r_{\theta+\phi}^{-1}-(r_\theta r_\phi)^{-1}}
              {\sinh\theta\sinh\phi}.
 \label{eq:supp-optimizer-defect-read}
\end{equation}
On a compact range where the observed ratios and
$\sinh\theta\sinh\phi$ are bounded away from zero, the right-hand side
has bounded first derivatives in its measured arguments.  It is consequently
locally Lipschitz there.  This is stability of a finite diagnostic, not
transverse attraction under the conservative blocking operation.
\end{proof}

\subsection{The nonlinear two-polarization regulator}
\label{supp:gowdy}

This appendix proves \Cref{thm:main-gowdy-regulator}.  The purpose is not to
establish unrestricted four-dimensional global regularity, but to show that the
finite renewal evolution, constraint, readout, and independent action tests can
be carried through a genuinely nonlinear spatially inhomogeneous vacuum sector.
We use the unpolarized Gowdy $T^3$ class in areal coordinates
\cite{Moncrief1981}; the lapse and shift are nevertheless retained as
independent variables until after the action variation.

\subsubsection{First-order frame system and constraints}

On $[t_0,t_1]\times\mathbb T^3$, $t_0>0$, consider
\begin{equation}
 g=t^{-1/2}e^{\lambda/2}(-\dd t^2+\dd\theta^2)
 +t\left[e^P(\dd x+Q\dd y)^2+e^{-P}\dd y^2\right].
 \label{eq:supp-gowdy-metric}
\end{equation}
With
$a=P_t$, $b=P_\theta$, $c=e^PQ_t$, $d=e^PQ_\theta$ and
\begin{equation}
 e=\frac12(a^2+b^2+c^2+d^2),\qquad f=ab+cd,
 \label{eq:supp-gowdy-ef}
\end{equation}
the vacuum equations take the first-order form
\begin{align}
 a_t&=b_\theta-a/t+c^2-d^2,&b_t&=a_\theta,\nonumber\\
 c_t&=d_\theta-c/t-ac+bd,&d_t&=c_\theta+ad-bc,\nonumber\\
 P_t&=a,&Q_t&=e^{-P}c,&\lambda_t&=2te,
 \label{eq:supp-gowdy-frame-system}
\end{align}
with
\begin{equation}
 P_\theta=b,\qquad e^PQ_\theta=d,\qquad
 \lambda_\theta=2tf.
 \label{eq:supp-gowdy-coordinate-constraints}
\end{equation}
A direct differentiation using \eqref{eq:supp-gowdy-frame-system} gives
\begin{equation}
 e_t-f_\theta=-\frac{a^2+c^2}{t},\qquad
 f_t-e_\theta=-\frac ft.
 \label{eq:supp-gowdy-stress-identities}
\end{equation}
Consequently
$\partial_t(\lambda_\theta-2tf)=0$ whenever the displayed equations hold, and
$\partial_t\int_{\mathbb T^1}tf\,\dd\theta=0$.  Thus the momentum constraint
and the periodic zero-mean condition for $\lambda$ propagate.

\subsubsection{Independent lapse and shift variation before areal gauge}

Let $(R,\lambda,P,Q,N,\beta)$ be independent fields with $R,N>0$, put
$v_u=\dot u-\beta u'$, and define
\begin{align}
 \mathscr L_{\rm G}={}&
 -\frac{v_R(v_\lambda-4\beta')}{2N}
 +\frac R{2N}(v_P^2+e^{2P}v_Q^2)\nonumber\\
 &+N\left[-2R''+\frac{R'\lambda'}2
 -\frac R2\{(P')^2+e^{2P}(Q')^2\}\right].
 \label{eq:supp-gowdy-adm-density}
\end{align}
The canonical momenta are
\begin{equation}
 p_R=-\frac{v_\lambda-4\beta'}{2N},\qquad
 p_\lambda=-\frac{v_R}{2N},\qquad
 \pi=\frac{Rv_P}{N},\qquad
 \varpi=\frac{Re^{2P}v_Q}{N}.
 \label{eq:supp-gowdy-momenta}
\end{equation}
Lapse and shift variation give respectively
\begin{align}
 \mathcal C={}&-2p_Rp_\lambda
 +\frac{\pi^2+e^{-2P}\varpi^2}{2R}
 +\frac R2\{(P')^2+e^{2P}(Q')^2\}
 +2R''-\frac12R'\lambda',\label{eq:supp-gowdy-H}\\
 \mathcal D={}&p_RR'+p_\lambda\lambda'+\pi P'+\varpi Q'-4p_\lambda'.
 \label{eq:supp-gowdy-D}
\end{align}
On a solution of \eqref{eq:supp-gowdy-frame-system}, set
\begin{equation}
 R=t,\quad N=1,\quad\beta=0,\quad
 p_\lambda=-\frac12,\quad p_R=-te,\quad
 \pi=ta,\quad\varpi=te^Pc.
 \label{eq:supp-gowdy-areal-lift}
\end{equation}
Substitution of \eqref{eq:supp-gowdy-coordinate-constraints} shows
$\mathcal C=0$ and $\mathcal D=0$; the remaining area equation follows from
\eqref{eq:supp-gowdy-stress-identities}.  Hence the areal-gauge solution is the
restriction of a critical point of the unreduced symmetry-class action rather
than a solution obtained after deleting the lapse and shift equations.

\subsubsection{Local splitting and exact staggered momentum}

Set
\begin{equation}
 U=(u_1,u_2,u_3,u_4)=\sqrt t\,(a,b,c,d),\qquad
 \mathcal J(U)=u_1u_2+u_3u_4,
 \qquad n(U)=|U|^2.
 \label{eq:supp-gowdy-rescaled-frame}
\end{equation}
The local source flow, with the evolution parameter denoted by $s$, is
\begin{align}
 \frac{\dd u_1}{\dd s}&=-\frac{u_1}{2t}
 +\frac{u_3^2-u_4^2}{\sqrt t},&
 \frac{\dd u_2}{\dd s}&=\frac{u_2}{2t},\nonumber\\
 \frac{\dd u_3}{\dd s}&=-\frac{u_3}{2t}
 +\frac{-u_1u_3+u_2u_4}{\sqrt t},&
 \frac{\dd u_4}{\dd s}&=\frac{u_4}{2t}
 +\frac{u_1u_4-u_2u_3}{\sqrt t},\nonumber\\
 \frac{\dd P}{\dd s}&=\frac{u_1}{\sqrt t},&
 \frac{\dd Q}{\dd s}&=e^{-P}\frac{u_3}{\sqrt t},&
 \frac{\dd t}{\dd s}&=1,
 \label{eq:supp-gowdy-source-flow}
\end{align}
with $\dd\lambda/\dd s=0$.  Let $\mathsf S_\sigma$ be its near-identity
implicit-midpoint map over source duration $\sigma$.

On a periodic grid of spacing $\ell$, let $S_\pm$ be the two shifts,
$D_+=(S_+-I)/\ell$, and $\mathsf A_+=(I+S_+)/2$.  Put
\begin{equation}
 w^+=\frac1{\sqrt2}(u_1+u_2,u_3+u_4),\qquad
 w^-=\frac1{\sqrt2}(u_1-u_2,u_3-u_4),\qquad
 g^\pm=|w^\pm|^2.
 \label{eq:supp-gowdy-characteristics}
\end{equation}
At frozen physical time define the transport map by
\begin{equation}
 (w^+)^+=S_+w^+,
 \qquad (w^-)^+=S_-w^-,
 \qquad
 \lambda^+=\lambda+\frac\ell2
 [(I+S_+)g^++(I+S_-)g^-],
 \label{eq:supp-gowdy-transport}
\end{equation}
leaving $P,Q$ fixed.

\begin{theorem}[Exact staggered momentum under symmetric splitting]
\label{thm:supp-gowdy-momentum}
With synchronized full step $h=\ell$, the source-half/transport/source-half
composition preserves exactly
\begin{equation}
 \mathcal K_\ell
 =D_+\lambda-2\mathsf A_+\mathcal J(U).
 \label{eq:supp-gowdy-staggered-momentum}
\end{equation}
The update is nearest neighbour, its background increment is a nonnegative sum
of squared characteristic amplitudes, and it is second-order consistent on
bounded smooth reference charts with $t\ge t_0>0$.
\end{theorem}

\begin{proof}
For the source flow, direct substitution in
\eqref{eq:supp-gowdy-source-flow} gives
$\nabla\mathcal J(U)\cdot\dot U=0$.  The difference of a quadratic function at
two endpoints equals its derivative at the arithmetic midpoint paired with the
increment, so implicit midpoint preserves $\mathcal J$ exactly while leaving
$\lambda$ fixed.  For the transport step,
\begin{equation}
 \Delta\mathcal J
 =\frac12[(S_+-I)g^+-(S_--I)g^-],\qquad
 D_+\Delta\lambda=2\mathsf A_+\Delta\mathcal J,
 \label{eq:supp-gowdy-transport-identity}
\end{equation}
which proves exact preservation without invoking an assumed discrete product
rule.  The $\lambda$ increment in \eqref{eq:supp-gowdy-transport} is a positive
sum of squared norms.  Characteristic translation is exact for the transport
fields, while the background increment is a trapezoidal approximation to the
integral of the translated squares with local error $O(\ell^3)$.  Implicit
midpoint has the same local order for the smooth source, and symmetric
composition therefore gives second-order consistency.
\end{proof}

\subsubsection{Residual, constraint, metric, and action budgets}

Let
$X_j=(U_j,P_j,Q_j,\lambda_j)$ satisfy
\begin{equation}
 X_{j+1}=\Phi_h(t_j,X_j)+r_j,
 \qquad h=\ell,
 \label{eq:supp-gowdy-residual-update}
\end{equation}
and define
\begin{equation}
 \mathcal E_h=\sum_j\frac{\|r_j\|_{2,\ell}^2}{2h},\qquad
 \mathcal F_h=\ell^{-(2r+1)}\mathcal E_h.
 \label{eq:supp-gowdy-residual-energy}
\end{equation}
Here $\mathsf S_\ell$ denotes grid sampling.

\begin{proposition}[Residual and constraint control]
\label{prop:supp-gowdy-residual-control}
On a fixed smooth reference chart, with initial $C^r_\ell$ error $\epsilon_0$
and either a smallness bootstrap or an enforced offset envelope,
\begin{equation}
 \max_j\|X_j-\mathsf S_\ell X_*(t_j)\|_{r,\infty,\ell}
 \le C(\epsilon_0+h^2+\sqrt{\mathcal F_h}).
 \label{eq:supp-gowdy-state-error}
\end{equation}
Writing $\widehat U_j$ for the deterministic predicted frame, the exact
constraint increment is
\begin{align}
 \mathcal K_\ell(X_{j+1})-\mathcal K_\ell(X_j)
 ={}&D_+r_j^\lambda\nonumber\\
 &-2\mathsf A_+
 [\nabla\mathcal J(\widehat U_j)\cdot r_j^U
   +\mathcal J(r_j^U)].
 \label{eq:supp-gowdy-constraint-increment}
\end{align}
For $r\ge1$, its accumulated norm is bounded by the initial residual plus
$C\sqrt{\mathcal F_h}$.
\end{proposition}

\begin{proof}
The deterministic split map has smooth-chart Lipschitz constant $1+Ch$ and
consistency error $Ch^3$.  Discrete Gronwall therefore bounds the global error
by
$C(\epsilon_0+h^2+\sum_j\|r_j\|_{r,\infty,\ell})$.
The one-dimensional inverse grid estimate and Cauchy--Schwarz bound the last
sum by $C\sqrt{\mathcal F_h}$, proving
\eqref{eq:supp-gowdy-state-error}.  Expanding the quadratic $\mathcal J$ about
$\widehat U_j$ and using exact preservation
\eqref{eq:supp-gowdy-staggered-momentum} gives
\eqref{eq:supp-gowdy-constraint-increment}; the same inverse estimate controls
its accumulated norm.
\end{proof}

For the pathwise curvature statement specialize the finite offset alphabet so
that every accepted offset obeys
\begin{equation}
 \|\epsilon_j\|_\infty\le c_0h^4,
 \label{eq:supp-gowdy-offset-cap}
\end{equation}
and take initial $C_h^1$ error $O(h^2)$.  The retained source implementation may
add a solver error $O(h^5)$ in $C_h^1$.  Since
$\|D_hv\|_\infty\le2h^{-1}\|v\|_\infty$, the offset has $C_h^1$ size
$O(h^3)$.  The one-step Lipschitz factor $1+Lh$, second-order consistency, and
discrete Gronwall then give, pathwise,
\begin{equation}
 \max_j\|X_j-\mathsf S_hX_*(t_j)\|_{1,\infty,h}\le Ch^2.
 \label{eq:supp-gowdy-pathwise-c1}
\end{equation}
For $f=P,Q,\lambda$, subtraction of the reference step and division by $h$
also gives an $O(h^2)$ bound for the nodal time differences of the errors.

Use the unchanged derivative records
\begin{align}
 u&=\left(u_1/\sqrt t,\ e^{-P}u_3/\sqrt t,\ |U|^2\right),\nonumber\\
 v&=\left(u_2/\sqrt t,\ e^{-P}u_4/\sqrt t,\ 2\mathcal J(U)\right),
 \qquad w=D_0u,
 \label{eq:supp-gowdy-hermite-jets}
\end{align}
where $D_0=(S_+-S_-)/(2h)$.  These are respectively the time, space, and mixed
jet records for $(P,Q,\lambda)$.  On each spacetime square let
$\mathcal H_h(f,u,v,w)$ be the tensor-product cubic Hermite interpolant using
\begin{align*}
 H_0(q)&=2q^3-3q^2+1,&H_1(q)&=-2q^3+3q^2,\\
 G_0(q)&=q^3-2q^2+q,&G_1(q)&=q^3-q^2.
\end{align*}
The identities $H_0+H_1=1$, $H_0'+H_1'=0$, and
$H_0''+H_1''=0$ imply the stability estimate
\begin{equation}
 \|F\|_\infty+\|DF\|_\infty\le C\eta,
 \qquad \|D^2F\|_\infty\le C\eta/h,
 \label{eq:supp-gowdy-hermite-second-jet}
\end{equation}
whenever the two nodal jet banks and the first nodal differences of their value
records differ by at most $\eta$.  Shared corner jets make the interpolant
$C^1$ across cell interfaces, so its weak second derivative is the cellwise
second derivative and has no interface measure.

\begin{theorem}[Pathwise full-curvature control on the Gowdy alphabet]
\label{thm:supp-gowdy-full-curvature}
Assume the smooth reference has bounded derivatives through the order required
by the Hermite remainder, \eqref{eq:supp-gowdy-offset-cap}, and
\eqref{eq:supp-gowdy-pathwise-c1}.  For the shared Hermite reconstruction of
$f=(P,Q,\lambda)$,
\begin{equation}
 \|f_h-f_*\|_{W^{1,\infty}}\le Ch^2,
 \qquad
 \|D^2f_h-D^2f_*\|_{L^\infty}\le Ch.
 \label{eq:supp-gowdy-second-jet}
\end{equation}
For the reconstructed metric \eqref{eq:supp-gowdy-metric},
\begin{align}
 \|g_h-g_*\|_{W^{1,\infty}}&\le Ch^2,\nonumber\\
 \|D^2g_h-D^2g_*\|_{L^\infty}
 +\|\operatorname{Riem}(g_h)-\operatorname{Riem}(g_*)\|_{L^\infty}
 &\le Ch.
 \label{eq:supp-gowdy-full-curvature}
\end{align}
Hence on every compact positive subslab
\begin{equation}
 \sup_h\|\operatorname{Riem}(g_h)\|_{L^2}<\infty,
 \qquad \|\operatorname{Ric}(g_h)\|_{L^\infty}\le Ch.
 \label{eq:supp-gowdy-curvature-certificate}
\end{equation}
All estimates are pathwise and the curvature limit is the Levi--Civita
curvature of $g_*$.  No all-time curvature norm is assumed.
\end{theorem}
\begin{proof}
The right sides of \eqref{eq:supp-gowdy-hermite-jets} are smooth functions on
the standing chart.  Equation \eqref{eq:supp-gowdy-pathwise-c1}, the nodal time
difference estimate, and centered-difference consistency therefore make every
corner value/first/mixed jet error $O(h^2)$.  Apply
\eqref{eq:supp-gowdy-hermite-second-jet}; the exact-reference cubic Hermite
remainder is $O(h^{4-k})$ through derivative order $k=2$.  This proves
\eqref{eq:supp-gowdy-second-jet}.

On $t\ge t_0>0$ the Gowdy metric map is smooth on the bounded chart and has
determinant $-te^\lambda$, so its inverse stays uniformly bounded.  The metric
estimates follow by differentiating this smooth composition twice.  In
coordinates, $\Gamma=g^{-1}*Dg$, hence
$\|\Gamma_h-\Gamma_*\|_\infty\le Ch^2$ and
$\|D\Gamma_h-D\Gamma_*\|_\infty\le Ch$.  Because $g_h$ is $C^1$ across
interfaces there is no distributional surface term.  Substitution in
$\operatorname{Riem}=D\Gamma+\Gamma*\Gamma$ gives
\eqref{eq:supp-gowdy-full-curvature}.  The reference is vacuum, so contraction
with the uniformly bounded metric and inverse gives the Ricci estimate; finite
volume gives the $L^2$ curvature certificate.
\end{proof}

A finite positive realization is obtained by giving each predicted local state
a finite offset list containing zero and contained in
\eqref{eq:supp-gowdy-offset-cap}, a strictly positive reference row, and
acceptance probability $\exp[-E_j/\varepsilon]$ for a nonnegative deviation
cost.  If $p_{*,j}$ is the reference mass of the zero offset, the exact Gibbs
identity gives
\begin{equation}
 \mathbb E\sum_jE_j
 \le\varepsilon\left[
 \mathcal K_{\rm path}+\sum_j\log(1/p_{*,j})\right],
 \label{eq:supp-gowdy-gibbs-bound}
\end{equation}
where $\mathcal K_{\rm path}$ is the chain-rule relative entropy of the observed
history law against the declared accepted law.  Product reference masses count
all local trials; solver residuals, finite retry failures, and iteration errors
remain separate budgets.  The positive selection cost is not identified with
the signed ADM action.  Since the curvature estimates above are deterministic
and uniform over the permitted alphabet, every successful history inherits
\eqref{eq:supp-gowdy-curvature-certificate}.

\subsubsection{Measured proposal access and finite execution}

\paragraph{Measured access to the curvature-controlled histories.}
A finite alphabet with admissible offsets establishes realizability, but does
not by itself give a uniform probability of proposing those offsets.  The
following result separates this operational issue from the deterministic
curvature estimate.

\begin{proposition}[Finite proposal-access certificate]
\label{prop:main-safe-access}
At stage $j$, let $p_j$ be the measured finite proposal row and let
$\mathcal G_j$ contain only proposals satisfying the offset and solver bounds
of \Cref{thm:main-gowdy-regulator}.  Suppose
$p_j(\mathcal G_j)\ge\mu_j>0$ and $0\le E_j\le c_j$ on this set,
uniformly over the supported source histories.  Accept with the actual rule
$\one_{\mathcal G_j}\exp(-E_j/\varepsilon_j)$.  Then
\begin{equation}
 \mu_j e^{-c_j/\varepsilon_j}\le Z_j\le p_j(\mathcal G_j),
 \qquad
 D_{\rm KL}(Q_j\Vert p_j)\ge-\log p_j(\mathcal G_j).
 \label{eq:main-safe-proposal-bounds}
\end{equation}
With a genuine reset to the same source between retries, caps $R_j$ give
\begin{equation}
 \Pr(\hbox{some stage fails})
       \le\sum_j\exp[-R_j\mu_j e^{-c_j/\varepsilon_j}].
 \label{eq:main-safe-retry-bound}
\end{equation}
Every successful history inherits the pathwise curvature and independent
action-test estimates of \Cref{thm:main-gowdy-regulator}.  If the displayed
failure bounds are summable along a countable refinement, these conclusions
hold on every sufficiently fine successful refinement almost surely.
\end{proposition}

The proof in Appendix~\ref{supp:gowdy} also gives a sufficient proposal floor
from finite launch-share and completed-writer records.  A missing safe
proposal mass is a genuine obstruction which changing temperature cannot
remove.  These are bounds on trial counts; an imposed physical deadline must
also accommodate execution of the capped trials.  The selection rule and its
normalizer remain physical records, and its positive deviation cost is not
identified with the signed gravitational action.


The deterministic second-jet estimate is uniform over the stated offset
alphabet.  It therefore survives changes of the actual proposal weights,
provided the acceptance rule retains that support and successful proposals
really occur.  At a fixed stage and source history write
$\mu=p(\mathcal G)$ for the measured mass of the admissible set, and use
\[
 a(u)=\one_{\mathcal G}(u)e^{-E(u)/\varepsilon},\qquad
 Z=\sum_up(u)a(u),\qquad Q(u)=p(u)a(u)/Z.
\]
This is a declared selection rule, not an unrecorded replacement of the
original microscopic transition.

\begin{proof}[Proof of \Cref{prop:main-safe-access}]
Since $0\le E\le c$ on $\mathcal G$,
$\mu e^{-c/\varepsilon}\le Z\le\mu$.
For $\mu>0$, let $p_{\mathcal G}=p/\mu$ on $\mathcal G$.  Every
$Q\ll p$ supported on $\mathcal G$ obeys the exact identity
\[
 D_{\rm KL}(Q\Vert p)=D_{\rm KL}(Q\Vert p_{\mathcal G})-\log\mu
       \ge-\log\mu.
\]
In particular $\mu=0$ precludes such an accepted row altogether.
No choice of temperature can correct this support failure.

On the accepted support the complete trial table still reconstructs its cost:
$E(u)=-\varepsilon\log(ZQ(u)/p(u))$.  Averaging also gives
$\mathbb E_QE=-\varepsilon\log Z-
\varepsilon D_{\rm KL}(Q\Vert p)$.  These identities concern the recorded
positive deviation cost; the independently varied signed action is unchanged.

With independent reset retries conditional on the same source, the probability
that $R_j$ attempts fail is $(1-Z_j)^{R_j}\le e^{-R_jZ_j}$.
The same bound follows for history-dependent retries when the displayed lower
success bound holds conditionally after every failed attempt.  Taking the
supported lower bound for each stage and using a union bound proves
\eqref{eq:main-safe-retry-bound}.  For $J$ stages and total failure target
$0<\delta<1$, one sufficient finite choice is
\begin{equation}
 R_j=\left\lceil\mu_j^{-1}e^{c_j/\varepsilon_j}
                         \log(J/\delta)\right\rceil.
 \label{eq:supp-safe-retry-cap}
\end{equation}
Conditioning the entire path on success may reweight its intermediate source
states.  No equality with the unconditioned product of accepted rows is needed:
every successful path has the same deterministic admissible support, and hence
the same curvature and action bounds.  Summability of the complete failure
probabilities and the first Borel--Cantelli lemma give the final assertion,
without independence between refinements.
\end{proof}

\begin{lemma}[A finite execution certificate for the proposal floor]
\label{lem:supp-executed-proposal-floor}
At a retained source snapshot let a launch race have finitely many bins.
Conditional on no earlier launch, let $a_\ell$ be the probability of any
launch and $f_\ell$ the probability of a predictor-marked launch in bin
$\ell$.  Set $S_0=1$, $S_\ell=\prod_{i\le\ell}(1-a_i)$.
Suppose
\begin{equation}
 f_\ell\ge\rho_*a_\ell,\qquad
 -\log S_J\ge A_*>0,\qquad 0<\rho_*\le1.
 \label{eq:supp-launch-share}
\end{equation}
If, conditional on every predictor launch, the actually executed writer
returns a certified member of $\mathcal G$ before the attempt deadline with
probability at least $q_*>0$, then
\begin{equation}
 p(\mathcal G)\ge q_*\rho_*(1-e^{-A_*}).
 \label{eq:supp-executed-proposal-floor}
\end{equation}
\end{lemma}
\begin{proof}
The disjoint first-launch events give predictor launch probability
$\sum_{\ell=1}^J S_{\ell-1}f_\ell$.  Since
$S_{\ell-1}a_\ell=S_{\ell-1}-S_\ell$, this is at least
$\rho_*(1-S_J)\ge\rho_*(1-e^{-A_*})$.  On each such event the conditional
completed-safe-output probability is at least $q_*$.  Multiplication of
conditional probabilities and summation prove the claim.  Launch,
interruption, completion and output certificates belong to the actual
same-history record; a possible but unexecuted numerical predictor supplies
none of these probabilities.
\end{proof}

The execution probability and deadline in this lemma are physical inputs,
not consequences of a fast geometric root clock.  For example, one root clock
competing with $D$ clocks of equal rate has share $1/(D+1)$, regardless of
their common rate scale.  Moreover, a cap on the number of retries bounds their
physical duration only when each attempt has a certified duration bound.  With
attempt durations at most $d_j$, the allowed stage duration must accommodate
$R_jd_j$.  Thus the proposition removes the need to assume a separate
continuum proposal profile on a class with the displayed finite execution
budgets, without asserting a uniform computation or sampling cost for every
microscopic law.

For the general residual/action estimate below put
\begin{equation}
 \kappa_h:=\epsilon_0+h^2+\sqrt{\mathcal F_h}.
 \label{eq:supp-gowdy-kappa}
\end{equation}

Finally discretize \eqref{eq:supp-gowdy-adm-density} by midpoint time sampling
and centered spatial differences, keep $N$ and $\beta$ independent, and replace
$-2NR''$ by $2N'R'$ under the periodic boundary convention.  Evaluation at the
recorded areal-gauge fields gives, for every fixed smooth compactly supported
independent test $(\delta R,\delta\lambda,\delta P,\delta Q,\delta N,\delta\beta)$,
\begin{equation}
 |\delta S_h-\delta S_*|\le C\kappa_h.
 \label{eq:supp-gowdy-action-test}
\end{equation}
The estimate follows from the same first-jet comparison used above, and the
variation is taken before the areal-gauge restriction.  Hence
$\mathbb E\mathcal F_h=O(h^4)$ and $\epsilon_0=O(h^2)$ yield second-order mean
metric and action-test errors.  Exact preservation of
$\mathcal K_\ell$ is one constraint statement and is not substituted for the
independent lapse/shift action tests.

\begin{proof}[Proof of \Cref{thm:main-gowdy-regulator}]
Item~\textup{(G1)} is \Cref{thm:supp-gowdy-momentum}.  Item~\textup{(G2)} is
\Cref{prop:supp-gowdy-residual-control}.  The pathwise $C_h^1$ estimate and
\Cref{thm:supp-gowdy-full-curvature} give item~\textup{(G3)}.  The unreduced
variational calculation
\eqref{eq:supp-gowdy-adm-density}--\eqref{eq:supp-gowdy-areal-lift}, the action
test \eqref{eq:supp-gowdy-action-test}, and the Gibbs identity
\eqref{eq:supp-gowdy-gibbs-bound} give item~\textup{(G4)} while retaining the
lapse and shift variations before gauge restriction.
\end{proof}

\subsection{Flat and de Sitter regulator limits}
\label{supp:vacuum}

This section proves the spatial and Cartan-continuum parts of
\Cref{mt:vacuum}.  The temporal profiles are no longer posited here: they are
the stationary renewal solutions derived in
Appendix~\ref{supp:homogeneous-dynamics}.  The present appendix verifies that
the associated $A_3$ spatial generators have the required normalization and
compact screens, and that the reconstructed coframes and exact link transports
converge to the corresponding flat and de~Sitter Cartan geometries.

\paragraph{Stationary geometric profiles.}

For $\Lambda=0$, choose the constant stationary branch with $a_0=1$.
Then
\begin{equation}
 \mathcal B_h=I_3,
 \qquad \varrho_h=1,
 \qquad g_h=I_3,
 \qquad N_h=1,
 \qquad \beta_h=0,
 \label{eq:flat-adm-main}
\end{equation}
and the graph generators converge to $\tfrac12\Delta_{\mathbb T^3}$.  The
Lorentzian limit is
\begin{equation}
 (\R\times\mathbb T^3,
 -\dd t^2+\dd x_1^2+\dd x_2^2+\dd x_3^2).
 \label{eq:flat-metric-main}
\end{equation}
Its curvature and vacuum Einstein tensor vanish.

For $\Lambda>0$, choose the expanding stationary branch of
\Cref{cor:main-derived-desitter-rates} with $a_0=1$, $t_0=0$, and
$H=H_0=\sqrt{\Lambda/3}$.  Then
\begin{equation}
 \mathcal B_h(t)=e^{-2Ht}I_3,
 \qquad
 \varrho_h(t)=e^{3Ht},
 \qquad
 g_h(t)=e^{2Ht}I_3,
 \qquad
 N_h=1,
 \qquad
 \beta_h=0.
 \label{eq:desitter-adm-main}
\end{equation}
The continuum coframe and connection are
\begin{equation}
 \vartheta^0=\dd t,
 \qquad
 \vartheta^a=e^{Ht}\dd x^a,
 \qquad
 \omega^a{}_0=H\vartheta^a,
 \qquad
 \omega^a{}_b=0.
 \label{eq:desitter-cartan-main}
\end{equation}
Cartan's equations give zero torsion and constant sectional curvature $H^2$,
so
\begin{equation}
 g_H=-\dd t^2+e^{2Ht}
 (\dd x_1^2+\dd x_2^2+\dd x_3^2),
 \qquad
 G_{\mu\nu}+3H^2(g_H)_{\mu\nu}=0.
 \label{eq:desitter-einstein-main}
\end{equation}
The spatial quotient is globally hyperbolic and is the compact spatial quotient
of the flat de~Sitter slicing.

\begin{theorem}[Stationary flat and de Sitter renewal regulators]
\label{mt:vacuum}
For the operational predictive carrier and geometric branch of
\Cref{mt:grand-tensor,mt:adm}, the stationary homogeneous families above have
the following properties.
\begin{enumerate}[label=\textup{(\roman*)},leftmargin=2.8em]
\item their rate and density profiles solve the lapse-varied homogeneous renewal
      Euler system of \Cref{thm:main-homogeneous-friedmann} and arise as the
      continuum limits of the finite stationary system in
      \Cref{thm:main-finite-homogeneous-stationarity};
\item each spatial cutoff defines a normalized finite $A_3$ renewal law whose
      root moments and speed densities reconstruct the displayed ADM variables
      without a primitive metric;
\item the flat sequence converges to \eqref{eq:flat-metric-main}, with
      $\Lambda=0$, while the curved sequence converges to
      \eqref{eq:desitter-einstein-main}, with $\Lambda=3H^2$;
\item on every compact time slab the cut, Poincar\'e, low-energy counting, and
      compact-screen margins are uniform in $h$;
\item the flat curvature vanishes identically, while on each compact de~Sitter
      slab the reconstructed curvature and observer/transport coefficients are
      uniformly bounded; the initial-energy and integrated propagation/source
      budgets of \Cref{thm:main-curvature-propagation} therefore pass directly,
      and in fact the curvature is uniformly bounded in $L^\infty$;
\item along $h_n=2^{-n^2}$ with time mesh $h_n$ and screen $[-n,n]$, exact
      parallel transport of the reconstructed Levi--Civita connection gives
      metric-unitary links, while the normalized torsion, plaquette-curvature,
      and Palatini first-variation defects are $O_T(h_n)$ and summable.
\end{enumerate}
In addition, \Cref{thm:main-explicit-operational-family} supplies a nontrivial
finite positive accept/reject realization whose measured action and accepted
dynamics converge to the same expanding stationary branch.
\end{theorem}

The detailed homogeneous variational and operational proofs are in
Appendix~\ref{supp:homogeneous-dynamics}; the spatial generator, compact-screen,
Cartan-curvature, and exact-link estimates are in Appendix~\ref{supp:vacuum}.
The result is not an attractor theorem: it does not select $\Lambda$, the time
orientation, or the quadratic operational class from arbitrary microscopic
renewal data.


\subsubsection{The \texorpdfstring{$A_3$}{A3} identity and periodic regulators}

Let
\begin{equation}
 W_0=\left\{x\in\R^4:\sum_{i=1}^4x_i=0\right\},
 \qquad
 \Lambda_{A_3}=W_0\cap\Z^4,
 \label{eq:supp-a3-lattice}
\end{equation}
and let
\begin{equation}
 \Phi_{A_3}=\{e_i-e_j:1\le i\ne j\le4\}
 \label{eq:supp-a3-oriented-roots}
\end{equation}
be the twelve oriented roots.  The scalar product on $W_0$ is the restriction
of the standard Euclidean scalar product on $\R^4$.

\begin{lemma}[$A_3$ second-moment identity]
\label{lem:supp-a3-second-moment}
On $W_0$,
\begin{equation}
 \sum_{\alpha\in\Phi_{A_3}}
        \alpha\alpha^{\mathsf T}=8I_{W_0}.
 \label{eq:supp-a3-second-moment}
\end{equation}
Equivalently, for every $v\in W_0$,
\begin{equation}
 \sum_{i\ne j}(v_i-v_j)^2=8\norm{v}^2.
 \label{eq:supp-a3-quadratic-identity}
\end{equation}
\end{lemma}

\begin{proof}
For $v\in W_0$, expand
\[
 \sum_{i\ne j}(v_i-v_j)^2
 =\sum_{i\ne j}(v_i^2+v_j^2-2v_iv_j).
\]
Each $v_i^2$ occurs six times in the first two sums, while
\[
 \sum_{i\ne j}v_iv_j
 =\left(\sum_i v_i\right)^2-\sum_i v_i^2
 =-\sum_i v_i^2.
\]
Hence the displayed expression is
$6\sum_i v_i^2+2\sum_i v_i^2=8\norm{v}^2$.
Polarization gives \eqref{eq:supp-a3-second-moment}.
\end{proof}

For an integer $d\ge2$, put $h=d^{-1}$ and
\begin{equation}
 V_h=\Lambda_{A_3}/d\Lambda_{A_3}.
 \label{eq:supp-periodic-vertices}
\end{equation}
We regard $hx$, $x\in V_h$, as a periodic mesh of the flat torus
$\mathbb T^3_{A_3}=W_0/\Lambda_{A_3}$.  An oriented root $\alpha$ connects
$x$ to $x+\alpha$ in $V_h$.  Let $m_h>0$ be the mass of each vertex and let
$k_{\alpha,h}>0$ be the jump rate along the oriented root $\alpha$.  The
Markov generator on functions $f:V_h\to\C$ is
\begin{equation}
 (L_hf)(x)
 =\sum_{\alpha\in\Phi_{A_3}}k_{\alpha,h}
  \bigl[f(x+\alpha)-f(x)\bigr].
 \label{eq:supp-root-generator}
\end{equation}
Its positive Dirichlet form on
$L^2(V_h,m_h)$ is
\begin{equation}
 \mathcal E_h(f)
 =-\ip{f}{L_hf}_{L^2(m_h)}
 =\frac{m_h}{2}\sum_{x\in V_h}
   \sum_{\alpha\in\Phi_{A_3}}k_{\alpha,h}
   \abs{f(x+\alpha)-f(x)}^2,
 \label{eq:supp-root-form}
\end{equation}
whenever the rates are reversal symmetric.

\begin{lemma}[Fourier symbol and flat generator limit]
\label{lem:supp-flat-symbol}
Suppose
\begin{equation}
 k_{\alpha,h}=\frac1{8h^2}
 \qquad(\alpha\in\Phi_{A_3}).
 \label{eq:supp-flat-rates}
\end{equation}
For a dual-lattice mode $e_k(x)=e^{2\pi i h\langle k,x\rangle}$,
$-L_h$ has eigenvalue
\begin{equation}
 \lambda_h(k)
 =\frac1{4h^2}\sum_{a=1}^{6}
  \left[1-\cos\bigl(2\pi h\langle r_a,k\rangle\bigr)\right],
 \label{eq:supp-flat-symbol}
\end{equation}
where $r_1,\ldots,r_6$ is one orientation of the six root pairs.  For every
fixed $k$,
\begin{equation}
 \lambda_h(k)\longrightarrow2\pi^2\norm{k}^2,
 \label{eq:supp-flat-symbol-limit}
\end{equation}
and, under the natural piecewise-constant embeddings,
\begin{equation}
 L_h\longrightarrow\frac12\Delta_{\mathbb T^3_{A_3}}
 \label{eq:supp-flat-generator-limit}
\end{equation}
in strong resolvent sense and strongly at the semigroup level.
\end{lemma}

\begin{proof}
Pair the roots $\alpha$ and $-\alpha$ in
\eqref{eq:supp-root-generator}.  This gives
\eqref{eq:supp-flat-symbol}.  Taylor expansion yields
\[
 \lambda_h(k)
 =\frac{\pi^2}{2}\sum_{a=1}^6\langle r_a,k\rangle^2+O_k(h^2).
\]
Because the sum over all twelve roots is twice the sum over the six unoriented
roots, \Cref{lem:supp-a3-second-moment} gives
$\sum_{a=1}^6r_ar_a^{\mathsf T}=4I$, and therefore the leading term is
$2\pi^2\norm{k}^2$.

The trigonometric symbol is nonnegative and converges pointwise on the common
Fourier core.  On each fixed finite Fourier band the convergence is uniform.
The corresponding quadratic forms converge on trigonometric polynomials and
have the common lower bound zero.  Their closure gives Mosco convergence to
the form of $\tfrac12\Delta$; the resolvent and semigroup conclusions follow
from the Mosco--resolvent equivalence.  Equivalently, one may apply dominated
convergence to the Fourier multipliers and then use density.
\end{proof}

\begin{lemma}[Uniform spatial screens for the periodic $A_3$ graph]
\label{lem:supp-periodic-screens}
Let the vertex mass be $m_h=a^3h^3$ and the symmetric root rate be
$k_{\alpha,h}=a^{-2}/(8h^2)$, where $a$ ranges in a compact interval
$0<a_-\le a\le a_+<\infty$.  Then there are constants
$c_{\rm P},C_{\rm W},c_{\rm cut}>0$, depending only on
$a_-,a_+$ and the fixed torus, such that for all sufficiently small $h$:
\begin{enumerate}[label=\textup{(\roman*)},leftmargin=2.8em]
\item the first nonzero eigenvalue of $-L_h$ is at least $c_{\rm P}$;
\item the low-energy counting function satisfies
      \begin{equation}
       N_h(R):=\#\{\lambda\in\spec(-L_h):\lambda\le R\}
       \le C_{\rm W}(1+R^{3/2});
       \label{eq:supp-weyl-count}
      \end{equation}
\item with edge conductance
      $c_{e,h}=m_hk_{\alpha,h}=ah/8$, every vertex set $A$ with
      $0<\mu_h(A)\le\mu_h(V_h)/2$ obeys
      \begin{equation}
       h\,c_h(\partial A)
       \ge c_{\rm cut}\,\mu_h(A)^{2/3}.
       \label{eq:supp-periodic-cut}
      \end{equation}
\end{enumerate}
Consequently the spectral projections
$P_{h,R}=\one_{[0,R]}(-L_h)$ have uniformly bounded rank and
\begin{equation}
 \norm{(I-P_{h,R})f}_{L^2(m_h)}^2
 \le R^{-1}\mathcal E_h(f).
 \label{eq:supp-screen-tail}
\end{equation}
Thus they form a uniformly exhausted compact-screen family.
\end{lemma}

\begin{proof}
The scalar $a^{-2}$ multiplies the flat symbol
\eqref{eq:supp-flat-symbol}, and the compact interval for $a$ changes all
spectral comparisons only by fixed constants.  For $\abs{2\pi h\langle
r_a,k\rangle}\le\pi$, the inequalities
$2s^2/\pi^2\le1-\cos s\le s^2/2$, together with the root-frame identity, give
\begin{equation}
 c_1\norm{k}^2\le\lambda_h(k)\le c_2\norm{k}^2
 \label{eq:supp-symbol-comparison}
\end{equation}
for modes in a fundamental dual cell, with constants independent of $h$.
The first nonzero dual vector has fixed positive length, proving (i).
Counting dual-lattice points in a three-dimensional ball gives (ii).

For (iii), retain only the three basis-root directions
$\pm e_1,\pm e_2,\pm e_3$ from \eqref{eq:a3-roots-main}.  Their Cayley
subgraph is a periodic three-dimensional grid and each retained edge has
conductance $ah/8\ge a_-h/8$.  The discrete edge-isoperimetric inequality on
the $d^3$ periodic grid follows by slicing $A$ along the three coordinate
directions and applying the Loomis--Whitney inequality to the three coordinate
projections; in physical units it reads
$h^2\#\partial A\ge c_0(h^3\#A)^{2/3}$.  Since
$\mu_h(A)=a^3h^3\#A$ and $c_{e,h}=ah/8$, it follows that
$h c_h(\partial A)\ge[c_0/(8a_+)]\mu_h(A)^{2/3}$.
Adding the other $A_3$ edges only increases this boundary capacity, proving
\eqref{eq:supp-periodic-cut}.  Finally,
$-L_h\succeq R(I-P_{h,R})$ gives
\eqref{eq:supp-screen-tail} by the spectral theorem.
\end{proof}

\subsubsection{Flat renewal vacuum}

\begin{construction}[Flat periodic renewal regulator]
\label{con:supp-flat-regulator}
At mesh $h=d^{-1}$, take the persistent spatial state $Z_n\in V_h$, one point
in every nongeometric fibre, the twelve root branches
\eqref{eq:supp-flat-rates}, vertex mass
\begin{equation}
 m_h=h^3,
 \label{eq:supp-flat-mass}
\end{equation}
the identity root router, the identity slow row, and zero nonvacuum reward.
The event alphabet consists of the root-race outcomes together with the
protected orientation and one-point spectator fibres.  Its conditional law is
normalized directly, so the predictive carrier and scheduler records are
literal parts of the operational regulator.
\end{construction}

\begin{theorem}[Flat ADM and vacuum limit]
\label{thm:supp-flat-vacuum}
The regulator of \Cref{con:supp-flat-regulator} satisfies
\begin{equation}
 \mathcal B_h=I_3,\qquad
 \varrho_h=1,\qquad
 g_h=I_3,\qquad N_h=1,\qquad\beta_h=0.
 \label{eq:supp-flat-adm}
\end{equation}
Every undirected root edge has conductance $h/8$.  The spatial generators
converge to $\tfrac12\Delta_{\mathbb T^3_{A_3}}$, and the Lorentzian continuum
is
\begin{equation}
 (M_0,g_0)=
 \left(\R\times\mathbb T^3_{A_3},
 -\dd t^2+\dd x_1^2+\dd x_2^2+\dd x_3^2\right).
 \label{eq:supp-flat-lorentzian}
\end{equation}
Its Levi--Civita connection and curvature vanish, so
$G(g_0)=0$.  The finite torsion, curvature, and vacuum Palatini first-variation
residuals vanish identically.
\end{theorem}

\begin{proof}
The root identity and rates give
\[
 \mathcal B_h
 =h^2\sum_{\alpha\in\Phi_{A_3}}
   \frac1{8h^2}\alpha\alpha^{\mathsf T}=I_3.
\]
Reversal symmetry gives $\beta_h=0$.  The mass density per physical volume is
one, so $\varrho_h=1$.  Substitution in the ADM inversion theorem
\Cref{thm:supp-adm-frame} gives \eqref{eq:supp-flat-adm}.  Multiplication of
one directed rate by one vertex mass gives $h/8$.

Generator convergence is \Cref{lem:supp-flat-symbol}.  The coframe
$\vartheta^0=\dd t$, $\vartheta^a=\dd x^a$ is constant, and the exact identity
router gives the zero spin connection.  Hence
$T^A=\dd\vartheta^A+\omega^A{}_B\wedge\vartheta^B=0$ and
$\Omega^A{}_B=\dd\omega^A{}_B+\omega^A{}_C\wedge\omega^C{}_B=0$.
Therefore $\Ric=0$, $R=0$, and $G=0$.  Any face-based finite Cartan and
Palatini writer constructed from these exact constant data is identically zero,
as is its first variation.  Normalization and predictive persistence follow directly from the operational event law and \Cref{thm:supp-persistence-derived}.
\end{proof}

\subsubsection{Homogeneous de Sitter renewal vacuum}

Let $\Lambda>0$, $H=\sqrt{\Lambda/3}$, and choose the expanding proper-time
stationary branch of \Cref{cor:main-derived-desitter-rates} with
$a_0=1$ and $t_0=0$.  On a compact time slab the resulting renewal rates and
masses are
\begin{equation}
 k_{\alpha,h}(t)=\frac{e^{-2Ht}}{8h^2},
 \qquad
 m_h(t)=e^{3Ht}h^3,
 \label{eq:supp-desitter-rates-mass}
\end{equation}
with physical speed density $\varrho_h(t)=e^{3Ht}$.  These profiles are
stationary outputs of the homogeneous renewal action, not assumptions made in
this appendix.

\begin{theorem}[Reconstructed de Sitter ADM packet]
\label{thm:supp-desitter-adm}
For every $h$ and $t$,
\begin{equation}
 \mathcal B_h(t)=e^{-2Ht}I_3,
 \quad
 \varrho_h(t)=e^{3Ht},
 \quad
 g_h(t)=e^{2Ht}I_3,
 \quad
 N_h(t)=1,
 \quad
 \beta_h(t)=0.
 \label{eq:supp-desitter-adm}
\end{equation}
Each undirected edge has conductance
\begin{equation}
 c_{e,h}(t)=\frac{e^{Ht}}8h.
 \label{eq:supp-desitter-conductance}
\end{equation}
On every compact time slab, the cut, Poincar\'e, counting, and compact-screen
constants are uniform in $h$ and $t$.
\end{theorem}

\begin{proof}
By \Cref{lem:supp-a3-second-moment},
\[
 \mathcal B_h(t)
 =h^2\frac{e^{-2Ht}}{8h^2}
  \sum_{\alpha\in\Phi_{A_3}}\alpha\alpha^{\mathsf T}
 =e^{-2Ht}I_3.
\]
The rates are reversal symmetric, so $\beta_h=0$.  Formula
\eqref{eq:adm-inversion-main} gives
\[
 g_h=e^{2Ht}I_3,
 \qquad
 N_h=(e^{3Ht})^{1/3}(e^{-6Ht})^{1/6}=1.
\]
The conductance is the product of vertex mass and directed rate:
$e^{3Ht}h^3e^{-2Ht}/(8h^2)=e^{Ht}h/8$.  On $[-T,T]$, the scalar
$a(t)=e^{Ht}$ is bounded above and below by positive constants, so
\Cref{lem:supp-periodic-screens} applies uniformly.
\end{proof}

Let internal indices $A,B\in\{0,1,2,3\}$ be raised and lowered with
$\eta=\diag(-1,1,1,1)$.  Define
\begin{equation}
 \vartheta^0=\dd t,
 \qquad
 \vartheta^a=e^{Ht}\dd x^a,
 \qquad
 \omega^a{}_0=H\vartheta^a,
 \qquad
 \omega^a{}_b=0,
 \label{eq:supp-desitter-cartan}
\end{equation}
with $\omega_{AB}=-\omega_{BA}$.

\begin{theorem}[Cartan curvature and the vacuum Einstein equation]
\label{thm:supp-desitter-einstein}
The connection \eqref{eq:supp-desitter-cartan} is the torsion-free
Levi--Civita spin connection of
\begin{equation}
 g_H=-\dd t^2+e^{2Ht}\sum_{a=1}^3(\dd x^a)^2.
 \label{eq:supp-desitter-metric}
\end{equation}
Its curvature two-form is
\begin{equation}
 \Omega^A{}_B=H^2\vartheta^A\wedge\vartheta_B.
 \label{eq:supp-desitter-curvature}
\end{equation}
Consequently
\begin{equation}
 \Ric_{\mu\nu}=3H^2g_{\mu\nu},
 \qquad
 R=12H^2,
 \qquad
 G_{\mu\nu}+3H^2g_{\mu\nu}=0.
 \label{eq:supp-desitter-einstein}
\end{equation}
The slices $t=\mathrm{const}$ are compact Cauchy hypersurfaces, so the spatial
quotient is globally hyperbolic.
\end{theorem}

\begin{proof}
Since
$\dd\vartheta^a=H\vartheta^0\wedge\vartheta^a$,
\[
 T^a=\dd\vartheta^a+\omega^a{}_0\wedge\vartheta^0
 =H\vartheta^0\wedge\vartheta^a
  +H\vartheta^a\wedge\vartheta^0=0,
\]
and $T^0=0$.  The coframe is nondegenerate, so this metric-compatible
connection is Levi--Civita.

Cartan's second equation gives
\[
 \Omega^a{}_0=\dd(H\vartheta^a)=H^2\vartheta^0\wedge\vartheta^a,
 \qquad
 \Omega^a{}_b=\omega^a{}_0\wedge\omega^0{}_b
 =H^2\vartheta^a\wedge\vartheta^b,
\]
which is \eqref{eq:supp-desitter-curvature} in invariant notation.  Hence
$R_{ABCD}=H^2(\eta_{AC}\eta_{BD}-\eta_{AD}\eta_{BC})$,
so contraction gives $\Ric_{AB}=3H^2\eta_{AB}$ and $R=12H^2$.  Therefore
$G=\Ric-\tfrac12Rg=-3H^2g$, proving
\eqref{eq:supp-desitter-einstein}.

The function $t$ has timelike gradient and every inextendible causal curve
meets each compact slice $\{t\}\times\mathbb T^3_{A_3}$ exactly once; this is
the standard global-hyperbolicity criterion for a product metric with positive
scale factor.
\end{proof}

\subsubsection{Exact links and summable finite Cartan--Palatini defects}

Let $K_h$ be the spacetime cell complex obtained by combining the spatial root
mesh with a time mesh of size $h$.  For an oriented edge $e$, let
\begin{equation}
 U_{e,h}=\mathcal P\exp\left(-\int_e\omega\right)
 \label{eq:supp-exact-link}
\end{equation}
be exact parallel transport of the reconstructed continuum connection along
that edge.  Then
$U_{e,h}^{\mathsf T}\eta U_{e,h}=\eta$ exactly.  For an oriented face $p$ with
area $|p|$, define the normalized curvature defect
\begin{equation}
 \mathfrak r_{p,h}
 =|p|^{-1}\log U_{\partial p}+\Omega(x_p),
 \label{eq:supp-curvature-defect}
\end{equation}
where $x_p$ is the face midpoint and the logarithm is taken on the principal
chart for sufficiently small $h$.  Define the normalized torsion defect by the
covariant developed-edge closure
\begin{equation}
 \mathfrak t_{p,h}
 =|p|^{-1}\sum_{e\subset\partial p}
  U_{x_p\leftarrow s(e),h}
  \int_e\vartheta.
 \label{eq:supp-torsion-defect}
\end{equation}
Finally, let $\delta S_h(e_h,U_h)$ be the midpoint quadrature of the
Palatini--cosmological first variation and let
$\delta S(e,\omega)$ be its continuum value.

\begin{theorem}[First-order consistency of the exact-link regulator]
\label{thm:supp-finite-defects}
For every compact time slab $[-T,T]$ there is $C_T<\infty$, independent of
$h$, such that
\begin{equation}
 \sup_p\bigl(\norm{\mathfrak r_{p,h}}+
\norm{\mathfrak t_{p,h}}\bigr)
 +\sup_{\norm{v_h}_{C^1}\le1}
  \abs{\delta S_h[v_h]-\delta S[v_h]}
 \le C_Th.
 \label{eq:supp-finite-defect-bound}
\end{equation}
For the flat regulator the left-hand side is zero.  For the de Sitter
regulator, take
\begin{equation}
 h_n=2^{-n^2},
 \qquad
 T_n=n.
 \label{eq:supp-diagonal-exhaustion}
\end{equation}
Then every normalized defect is summable along the diagonal exhaustion:
\begin{equation}
 \sum_{n=1}^{\infty}C_{T_n}h_n<\infty.
 \label{eq:supp-defect-summability}
\end{equation}
\end{theorem}

\begin{proof}
On a compact slab, the coframe, connection, curvature, and the finitely many
derivatives needed below are bounded.  The nonabelian Stokes expansion for the
exact link gives
\[
 \log U_{\partial p}
 =-|p|\Omega(x_p)+O_T(h^3),
\]
because $|p|=O(h^2)$ and the first variation of the curvature across the face
is $O_T(h)$.  Division by $|p|$ gives the curvature part of
\eqref{eq:supp-finite-defect-bound}.  The covariant sum of developed edge
coframes is the integral of the torsion over the face plus an $O_T(h^3)$
quadrature remainder.  Since the continuum torsion vanishes, division by area
gives the torsion bound.

The Palatini--cosmological first variation is a finite sum of products of
coframe, curvature, and test-field coefficients.  Midpoint quadrature on each
cell has local error $O_T(h^5)$ for a four-dimensional cell of volume
$O(h^4)$; summing $O(h^{-4})$ cells on a fixed slab gives $O_T(h)$.
Equivalent estimates follow directly from Taylor expansion of the finite
wedge and holonomy writers.  In the flat case all coefficients are constant
and the exact links are identities, so the defects vanish.

For de Sitter, derivatives of order bounded independently of $n$ grow at most
exponentially on $[-n,n]$: $C_{T_n}\le C e^{cn}$ for fixed constants $C,c$.
Hence
$C_{T_n}h_n\le C\exp(cn-(\log2)n^2)$, whose series converges.
\end{proof}

\subsubsection{Proof of Theorem~\ref{mt:vacuum}}

\begin{proof}[Proof of Theorem~\ref{mt:vacuum}]
The predictive carrier at every cutoff is reconstructed by
\Cref{thm:supp-persistence-derived}.  The homogeneous temporal profiles and
their finite stationary origin are
\Cref{thm:main-homogeneous-friedmann,cor:main-derived-desitter-rates,thm:main-finite-homogeneous-stationarity};
the nonempty operational matching family is
\Cref{thm:main-explicit-operational-family}.  The spatial ADM reconstruction is
\Cref{thm:supp-flat-vacuum,thm:supp-desitter-adm}.  The strong-resolvent and
semigroup spatial limits follow from \Cref{lem:supp-flat-symbol}; the scalar
time-dependent factor is uniformly controlled on compact slabs.  Uniform cut,
Poincar\'e, counting, and compact-screen bounds are
\Cref{lem:supp-periodic-screens}.  The flat Einstein equation follows from the
zero connection in \Cref{thm:supp-flat-vacuum}.  The de Sitter Cartan and
Einstein calculation is \Cref{thm:supp-desitter-einstein}.  Exact metric-unitary
links and the $O_T(h)$, diagonally summable consistency estimates are
\Cref{thm:supp-finite-defects}.  These statements establish the six numbered clauses and the operational
realization asserted in the main theorem.
\end{proof}

\begin{remark}[Meaning of constructive vacuum emergence]
The temporal rate profiles used here are not chosen in order to reproduce a
target metric: on the homogeneous branch they solve the lapse-varied renewal
Euler system of Appendix~\ref{supp:homogeneous-dynamics}, and a nonempty finite
accept/reject class reconstructs the same reduced action from its measured
trial probabilities.  The present appendix verifies the downstream spatial and
Cartan geometry of those stationary profiles.  The result still does not claim
that an arbitrary renewal table selects the quadratic operational class, the
rank-one compatibility condition, the value of $\Lambda$, or the expanding
sign, nor that de~Sitter is a universal attractor under coarse graining.
\end{remark}


\begin{thebibliography}{99}

\bibitem{bombelli-sorkin}
L.~Bombelli, J.~Lee, D.~Meyer, and R.~D.~Sorkin,
\newblock Space-time as a causal set,
\newblock \emph{Physical Review Letters} \textbf{59} (1987), 521--524.

\bibitem{surya}
S.~Surya,
\newblock The causal set approach to quantum gravity,
\newblock \emph{Living Reviews in Relativity} \textbf{22} (2019), 5.

\bibitem{konopkamarkopouloseverini2008}
T.~Konopka, F.~Markopoulou, and S.~Severini,
\newblock Quantum graphity: a model of emergent locality,
\newblock \emph{Physical Review D} \textbf{77} (2008), 104029.

\bibitem{ambjornjurkiewiczloll2009}
J.~Ambj{\o}rn, J.~Jurkiewicz, and R.~Loll,
\newblock Quantum gravity, or the art of building spacetime,
\newblock in \emph{Approaches to Quantum Gravity},
D.~Oriti, ed.,
Cambridge University Press, Cambridge, 2009, pp.~341--359.

\bibitem{oriti2009}
D.~Oriti, ed.,
\newblock \emph{Approaches to Quantum Gravity: Toward a New Understanding of
Space, Time and Matter},
\newblock Cambridge University Press, Cambridge, 2009.

\bibitem{Connes1994}
A.~Connes,
\newblock \emph{Noncommutative Geometry},
\newblock Academic Press, San Diego, 1994.

\bibitem{jacobson1995}
T.~Jacobson,
\newblock Thermodynamics of spacetime: the Einstein equation of state,
\newblock \emph{Physical Review Letters} \textbf{75} (1995), 1260--1263.

\bibitem{verlinde2011}
E.~Verlinde,
\newblock On the origin of gravity and the laws of Newton,
\newblock \emph{Journal of High Energy Physics} \textbf{2011} (2011), no.~4,
029.

\bibitem{vanraamsdonk2010}
M.~Van Raamsdonk,
\newblock Building up spacetime with quantum entanglement,
\newblock \emph{General Relativity and Gravitation} \textbf{42} (2010),
2323--2329.

\bibitem{shalizicrutchfield}
C.~R.~Shalizi and J.~P.~Crutchfield,
\newblock Computational mechanics: pattern and prediction, structure and
simplicity,
\newblock \emph{Journal of Statistical Physics} \textbf{104} (2001),
817--879.

\bibitem{TraversCrutchfield2014}
N.~F.~Travers and J.~P.~Crutchfield,
\newblock Equivalence of history and generator $\epsilon$-machines,
\newblock \emph{Physical Review E} \textbf{89} (2014), 062107.

\bibitem{MarzenCrutchfield2017}
S.~E.~Marzen and J.~P.~Crutchfield,
\newblock Structure and randomness of continuous-time, discrete-event
processes,
\newblock \emph{Journal of Statistical Physics} \textbf{169} (2017),
303--315.

\bibitem{JurgensCrutchfield2021}
A.~M.~Jurgens and J.~P.~Crutchfield,
\newblock Divergent predictive states: the statistical complexity dimension
of stationary, ergodic hidden Markov processes,
\newblock \emph{Chaos} \textbf{31} (2021), 083114.

\bibitem{feller}
W.~Feller,
\newblock \emph{An Introduction to Probability Theory and Its Applications,
Vol.~II},
\newblock Wiley, New York, 1971.

\bibitem{garsialamperti1962}
A.~M.~Garsia and J.~Lamperti,
\newblock A discrete renewal theorem with infinite mean,
\newblock \emph{Commentarii Mathematici Helvetici} \textbf{37} (1962),
221--234.

\bibitem{erickson1970}
K.~B.~Erickson,
\newblock Strong renewal theorems with infinite mean,
\newblock \emph{Transactions of the American Mathematical Society}
\textbf{151} (1970), 263--291.

\bibitem{caravennadoney2019}
F.~Caravenna and R.~Doney,
\newblock Local large deviations and the strong renewal theorem,
\newblock \emph{Electronic Journal of Probability} \textbf{24} (2019),
Paper No.~72, 1--48.

\bibitem{Stinespring1955}
W.~F.~Stinespring,
\newblock Positive functions on $C^*$-algebras,
\newblock \emph{Proceedings of the American Mathematical Society}
\textbf{6} (1955), 211--216.

\bibitem{Choi1975}
M.-D.~Choi,
\newblock Completely positive linear maps on complex matrices,
\newblock \emph{Linear Algebra and its Applications} \textbf{10} (1975),
285--290.

\bibitem{strohmaier}
A.~Strohmaier,
\newblock On noncommutative and pseudo-Riemannian geometry,
\newblock \emph{Journal of Geometry and Physics} \textbf{56} (2006),
175--195.

\bibitem{franco}
N.~Franco,
\newblock Temporal Lorentzian spectral triples,
\newblock \emph{Reviews in Mathematical Physics} \textbf{26} (2014),
1430007.

\bibitem{franco-eckstein}
N.~Franco and M.~Eckstein,
\newblock An algebraic formulation of causality for noncommutative geometry,
\newblock \emph{Classical and Quantum Gravity} \textbf{30} (2013), 135007.

\bibitem{vandendungenpaschkerennie2013}
K.~van den Dungen, M.~Paschke, and A.~Rennie,
\newblock Pseudo-Riemannian spectral triples and the harmonic oscillator,
\newblock \emph{Journal of Geometry and Physics} \textbf{73} (2013),
37--55.

\bibitem{vandendungen2016}
K.~van den Dungen,
\newblock Krein spectral triples and the fermionic action,
\newblock \emph{Mathematical Physics, Analysis and Geometry}
\textbf{19} (2016), 4.

\bibitem{DevastatoFarnsworthLizziMartinetti2018}
A.~Devastato, S.~Farnsworth, F.~Lizzi, and P.~Martinetti,
\newblock Lorentz signature and twisted spectral triples,
\newblock \emph{Journal of High Energy Physics} \textbf{2018} (2018),
no.~3, 089.

\bibitem{DevastatoFilaciMartinettiSingh2020}
A.~Devastato, M.~Filaci, P.~Martinetti, and D.~Singh,
\newblock Actions for twisted spectral triple and the transition from the
Euclidean to the Lorentzian,
\newblock \emph{International Journal of Geometric Methods in Modern Physics}
\textbf{17} (2020), suppl.~1, 2030001.

\bibitem{nieuviarts2026a}
G.~Nieuviarts,
\newblock Emergence of time from a twisted spectral triple in
almost-commutative geometry,
\newblock \emph{Journal of Noncommutative Geometry} (2026), to appear;
arXiv:2512.15450.

\bibitem{nieuviarts2026b}
G.~Nieuviarts,
\newblock Emergence of Lorentz symmetry from an almost-commutative twisted
spectral triple,
\newblock \emph{International Journal of Geometric Methods in Modern Physics}
(2026), to appear; arXiv:2502.18105.

\bibitem{Bregman1967}
L.~M.~Bregman,
\newblock The relaxation method of finding the common point of convex sets and
its application to the solution of problems in convex programming,
\newblock \emph{USSR Computational Mathematics and Mathematical Physics}
\textbf{7} (1967), 200--217.

\bibitem{Chung1997}
F.~R.~K.~Chung,
\newblock \emph{Spectral Graph Theory},
\newblock CBMS Regional Conference Series in Mathematics, vol.~92,
American Mathematical Society, Providence, RI, 1997.

\bibitem{SaloffCoste2002}
L.~Saloff-Coste,
\newblock \emph{Aspects of Sobolev-Type Inequalities},
\newblock London Mathematical Society Lecture Note Series, vol.~289,
Cambridge University Press, Cambridge, 2002.

\bibitem{ACY1998}
A.~Anderson, Y.~Choquet-Bruhat, and J.~W.~York, Jr.,
\newblock Einstein--Bianchi hyperbolic system for general relativity,
\newblock arXiv:gr-qc/9710041, version 2 (1998).

\bibitem{KRS2015}
S.~Klainerman, I.~Rodnianski, and J.~Szeftel,
\newblock The bounded $L^2$ curvature conjecture,
\newblock \emph{Inventiones Mathematicae} \textbf{202} (2015), 91--216.

\bibitem{York1972}
J.~W.~York, Jr.,
\newblock Role of conformal three-geometry in the dynamics of gravitation,
\newblock \emph{Physical Review Letters} \textbf{28} (1972), 1082--1085.

\bibitem{GibbonsHawking1977}
G.~W.~Gibbons and S.~W.~Hawking,
\newblock Action integrals and partition functions in quantum gravity,
\newblock \emph{Physical Review D} \textbf{15} (1977), 2752--2756.

\bibitem{MarsdenWest2001}
J.~E.~Marsden and M.~West,
\newblock Discrete mechanics and variational integrators,
\newblock \emph{Acta Numerica} \textbf{10} (2001), 357--514.

\bibitem{Moncrief1981}
V.~Moncrief,
\newblock Global properties of Gowdy spacetimes with $T^3\times\mathbb R$ topology,
\newblock \emph{Annals of Physics} \textbf{132} (1981), 87--107.

\bibitem{Chiribella2009}
G.~Chiribella, G.~M.~D'Ariano, and P.~Perinotti,
\newblock Theoretical framework for quantum networks,
\newblock \emph{Physical Review A} \textbf{80} (2009), 022339.

\bibitem{Pollock2018}
F.~A.~Pollock, C.~Rodr\'{\i}guez-Rosario, T.~Frauenheim,
M.~Paternostro, and K.~Modi,
\newblock Non-Markovian quantum processes: Complete framework and efficient
characterization,
\newblock \emph{Physical Review A} \textbf{97} (2018), 012127.

\bibitem{BerstelReutenauer}
J.~Berstel and C.~Reutenauer,
\newblock \emph{Noncommutative Rational Series with Applications},
\newblock Encyclopedia of Mathematics and its Applications, vol.~137,
Cambridge University Press, Cambridge, 2010.

\bibitem{CorichiReview}
A.~Corichi, I.~Rubalcava-Garc\'{\i}a, and T.~Vuka\v{s}inac,
\newblock Actions, topological terms and boundaries in first order gravity:
A review,
\newblock arXiv:1604.07764 (2016).

\bibitem{SimonCompact}
J.~Simon,
\newblock Compact sets in the space $L^p(0,T;B)$,
\newblock \emph{Annali di Matematica Pura ed Applicata}
\textbf{146} (1986), 65--96.

\bibitem{Mattheij1985}
R.~M.~M.~Mattheij,
\newblock Decoupling and stability of algorithms for boundary value problems,
\newblock \emph{SIAM Review} \textbf{27} (1985), 1--44.

\bibitem{BatesLuZeng}
P.~W.~Bates, K.~Lu, and C.~Zeng,
\newblock Invariant foliations near normally hyperbolic invariant manifolds
for semiflows,
\newblock \emph{Transactions of the American Mathematical Society}
\textbf{352} (2000), 4641--4676.


\bibitem{Wimmer2012}
M.~Wimmer,
\newblock Algorithm 923: Efficient numerical computation of the Pfaffian for
 dense and banded skew-symmetric matrices,
\newblock \emph{ACM Transactions on Mathematical Software} \textbf{38}
 (2012), no.~4, Article~30, 1--17.

\bibitem{Bisio2011}
A.~Bisio, G.~Chiribella, G.~M.~D'Ariano, and P.~Perinotti,
\newblock Minimal computational-space implementation of multi-round quantum protocols,
\newblock \emph{Physical Review A} \textbf{83} (2011), 022325.

\bibitem{Lindblom2006}
L.~Lindblom, M.~A.~Scheel, L.~E.~Kidder, R.~Owen, and O.~Rinne,
\newblock A new generalized harmonic evolution system,
\newblock \emph{Classical and Quantum Gravity} \textbf{23} (2006), S447--S462.

\end{thebibliography}

\addtocontents{toc}{\protect\setcounter{tocdepth}{1}}
\end{document}
